% Options for packages loaded elsewhere
\PassOptionsToPackage{unicode}{hyperref}
\PassOptionsToPackage{hyphens}{url}
\PassOptionsToPackage{dvipsnames,svgnames,x11names}{xcolor}
\documentclass[
  spanish,
  12pt,
]{book}
\usepackage{xcolor}
\usepackage[margin = 2cm]{geometry}
\usepackage{amsmath,amssymb}
\setcounter{secnumdepth}{5}
\usepackage{iftex}
\ifPDFTeX
  \usepackage[T1]{fontenc}
  \usepackage[utf8]{inputenc}
  \usepackage{textcomp} % provide euro and other symbols
\else % if luatex or xetex
  \usepackage{unicode-math} % this also loads fontspec
  \defaultfontfeatures{Scale=MatchLowercase}
  \defaultfontfeatures[\rmfamily]{Ligatures=TeX,Scale=1}
\fi
\usepackage{lmodern}
\ifPDFTeX\else
  % xetex/luatex font selection
\fi
% Use upquote if available, for straight quotes in verbatim environments
\IfFileExists{upquote.sty}{\usepackage{upquote}}{}
\IfFileExists{microtype.sty}{% use microtype if available
  \usepackage[]{microtype}
  \UseMicrotypeSet[protrusion]{basicmath} % disable protrusion for tt fonts
}{}
\makeatletter
\@ifundefined{KOMAClassName}{% if non-KOMA class
  \IfFileExists{parskip.sty}{%
    \usepackage{parskip}
  }{% else
    \setlength{\parindent}{0pt}
    \setlength{\parskip}{6pt plus 2pt minus 1pt}}
}{% if KOMA class
  \KOMAoptions{parskip=half}}
\makeatother
\usepackage{color}
\usepackage{fancyvrb}
\newcommand{\VerbBar}{|}
\newcommand{\VERB}{\Verb[commandchars=\\\{\}]}
\DefineVerbatimEnvironment{Highlighting}{Verbatim}{commandchars=\\\{\}}
% Add ',fontsize=\small' for more characters per line
\usepackage{framed}
\definecolor{shadecolor}{RGB}{248,248,248}
\newenvironment{Shaded}{\begin{snugshade}}{\end{snugshade}}
\newcommand{\AlertTok}[1]{\textcolor[rgb]{0.94,0.16,0.16}{#1}}
\newcommand{\AnnotationTok}[1]{\textcolor[rgb]{0.56,0.35,0.01}{\textbf{\textit{#1}}}}
\newcommand{\AttributeTok}[1]{\textcolor[rgb]{0.13,0.29,0.53}{#1}}
\newcommand{\BaseNTok}[1]{\textcolor[rgb]{0.00,0.00,0.81}{#1}}
\newcommand{\BuiltInTok}[1]{#1}
\newcommand{\CharTok}[1]{\textcolor[rgb]{0.31,0.60,0.02}{#1}}
\newcommand{\CommentTok}[1]{\textcolor[rgb]{0.56,0.35,0.01}{\textit{#1}}}
\newcommand{\CommentVarTok}[1]{\textcolor[rgb]{0.56,0.35,0.01}{\textbf{\textit{#1}}}}
\newcommand{\ConstantTok}[1]{\textcolor[rgb]{0.56,0.35,0.01}{#1}}
\newcommand{\ControlFlowTok}[1]{\textcolor[rgb]{0.13,0.29,0.53}{\textbf{#1}}}
\newcommand{\DataTypeTok}[1]{\textcolor[rgb]{0.13,0.29,0.53}{#1}}
\newcommand{\DecValTok}[1]{\textcolor[rgb]{0.00,0.00,0.81}{#1}}
\newcommand{\DocumentationTok}[1]{\textcolor[rgb]{0.56,0.35,0.01}{\textbf{\textit{#1}}}}
\newcommand{\ErrorTok}[1]{\textcolor[rgb]{0.64,0.00,0.00}{\textbf{#1}}}
\newcommand{\ExtensionTok}[1]{#1}
\newcommand{\FloatTok}[1]{\textcolor[rgb]{0.00,0.00,0.81}{#1}}
\newcommand{\FunctionTok}[1]{\textcolor[rgb]{0.13,0.29,0.53}{\textbf{#1}}}
\newcommand{\ImportTok}[1]{#1}
\newcommand{\InformationTok}[1]{\textcolor[rgb]{0.56,0.35,0.01}{\textbf{\textit{#1}}}}
\newcommand{\KeywordTok}[1]{\textcolor[rgb]{0.13,0.29,0.53}{\textbf{#1}}}
\newcommand{\NormalTok}[1]{#1}
\newcommand{\OperatorTok}[1]{\textcolor[rgb]{0.81,0.36,0.00}{\textbf{#1}}}
\newcommand{\OtherTok}[1]{\textcolor[rgb]{0.56,0.35,0.01}{#1}}
\newcommand{\PreprocessorTok}[1]{\textcolor[rgb]{0.56,0.35,0.01}{\textit{#1}}}
\newcommand{\RegionMarkerTok}[1]{#1}
\newcommand{\SpecialCharTok}[1]{\textcolor[rgb]{0.81,0.36,0.00}{\textbf{#1}}}
\newcommand{\SpecialStringTok}[1]{\textcolor[rgb]{0.31,0.60,0.02}{#1}}
\newcommand{\StringTok}[1]{\textcolor[rgb]{0.31,0.60,0.02}{#1}}
\newcommand{\VariableTok}[1]{\textcolor[rgb]{0.00,0.00,0.00}{#1}}
\newcommand{\VerbatimStringTok}[1]{\textcolor[rgb]{0.31,0.60,0.02}{#1}}
\newcommand{\WarningTok}[1]{\textcolor[rgb]{0.56,0.35,0.01}{\textbf{\textit{#1}}}}
\usepackage{longtable,booktabs,array}
\newcounter{none} % for unnumbered tables
\usepackage{calc} % for calculating minipage widths
% Correct order of tables after \paragraph or \subparagraph
\usepackage{etoolbox}
\makeatletter
\patchcmd\longtable{\par}{\if@noskipsec\mbox{}\fi\par}{}{}
\makeatother
% Allow footnotes in longtable head/foot
\IfFileExists{footnotehyper.sty}{\usepackage{footnotehyper}}{\usepackage{footnote}}
\makesavenoteenv{longtable}
\usepackage{graphicx}
\makeatletter
\newsavebox\pandoc@box
\newcommand*\pandocbounded[1]{% scales image to fit in text height/width
  \sbox\pandoc@box{#1}%
  \Gscale@div\@tempa{\textheight}{\dimexpr\ht\pandoc@box+\dp\pandoc@box\relax}%
  \Gscale@div\@tempb{\linewidth}{\wd\pandoc@box}%
  \ifdim\@tempb\p@<\@tempa\p@\let\@tempa\@tempb\fi% select the smaller of both
  \ifdim\@tempa\p@<\p@\scalebox{\@tempa}{\usebox\pandoc@box}%
  \else\usebox{\pandoc@box}%
  \fi%
}
% Set default figure placement to htbp
\def\fps@figure{htbp}
\makeatother
\ifLuaTeX
\usepackage[bidi=basic,shorthands=off]{babel}
\else
\usepackage[bidi=default,shorthands=off]{babel}
\fi
\ifLuaTeX
  \usepackage{selnolig} % disable illegal ligatures
\fi
\setlength{\emergencystretch}{3em} % prevent overfull lines
\providecommand{\tightlist}{%
  \setlength{\itemsep}{0pt}\setlength{\parskip}{0pt}}
\usepackage[]{natbib}
\bibliographystyle{plainnat}
\usepackage{setspace}
%\linespread{1.5}
%\usepackage{tikz}
%\usetikzlibrary{shapes,arrows}
\usepackage{pdflscape}
\newcommand{\blandscape}{\begin{landscape}}
\newcommand{\elandscape}{\end{landscape}}
\renewcommand{\thesection}{\Roman{section}}
\renewcommand{\thesubsection}{\Alph{subsection}}
\usepackage{booktabs}
\usepackage{amsthm}
\makeatletter
\def\thm@space@setup{%
  \thm@preskip=8pt plus 2pt minus 4pt
  \thm@postskip=\thm@preskip
}
\makeatother
\usepackage{amsmath}
\usepackage[ruled,vlined,linesnumbered]{algorithm2e}
\usepackage{hyperref}
\usepackage{bookmark}
\IfFileExists{xurl.sty}{\usepackage{xurl}}{} % add URL line breaks if available
\urlstyle{same}
\hypersetup{
  pdftitle={Análisis de encuestas de hogares con R},
  pdfauthor={Andrés Gutiérrez, Stalyn Guerrero, Giovany Babativa-Márquez, Cristian Téllez},
  pdflang={es},
  colorlinks=true,
  linkcolor={blue},
  filecolor={Maroon},
  citecolor={Blue},
  urlcolor={Blue},
  pdfcreator={LaTeX via pandoc}}

\title{Análisis de encuestas de hogares con R}
\author{Andrés Gutiérrez\footnote{Asesor Regional en Estadísticas Sociales - Comisión Económica para América Latina y el Caribe (CEPAL) - \href{mailto:andres.gutierrez@cepal.org}{\nolinkurl{andres.gutierrez@cepal.org}}}, Stalyn Guerrero\footnote{Consultor - Comisión Económica para América Latina y el Caribe (CEPAL), \href{mailto:guerrerostalyn@gmail.com}{\nolinkurl{guerrerostalyn@gmail.com}}}, Giovany Babativa-Márquez\footnote{Profesor asociado - Departamento de Estadística, Facultad de Ciencias, Universidad Nacional de Colombia, Sede Bogotá - \href{mailto:jgbabativam@unal.edu.co}{\nolinkurl{jgbabativam@unal.edu.co}}}, Cristian Téllez\footnote{Consultor - Comisión Económica para América Latina y el Caribe (CEPAL), \href{mailto:cftellezp@unal.edu.co}{\nolinkurl{cftellezp@unal.edu.co}}}}
\date{2026-06-18}

\begin{document}
\maketitle

\renewcommand*\contentsname{Tabla de Contenido}
{
\setcounter{tocdepth}{1}
\tableofcontents
}
\listoffigures
\listoftables
\chapter*{Prefacio}\label{prefacio}
\addcontentsline{toc}{chapter}{Prefacio}

La versión online de este libro está licenciada bajo una \href{http://creativecommons.org/licenses/by-nc-sa/4.0/}{Licencia Internacional de Creative Commons para compartir con atribución no comercial 4.0}.

También puede descargar la versión en PDF aquí:
Descargar libro en PDF.

Este libro es el resultado de un compendio de las experiencias internacionales prácticas adquiridas por los autores en su asesoría a los países de América Latina y el Caribe desde la División de Estadísticas de la CEPAL.

\chapter*{Resumen}\label{resumen}
\addcontentsline{toc}{chapter}{Resumen}

Las encuestas de hogares son la principal fuente de información cuantitativa para medir indicadores sociales y económicos (pobreza, desigualdad, empleo, educación, entreo muchos otros) en América Latina y en el mundo. A partir de ellas se estiman indicadores que orientan el diseño, la evaluación y el seguimiento de políticas públicas, y también constituyen una fuente principal de información para el monitoreo de los Objetivos de Desarrollo Sostenible en el marco de la Agenda 2030 de las Naciones Unidas. La Comisión Económica para América Latina y el Caribe (CEPAL) gestionan repositorios de microdatos, entre ellos el Banco de Datos de Encuestas de Hogares (BADEHOG), que reúnen, armonizan y documentan encuestas de más de 18 países de la región desde el año 2000, precisamente porque la información que producen estas fuentes es de vital importancia para la toma de decisiones basada en evidencia.

Sin embargo, detrás de este aparente acceso democratizado a los datos existe un problema técnico estructural que con frecuencia pasa inadvertido, y es que la mayoría de las encuestas de hogares no proviene de muestras simples. Por el contrario, están sustentadas en diseños de muestreo complejos que incorporan estratificación, múltiples etapas, probabilidades de selección desiguales y ajustes de ponderación derivados de la calibración. Cuando un analista aplica los métodos estadísticos tradicionales diseñados para muestras independientes e idénticamente distribuidas sobre datos que no poseen esas características, las consecuencias son estimaciones sesgadas, errores estándar subestimados, intervalos de confianza demasiado estrechos y pruebas de hipótesis con tasas de error de tipo infladas.

A manera de ejemplo ilustrativo, suponga una población con dos regiones de tamaños muy distintos y con ingresos medios muy diferentes. Una encuesta que selecciona igual número de personas en cada región y analiza los datos sin ponderación producirá un estimador del ingreso medio nacional con un sesgo estructural, porque las unidades de la región pequeña tienen un peso empírico en la muestra superior a su peso real en la población. Este error se propaga también hacia la varianza. Un analista que calcule la varianza ignorando la estructura por conglomerados tratará a hogares del mismo sector censal como observaciones independientes, cuando en realidad comparten características comunes que reducen la información efectiva de la muestra.

Este documento busca cerrar esa brecha y proporcionar al investigador los fundamentos teóricos y las herramientas computacionales necesarias para trabajar adecuadamente con encuestas provenientes de diseños complejos. A lo largo del texto se abordan de manera integrada aspectos prácticos que van desde la definición del diseño de muestreo hasta la estimación de indicadores ---como totales, tamaños, medias y proporciones---, la construcción de sus respectivos intervalos de confianza, la formulación de pruebas de hipótesis, la modelación mediante técnicas de regresión avanzadas y el tratamiento de problemas habituales de imputación.

\chapter{Introducción}\label{cap1}

El análisis de encuestas con diseños complejos cuenta con una trayectoria consolidada en la literatura estadística, sustentada en el enfoque de inferencia basado en el diseño de muestreo. Como señalan \citet{sarndal2003model} y \citet{Gutierrez2016}, este enfoque reconoce a la distribución de probabilidad inducida por el proceso de selección de la muestra como el único mecanismo probabilístico que rige la inferencia. Bajo este paradigma, las propiedades estadísticas de los estimadores ---como el insesgamiento, la precisión y la consistencia--- se evalúan con respecto a la distribución de todas las posibles muestras que el diseño puede generar, sin imponer supuestos distribucionales sobre la variable de interés.

En este marco, la estimación de varianzas ha ocupado un lugar central y ha evolucionado en torno a tres grandes familias metodológicas complementarias. La primera, el método del último conglomerado \citep{hansen1953sample}, simplifica la estructura multinivel al concentrar la variabilidad en las unidades primarias de muestreo (UPM), tratando la selección como si se realizara con reemplazo. La segunda, la linealización de Taylor \citep{binder1983variances}, extiende este principio a parámetros no lineales ---como medias, proporciones y razones--- mediante una aproximación de primer orden que reduce el problema de estimación de la varianza al de un total. Finalmente, la tercera familia, los métodos de replicación (jackknife, bootstrap y balanced repeated replication), estima la varianza a partir de la dispersión entre estimaciones obtenidas mediante réplicas sistemáticas de la muestra.

Paralelamente, en el ámbito de la modelización estadística, contribuciones fundamentales como las de \citet{kish1974inference}, \citet{fuller1975regression} y \citet{binder1983variances} sentaron las bases de la regresión ponderada y de la estimación de varianzas de coeficientes en diseños estratificados y de múltiples etapas. Sobre este desarrollo, \citet{pfeffermann1993role} formalizó el método de máxima pseudo-verosimilitud (MPV), que incorpora los pesos muestrales en la función de log-verosimilitud para producir estimadores consistentes bajo el diseño, extendiendo así los modelos lineales generalizados de \citet{nelder1972generalized} al contexto de encuestas con diseño complejo.

Más recientemente, \citet{heeringa2017applied} han sintetizado estas contribuciones en un marco analítico unificado que integra de manera coherente el diseño muestral, la ponderación y la modelización, constituyéndose en una referencia estándar para la práctica contemporánea del análisis de encuestas complejas.

Este documento está dirigido a lectores con formación en estadística o ciencia de datos, o con experiencia profesional equivalente. Está especialmente pensado para quienes cuentan con conocimientos de los fundamentos de la teoría de muestreo ---como probabilidades de inclusión, estimadores e intervalos de confianza---, de modelos de regresión lineal y logística, y con un manejo introductorio del lenguaje de programación R.

En este contexto, el documento se posiciona como un texto de nivel intermedio-avanzado, orientado a profundizar y consolidar estos conocimientos, facilitando su aplicación rigurosa en contextos complejos y de alta exigencia técnica. Su propósito es acompañar al investigador en la transición desde los fundamentos hacia su implementación adecuada en la producción de estadísticas oficiales, la investigación académica en ciencias sociales, la evaluación de políticas públicas y el análisis de datos bajo estándares internacionales de calidad.

\section{Justificación}\label{justificacion}

A pesar de la riqueza metodológica disponible en la literatura especializada, persiste una brecha significativa entre el estado del arte y la práctica cotidiana del análisis de encuestas en la región. Una de las principales manifestaciones de esta brecha se relaciona con el acceso y uso del software estadístico. Algunos textos de referencia ampliamente reconocidos en análisis de encuestas, como \citet{heeringa2017applied}, \citet{cochran1977sampling} y \citet{sarndal2003model}, constituyen fuentes fundamentales desde el punto de vista teórico; sin embargo, están orientados principalmente hacia software con licencia, como \texttt{SAS}, \texttt{Stata} o \texttt{SPSS}, o bien no logran integrarse plenamente con el ecosistema computacional de \texttt{R}. En contraste, muchos tutoriales disponibles en línea sobre análisis de encuestas en \texttt{R} suelen privilegiar los aspectos operativos del software, pero rara vez alcanzan el nivel de profundidad conceptual necesario para un uso adecuado y responsable en la producción de estadísticas oficiales.

Adicionalmente, existe una fragmentación metodológica importante entre los distintos componentes que intervienen en el análisis de encuestas complejas. Con frecuencia, los fundamentos del diseño muestral, la gestión computacional de los datos, la estimación de parámetros descriptivos e inferenciales con corrección por diseño y el tratamiento de la no respuesta son abordados de manera separada, dificultando la construcción de flujos de trabajo integrales y reproducibles. En este contexto, el presente texto propone una integración coherente de estos elementos dentro de un único marco analítico implementado completamente en \texttt{R}.

Otro aspecto relevante corresponde a la contextualización regional de los ejemplos y aplicaciones. Una proporción considerable de los textos de referencia utiliza bases de datos y ejemplos provenientes de encuestas de países anglosajones, lo que limita parcialmente su aplicabilidad inmediata en América Latina. Con el propósito de acercar los contenidos a las realidades estadísticas de la región, este documento emplea una base de datos nativa de \texttt{R}, diseñada para replicar la estructura tradicional de las encuestas de hogares latinoamericanas y, por tanto, proporcionar ejemplos más cercanos a los problemas enfrentados por las oficinas nacionales de estadística y los equipos técnicos de la región.

Finalmente, este documento constituye en sí mismo un ejercicio de investigación reproducible. Todo su contenido, incluyendo tablas, figuras, ecuaciones y resultados numéricos, es generado íntegramente mediante código en R incorporado directamente en los capítulos. De esta manera, el lector no solo puede estudiar los métodos presentados, sino también replicarlos, modificarlos y adaptarlos a sus propios contextos y bases de datos, fortaleciendo así las capacidades analíticas y la transparencia en la producción estadística.

\section{Uso del software R para el análisis de encuestas}\label{uso-del-software-r-para-el-analisis-de-encuestas}

La elección de R como entorno computacional para este documento no es arbitraria. R reúne al menos tres características que lo convierten en una opción idónea para el análisis de encuestas de hogares en el contexto latinoamericano.

En primer lugar, su condición de acceso libre. A diferencia de otros programas con costos de licencia que pueden resultar prohibitivos para muchas oficinas nacionales de estadística en países en desarrollo, R se distribuye bajo licencia GNU, lo que garantiza que los métodos presentados en este documento sean plenamente accesibles y replicables, independientemente de las restricciones presupuestarias de las instituciones.

En segundo lugar, destaca su ecosistema especializado. El paquete \texttt{survey} \citep{Lumley2024}, disponible desde 2004 y en constante evolución, implementa una amplia gama de métodos de estimación para diseños complejos: estimadores de Horvitz--Thompson, técnicas de calibración, linealización de Taylor, métodos de replicación, modelos de regresión ponderados y pruebas de hipótesis ajustadas al diseño (como Rao-Scott), entre otros. A su vez, el paquete \texttt{srvyr} \citep{Freedman2024} extiende estas capacidades al incorporar la sintaxis del ecosistema \texttt{tidyverse}, facilitando una transición fluida entre el análisis exploratorio y el inferencial. Paquetes complementarios como \texttt{TeachingSampling} \citep{GutierrezT2020} y \texttt{convey} \citep{convey2024} consolidan un entorno analítico difícil de igualar en otros lenguajes.

Finalmente, R ofrece una sólida integración con herramientas de reproducibilidad científica. En particular, su combinación con R Markdown y el sistema \texttt{bookdown} permite generar documentos en los que resultados, tablas y figuras se producen directamente a partir del código fuente. Esta capacidad resulta especialmente relevante en el ámbito de la estadística oficial, donde la trazabilidad y la transparencia metodológica son requisitos fundamentales.

\section{Estructura del documento}\label{estructura-del-documento}

Este documento está organizado siguiendo una progresión lógica que va desde los fundamentos conceptuales del muestreo hasta las aplicaciones más avanzadas de modelamiento estadístico y tratamiento de la no respuesta. Esta secuencia no es arbitraria: cada nivel de conocimiento constituye un requisito necesario para comprender el siguiente, de modo que el lector pueda transitar de los principios básicos hacia herramientas analíticas más sofisticadas de manera coherente.

En primer lugar, se abordan los fundamentos del diseño muestral. El análisis adecuado de una encuesta comienza necesariamente con la comprensión de su diseño probabilístico. Antes de estimar cualquier parámetro, el analista debe familiarizarse con tres elementos constitutivos: la población objetivo, entendida como el conjunto de unidades sobre las cuales se desean hacer inferencias; el marco de muestreo, que corresponde a la enumeración operacional de dichas unidades; y las unidades de muestreo en cada etapa del diseño. Estos componentes determinan tanto la validez como el alcance de las inferencias posteriores.

A continuación, se introducen los elementos computacionales que permiten materializar estos conceptos en un entorno aplicado. El documento incorpora el manejo y transformación de variables mediante un ambiente de programación adecuado para la incorporación del diseño de muestreo, que constituye el eje central del flujo de análisis y garantiza que todas las operaciones estadísticas respeten dicha estructura.

Con el diseño correctamente especificado, el documento avanza hacia la estimación de parámetros descriptivos. Esta sección cubre los estimadores más utilizados en el análisis de encuestas para variables continuas, tales como ingresos, gastos o activos, incluyendo totales, medias y razones. Asimismo, se presentan sus correspondientes medidas de precisión (varianzas, errores estándar e intervalos de confianza) ajustadas a la complejidad del diseño muestral.

Posteriormente, se desarrolla el modelamiento estadístico, extendiendo el análisis más allá de la estadística descriptiva hacia un enfoque analítico e inferencial. En particular, se abordan los modelos de regresión lineal bajo diseños complejos, donde la incorporación explícita de los pesos muestrales resulta fundamental para obtener estimadores insesgados respecto al diseño. Este contexto pone de manifiesto la discusión entre el enfoque basado en el diseño, que prioriza las propiedades estadísticas a nivel poblacional, y el enfoque basado en los modelos, que depende de la correcta especificación funcional. En consecuencia, se enfatiza la importancia de evaluar la sensibilidad de los resultados mediante comparaciones entre estimaciones con y sin ponderación. Adicionalmente, se introducen los modelos lineales generalizados y los modelos multinivel, que amplían el alcance del análisis a variables respuesta no normales, como variables binarias (regresión logística), conteos (regresión de Poisson) y proporciones.

La organización de los capítulos sigue la arquitectura conceptual descrita anteriormente, avanzando de los fundamentos hacia las aplicaciones más complejas. De esta forma, el capítulo 2 presenta los conceptos básicos para el análisis de encuestas de hogares, destacando que la validez de las estimaciones depende de incorporar adecuadamente el diseño muestral. Se explican los elementos centrales de una encuesta, como la población objetivo, el marco muestral, las unidades de muestreo, la estratificación, la conglomeración y los pesos de muestreo. Además, el capítulo introduce los principales estimadores de muestreo para totales, medias y otros parámetros poblacionales y desarrolla los fundamentos para estimar la varianza e intervalos de confianza.

En el capítulo 3 se aborda la estimación rigurosa de parámetros descriptivos para variables numéricas bajo diseños complejos y se desarrollan técnicas para estimar totales y medias, junto con la estimación de varianzas mediante el método del último conglomerado y la construcción de intervalos de confianza con grados de libertad ajustados. Asimismo, se presentan enfoques alternativos para la estimación de la varianza, como la linealización de Taylor, las ecuaciones de estimación y los métodos de replicación, incluyendo bootstrap y jackknife. El capítulo también cubre la estimación de cuantiles ponderados, el análisis de desigualdad mediante indicadores como el índice de Gini y la curva de Lorenz, la estimación en dominios y subpoblaciones, y la visualización de resultados con intervalos de confianza.

En el capítulo 4 se aborda el análisis de variables categóricas bajo diseños complejos, incluyendo la evaluación de asociaciones entre categorías mediante métodos que incorporan adecuadamente el diseño muestral. Se desarrollan técnicas para estimar tamaños poblacionales y proporciones a partir de ponderaciones, junto con la estimación de su variabilidad mediante linealización y la construcción de intervalos de confianza, incluyendo enfoques robustos para proporciones extremas. Asimismo, se presentan tablas de contingencia ponderadas con sus respectivas frecuencias y proporciones, así como pruebas de independencia que ajustan la estadística clásica para reflejar el efecto del diseño. El capítulo también cubre la estimación e interpretación de medidas de asociación como la razón de odds, el análisis de diferencias de proporciones entre subpoblaciones mediante contrastes, y diversas estrategias de visualización, incluyendo gráficos de barras con intervalos de confianza y mapas.

El capítulo 5 se enfoca en analizar relaciones entre variables y construir modelos predictivos o explicativos válidos bajo diseños complejos, cuando los supuestos clásicos de independencia e idéntica distribución no se cumplen. Se presenta la evolución del problema desde los aportes iniciales hasta los desarrollos más recientes, y discute las implicaciones de trabajar con datos de encuestas en el contexto de los modelos de regresión, incluyendo la violación de los supuestos del modelo lineal clásico. Se contrasta el enfoque basado en el diseño con el basado en los modelos. El capítulo también aborda la realización de pruebas de hipótesis con ajustes apropiados, el diagnóstico de residuales y la identificación de observaciones influyentes.

En el capítulo 6 se extiende la modelación estadística a variables que no siguen una distribución normal (como variables binarias, multinomiales y estirctamente positivas) en el contexto de diseños de muestreo complejo. Para ello, se construye un marco unificado basado en los modelos lineales generalizados, incorporando sus componentes fundamentales y se discute la dualidad inferencial en poblaciones finitas, donde coexisten la medida de probabilidad del modelo y la del diseño muestral. Asimismo, se contrasta el enfoque clásico de máxima verosimilitud con el de máxima pseudo-verosimilitud, basado en ecuaciones de estimación ponderadas que incorporan los pesos muestrales, y se estudia la estimación de la varianza mediante técnicas de linealización.

En el capítulo 7 se revisan los modelos multinivel, también llamados jerárquicos o de efectos mixtos, para analizar datos con estructura anidada incorporando simultáneamente variables a nivel individual y contextual. A través de ejemplos de regresión (con interceptos y pendientes que varían por estrato), se muestra que ignorar la estructura jerárquica puede llevar a inferencias incorrectas, mientras que modelarla explícitamente permite capturar heterogeneidad entre grupos y aprovechar información compartida mediante. El enfoque se formaliza con el modelo nulo, que descompone la varianza en componentes dentro y entre estratos, y permite calcular la correlación intraclase (ICC), indicador clave para cuantificar la dependencia intragrupo.

Adicionalmente, el primer apéndice introduce el manejo de bases de datos en \texttt{R} mediante el entorno \texttt{tidyverse}, con énfasis en el paquete \texttt{dplyr}. A partir de la base de datos \texttt{BigCity}, se presenta un flujo básico de trabajo que incluye la carga de librerías, la inspección inicial de la base, el uso del operador de encadenamiento \texttt{\%\textgreater{}\%} y la aplicación de los principales verbos de \texttt{dplyr}: \texttt{filter()}, \texttt{select()}, \texttt{arrange()}, \texttt{mutate()}, \texttt{summarise()} y \texttt{group\_by()}. El capítulo muestra cómo seleccionar registros y variables, ordenar datos, crear nuevas variables y producir resúmenes descriptivos, explicando tanto la lógica de los códigos como la interpretación de sus salidas.

Finalmente, el segundo apéndice aborda la inferencia en poblaciones finitas y la diferencia entre la inferencia basada en modelos y la inferencia basada en el diseño muestral. A partir de ejemplos de simulación, muestra que en encuestas complejas la aleatoriedad proviene principalmente de la selección de la muestra y que, por tanto, las estimaciones deben incorporar las probabilidades de inclusión y los pesos muestrales. Se presenta el método de máxima verosimilitud como punto de partida para modelos estadísticos bajo supuestos de observaciones independientes, y luego presenta la máxima pseudo-verosimilitud como una extensión adecuada para datos de encuestas, donde cada contribución individual se pondera según el diseño.

El lector que los recorra en orden los capítulos habrá construido la capacidad de tomar cualquier encuesta de hogares latinoamericana, entender su diseño, procesar sus datos, estimar los indicadores de interés con la precisión estadística correcta, modelar las relaciones entre variables de manera válida bajo el diseño complejo, y manejar la no respuesta con métodos apropiados. Esa capacidad es, en esencia, lo que la producción rigurosa de estadísticas sociales exige.

\chapter{Conceptos básicos en encuestas de hogares}\label{cap2}

Las encuestas de hogares constituyen una de las principales herramientas para comprender la realidad social y económica de una población. A partir de ellas se obtiene información esencial sobre las condiciones de vida, el empleo, la educación, la salud y otros aspectos que orientan la formulación de políticas públicas y la toma de decisiones basadas en evidencia. El análisis riguroso de estas encuestas es fundamental para estudiar la realidad social, económica y demográfica de una población. Sin embargo, la validez de los resultados no depende únicamente del tamaño de la muestra o de la calidad de la recolección de datos, sino también de la manera en que la encuesta fue diseñada y de cómo dicho diseño se incorpora en el análisis estadístico. Ignorar el diseño muestral y aplicar métodos de análisis tradicionales que asumen una muestra aleatoria simple puede conducir a estimaciones sesgadas, errores en la inferencia y conclusiones equivocadas. Por esta razón, el análisis de las encuestas de hogares no puede desligarse de su diseño de selección, ya que este constituye el fundamento que sustenta la validez y representatividad de la información producida.

En la práctica, las encuestas de hogares suelen utilizar diseños complejos que combinan la estratificación con la selección de informantes en varias etapas. Este tipo de diseños busca optimizar recursos y mejorar la precisión de las estimaciones, aunque puede generar probabilidades de inclusión desiguales. Ignorar estas características puede conducir a estimaciones sesgadas, errores estándar incorrectos y conclusiones erróneas. Como señalan \citet{sarndal2003model} y \citet{Gutierrez2016}, toda encuesta parte de tres conceptos fundamentales: la población objetivo, el marco muestral y la muestra seleccionada. Estos elementos constituyen la base del diseño de la encuesta y son esenciales para garantizar la validez de las inferencias y la representatividad de los resultados.

En este contexto, un aspecto fundamental es la adecuada consideración del diseño muestral. Para obtener conclusiones válidas sobre la población, es necesario adoptar un enfoque de inferencia basado en el diseño, el cual reconoce que la muestra no proviene de una selección aleatoria cualquiera, sino de un plan probabilístico cuidadosamente definido. En este esquema, cada unidad de la población posee una probabilidad conocida y distinta de cero de ser seleccionada, lo que constituye la principal garantía de que los resultados puedan generalizarse a la población de referencia.

Bajo este enfoque, es posible demostrar que las estimaciones son insesgadas (o aproximadamente insesgadas), con respecto al esquema de muestreo, sin necesidad de asumir distribuciones específicas para la variable de interés. Esta característica confiere a la inferencia basada en el diseño un carácter robusto y ampliamente aceptado en el análisis de encuestas de hogares. Un componente central de este proceso son los pesos de muestreo, los cuales indican cuántas unidades de la población representa cada unidad seleccionada. Estos pesos permiten ajustar las estimaciones a las particularidades del diseño, garantizando que los resultados reflejen adecuadamente la estructura de la población.

Como muestran \citet{korn1995analysis}, las estimaciones ponderadas y no ponderadas pueden diferir sustancialmente, lo que evidencia la importancia de utilizar métodos de análisis coherentes con el diseño de la encuesta. En consecuencia, considerar adecuadamente los pesos de muestreo, la estratificación y las unidades de selección no constituye un aspecto técnico secundario, sino un requisito esencial para garantizar la credibilidad de los resultados.

\citet{UNhandbook} menciona el siguiente ejemplo. Suponga un país conformado por dos regiones: la Región A, con 100 habitantes e ingreso promedio de \$10.000, y la Región B, con 900 habitantes e ingreso promedio de \$2.000. El ingreso promedio verdadero de la población es:

\[
\theta = \frac{(100 \times 10.000) + (900 \times 2.000)}{100 + 900} = 2.800
\]

Si se seleccionan 50 personas en cada región y se ignora el diseño de muestreo, asignando el mismo peso a todas las observaciones, el promedio estimado sería:

\[
\hat{\theta} = \frac{(50 \times 10.000) + (50 \times 2.000)}{100} = 6.000
\]

En este caso, la estimación sobrestima considerablemente el ingreso nacional, ya que la Región A, que representa únicamente el 10\% de la población, termina influyendo tanto como la Región B, que concentra el 90\%. En cambio, al aplicar pesos proporcionales al tamaño poblacional (1 para la Región A y 9 para la Región B), se obtiene la siguiente estimación:

\[
\hat{\theta} = \frac{(1 \times 50 \times 10.000) + (9 \times 50 \times 2.000)}{(1 \times 50) + (9 \times 50)} = 2.800
\]

Este resultado coincide con el valor verdadero de la población y muestra cómo el uso adecuado de los pesos corrige el sesgo introducido por un análisis que ignora el diseño muestral.

\section{Población objetivo y unidades de muestreo}\label{poblacion-objetivo-y-unidades-de-muestreo}

\citet{gambino2009sampling} menciona que, desde la popularización de las encuestas de hogares en la década de 1940, se han evidenciado diversas tendencias asociadas a los avances tecnológicos tanto en las agencias estadísticas como en la sociedad en general, las cuales se aceleraron con la introducción del computador. En este contexto, el muestreo surge como una respuesta a la necesidad de obtener información estadística precisa sobre una población objetivo sin recurrir a la realización de un censo completo. Como señala \citet{Gutierrez2016}, el muestreo consiste en realizar investigaciones parciales sobre una población con el propósito de inferir resultados para el conjunto total.

En las últimas décadas, esta metodología se ha consolidado en distintos campos, especialmente en el sector gubernamental, mediante la producción de estadísticas oficiales que permiten el seguimiento de políticas públicas y de los Objetivos de Desarrollo Sostenible. Asimismo, el uso del muestreo se ha extendido al ámbito académico, al sector privado y a los medios de comunicación, donde constituye una herramienta fundamental para la generación y análisis de información.

Toda encuesta se encuentra asociada a una población finita compuesta por individuos o elementos sobre los cuales se desea obtener información. El conjunto de unidades sobre las que se producirán estimaciones y resultados recibe el nombre de población objetivo. Asimismo, las encuestas definen unidades de análisis, que corresponden a los distintos niveles de desagregación para los cuales se presentan estadísticas de interés, como personas, hogares o viviendas.

En las encuestas de hogares de América Latina es común utilizar diseños de muestreo en varias etapas. Para ello, se seleccionan diferentes unidades de muestreo que permiten llegar finalmente a los hogares que harán parte de la muestra. Por ejemplo, las unidades primarias de muestreo (UPM) pueden corresponder a sectores censales construidos a partir del último censo demográfico, mientras que las unidades secundarias de muestreo (USM) pueden ser las viviendas ubicadas dentro de dichos sectores.

Para realizar un proceso de selección sistemática de hogares es indispensable contar con un marco muestral que sirva como enlace entre las unidades de muestreo y los hogares que conforman la población objetivo. Este marco corresponde al conjunto en el que se identifican todos los elementos que integran la población de estudio y a partir del cual se selecciona la muestra. En las encuestas de hogares con diseños muestrales complejos, los marcos de muestreo suelen estar basados en áreas geográficas. Es decir, se construyen a partir de estructuras espaciales que vinculan a los hogares o personas con zonas delimitadas del territorio. Este tipo de marco facilita la organización territorial en unidades más manejables y permite implementar procesos de selección en varias etapas, manteniendo los principios del muestreo probabilístico.

Otra característica fundamental de las encuestas de hogares es la conglomeración. En la práctica, muchas encuestas seleccionan primero unidades primarias de muestreo, como sectores censales o áreas de enumeración, y posteriormente seleccionan hogares dentro de cada una de ellas. Este tipo de diseño permite reducir costos operativos y facilitar el trabajo de campo; sin embargo, también tiene implicaciones importantes sobre la precisión de las estimaciones. Cuando los hogares pertenecientes a un mismo conglomerado comparten características similares, la información adicional que aporta cada nueva observación dentro del conglomerado disminuye.

Por ejemplo, supóngase una encuesta que selecciona 100 conglomerados y, dentro de cada uno, 10 hogares, obteniendo así una muestra total de 1.000 hogares. Si los hogares de un mismo conglomerado presentan comportamientos muy parecidos (por ejemplo, si todos cuentan con acceso a electricidad), la variabilidad efectiva de la muestra se reduce considerablemente. En términos prácticos, la precisión obtenida podría ser equivalente a la de una muestra aleatoria simple de apenas 100 hogares. Analizar los 1.000 hogares como si fueran observaciones completamente independientes, ignorando la estructura de conglomeración, conduce a una subestimación de los errores estándar y, en consecuencia, a intervalos de confianza artificialmente estrechos y conclusiones potencialmente erróneas sobre la significancia estadística de los resultados.

\section{Estimadores de muestreo}\label{estimadores-de-muestreo}

Al analizar datos de encuestas, también es imprescindible definir el parámetro de interés: un valor numérico fijo que describe una característica de toda la población, denotada como \(U\). Entre los parámetros más comunes se encuentran las proporciones, los tamaños, los totales, las medias y las razones, entre otros. Dado que, en la práctica, no es posible observar a toda la población, las encuestas por muestreo permiten inferir estos parámetros a partir de una muestra, denotada como \(s\).

En el muestreo probabilístico, cada unidad de la población tiene una probabilidad de inclusión conocida y mayor que cero de ser seleccionada en la muestra. Estas probabilidades constituyen la base para el cálculo de los pesos básicos de muestreo, los cuales se utilizan posteriormente para estimar los parámetros poblacionales mediante sumas ponderadas de los datos recolectados en la encuesta. Cuando el diseño y la selección se implementan correctamente, las estimaciones obtenidas son insesgadas, es decir, en promedio coinciden con el valor real de la población si el proceso de muestreo se repitiera bajo las mismas condiciones.

No obstante, los pesos básicos de muestreo suelen requerir ajustes adicionales con el fin de mejorar la precisión y la robustez de las estimaciones. Uno de los ajustes más comunes corresponde al tratamiento de la no respuesta, mediante el cual los pesos de las unidades que sí respondieron se incrementan para representar también a aquellas unidades seleccionadas que no participaron en la encuesta, contribuyendo así a reducir posibles sesgos. Otro ajuste ampliamente utilizado es la calibración, que consiste en modificar los pesos para garantizar que las sumas ponderadas de ciertas variables auxiliares, como la edad o el sexo, coincidan con totales poblacionales conocidos provenientes de censos o proyecciones demográficas. Además de mejorar la consistencia de las estimaciones, la calibración constituye una herramienta útil para detectar problemas de cobertura o patrones de no respuesta.

Los ajustes de los pesos, y en particular la calibración, desempeñan un papel fundamental en el análisis de encuestas con diseños complejos. Estos procedimientos no solo corrigen las imperfecciones derivadas de la no respuesta o la cobertura incompleta, sino que también fortalecen la coherencia entre la muestra y la población objetivo, garantizando que las inferencias sean estadísticamente válidas y comparables con otras fuentes oficiales \citep{kalton2003weighting}.

La calibración es especialmente relevante porque permite confrontar las estimaciones iniciales de la encuesta con valores de referencia externos, generalmente obtenidos de censos, registros administrativos o proyecciones demográficas. Esta comparación constituye una herramienta poderosa para detectar inconsistencias en la composición de la muestra, identificar sesgos potenciales asociados a la no respuesta y evaluar la calidad del trabajo de campo.

Además, los pesos calibrados tienden a reducir la varianza de las estimaciones, mejorando su precisión sin alterar los totales poblacionales conocidos \citep{sarndal2005nonresponse}. En consecuencia, la calibración no solo corrige, sino que también optimiza el uso de la información disponible, permitiendo obtener resultados más precisos.

La estimación de totales en encuestas constituye un paso central en el análisis estadístico aplicado a poblaciones finitas. Gran parte de los indicadores de interés para la formulación de políticas públicas (como por ejemplo, el número de personas en situación de pobreza, el total de ocupados o el gasto agregado de los hogares) se derivan de un total poblacional. Por esta razón, comprender cómo se definen y estiman los totales resulta fundamental para garantizar la calidad y pertinencia de la información producida.

En términos formales, si \(y_k\) denota el valor de una variable de interés para la unidad \(k \in U\), el total poblacional se define como

\[
t_y = \sum_{U} y_k,
\]

Siendo \(N\) el tamano de la población, la media poblacional se define como \(\bar{y} = \frac{t_y}{N}\). Dado que en la práctica solo se observa una muestra \(s \subset U\), es necesario recurrir a estimadores que incorporen el diseño de muestreo. Por ejemplo, el estimador de Horvitz--Thompson (HT) se expresa como

\[
\hat{t}_{y} = \sum_{s} d_k y_k
\]

donde \(d_k = 1/\pi_k\) son los pesos básicos provenientes del diseño de muestreo y \(\pi_k = \Pr(k \in s)\) son las probabilidades de inclusión de primer orden.

En la práctica, los pesos de diseño suelen modificarse para reflejar procesos adicionales como el ajuste por no respuesta o la calibración a totales poblacionales conocidos, obteniendo así nuevos pesos ajustados, denotados como \(w_k\) para todo \(k \in s\). El reemplazo de \(d_k\) por \(w_k\) permite mejorar la precisión y reducir sesgos en las estimaciones, especialmente cuando existen fuentes auxiliares de información confiables.

No obstante, toda estimación a partir de una muestra conlleva incertidumbre. Incluso cuando el estimador es insesgado, los resultados varían de una muestra a otra debido a la propiedad de aleatoriedad del diseño. Esta variabilidad se cuantifica mediante la varianza del estimador con la cual se puede calcular el error estándar (\(se\)) o el coeficiente de variación (\(cv\)). Estos indicadores son herramientas indispensables para evaluar la precisión de los totales estimados y, por tanto, para interpretar de manera adecuada la información estadística. Bajo el enfoque de diseño, la varianza insesgada del estimador de Horvitz--Thompson puede expresarse como:

\[
\widehat{Var}_p(\hat{t}_{y}) = \sum_{k \in s} \sum_{l \in s} \bigl( d_k d_l - d_{kl} \bigr) y_k y_l,
\]

donde \(d_{kl} = 1/\pi_{kl}\) y \(\pi_{kl} = \Pr(k,l \in s)\) representan las probabilidades conjuntas de inclusión. Esta expresión requiere que el diseño de muestreo cumpla que \(\pi_{kl} > 0\) para todo par de unidades \(k,l \in U\).

Para comprender de manera más concreta la relevancia de considerar el diseño muestral en la estimación de totales y de sus varianzas, analicemos un ejemplo de \citet{UNhandbook} que permite visualizar cómo el método de estimación se ajusta al esquema de selección adoptado. Suponga una población finita de tamaño \(N=6\), de la cual se selecciona una muestra aleatoria simple (MAS) sin reemplazo de tamaño \(n=3\). En dicha muestra se observan los valores \((y_1=10, y_2=14, y_3=18)\). Bajo este diseño, el estimador de Horvitz--Thompson del total poblacional se define como la suma de los valores observados divididos por sus probabilidades de inclusión, y su varianza estimada se obtiene a partir de las covarianzas inducidas por el proceso de selección. De esta forma, la varianza estimada del estimador de Horvitz--Thompson se calcula como:

\[
\widehat{Var}_{MAS}(\hat{t}_{y}) = \frac{N^2}{n}\left(1-\frac{n}{N}\right)S_{y_s}^2
\]

donde \(S_{y_s}^2\) corresponde a la varianza muestral de los valores observados. Sustituyendo en la expresión, se obtiene

\[
\widehat{Var}_p(\hat{t}_{y}) = \frac{36}{3}\left(1-\frac{3}{6}\right)16 = 96
\]

En contraste, si se ignora el diseño de muestreo, un analista inexperto podría calcular erróneamente la varianza mediante la fórmula simplificada:

\[
\frac{N^2}{n}S_{y_s}^2 = 192,
\]

lo que conduciría a una sobreestimación de la varianza por no considerar las características del diseño muestral. Asimismo, la estimación del total poblacional es \(\hat{t}_{y}=84\). El error estándar calculado según el diseño de muestreo, es

\[
\sqrt{\widehat{Var}_p(\hat{t}_{y})} = \sqrt{96} \approx 9.80.
\]

En cambio, si la varianza se estima usando un método ingenuo que ignore el diseño de muestreo, el intervalo de confianza sería más amplio con evidentes sesgos, lo que podría conducir a inferencias erróneas.

Por último, de acuerdo con \citet[p.~258]{kish1965survey}, el efecto del diseño (DEFF) se define como la razón entre la varianza de un estimador obtenida bajo un diseño de muestreo complejo y la varianza del mismo estimador bajo un muestreo aleatorio simple (MAS) de igual tamaño de muestra. Su estimación se expresa como:

\[
\widehat{{DEFF}} = \frac{\widehat{Var}_{p}(\hat{\theta})}{\widehat{Var}_{MAS}(\hat{\theta})}
\]

En donde, \(\widehat{Var}_{p}(\hat{\theta})\) corresponde a la varianza estimada de \(\hat{\theta}\) bajo el diseño complejo \(p(\cdot)\), mientras que \(\widehat{Var}_{MAS}(\hat{\theta})\)
representa la varianza estimada del mismo estimador bajo un diseño MAS con igual tamaño muestral.

Este indicador permite cuantificar cuánto aumenta la varianza debido a la conglomeración y otras características propias de los diseños complejos en comparación con un muestreo simple. Según \citet[p.~40]{ONUDeff}, el DEFF puede entenderse de tres maneras: como el factor de incremento de la varianza frente al MAS, como una medida de la pérdida relativa de precisión o como una indicación del aumento en el tamaño de la muestra que sería necesario en un diseño complejo para alcanzar el mismo nivel de varianza que en un MAS. Según \citet{Park2003Design}, el efecto del diseño de una encuesta puede deberse a los siguientes tres factores:

\begin{itemize}
\tightlist
\item
  Ponderación desigual: la presencia de pesos muestrales no uniformes suele incrementar ligeramente la varianza; por ello, el uso de pesos uniformes resulta ventajoso y explica por qué los diseños auto-ponderados son preferidos en encuestas de hogares.
\item
  Estratificación: cuando se aplica correctamente, puede disminuir la varianza, aunque en la práctica su efecto reductor suele ser moderado.
\item
  Muestreo en varias etapas: generalmente incrementa la varianza, ya que las unidades dentro de un mismo conglomerado tienden a ser más homogéneas entre sí que en comparación con las de otros conglomerados.
\end{itemize}

En el análisis de encuestas, el efecto del diseño (DEFF) constituye un indicador esencial para evaluar la precisión y eficiencia de las estimaciones, así como para orientar la planificación de futuros estudios. Un valor elevado evidencia que el diseño complejo introduce un incremento de la varianza, reduciendo la precisión de los resultados. Por el contrario, un valor cercano a uno indica que el diseño tiene un efecto mínimo sobre la varianza. Esta información permite a los investigadores identificar si es necesario ajustar la ponderación, optimizar la estratificación o modificar el tamaño del submuestreo para aumentar la eficiencia en levantamientos posteriores.

La interpretación de un DEFF alto debe hacerse con precaución, ya que no siempre implica que el diseño muestral sea inadecuado. Es fundamental considerar el contexto de la encuesta: un valor superior a tres podría parecer alarmante, pero a menudo se debe a limitaciones prácticas como restricciones presupuestales, dificultades logísticas o la necesidad de garantizar la participación de los entrevistados. En ciertos levantamientos de hogares, puede ser indispensable seleccionar solo una fracción de los individuos elegibles en cada hogar. Además, problemas de cobertura o tasas de no respuesta pueden aumentar la variabilidad de los pesos muestrales y, en consecuencia, elevar los valores del DEFF. \citet{GutierrezBabativa2023} realizan un análisis detallado de los efectos de diseño en encuestas de hogares en América Latina usando los datos de BADEHOG.

\section{Fundamentos teóricos}\label{fundamentos-teoricos}

Tal como destacan \citet{heeringa2017applied}, el cálculo de totales y medias poblacionales, junto con sus varianzas, ha sido esencial para el desarrollo de la teoría del muestreo probabilístico y la interpretación adecuada de los resultados de encuestas de hogares. La determinación de los totales poblacionales constituye uno de los pilares fundamentales del análisis de encuestas. Tanto las medias como las proporciones y las razones se derivan de estos totales. Un total se define como la suma de una variable específica (por ejemplo, ingreso o gasto) a nivel de toda la población.

Con encuestas de hogares, el análisis de datos numéricoséricos implica con frecuencia calcular estadísticas descriptivas como medias, totales y razones, ya que estas permiten sintetizar las principales características de la población y sirven de base para la toma de decisiones. Dichas estimaciones pueden calcularse para la población en su conjunto o para subgrupos específicos, según los objetivos de la investigación.

\subsection{Estimación puntual}\label{estimacion-puntual}

Una vez definido el objeto del diseño muestral, se procede a realizar los procesos de estimación de los parámetros de interés. Para efectos de este documento, se iniciará con la estimación del total de ingresos de los hogares. Para la estimación de totales en diseños muestrales complejos, se denotará a los estratos con la letra \(h\); a las unidades primarias de muestreo, las cuales están contenidas dentro de los estratos, con la letra \(i\); mientras que las unidades observadas (hogares o personas), se denotarán con la letra \(k\). De esta forma, el estimador del total puede expresarse como:

\[
\hat{t}_y = \sum_{h} \sum_{i} \sum_{k} w_{hik} \, y_{hik}
\]

En donde, \(w_{hik}\) corresponde al factor de expansión de la unidad \(k\), mientras que \(y_{hik}\) corresponde a la observación de la variable de interés en esta misma unidad. El cálculo de la estimación del total y de su varianza asociada puede resultar complejo, por lo que es necesario considerar distintos enfoques metodológicos.

En términos de notación y con el propósito de simplificar la escritura a lo largo de este documento, a menos que se indique lo contrario, \(\sum_{h}\) denotará la suma sobre todos los estratos de la población; \(\sum_{i}\) representará la suma sobre todas las unidades primarias de muestreo seleccionadas en la muestra dentro del estrato \(h\); y \(\sum_{k}\) corresponderá a la suma sobre todos los elementos observados en la muestra de la UPM \(i\), perteneciente al estrato \(h\) de la población de interés.

\subsection{Estimación de la varianza}\label{estimacion-de-la-varianza}

Al trabajar con encuestas de hogares es fundamental no solo obtener estimaciones puntuales, sino también cuantificar la incertidumbre asociada a dichas estimaciones. Comprender y estimar esta incertidumbre constituye una parte esencial del análisis de los datos de encuestas de hogares. Mediante la aplicación de métodos apropiados, los usuarios pueden evaluar la precisión de sus estimaciones. Existen diversos métodos para estimar dicha precisión y, con el apoyo de software moderno, estos enfoques pueden implementarse de manera eficiente para respaldar análisis rigurosos y confiables. Entre los principales métodos se encuentran:

\begin{itemize}
\item
  Ecuaciones de estimación: proporcionan un marco flexible para estimar totales, medias, razones y otros parámetros, así como sus varianzas correspondientes, integrando una visión unificada de la teoría del muestreo \citep{binder1983variances}.
\item
  Linealización de Taylor: consiste en aproximar estadísticas no lineales complejas mediante expresiones lineales, y posteriormente estimar la varianza de dicha aproximación.
\item
  Método del último conglomerado: utilizado con frecuencia en encuestas basadas en muestreo estratificado en múltiples etapas. Se fundamenta en el cálculo de la varianza a partir de las diferencias entre las estimaciones obtenidas a nivel de las unidades primarias de muestreo (UPM). Este método suele combinarse con la linealización de Taylor para estimar la varianza de estadísticas no lineales, como medias o razones.
\item
  Bootstrap y otros métodos de replicación: se fundamentan en tomar repetidas submuestras del conjunto de datos observado, calcular estimaciones para cada réplica y luego utilizar la variabilidad entre estas estimaciones replicadas para inferir la varianza del estimador principal.
\end{itemize}

En el contexto de encuestas de hogares, muchos parámetros complejos pueden formularse como combinaciones lineales de otros parámetros \(f(\theta_1,\ldots,\theta_J)\). En efecto, sea:

\[
f=f(\theta_1,\ldots,\theta_J) = \sum_{j=1}^J a_j \theta_j,
\]

donde los coeficientes \(a_j\) son constantes conocidas y \(\theta_j\) representan parámetros poblacionales. Si, \(\hat\theta_j\) es un estimador de \(\theta_j\), entonces, un estimador muestral de \(f\) se define como:

\[
\hat{f} = \sum_{j=1}^J a_j \hat{\theta}_j,
\]

De esta manera, su varianza puede expresarse como:

\[
Var(\hat{f}) = \sum_{j=1}^J a_j^2 Var(\hat{\theta}_j)
\;+\;
2\sum_{j=1}^{J-1}\sum_{k>j}^J
a_j a_k \, Cov(\hat{\theta}_j,\hat{\theta}_k).
\]

Este resultado es particularmente importante porque incluye numerosos casos de interés práctico que se desarrollarán a lo largo de este documento.

\subsubsection{Ecuaciones de estimación}\label{ecuaciones-de-estimacion}

Muchos parámetros poblacionales pueden expresarse como soluciones de ecuaciones de estimación que involucran totales poblacionales. Aunque los detalles técnicos pueden ser complejos, la idea fundamental es que los mismos principios utilizados para estimar totales pueden aplicarse también para la estimación de varianzas. Este marco general hace que el método sea sencillo y flexible, facilitando su implementación en software estadístico especializado. Una ecuación de estimación poblacional genérica puede expresarse de la siguiente manera:

\[
\sum_{k\in U} z_k(\theta)=0,
\]

donde \(z_k(\cdot)\) es una función de estimación evaluada para la unidad \(k\) y \(\theta\) representa el parámetro poblacional de interés. Estas ecuaciones proporcionan un marco general para definir y calcular diversos parámetros de la población, como totales, medias y razones. Por ejemplo, para el total poblacional, se define \(z_k(\theta)=y_k-\theta/N\), de manera que la ecuación de estimación es \(\sum_{k\in U}(y_k-\theta/N)=0\). La solución de esta ecuación conduce a \(\theta=\sum_{k\in U} y_k = Y\), es decir, al total poblacional.

De forma similar, para la media poblacional se utiliza \(z_k(\theta)=y_k-\theta\), y la ecuación \(\sum_{k\in U}(y_k-\theta)=0\) tiene como solución \(\theta=\left(\sum_{k\in U} y_k\right)/N = \overline{Y}\), correspondiente a la media de la población. Asimismo, en el caso de las razones de totales, se define \(z_k(\theta)=y_k-\theta x_k\). Así, la ecuación \(\sum_{k\in U}(y_k-\theta x_k)=0\) conduce a la solución \(\theta=\dfrac{\sum_{k\in U} y_k}{\sum_{k\in U} x_k} = R\), que corresponde a la razón poblacional entre los totales de las variables \(y\) y \(x\). Estas formulaciones muestran cómo distintos parámetros de interés pueden expresarse mediante ecuaciones de estimación que comparten una estructura común.

La idea de definir parámetros poblacionales como soluciones de ecuaciones de estimación a nivel de población conduce naturalmente a un método general para obtener los estimadores muestrales. En este caso, se utilizan ecuaciones de la forma:

\[
\sum_{k\in s} w_k\, z_k(\theta)=0,
\]

donde \(w_k\) son los factores de expansión y \(z_k(\theta)\) la función de estimación evaluada para cada unidad de la muestra. Bajo un muestreo probabilístico y asumiendo respuesta completa, la suma muestral \(\sum_{k\in s} d_k\, z_k(\theta)\) es insesgada respecto a su análoga poblacional, lo que garantiza que las soluciones de estas ecuaciones sean estimadores consistentes de los parámetros poblacionales.

\subsubsection{Linealización de Taylor}\label{linealizacion-de-taylor}

Esta técnica permite aproximar la varianza de estimadores no lineales. El procedimiento consiste en aplicar una expansión de Taylor de primer orden alrededor del parámetro estimado, con el fin de reemplazar el estimador no lineal por una expresión lineal. De esta manera, se facilita el cálculo de varianzas en situaciones donde no existen fórmulas exactas o su derivación resulta demasiado compleja. Un estimador consistente de la varianza, derivado mediante linealización de Taylor para soluciones de ecuaciones de estimación muestrales, puede expresarse como:

\[
\widehat{Var}(\hat{\theta}) \approx [\hat{J}(\hat{\theta})]^{-1} \, \widehat{Var}_p \Bigg[\sum_{k\in s} w_k\, z_k(\hat{\theta})\Bigg] \, [\hat{J}(\hat{\theta})]^{-1}
\]

donde \(\hat{J}(\hat{\theta}) = \sum_{k\in s} w_k \left[ \frac{\partial z_k(\theta)}{\partial \theta} \right]_{\theta=\hat{\theta}}\). Este resultado muestra cómo la linealización de Taylor convierte la estimación de varianzas de parámetros complejos en un problema de estimación de totales, lo que explica su amplia adopción en software especializado para el análisis de encuestas.

\subsubsection{Último conglomerado}\label{ultimo-conglomerado}

El método del último conglomerado constituye un enfoque directo y robusto para estimar la varianza de totales en encuestas que emplean diseños de muestreo por conglomerados estratificados en múltiples etapas. Propuesto por \citet{hansen1953sample}, este método simplifica la complejidad de los diseños multinivel al centrarse únicamente en la variación entre las Unidades Primarias de Muestreo (UPM). Se asume que, dentro de cada estrato de muestreo, las UPM fueron seleccionadas de manera independiente con reemplazo (posiblemente con probabilidades desiguales), aunque en la práctica la selección suele realizarse sin reemplazo.

El método se basa en la variación entre las estadísticas calculadas a nivel de UPM. Cuando se aplica correctamente, refleja implícitamente cualquier submuestreo realizado dentro de las UPM, permitiendo estimaciones de varianza más simples pero confiables. Es especialmente útil en diseños complejos que incluyen estratificación y probabilidades desiguales de selección tanto de UPM como de unidades de niveles inferiores (hogares e individuos). Los requisitos para aplicar este método incluyen la disponibilidad de estimaciones insesgadas de totales para las variables de interés en cada UPM seleccionada. Además, si la muestra se estratifica en la primera etapa, es necesario contar con al menos dos UPM disponibles por estrato, ya que esto permite estimar adecuadamente la variabilidad entre conglomerados dentro de cada estrato.

De esta forma, considere un diseño de muestreo en múltiples etapas donde se seleccionan \(n_h\) UPM en el estrato \(h\), (\(h=1,\dots,H\)). Una estimación del total de la variable de interés en la UPM \(i\) del estrato \(h\) es:

\[
\hat{t}_{y_i} = \sum_{k\in s_{hi}} w_{hik} \ y_{hik}
\]

En donde \(s_{hi}\) corresponde a la muestra de unidades en la UPM \(i\) del estrato \(h\); por tanto, un estimador insesgado del total poblacional se expresaría como

\[
\hat{t}_{y}  = \sum_{h}\sum_{i} \hat{t}_{y_i}
\]

Según el método del último conglomerado, suponiendo que \(\hat{\bar t}_{y_h}=(1/{n_h}) \sum_{i} \hat{t}_{y_i}\) es la media muestral de los totales estimados de las UPM en el estrato \(h\), un estimador aproximado de la varianza de \(\hat t_y\) se obtiene mediante al siguiente expresión:

\[
\widehat{Var}(\hat{t}_y) = \sum_{h} \frac{n_h}{n_h-1} \sum_{i} (\hat{t}_{y_i} - \hat{\bar t}_{y_h})^2
\]

Para más detalles, véase Hansen, \citet[p.~258]{hansen1953sample} o \citet{wolter2007}. Aunque este método se propuso originalmente para calcular varianzas de estimadores de totales, el método puede combinarse con linealización de Taylor o ecuaciones de estimación para derivar varianzas de otros parámetros poblacionales formulables como soluciones de ecuaciones de estimación. Esta flexibilidad hace que el método sea aplicable a diversos contextos de análisis de encuestas de hogares.

Un supuesto fundamental de esta técnica es que, dentro de cada estrato, las UPM se eligen de forma independiente y con reemplazo. En la práctica, la mayoría de las encuestas selecciona UPM sin reemplazo, generando diseños más eficientes. En consecuencia, las varianzas calculadas bajo la hipótesis de independencia constituyen aproximaciones de las verdaderas varianzas de muestreo. Cuando la fracción muestral es pequeña (por ejemplo, inferior al \(5\%\)), dichas aproximaciones suelen ser suficientemente precisas para su utilización por parte de las oficinas nacionales de estadística o de otros analistas.

Este método destaca por su simplicidad y robustez, lo que lo hace muy atractivo en la práctica. Aunque los métodos más sofisticados que consideran todas las etapas del diseño pueden ofrecer estimaciones de varianza ligeramente más precisas, su aplicación requiere información más detallada y mayor complejidad computacional. Por el contrario, el método del último conglomerado proporciona una aproximación confiable y eficiente, especialmente útil al estimar totales o medias en encuestas de hogares. Para un análisis detallado sobre la precisión de esta aproximación y posibles alternativas, véase \citet{sarndal2003model}.

\subsubsection{Bootstrap}\label{bootstrap}

En muchos casos, los microdatos de encuestas públicas omiten información esencial del diseño, como identificadores de estratos o de unidades primarias de muestreo (UPM), para proteger la confidencialidad de los encuestados. Esta omisión limita la capacidad de los usuarios para calcular varianzas válidas. En tales situaciones, se recomienda que las oficinas nacionales de estadística proporcionen pesos de replicación, lo que permite a los analistas estimar errores estándar de manera correcta. Sin estos datos, los usuarios secundarios no pueden reproducir los errores estándar publicados ni considerar adecuadamente el diseño complejo de la encuesta.

Los métodos de replicación estiman la varianza generando subconjuntos de la muestra original, calculando estimaciones para cada uno y utilizando la variabilidad observada entre estas estimaciones para aproximar la varianza del estimador principal. Son particularmente útiles cuando no se dispone de información sobre estratos o UPM, situación en la que no se puede aplicar el método de Ultimate Cluster.

Por ejemplo, el Bootstrap es una herramienta de replicación robusta y versátil. Originalmente introducido por \citet{efron1979bootstrap} para datos que no provenían de encuestas, su adaptación más usada para encuestas de hogares es el Bootstrap de Reescalamiento Rao-Wu-Yue \citep{rao1992resampling}. Este método se ajusta de manera óptima a diseños de muestreo estratificados y multietápicos, y es ampliamente empleado para la estimación de varianzas en encuestas complejas.

El procedimiento consiste en generar muchas réplicas de la muestra original, simulando extracciones repetidas de la población. Cada réplica se construye mediante la creación de columnas adicionales de pesos de replicación en la base de datos, siguiendo este proceso:

\begin{itemize}
\tightlist
\item
  Para cada estrato, se seleccionan aleatoriamente las UPM con reemplazo; algunas pueden repetirse y otras no aparecer. Cada UPM elegida se incorpora con todas sus observaciones. Si el tamaño de la muestra de primera etapa en el estrato \(h\) es mayor que dos (\(n_h > 2\)), el número de UPM seleccionadas por réplica es \(n_h - 1\).
\item
  Este proceso se repite muchas veces, habitualmente cientos, generando un gran número de réplicas. La cantidad de veces que una UPM \(i\) del estrato \(h\) aparece en la réplica \(r\) se denota \(n_{hi}^{(r)}\), variando entre 0 y \(n_h - 1\).
\item
  A partir de cada réplica se calculan nuevos pesos bootstrap para todas las unidades, reflejando cuántas veces fue seleccionada su UPM. El peso de la unidad \(k\) en la réplica \(r\) se calcula como:
\end{itemize}

\[
w_{hik}^{(r)} = w_{hik} \times \frac{n_h}{n_h - 1} \times n_{hi}^{(r)}
\]

Si los pesos originales incluyen ajustes por no respuesta o calibración, estos deben aplicarse también a cada conjunto de pesos bootstrap.

Cuando se proporciona únicamente pesos de replicación Bootstrap, los analistas pueden estimar errores estándar correctamente, aun sin identificadores de estratos o UPM. Para cada réplica \(r\), se calcula el parámetro de interés \(\hat{\theta}^{(r)}\) usando los pesos bootstrap \(w_{hik}^{(r)}\). La varianza del estimador original se aproxima mediante la variabilidad entre todas las réplicas:

\[
\widehat{Var}_B(\hat{\theta}) = \frac{1}{R} \sum_{r=1}^{R} \left(\hat{\theta}^{(r)} - \tilde{\theta}\right)^2, \quad
\tilde{\theta} = \frac{1}{R} \sum_{r=1}^{R} \hat{\theta}^{(r)}
\]

Este enfoque asegura que la dispersión entre réplicas capture fielmente la incertidumbre del parámetro. El Bootstrap ofrece múltiples ventajas. A pesar de requerir un mayor procesamiento computacional, es eficaz para diseños de encuesta complejos y permite estimar parámetros difíciles de calcular con métodos tradicionales, como medianas u otras estadísticas no lineales. Es especialmente útil para analistas que trabajan con bases de datos sin identificadores de estratos y UPM, pero con pesos de replicación.

La simplicidad del método facilita su aplicación incluso sin software estadístico especializado. Sin embargo, la mayoría de los paquetes estadísticos modernos ya incluyen procedimientos para aplicar Bootstrap y calcular varianzas, ampliando su disponibilidad y robustez. No obstante, su uso no es recomendable en encuestas repetidas con muestras superpuestas ni en situaciones con fracciones de muestreo grandes y tamaños de muestra pequeños \citep{bruch2011variance}.

\section{Intervalos de confianza}\label{intervalos-de-confianza}

Además de producir estimaciones puntuales, uno de los objetivos fundamentales de una encuesta es cuantificar la incertidumbre asociada a dichas estimaciones. En este contexto, los intervalos de confianza permiten construir un rango de valores plausibles para el parámetro poblacional de interés, incorporando explícitamente la variabilidad introducida por el diseño muestral. El intervalo de confianza de nivel \((1-\alpha)100\%\) para el total poblacional \(Y\) se calcula como:

\[
\hat{\theta} \pm t_{1-\alpha/2, df} \times \sqrt{\widehat{Var}(\hat{\theta})}
\]

donde \(\hat{\theta}\) corresponde al estimador por muestreo del parámetro poblacional \(\theta\); \(\widehat{Var}(\hat{\theta})\) representa el estimador de varianza bajo el diseño complejo de la encuesta; y \(t_{1-\alpha/2, df}\) es el cuantil de la distribución \(t\) de Student con \(df\) grados de libertad. En encuestas complejas, los grados de libertad suelen estar relacionados con la cantidad de unidades primarias de muestreo (UPM) y los estratos considerados en el diseño. Una aproximación común consiste en calcularlos como:

\[
df = m - H
\]

donde \(m\) representa el número total de UPM observadas en la muestra y \(H\) el número de estratos. Esta aproximación refleja la cantidad efectiva de información independiente disponible para estimar la variabilidad muestral. A medida que los grados de libertad aumentan, la distribución \(t\) de Student converge hacia la distribución normal estándar. Esta propiedad explica por qué muchas veces se utilizan aproximaciones normales para reportar intervalos de confianza, especialmente en encuestas de hogares de gran escala. Sin embargo, cuando el número de conglomerados es reducido o existen estratos con pocas UPM, la aproximación normal puede subestimar la incertidumbre y producir intervalos excesivamente estrechos.

\chapter{Análisis de las variables continuas}\label{cap3}

El análisis estadístico de encuestas ha evolucionado de manera constante, impulsado por el desarrollo de nuevas metodologías y enfoques que buscan mejorar la precisión y la representatividad de los resultados. En general, estos avances surgen en el ámbito académico y, con el tiempo, son adoptados por instituciones públicas, empresas privadas y organismos estadísticos, que demandan su incorporación en los principales programas de análisis de datos.

Este capítulo presenta los principales procedimientos para el análisis estadístico de variables continuas utilizando \texttt{R}, incorporando adecuadamente las características del diseño muestral, tales como los pesos de expansión, la estratificación y la conglomeración. A lo largo del capítulo se ilustran métodos para estimar parámetros poblacionales, medir su precisión y realizar inferencias válidas a partir de datos provenientes de encuestas complejas. Asimismo, se introducen funciones y herramientas del paquete \texttt{survey} \citep{Lumley2024}, uno de los recursos más utilizados en \texttt{R} para el análisis de encuestas.

Las bases de datos utilizadas en el análisis de encuestas pueden encontrarse en una amplia variedad de formatos (.xlsx, .dat, .csv, .parquet, .sas7bdat, .sav, .txt, entre otros). Sin embargo, una buena práctica consiste en importar la información desde cualquiera de estos formatos y guardarla inmediatamente en un archivo con extensión \texttt{.rds}, que es el formato nativo de \texttt{R}.

Los archivos \texttt{.rds} permiten almacenar cualquier objeto o información en \texttt{R}, como marcos de datos, vectores, matrices o listas, y se caracterizan por su compatibilidad completa con el entorno de trabajo R. Además, facilitan la lectura y escritura eficiente de los datos, reduciendo los tiempos de carga y evitando problemas de compatibilidad entre sesiones. Otro formato propio de \texttt{R} es \texttt{.Rdata}, el cual permite guardar múltiples objetos en un solo archivo. En contraste, los archivos \texttt{.rds} almacenan únicamente un objeto, lo que hace su uso más controlado y reproducible. Por esta razón, se recomienda trabajar preferentemente con el formato \texttt{.rds} al documentar proyectos o reproducir análisis.

Para ilustrar la sintaxis computacional que se utilizará a lo largo del capítulo, se empleará la misma base del capítulo anterior, que contiene una muestra de 2.605 registros que proviene de un muestreo complejo. A continuación, se muestra cómo cargar un archivo con extensión \texttt{.rds}:

\begin{Shaded}
\begin{Highlighting}[]
\FunctionTok{library}\NormalTok{(tidyverse)}

\NormalTok{survey\_data }\OtherTok{\textless{}{-}} \FunctionTok{readRDS}\NormalTok{(}\StringTok{"Data/encuesta.rds"}\NormalTok{)}
\end{Highlighting}
\end{Shaded}

De acuerdo con \citet[sección 7.8]{ONUDeff}, es fundamental que la estructura del diseño muestral sea tenida en cuenta en el proceso de inferencia al estimar estadísticas oficiales basadas en encuestas de hogares. Como se ha insistido en este docuemnto, ignorar este aspecto puede generar estimaciones sesgadas y errores de muestreo subestimados. En este sentido, los programas estadísticos incorporan funcionalidades específicas para manejar datos provenientes de encuestas con diseños complejos \citep[apéndice A]{heeringa2017applied}.

Una vez cargada la muestra de hogares en \texttt{R}, el siguiente paso es definir el diseño muestral del cual proviene dicha muestra. Para ello se utilizará el paquete \texttt{srvyr} \citep{Freedman2024}, que surge como un complemento del paquete \texttt{survey} \citep{Lumley2024}. Estas librerías permiten definir objetos tipo \texttt{survey.design}, a los que se aplican las funciones de estimación y análisis de encuestas, y que pueden ser combinadas con la programación del paquete \texttt{tidyverse} \citep{tidyverse2026}. A continuación, se muestra un ejemplo en el que se define un objeto de tipo \texttt{survey\_design} para la encuesta:

\begin{Shaded}
\begin{Highlighting}[]
\FunctionTok{options}\NormalTok{(}\AttributeTok{survey.lonely.UPM =} \StringTok{"adjust"}\NormalTok{)}

\FunctionTok{library}\NormalTok{(survey)}
\FunctionTok{library}\NormalTok{(srvyr)}

\NormalTok{survey\_design }\OtherTok{\textless{}{-}}\NormalTok{ survey\_data }\SpecialCharTok{\%\textgreater{}\%}
  \FunctionTok{as\_survey\_design}\NormalTok{(}
    \AttributeTok{strata =}\NormalTok{ Stratum,}
    \AttributeTok{ids =}\NormalTok{ PSU,}
    \AttributeTok{weights =}\NormalTok{ wk,}
    \AttributeTok{nest =} \ConstantTok{TRUE}
\NormalTok{  )}
\end{Highlighting}
\end{Shaded}

En este código, \texttt{survey\_data} corresponde a la base de datos que contiene la muestra seleccionada. A partir de estos datos, la función \texttt{as\_survey\_design()} convierte la información en un objeto de clase \texttt{survey\_design}, el cual almacena todos los elementos necesarios para describir el diseño muestral de la encuesta. Entre ellos se incluyen la variable de estratificación (\texttt{strata}), que identifica los estratos definidos en el diseño, y la variable especificada en \texttt{ids}, que representa las unidades primarias de muestreo (UPM) o conglomerados seleccionados en la primera etapa del muestreo. Asimismo, el argumento \texttt{weights} define la variable que contiene los factores de expansión asociados a cada observación. Finalmente, la opción \texttt{nest\ =\ TRUE} indica que las UPM están anidadas dentro de los estratos, una característica necesaria para representar correctamente la estructura jerárquica del diseño de muestreo.

\section{Estimación de parámetros descriptivos}\label{estimacion-de-parametros-descriptivos}

En las encuestas de hogares es común estimar parámetros descriptivos asociados a variables numéricas, tales como totales, medias, razones y desviaciones estándar. Estas medidas permiten caracterizar distintos fenómenos sociales y económicos de la población, por ejemplo el ingreso, el gasto, las horas trabajadas o el tamaño de los hogares. En esta sección se presentan la lógica detrás de la sintaxis utilizada para este tipo de estimación, junto con sus correspondientes medidas de precisión.

\subsection{Estimación del total}\label{estimacion-del-total}

En \texttt{R}, la función \texttt{survey\_total()} se utiliza para estimar totales a partir del objeto de diseño. Esta función incorpora automáticamente los factores de expansión y la estructura del diseño, garantizando que las estimaciones reflejen correctamente el modelo de inferencia.

\begin{Shaded}
\begin{Highlighting}[]
\FunctionTok{survey\_total}\NormalTok{(}
\NormalTok{  x,}
  \AttributeTok{na.rm =} \ConstantTok{FALSE}\NormalTok{,}
  \AttributeTok{vartype =} \FunctionTok{c}\NormalTok{(}\StringTok{"se"}\NormalTok{, }\StringTok{"ci"}\NormalTok{, }\StringTok{"var"}\NormalTok{, }\StringTok{"cv"}\NormalTok{),}
  \AttributeTok{level =} \FloatTok{0.95}\NormalTok{,}
  \AttributeTok{deff =} \ConstantTok{FALSE}
\NormalTok{)}
\end{Highlighting}
\end{Shaded}

Los argumentos principales de la función \texttt{survey\_total()} son \texttt{x}, \texttt{na.rm}, \texttt{vartype}, \texttt{level} y \texttt{deff}. El argumento \texttt{x} corresponde a la variable sobre la cual se desea calcular el total; también puede dejarse vacío cuando se busca obtener únicamente el total ponderado de casos de la muestra. Por su parte, \texttt{na.rm} indica si los valores faltantes (\texttt{NA}) deben eliminarse antes de realizar el cálculo, siendo \texttt{FALSE} el valor predeterminado.

El argumento \texttt{vartype} permite especificar el tipo de medida de incertidumbre que se desea estimar, pudiendo solicitarse uno o varios resultados simultáneamente, tales como el error estándar (\texttt{"se"}), el intervalo de confianza (\texttt{"ci"}), la varianza (\texttt{"var"}) o el coeficiente de variación (\texttt{"cv"}). Asimismo, \texttt{level} define el nivel de confianza utilizado para el cálculo de los intervalos de confianza, cuyo valor por defecto es 0.95. Por otro lado, el intervalo de confianza puede obtenerse directamente al incluir el argumento \texttt{vartype\ =\ "ci"} dentro de \texttt{survey\_total()}, o bien utilizando la función \texttt{confint()} sobre el objeto resultante, especificando el nivel de confianza requerido. Finalmente, \texttt{deff} es un valor lógico que indica si se debe retornar el efecto del diseño (DEFF).

El siguiente ejemplo presenta cómo estimar totales y sus intervalos de confianza para diferentes variables de interés en \texttt{R}, utilizando la función \texttt{survey\_total()}, y presentada en la Tabla \ref{tab:tab-total-ingreso}:

\begin{Shaded}
\begin{Highlighting}[]
\NormalTok{survey\_design }\SpecialCharTok{\%\textgreater{}\%}
  \FunctionTok{summarise}\NormalTok{(}\AttributeTok{total =} \FunctionTok{survey\_total}\NormalTok{(Income, }\AttributeTok{vartype =} \FunctionTok{c}\NormalTok{(}\StringTok{"se"}\NormalTok{, }\StringTok{"ci"}\NormalTok{), }\AttributeTok{deff =} \ConstantTok{TRUE}\NormalTok{))}
\end{Highlighting}
\end{Shaded}

\begin{longtable}[]{@{}lcccc@{}}
\caption{\label{tab:tab-total-ingreso}Estimación del total de ingresos de los hogares con intervalo de confianza al 95\%}\tabularnewline
\toprule\noalign{}
total & total\_se & total\_low & total\_upp & total\_deff \\
\midrule\noalign{}
\endfirsthead
\toprule\noalign{}
total & total\_se & total\_low & total\_upp & total\_deff \\
\midrule\noalign{}
\endhead
\bottomrule\noalign{}
\endlastfoot
85793667 & 4778674 & 76331414 & 95255920 & 11 \\
\end{longtable}

De la tabla, se obtiene que la estimación del total poblacional para el ingreso total es de 85.793.667. Este valor corresponde al estimador puntual obtenido a partir de la encuesta, incorporando los factores de expansión y las características del diseño muestral. El error estándar asociado a esta estimación es de 4.778.674, lo que refleja la variabilidad muestral del estimador. A partir de este valor se construye un intervalo de confianza del 95\%, cuyos límites inferior y superior son 76.331.414 y 95.255.920, respectivamente. El efecto de diseño (\texttt{deff}) es igual a 11, lo cual sugiere una pérdida importante de precisión asociada principalmente a las características del diseño complejo. En consecuencia, aunque el tamaño muestral nominal pueda ser grande, el tamaño muestral efectivo resultaría considerablemente menor debido a la correlación entre observaciones dentro de los conglomerados.

Para continuar ilustrando el uso de la función \texttt{survey\_total()} dentro del enfoque de \texttt{srvyr}, estimemos el total de gastos de los hogares, pero ahora calculando el intervalo de confianza al 90\%. El siguiente código realiza esta estimación, presentada en la Tabla \ref{tab:tab-total-gasto}:

\begin{Shaded}
\begin{Highlighting}[]
\NormalTok{survey\_design }\SpecialCharTok{\%\textgreater{}\%}
  \FunctionTok{summarise}\NormalTok{(}\AttributeTok{total =} \FunctionTok{survey\_total}\NormalTok{(Expenditure,}
                              \AttributeTok{vartype =} \StringTok{"ci"}\NormalTok{, }\AttributeTok{level =} \FloatTok{0.9}\NormalTok{,}
                              \AttributeTok{deff =} \ConstantTok{TRUE}\NormalTok{))}
\end{Highlighting}
\end{Shaded}

\begin{longtable}[]{@{}lccc@{}}
\caption{\label{tab:tab-total-gasto}Estimación del total de gastos de los hogares con intervalo de confianza al 90\%}\tabularnewline
\toprule\noalign{}
total & total\_low & total\_upp & total\_deff \\
\midrule\noalign{}
\endfirsthead
\toprule\noalign{}
total & total\_low & total\_upp & total\_deff \\
\midrule\noalign{}
\endhead
\bottomrule\noalign{}
\endlastfoot
55677504 & 51360469 & 59994539 & 10.2 \\
\end{longtable}

Si el objetivo es estimar el total de ingresos de los hogares, pero desagregado por sexo, se puede utilizar la función \texttt{group\_by()} para agrupar por la variable de interés, junto con \texttt{cascade()} de la librería \texttt{srvyr}, que permite agregar una fila con el total general al final de la tabla. De esta manera, se pueden obtener fácilmente los totales por categoría y el total global dentro del mismo marco de resumen, como se muestra en la Tabla \ref{tab:tab-total-ingreso-sexo}.

\begin{Shaded}
\begin{Highlighting}[]
\NormalTok{survey\_design }\SpecialCharTok{\%\textgreater{}\%}
  \FunctionTok{group\_by}\NormalTok{(Sex) }\SpecialCharTok{\%\textgreater{}\%}
  \FunctionTok{cascade}\NormalTok{(}\AttributeTok{total =} \FunctionTok{survey\_total}\NormalTok{(Income,}
                               \AttributeTok{vartype =} \FunctionTok{c}\NormalTok{(}\StringTok{"se"}\NormalTok{, }\StringTok{"ci"}\NormalTok{),}
                               \AttributeTok{level =} \FloatTok{0.95}\NormalTok{),}
          \AttributeTok{.fill =} \StringTok{"Total income"}\NormalTok{)}
\end{Highlighting}
\end{Shaded}

\begin{longtable}[]{@{}lcccc@{}}
\caption{\label{tab:tab-total-ingreso-sexo}Total de ingresos de los hogares desagregado por sexo}\tabularnewline
\toprule\noalign{}
Sex & total & total\_se & total\_low & total\_upp \\
\midrule\noalign{}
\endfirsthead
\toprule\noalign{}
Sex & total & total\_se & total\_low & total\_upp \\
\midrule\noalign{}
\endhead
\bottomrule\noalign{}
\endlastfoot
Female & 44153820 & 2324452 & 39551172 & 48756467 \\
Male & 41639847 & 2870194 & 35956576 & 47323118 \\
Total income & 85793667 & 4778674 & 76331414 & 95255920 \\
\end{longtable}

Como se observa en los códigos anteriores, una forma practica de obtener las estimaciones del total, su error estándar y su intervalo de confianza es mediante el argumento \texttt{vartype}, especificando las opciones \texttt{"se"} y \texttt{"ci"} respectivamente.

\subsection{Estimación de la media}\label{estimacion-de-la-media}

La estimación de la media poblacional ocupa un lugar central en las encuestas de hogares, ya que muchos de los indicadores utilizados para el análisis socioeconómico corresponden a valores promedio, como el ingreso medio del hogar, el gasto per cápita o las horas promedio trabajadas. Este tipo de parámetros permite resumir el comportamiento general de las variables de interés y describir las tendencias centrales de la población objetivo. De acuerdo con \citet{Gutierrez2016}, el estimador de la media poblacional puede expresarse como una razón no lineal entre dos totales poblacionales, los cuales se estiman de la siguiente manera:

\[
\hat{\bar{y}}=  \frac{\hat{t}_y}{\hat{N}}  =  \frac{\sum_{h}\sum_{i}\sum_{k} w_{h i k} \ y_{hik}}{\sum_{h}\sum_{i}\sum_{k} w_{h i k}}
\]

Como \(\hat{\bar{y}}\) no es una estadística lineal, no existe una fórmula cerrada para la varianza de este estimador. Por esta razón, es necesario recurrir a métodos de remuestreo o a la linealización de Taylor para aproximar su varianza. En este caso particular, utilizando series de Taylor, la varianza se define como:

\[
\widehat{Var}\left(\hat{\bar{y}}\right)  \approx  \frac{\widehat{Var}\left(\hat{t}_y\right)+\hat{\bar{y}}^{2}\ \widehat{Var}\left(\hat{N}\right)-2\ \hat{\bar{y}} \  \widehat{Cov}\left(\hat{t}_y,\hat{N}\right)}{\hat{N}^{2}}
\]

Como se puede observar, el cálculo de la varianza de la media poblacional implica componentes analíticos complejos, como la covarianza entre el total estimado y el tamaño poblacional estimado. Sin embargo, el paquete \texttt{srvyr} en \texttt{R} facilita estos cálculos al incorporar funciones que estiman directamente la media, su error estándar, el intervalo de confianza y el efecto del diseño. A continuación, se presenta la sintaxis para estimar la media de los ingresos de los hogares, junto con su intervalo de confianza al 95\%, presentada en la Tabla \ref{tab:tab-media-ingreso}:

\begin{Shaded}
\begin{Highlighting}[]
\NormalTok{survey\_design }\SpecialCharTok{\%\textgreater{}\%}
  \FunctionTok{summarise}\NormalTok{(}\AttributeTok{mean =} \FunctionTok{survey\_mean}\NormalTok{(Income,}
                                \AttributeTok{vartype =} \FunctionTok{c}\NormalTok{(}\StringTok{"se"}\NormalTok{, }\StringTok{"ci"}\NormalTok{),}
                                \AttributeTok{level =} \FloatTok{0.95}\NormalTok{,}
                                \AttributeTok{deff =} \ConstantTok{TRUE}\NormalTok{))}
\end{Highlighting}
\end{Shaded}

\begin{longtable}[]{@{}lcccc@{}}
\caption{\label{tab:tab-media-ingreso}Estimación de la media de ingresos de los hogares}\tabularnewline
\toprule\noalign{}
mean & mean\_se & mean\_low & mean\_upp & mean\_deff \\
\midrule\noalign{}
\endfirsthead
\toprule\noalign{}
mean & mean\_se & mean\_low & mean\_upp & mean\_deff \\
\midrule\noalign{}
\endhead
\bottomrule\noalign{}
\endlastfoot
571 & 28.5 & 515 & 627 & 8.82 \\
\end{longtable}

Como se aprecia, los argumentos utilizados en la función \texttt{survey\_mean()} son similares a los de \texttt{survey\_total()}. En este caso, \texttt{vartype} permite obtener el error estándar (\texttt{"se"}) y el intervalo de confianza (\texttt{"ci"}), mientras que el argumento \texttt{deff\ =\ TRUE} solicita el cálculo del efecto del diseño. De forma análoga, es posible estimar la media de los gastos de los hogares, empleando la misma estructura de código, presentada en la Tabla \ref{tab:tab-media-gasto}:

\begin{Shaded}
\begin{Highlighting}[]
\NormalTok{survey\_design }\SpecialCharTok{\%\textgreater{}\%}
  \FunctionTok{summarise}\NormalTok{(}\AttributeTok{mean =} \FunctionTok{survey\_mean}\NormalTok{(Expenditure,}
                                    \AttributeTok{vartype =} \FunctionTok{c}\NormalTok{(}\StringTok{"se"}\NormalTok{, }\StringTok{"ci"}\NormalTok{),}
                                    \AttributeTok{level =} \FloatTok{0.95}\NormalTok{,}
                                    \AttributeTok{deff =} \ConstantTok{TRUE}\NormalTok{))}
\end{Highlighting}
\end{Shaded}

\begin{longtable}[]{@{}lcccc@{}}
\caption{\label{tab:tab-media-gasto}Estimación de la media de gastos de los hogares}\tabularnewline
\toprule\noalign{}
mean & mean\_se & mean\_low & mean\_upp & mean\_deff \\
\midrule\noalign{}
\endfirsthead
\toprule\noalign{}
mean & mean\_se & mean\_low & mean\_upp & mean\_deff \\
\midrule\noalign{}
\endhead
\bottomrule\noalign{}
\endlastfoot
371 & 13.3 & 344 & 397 & 6.02 \\
\end{longtable}

También se pueden realizar estimaciones de la media por subgrupos siguiendo el mismo esquema mostrado previamente. Particularmente, los gastos de los hogares discriminados por sexo se presentan en la Tabla \ref{tab:tab-media-gasto-sexo}:

\begin{Shaded}
\begin{Highlighting}[]
\NormalTok{survey\_design }\SpecialCharTok{\%\textgreater{}\%}
  \FunctionTok{group\_by}\NormalTok{(Sex) }\SpecialCharTok{\%\textgreater{}\%}
  \FunctionTok{cascade}\NormalTok{(}\AttributeTok{mean =} \FunctionTok{survey\_mean}\NormalTok{(Expenditure,}
                              \AttributeTok{level =} \FloatTok{0.95}\NormalTok{,}
                              \AttributeTok{vartype =} \FunctionTok{c}\NormalTok{(}\StringTok{"se"}\NormalTok{, }\StringTok{"ci"}\NormalTok{)),}
        \AttributeTok{.fill =} \StringTok{"Mean expenditure"}\NormalTok{) }\SpecialCharTok{\%\textgreater{}\%}
  \FunctionTok{arrange}\NormalTok{(}\FunctionTok{desc}\NormalTok{(Sex))}
\end{Highlighting}
\end{Shaded}

\begin{longtable}[]{@{}lcccc@{}}
\caption{\label{tab:tab-media-gasto-sexo}Media de gastos de los hogares desagregada por sexo}\tabularnewline
\toprule\noalign{}
Sex & mean & mean\_se & mean\_low & mean\_upp \\
\midrule\noalign{}
\endfirsthead
\toprule\noalign{}
Sex & mean & mean\_se & mean\_low & mean\_upp \\
\midrule\noalign{}
\endhead
\bottomrule\noalign{}
\endlastfoot
Mean expenditure & 371 & 13.3 & 344 & 397 \\
Male & 374 & 16.1 & 343 & 406 \\
Female & 367 & 12.3 & 343 & 391 \\
\end{longtable}

De manera análoga, se pueden calcular las medias del gasto por zona, las cuales son presentadas en la Tabla \ref{tab:tab-media-gasto-zona}:

\begin{Shaded}
\begin{Highlighting}[]
\NormalTok{survey\_design }\SpecialCharTok{\%\textgreater{}\%}
  \FunctionTok{group\_by}\NormalTok{(Zone) }\SpecialCharTok{\%\textgreater{}\%}
  \FunctionTok{cascade}\NormalTok{(}\AttributeTok{mean =} \FunctionTok{survey\_mean}\NormalTok{(Expenditure,}
                              \AttributeTok{level =} \FloatTok{0.95}\NormalTok{,}
                              \AttributeTok{vartype =} \FunctionTok{c}\NormalTok{(}\StringTok{"se"}\NormalTok{, }\StringTok{"ci"}\NormalTok{)),}
        \AttributeTok{.fill =} \StringTok{"Mean expenditure"}\NormalTok{) }\SpecialCharTok{\%\textgreater{}\%}
  \FunctionTok{arrange}\NormalTok{(}\FunctionTok{desc}\NormalTok{(Zone))}
\end{Highlighting}
\end{Shaded}

\begin{longtable}[]{@{}lcccc@{}}
\caption{\label{tab:tab-media-gasto-zona}Media de gastos de los hogares desagregada por zona}\tabularnewline
\toprule\noalign{}
Zone & mean & mean\_se & mean\_low & mean\_upp \\
\midrule\noalign{}
\endfirsthead
\toprule\noalign{}
Zone & mean & mean\_se & mean\_low & mean\_upp \\
\midrule\noalign{}
\endhead
\bottomrule\noalign{}
\endlastfoot
Urban & 460 & 22.2 & 416 & 504 \\
Rural & 274 & 10.3 & 254 & 294 \\
Mean expenditure & 371 & 13.3 & 344 & 397 \\
\end{longtable}

Asimismo, es posible realizar una estimación combinada por sexo y zona, presentada en la Tabla \ref{tab:tab-media-gasto-zona-sexo}:

\begin{Shaded}
\begin{Highlighting}[]
\NormalTok{survey\_design }\SpecialCharTok{\%\textgreater{}\%}
  \FunctionTok{group\_by}\NormalTok{(Zone, Sex) }\SpecialCharTok{\%\textgreater{}\%}
  \FunctionTok{cascade}\NormalTok{(}\AttributeTok{mean =} \FunctionTok{survey\_mean}\NormalTok{(Expenditure,}
                              \AttributeTok{level =} \FloatTok{0.95}\NormalTok{,}
                              \AttributeTok{vartype =} \FunctionTok{c}\NormalTok{(}\StringTok{"se"}\NormalTok{, }\StringTok{"ci"}\NormalTok{)),}
        \AttributeTok{.fill =} \StringTok{"Mean expenditure"}\NormalTok{) }\SpecialCharTok{\%\textgreater{}\%}
  \FunctionTok{arrange}\NormalTok{(}\FunctionTok{desc}\NormalTok{(Zone), }\FunctionTok{desc}\NormalTok{(Sex))}
\end{Highlighting}
\end{Shaded}

\begin{longtable}[]{@{}
  >{\raggedright\arraybackslash}p{(\linewidth - 10\tabcolsep) * \real{0.2429}}
  >{\centering\arraybackslash}p{(\linewidth - 10\tabcolsep) * \real{0.2571}}
  >{\centering\arraybackslash}p{(\linewidth - 10\tabcolsep) * \real{0.0857}}
  >{\centering\arraybackslash}p{(\linewidth - 10\tabcolsep) * \real{0.1286}}
  >{\centering\arraybackslash}p{(\linewidth - 10\tabcolsep) * \real{0.1429}}
  >{\centering\arraybackslash}p{(\linewidth - 10\tabcolsep) * \real{0.1429}}@{}}
\caption{\label{tab:tab-media-gasto-zona-sexo}Media de gastos de los hogares desagregada por zona y sexo}\tabularnewline
\toprule\noalign{}
\begin{minipage}[b]{\linewidth}\raggedright
Zone
\end{minipage} & \begin{minipage}[b]{\linewidth}\centering
Sex
\end{minipage} & \begin{minipage}[b]{\linewidth}\centering
mean
\end{minipage} & \begin{minipage}[b]{\linewidth}\centering
mean\_se
\end{minipage} & \begin{minipage}[b]{\linewidth}\centering
mean\_low
\end{minipage} & \begin{minipage}[b]{\linewidth}\centering
mean\_upp
\end{minipage} \\
\midrule\noalign{}
\endfirsthead
\toprule\noalign{}
\begin{minipage}[b]{\linewidth}\raggedright
Zone
\end{minipage} & \begin{minipage}[b]{\linewidth}\centering
Sex
\end{minipage} & \begin{minipage}[b]{\linewidth}\centering
mean
\end{minipage} & \begin{minipage}[b]{\linewidth}\centering
mean\_se
\end{minipage} & \begin{minipage}[b]{\linewidth}\centering
mean\_low
\end{minipage} & \begin{minipage}[b]{\linewidth}\centering
mean\_upp
\end{minipage} \\
\midrule\noalign{}
\endhead
\bottomrule\noalign{}
\endlastfoot
Urban & Mean expenditure & 460 & 22.2 & 416 & 504 \\
Urban & Male & 470 & 27.0 & 416 & 523 \\
Urban & Female & 451 & 20.1 & 411 & 491 \\
Rural & Mean expenditure & 274 & 10.3 & 254 & 294 \\
Rural & Male & 275 & 10.2 & 255 & 296 \\
Rural & Female & 273 & 11.6 & 250 & 296 \\
Mean expenditure & Mean expenditure & 371 & 13.3 & 344 & 397 \\
\end{longtable}

\subsection{Estimación de razones}\label{estimacion-de-razones}

Un caso particular de funciones no lineales de totales es la razón poblacional, definida como el cociente entre dos totales poblacionales asociados a características de interés. Las razones permiten expresar la relación entre dos variables, lo cual resulta especialmente útil para construir indicadores comparativos y de seguimiento.

Este tipo de parámetro es ampliamente utilizado en encuestas de hogares para construir indicadores relativos, por ejemplo, la cantidad de hombres por cada mujer, la proporción de personas ocupadas respecto a la población en edad de trabajar o el número promedio de mascotas por hogar. Las razones también se utilizan en marcos internacionales. Por ejemplo, el Indicador 2.1.1 de los Objetivos de Desarrollo Sostenible (ODS) (prevalencia de subalimentación) se calcula a partir de la razón entre el consumo de alimentos, medido en calorías ingeridas, y los requerimientos energéticos mínimos de la dieta, determinados según edad, sexo y nivel de actividad física.

Dado que la razón corresponde al cociente entre dos parámetros poblacionales desconocidos, tanto el numerador como el denominador deben ser estimados a partir de la muestra \citep{bautista1998}. Siendo \(Y\) el total de la varaible \(y\) y \(X\) el total de la varaible \(x\), la razón poblacional se define como \(R = \frac{t_y}{t_x}\) y su estimador puntual en el marco de un muestreo complejo se expresa como:

\[
\hat{R} = \frac{\hat{t}_y}{\hat{t}_x}
= \frac{\sum_{h}\sum_{i}\sum_{k} w_{hik} \ y_{hik}}
       {\sum_{h}\sum_{i}\sum_{k} w_{hik} \ x_{hik}}
\]

Sin embargo, dado que \(\hat{R}\) es un cociente entre dos estimadores, es decir, dos variables aleatorias, el cálculo de su varianza no es trivial. Para ello se emplea la linealización de Taylor, como lo muestra \citet{Gutierrez2016}, o bien métodos de remuestreo. En particular, la función \texttt{survey\_ratio} implementa la estimación de razones y sus varianzas en el marco de encuestas complejas.

Un ejemplo directo en encuestas de hogares es la razón \textbf{gasto/ingreso}, que ayuda a identificar patrones de consumo y niveles de sostenibilidad económica de los hogares. A continuación, se ilustra cómo estimar la razón entre el gasto y el ingreso de los hogares, presentada en la Tabla \ref{tab:tab-razon-gasto-ingreso}:

\begin{Shaded}
\begin{Highlighting}[]
\NormalTok{survey\_design }\SpecialCharTok{\%\textgreater{}\%}
  \FunctionTok{summarise}\NormalTok{(}\AttributeTok{ratio =} \FunctionTok{survey\_ratio}\NormalTok{(}\AttributeTok{numerator =}\NormalTok{ Expenditure,}
                                 \AttributeTok{denominator =}\NormalTok{ Income,}
                                 \AttributeTok{level =} \FloatTok{0.95}\NormalTok{,}
                                 \AttributeTok{vartype =} \FunctionTok{c}\NormalTok{(}\StringTok{"se"}\NormalTok{, }\StringTok{"ci"}\NormalTok{))}
\NormalTok{            )}
\end{Highlighting}
\end{Shaded}

\begin{longtable}[]{@{}lccc@{}}
\caption{\label{tab:tab-razon-gasto-ingreso}Razón gasto/ingreso de los hogares}\tabularnewline
\toprule\noalign{}
ratio & ratio\_se & ratio\_low & ratio\_upp \\
\midrule\noalign{}
\endfirsthead
\toprule\noalign{}
ratio & ratio\_se & ratio\_low & ratio\_upp \\
\midrule\noalign{}
\endhead
\bottomrule\noalign{}
\endlastfoot
0.649 & 0.023 & 0.603 & 0.695 \\
\end{longtable}

En este caso, se especificaron la variable de gasto como numerador (\texttt{numerator}), la variable de ingreso como denominador (\texttt{denominator}), el nivel de confianza (\texttt{level}) utilizado para la construcción de los intervalos de confianza y las medidas de precisión requeridas mediante el argumento (\texttt{vartype}). Como se puede observar, la razón entre el gasto y el ingreso es, aproximando, 0.65. Lo que implica que por cada unidad 100 unidades monetarias que le ingrese al hogar, se gastan 65 unidades, consiguiendo un intervalo de confianza al 95\% de 0.60 y 0.69.

Si ahora el objetivo es estimar la razón entre mujeres y hombres en la base de ejemplo, se realiza de la siguiente manera, presentada en la Tabla \ref{tab:tab-razon-sexo}:

\begin{Shaded}
\begin{Highlighting}[]
\NormalTok{survey\_design }\SpecialCharTok{\%\textgreater{}\%}
  \FunctionTok{summarise}\NormalTok{(}\AttributeTok{ratio =} \FunctionTok{survey\_ratio}\NormalTok{(}\AttributeTok{numerator =}\NormalTok{ (Sex }\SpecialCharTok{==} \StringTok{"Female"}\NormalTok{),}
                                 \AttributeTok{denominator =}\NormalTok{ (Sex }\SpecialCharTok{==} \StringTok{"Male"}\NormalTok{),}
                                 \AttributeTok{level =} \FloatTok{0.95}\NormalTok{,}
                                 \AttributeTok{vartype =} \FunctionTok{c}\NormalTok{(}\StringTok{"se"}\NormalTok{, }\StringTok{"ci"}\NormalTok{))}
\NormalTok{            )}
\end{Highlighting}
\end{Shaded}

\begin{longtable}[]{@{}lccc@{}}
\caption{\label{tab:tab-razon-sexo}Razón de mujeres por hombre en la población}\tabularnewline
\toprule\noalign{}
ratio & ratio\_se & ratio\_low & ratio\_upp \\
\midrule\noalign{}
\endfirsthead
\toprule\noalign{}
ratio & ratio\_se & ratio\_low & ratio\_upp \\
\midrule\noalign{}
\endhead
\bottomrule\noalign{}
\endlastfoot
1.11 & 0.035 & 1.04 & 1.18 \\
\end{longtable}

Dado que la variable sexo en la base de datos es categórica, fue necesario acudir a variables indicadoras para realizar el cálculo de la razón, utilizando \texttt{Sex\ ==\ "Female"} para identificar a las mujeres y \texttt{Sex\ ==\ "Male"} para los hombres. Los resultados muestran que en la población analizada hay una mayor cantidad de mujeres que de hombres, obteniéndose una razón estimada de 1.11. Esto indica que, por cada 100 hombres, existen aproximadamente 111 mujeres. Además, el intervalo de confianza al 95\% sugiere que esta razón podría variar entre 1.05 y 1.18. Si se desea hacer la razón de mujeres y hombres pero en la zona rural, se haría de la siguiente manera, presentada en la Tabla \ref{tab:tab-razon-sexo-rural}:

\begin{Shaded}
\begin{Highlighting}[]
\NormalTok{rural\_subset }\OtherTok{\textless{}{-}}\NormalTok{ survey\_design }\SpecialCharTok{\%\textgreater{}\%}
  \FunctionTok{filter}\NormalTok{(Zone }\SpecialCharTok{==} \StringTok{"Rural"}\NormalTok{)}
\NormalTok{rural\_subset }\SpecialCharTok{\%\textgreater{}\%}
  \FunctionTok{summarise}\NormalTok{(}\AttributeTok{ratio =} \FunctionTok{survey\_ratio}\NormalTok{(}\AttributeTok{numerator =}\NormalTok{ (Sex }\SpecialCharTok{==} \StringTok{"Female"}\NormalTok{),}
                                  \AttributeTok{denominator =}\NormalTok{ (Sex }\SpecialCharTok{==} \StringTok{"Male"}\NormalTok{),}
                                  \AttributeTok{level =} \FloatTok{0.95}\NormalTok{,}
                                  \AttributeTok{vartype =} \FunctionTok{c}\NormalTok{(}\StringTok{"se"}\NormalTok{, }\StringTok{"ci"}\NormalTok{))}
\NormalTok{            )}
\end{Highlighting}
\end{Shaded}

\begin{longtable}[]{@{}lccc@{}}
\caption{\label{tab:tab-razon-sexo-rural}Razón de mujeres por hombre en zona rural}\tabularnewline
\toprule\noalign{}
ratio & ratio\_se & ratio\_low & ratio\_upp \\
\midrule\noalign{}
\endfirsthead
\toprule\noalign{}
ratio & ratio\_se & ratio\_low & ratio\_upp \\
\midrule\noalign{}
\endhead
\bottomrule\noalign{}
\endlastfoot
1.07 & 0.035 & 0.997 & 1.14 \\
\end{longtable}

Ahora bien, otro análisis de interés es estimar la razón de gastos pero solo en la población femenina. A continuación, se presentan los códigos computacionales, como se muestra en la Tabla \ref{tab:tab-razon-gasto-ingreso-mujer}.

\begin{Shaded}
\begin{Highlighting}[]
\NormalTok{female\_subset }\OtherTok{\textless{}{-}}\NormalTok{ survey\_design }\SpecialCharTok{\%\textgreater{}\%}
  \FunctionTok{filter}\NormalTok{(Sex }\SpecialCharTok{==} \StringTok{"Female"}\NormalTok{)}

\NormalTok{female\_subset }\SpecialCharTok{\%\textgreater{}\%}
  \FunctionTok{summarise}\NormalTok{(}\AttributeTok{ratio =} \FunctionTok{survey\_ratio}\NormalTok{(}\AttributeTok{numerator =}\NormalTok{ Expenditure,}
                                 \AttributeTok{denominator =}\NormalTok{ Income,}
                                 \AttributeTok{level =} \FloatTok{0.95}\NormalTok{,}
                                 \AttributeTok{vartype =} \FunctionTok{c}\NormalTok{(}\StringTok{"se"}\NormalTok{, }\StringTok{"ci"}\NormalTok{))}
\NormalTok{            )}
\end{Highlighting}
\end{Shaded}

\begin{longtable}[]{@{}lccc@{}}
\caption{\label{tab:tab-razon-gasto-ingreso-mujer}Razón gasto/ingreso en hogares con jefatura femenina}\tabularnewline
\toprule\noalign{}
ratio & ratio\_se & ratio\_low & ratio\_upp \\
\midrule\noalign{}
\endfirsthead
\toprule\noalign{}
ratio & ratio\_se & ratio\_low & ratio\_upp \\
\midrule\noalign{}
\endhead
\bottomrule\noalign{}
\endlastfoot
0.658 & 0.02 & 0.619 & 0.698 \\
\end{longtable}

Dando como resultado que por cada 100 unidades monetarias que le ingresan a las mujeres se gastan 66 con un intervalo de confianza entre 0.62 y 0.69. Por último, análogamente para los hombres, la razón de gastos resulta muy similar que para las mujeres, como se muestra en la Tabla \ref{tab:tab-razon-gasto-ingreso-hombre}.

\begin{Shaded}
\begin{Highlighting}[]
\NormalTok{male\_subset }\OtherTok{\textless{}{-}}\NormalTok{ survey\_design }\SpecialCharTok{\%\textgreater{}\%}
  \FunctionTok{filter}\NormalTok{(Sex }\SpecialCharTok{==} \StringTok{"Male"}\NormalTok{)}

\NormalTok{male\_subset }\SpecialCharTok{\%\textgreater{}\%}
  \FunctionTok{summarise}\NormalTok{(}\AttributeTok{ratio =} \FunctionTok{survey\_ratio}\NormalTok{(}\AttributeTok{numerator =}\NormalTok{ Expenditure,}
                                 \AttributeTok{denominator =}\NormalTok{ Income,}
                                 \AttributeTok{level =} \FloatTok{0.95}\NormalTok{,}
                                 \AttributeTok{vartype =} \FunctionTok{c}\NormalTok{(}\StringTok{"se"}\NormalTok{, }\StringTok{"ci"}\NormalTok{))}
\NormalTok{            )}
\end{Highlighting}
\end{Shaded}

\begin{longtable}[]{@{}lccc@{}}
\caption{\label{tab:tab-razon-gasto-ingreso-hombre}Razón gasto/ingreso en hogares con jefatura masculina}\tabularnewline
\toprule\noalign{}
ratio & ratio\_se & ratio\_low & ratio\_upp \\
\midrule\noalign{}
\endfirsthead
\toprule\noalign{}
ratio & ratio\_se & ratio\_low & ratio\_upp \\
\midrule\noalign{}
\endhead
\bottomrule\noalign{}
\endlastfoot
0.639 & 0.029 & 0.582 & 0.696 \\
\end{longtable}

\section{Estimación de parámetros de distribución y desigualdad}\label{estimacion-de-parametros-de-distribucion-y-desigualdad}

\subsection{Estimación de la dispersión}\label{estimacion-de-la-dispersion}

En las encuestas de hogares también es fundamental estimar medidas de dispersión de las variables estudiadas, ya que estas permiten evaluar el grado de heterogeneidad existente en la población. Por ejemplo, conocer qué tan dispersos son los ingresos de los hogares resulta clave para analizar desigualdades económicas y orientar el diseño de políticas públicas. Entre las medidas de dispersión más utilizadas se encuentra la desviación estándar, la cual cuantifica qué tan alejados se encuentran los valores observados respecto a su media. A continuación, se presenta el estimador de este parámetro:

\[
\hat s_y  =  \sqrt{\frac{\sum_{h}\sum_{i}\sum_{k} w_{hik} \ \left(y_{hik}-\hat{\bar{y}}\right)^{2}}{\sum_{h}\sum_{i}\sum_{k} w_{hik}-1}}
\]
Para estimar la desviación estándar de variables numéricas en encuestas de hogares utilizando \texttt{R}, puede emplearse la función \texttt{survey\_var}. A continuación, se presenta un ejemplo de su uso aplicado a la estimación de la desviación estándar del ingreso, cuyos resultados se muestran en la Tabla \ref{tab:tab-sd-ingreso-zona}.

\begin{Shaded}
\begin{Highlighting}[]
\NormalTok{survey\_design }\SpecialCharTok{\%\textgreater{}\%}
  \FunctionTok{group\_by}\NormalTok{(Zone) }\SpecialCharTok{\%\textgreater{}\%}
  \FunctionTok{summarise}\NormalTok{(}
    \AttributeTok{Var =} \FunctionTok{survey\_var}\NormalTok{(}
\NormalTok{      Income,}
      \AttributeTok{level =} \FloatTok{0.95}\NormalTok{,}
      \AttributeTok{vartype =} \FunctionTok{c}\NormalTok{(}\StringTok{"se"}\NormalTok{, }\StringTok{"ci"}\NormalTok{),}
      \AttributeTok{na.rm =} \ConstantTok{TRUE}
\NormalTok{    )}
\NormalTok{  ) }\SpecialCharTok{\%\textgreater{}\%}
  \FunctionTok{mutate}\NormalTok{(}
    \AttributeTok{Sd     =} \FunctionTok{sqrt}\NormalTok{(Var),}
    \AttributeTok{Sd\_low =} \FunctionTok{sqrt}\NormalTok{(Var\_low),}
    \AttributeTok{Sd\_upp =} \FunctionTok{sqrt}\NormalTok{(Var\_upp)}
\NormalTok{  )}
\end{Highlighting}
\end{Shaded}

\begin{longtable}[]{@{}lccccccc@{}}
\caption{\label{tab:tab-sd-ingreso-zona}Desviación estándar de los ingresos desagregada por zona}\tabularnewline
\toprule\noalign{}
Zone & Var & Var\_se & Var\_low & Var\_upp & Sd & Sd\_low & Sd\_upp \\
\midrule\noalign{}
\endfirsthead
\toprule\noalign{}
Zone & Var & Var\_se & Var\_low & Var\_upp & Sd & Sd\_low & Sd\_upp \\
\midrule\noalign{}
\endhead
\bottomrule\noalign{}
\endlastfoot
Rural & 96274 & 13794 & 68960 & 123588 & 310 & 263 & 352 \\
Urban & 338628 & 81241 & 177763 & 499494 & 582 & 422 & 707 \\
\end{longtable}

Como se observó en el ejemplo anterior, se estimó la desviación estándar de los ingresos por zona, reportando además su correspondiente intervalo de confianza al 95\%. Los argumentos de la función \texttt{survey\_var()} son equivalentes a los utilizados previamente para la estimación de medias y totales. Si el interés es estimar la desviación estándar de los ingresos desagregando simultáneamente por sexo y zona, los códigos computacionales correspondientes se presentan en la Tabla \ref{tab:tab-sd-ingreso-zona-sexo}.

\begin{Shaded}
\begin{Highlighting}[]
\NormalTok{survey\_design }\SpecialCharTok{\%\textgreater{}\%}
  \FunctionTok{group\_by}\NormalTok{(Zone, Sex) }\SpecialCharTok{\%\textgreater{}\%}
  \FunctionTok{summarise}\NormalTok{(}
    \AttributeTok{Var =} \FunctionTok{survey\_var}\NormalTok{(}
\NormalTok{      Income,}
      \AttributeTok{level =} \FloatTok{0.95}\NormalTok{,}
      \AttributeTok{vartype =} \FunctionTok{c}\NormalTok{(}\StringTok{"se"}\NormalTok{, }\StringTok{"ci"}\NormalTok{),}
      \AttributeTok{na.rm =} \ConstantTok{TRUE}
\NormalTok{    )}
\NormalTok{  ) }\SpecialCharTok{\%\textgreater{}\%}
  \FunctionTok{mutate}\NormalTok{(}
    \AttributeTok{Sd     =} \FunctionTok{sqrt}\NormalTok{(Var),}
    \AttributeTok{Sd\_low =} \FunctionTok{sqrt}\NormalTok{(Var\_low),}
    \AttributeTok{Sd\_upp =} \FunctionTok{sqrt}\NormalTok{(Var\_upp)}
\NormalTok{  )}
\end{Highlighting}
\end{Shaded}

\begin{longtable}[]{@{}
  >{\raggedright\arraybackslash}p{(\linewidth - 16\tabcolsep) * \real{0.0870}}
  >{\centering\arraybackslash}p{(\linewidth - 16\tabcolsep) * \real{0.1159}}
  >{\centering\arraybackslash}p{(\linewidth - 16\tabcolsep) * \real{0.1159}}
  >{\centering\arraybackslash}p{(\linewidth - 16\tabcolsep) * \real{0.1159}}
  >{\centering\arraybackslash}p{(\linewidth - 16\tabcolsep) * \real{0.1304}}
  >{\centering\arraybackslash}p{(\linewidth - 16\tabcolsep) * \real{0.1304}}
  >{\centering\arraybackslash}p{(\linewidth - 16\tabcolsep) * \real{0.0725}}
  >{\centering\arraybackslash}p{(\linewidth - 16\tabcolsep) * \real{0.1159}}
  >{\centering\arraybackslash}p{(\linewidth - 16\tabcolsep) * \real{0.1159}}@{}}
\caption{\label{tab:tab-sd-ingreso-zona-sexo}Desviación estándar de los ingresos desagregada por zona y sexo}\tabularnewline
\toprule\noalign{}
\begin{minipage}[b]{\linewidth}\raggedright
Zone
\end{minipage} & \begin{minipage}[b]{\linewidth}\centering
Sex
\end{minipage} & \begin{minipage}[b]{\linewidth}\centering
Var
\end{minipage} & \begin{minipage}[b]{\linewidth}\centering
Var\_se
\end{minipage} & \begin{minipage}[b]{\linewidth}\centering
Var\_low
\end{minipage} & \begin{minipage}[b]{\linewidth}\centering
Var\_upp
\end{minipage} & \begin{minipage}[b]{\linewidth}\centering
Sd
\end{minipage} & \begin{minipage}[b]{\linewidth}\centering
Sd\_low
\end{minipage} & \begin{minipage}[b]{\linewidth}\centering
Sd\_upp
\end{minipage} \\
\midrule\noalign{}
\endfirsthead
\toprule\noalign{}
\begin{minipage}[b]{\linewidth}\raggedright
Zone
\end{minipage} & \begin{minipage}[b]{\linewidth}\centering
Sex
\end{minipage} & \begin{minipage}[b]{\linewidth}\centering
Var
\end{minipage} & \begin{minipage}[b]{\linewidth}\centering
Var\_se
\end{minipage} & \begin{minipage}[b]{\linewidth}\centering
Var\_low
\end{minipage} & \begin{minipage}[b]{\linewidth}\centering
Var\_upp
\end{minipage} & \begin{minipage}[b]{\linewidth}\centering
Sd
\end{minipage} & \begin{minipage}[b]{\linewidth}\centering
Sd\_low
\end{minipage} & \begin{minipage}[b]{\linewidth}\centering
Sd\_upp
\end{minipage} \\
\midrule\noalign{}
\endhead
\bottomrule\noalign{}
\endlastfoot
Rural & Female & 86947 & 12459 & 62277 & 111618 & 295 & 250 & 334 \\
Rural & Male & 106119 & 15616 & 75197 & 137040 & 326 & 274 & 370 \\
Urban & Female & 323069 & 82058 & 160585 & 485553 & 568 & 401 & 697 \\
Urban & Male & 356141 & 83488 & 190826 & 521456 & 597 & 437 & 722 \\
\end{longtable}

\subsection{Estimación de percentiles}\label{estimacion-de-percentiles}

Las medidas de posición no central, como los percentiles y la mediana, permiten identificar puntos específicos de la distribución de una variable de interés distintos de su valor promedio. En las encuestas de hogares, este tipo de medidas resulta especialmente útil para caracterizar la distribución de variables económicas y sociales, como el ingreso, el gasto o las horas trabajadas. La mediana, por ejemplo, corresponde al valor que divide a la población en dos partes iguales: el 50\% de las observaciones se encuentra por debajo de dicho valor y el otro 50\% por encima. A diferencia de la media, la mediana es poco sensible a valores extremos, razón por la cual suele considerarse una medida robusta de tendencia central.

De manera similar, los percentiles permiten identificar segmentos específicos de la distribución poblacional. Por ejemplo, los percentiles de ingreso pueden utilizarse para identificar a la población perteneciente al 10\% superior de la distribución con fines tributarios, o para focalizar subsidios dirigidos a los hogares ubicados en los percentiles más bajos. Según \citet{loomis2005}, la estimación de cuantiles en encuestas complejas se basa en el uso de estimadores ponderados y en la estimación de la función de distribución acumulada (FDA) de la población. Para una población finita de tamaño \(N\), un estimador de esta distribución está dado por:

\[
\hat{F}\left(x\right) = \frac{\sum_{h} \sum_{i} \sum_{k} w_{hik} \ I\left(y_{k}\leq x\right)}{\sum_{h}\sum_{i}\sum_{k} w_{hik}}
\]

en donde la función \(I\left(y_{k}\leq x\right)\) es una variable indicadora que toma el valor uno si \(y_i \leq x\) y cero en otro caso.

Una vez estimada la FDA utilizando los pesos del diseño muestral, el cuantil \(q\)-ésimo de una variable \(y\) se define como el menor valor de \(y\) para el cual la FDA es mayor o igual que \(q\). En particular, la mediana corresponde al valor para el cual la FDA alcanza o supera 0.5; por tanto, la mediana estimada es el valor donde la FDA estimada es mayor o igual a 0.5. De esta forma, siguiendo las recomendaciones de \citet{heeringa2017applied} para la estimación de cuantiles en encuestas complejas, se ordenan las observaciones en orden ascendente (estadísticas de orden) de forma que \(y_{(1)} \geq y_{(2)} \geq  \ldots \geq y_{(n)}\) y se identifica el valor de \(j\) \((j=1,\ldots,n)\) tal que:

\[
\hat{F}\left(y_{(j)}\right)\leq q\leq\hat{F}\left(y_{(j+1)}\right)
\]

Bajo esta condición, el estimador del cuantil \(q\)-ésimo se obtiene mediante interpolación lineal, la cual permite obtener estimaciones más estables y precisas de los cuantiles, especialmente cuando la FDA presenta saltos importantes asociados al uso de pesos muestrales. En cuanto a la estimación de la varianza de percentiles y medianas, así como a la construcción de intervalos de confianza, \citet{kovar1988} realizaron un estudio de simulación en el contexto de diseños complejos y recomiendan el uso de métodos de replicación debido a su buen desempeño para este tipo de parámetros no lineales.

Tanto los estimadores de cuantiles como los métodos para estimar sus varianzas se encuentran implementados en \texttt{R}. En particular, la función \texttt{survey\_median()} permite obtener la mediana junto con su error estándar e intervalo de confianza. A continuación, se presenta la sintaxis para estimar la mediana del gasto de los hogares utilizando la base de datos de ejemplo, cuyos resultados se muestran en la Tabla \ref{tab:tab-mediana-gasto}.

\begin{Shaded}
\begin{Highlighting}[]
\NormalTok{survey\_design }\SpecialCharTok{\%\textgreater{}\%}
  \FunctionTok{summarise}\NormalTok{(}\AttributeTok{median =} \FunctionTok{survey\_median}\NormalTok{(Expenditure,}
                                    \AttributeTok{level =} \FloatTok{0.95}\NormalTok{,}
                                    \AttributeTok{vartype =} \FunctionTok{c}\NormalTok{(}\StringTok{"se"}\NormalTok{, }\StringTok{"ci"}\NormalTok{)}
\NormalTok{                                    )}
\NormalTok{            )}
\end{Highlighting}
\end{Shaded}

\begin{longtable}[]{@{}lccc@{}}
\caption{\label{tab:tab-mediana-gasto}Estimación de la mediana de los gastos de los hogares}\tabularnewline
\toprule\noalign{}
median & median\_se & median\_low & median\_upp \\
\midrule\noalign{}
\endfirsthead
\toprule\noalign{}
median & median\_se & median\_low & median\_upp \\
\midrule\noalign{}
\endhead
\bottomrule\noalign{}
\endlastfoot
298 & 8.82 & 282 & 317 \\
\end{longtable}

Como puede observarse, los argumentos de la función \texttt{survey\_median()} son equivalentes a los utilizados previamente para estimar totales y medias. Al igual que con otros parámetros poblacionales, la mediana también puede estimarse para distintos dominios de estudio, como zona geográfica, sexo o grupos de edad. La Tabla \ref{tab:tab-mediana-gasto-zona} muestra los resultados por zona geográfica.

\begin{Shaded}
\begin{Highlighting}[]
\NormalTok{survey\_design }\SpecialCharTok{\%\textgreater{}\%}
  \FunctionTok{group\_by}\NormalTok{(Zone) }\SpecialCharTok{\%\textgreater{}\%}
  \FunctionTok{summarise}\NormalTok{(}\AttributeTok{median =} \FunctionTok{survey\_median}\NormalTok{(Expenditure,}
                                    \AttributeTok{level =} \FloatTok{0.95}\NormalTok{,}
                                    \AttributeTok{vartype =} \FunctionTok{c}\NormalTok{(}\StringTok{"se"}\NormalTok{, }\StringTok{"ci"}\NormalTok{)}
\NormalTok{                                    )}
\NormalTok{            )}
\end{Highlighting}
\end{Shaded}

\begin{longtable}[]{@{}lcccc@{}}
\caption{\label{tab:tab-mediana-gasto-zona}Estimación de la mediana de los gastos desagregada por zona}\tabularnewline
\toprule\noalign{}
Zone & median & median\_se & median\_low & median\_upp \\
\midrule\noalign{}
\endfirsthead
\toprule\noalign{}
Zone & median & median\_se & median\_low & median\_upp \\
\midrule\noalign{}
\endhead
\bottomrule\noalign{}
\endlastfoot
Rural & 241 & 11.0 & 214 & 258 \\
Urban & 381 & 19.8 & 337 & 416 \\
\end{longtable}

Si el objetivo es estimar un percentil específico diferente de la mediana, por ejemplo el percentil 25 (primer cuartil), puede utilizarse la función \texttt{survey\_quantile()} de la siguiente manera. Los resultados correspondientes se presentan en la Tabla \ref{tab:tab-cuantil-gasto}.

\begin{Shaded}
\begin{Highlighting}[]
\NormalTok{survey\_design }\SpecialCharTok{\%\textgreater{}\%}
  \FunctionTok{summarise}\NormalTok{(}\AttributeTok{quantile =} \FunctionTok{survey\_quantile}\NormalTok{(Expenditure,}
                                \AttributeTok{quantiles =} \FloatTok{0.25}\NormalTok{,}
                                \AttributeTok{level =} \FloatTok{0.95}\NormalTok{,}
                                \AttributeTok{vartype =} \FunctionTok{c}\NormalTok{(}\StringTok{"se"}\NormalTok{, }\StringTok{"ci"}\NormalTok{),}
                                \AttributeTok{interval\_type =} \StringTok{"score"}
\NormalTok{                                )}
\NormalTok{            )}
\end{Highlighting}
\end{Shaded}

\begin{longtable}[]{@{}lccc@{}}
\caption{\label{tab:tab-cuantil-gasto}Estimación del primer cuartil de los gastos de los hogares}\tabularnewline
\toprule\noalign{}
quantile\_q25 & quantile\_q25\_se & quantile\_q25\_low & quantile\_q25\_upp \\
\midrule\noalign{}
\endfirsthead
\toprule\noalign{}
quantile\_q25 & quantile\_q25\_se & quantile\_q25\_low & quantile\_q25\_upp \\
\midrule\noalign{}
\endhead
\bottomrule\noalign{}
\endlastfoot
200 & 13.8 & 163 & 218 \\
\end{longtable}

Observe que se incorporó el argumento \texttt{interval\_type\ =\ "score"}, el cual indica que los intervalos de confianza se construyen mediante el método \emph{score} basado en la función de influencia del percentil. Este enfoque es recomendado por \citet{Lumley2010} para el análisis de datos provenientes de diseños muestrales complejos. De manera análoga, la estimación del percentil 25 puede obtenerse para distintos subgrupos de la población, como sexo o zona geográfica, utilizando \texttt{group\_by()}, tal como se muestra en la Tabla \ref{tab:tab-cuantil-gasto-zona-sexo}.

\begin{Shaded}
\begin{Highlighting}[]
\NormalTok{survey\_design }\SpecialCharTok{\%\textgreater{}\%}
  \FunctionTok{group\_by}\NormalTok{(Zone, Sex) }\SpecialCharTok{\%\textgreater{}\%}
  \FunctionTok{summarise}\NormalTok{(}\AttributeTok{quantile =} \FunctionTok{survey\_quantile}\NormalTok{(Expenditure,}
                                \AttributeTok{quantiles =} \FloatTok{0.25}\NormalTok{,}
                                \AttributeTok{level =} \FloatTok{0.95}\NormalTok{,}
                                \AttributeTok{vartype =} \FunctionTok{c}\NormalTok{(}\StringTok{"se"}\NormalTok{, }\StringTok{"ci"}\NormalTok{),}
                                \AttributeTok{interval\_type =} \StringTok{"score"}
\NormalTok{                                )}
\NormalTok{            )}
\end{Highlighting}
\end{Shaded}

\begin{longtable}[]{@{}
  >{\raggedright\arraybackslash}p{(\linewidth - 10\tabcolsep) * \real{0.0741}}
  >{\centering\arraybackslash}p{(\linewidth - 10\tabcolsep) * \real{0.0988}}
  >{\centering\arraybackslash}p{(\linewidth - 10\tabcolsep) * \real{0.1728}}
  >{\centering\arraybackslash}p{(\linewidth - 10\tabcolsep) * \real{0.2099}}
  >{\centering\arraybackslash}p{(\linewidth - 10\tabcolsep) * \real{0.2222}}
  >{\centering\arraybackslash}p{(\linewidth - 10\tabcolsep) * \real{0.2222}}@{}}
\caption{\label{tab:tab-cuantil-gasto-zona-sexo}Primer cuartil de los gastos desagregado por zona y sexo}\tabularnewline
\toprule\noalign{}
\begin{minipage}[b]{\linewidth}\raggedright
Zone
\end{minipage} & \begin{minipage}[b]{\linewidth}\centering
Sex
\end{minipage} & \begin{minipage}[b]{\linewidth}\centering
quantile\_q25
\end{minipage} & \begin{minipage}[b]{\linewidth}\centering
quantile\_q25\_se
\end{minipage} & \begin{minipage}[b]{\linewidth}\centering
quantile\_q25\_low
\end{minipage} & \begin{minipage}[b]{\linewidth}\centering
quantile\_q25\_upp
\end{minipage} \\
\midrule\noalign{}
\endfirsthead
\toprule\noalign{}
\begin{minipage}[b]{\linewidth}\raggedright
Zone
\end{minipage} & \begin{minipage}[b]{\linewidth}\centering
Sex
\end{minipage} & \begin{minipage}[b]{\linewidth}\centering
quantile\_q25
\end{minipage} & \begin{minipage}[b]{\linewidth}\centering
quantile\_q25\_se
\end{minipage} & \begin{minipage}[b]{\linewidth}\centering
quantile\_q25\_low
\end{minipage} & \begin{minipage}[b]{\linewidth}\centering
quantile\_q25\_upp
\end{minipage} \\
\midrule\noalign{}
\endhead
\bottomrule\noalign{}
\endlastfoot
Rural & Female & 160 & 7.18 & 135 & 163 \\
Rural & Male & 163 & 3.45 & 150 & 163 \\
Urban & Female & 258 & 10.40 & 249 & 291 \\
Urban & Male & 259 & 10.03 & 257 & 297 \\
\end{longtable}

\subsection{Estimación del coeficiente de Gini}\label{estimacion-del-coeficiente-de-gini}

El tema sobre desigualdad económica trasciende el ámbito estrictamente estadístico y constituye uno de los temas centrales del análisis social contemporáneo. Existe un amplio consenso respecto a la importancia de comprender cómo las desigualdades económicas condicionan las oportunidades, el bienestar y la calidad de vida de las personas. En este sentido, la medición rigurosa de la desigualdad representa un insumo fundamental para el diseño, seguimiento y evaluación de políticas públicas orientadas a la equidad social.

Entre los indicadores más ampliamente utilizados para cuantificar la desigualdad económica se encuentra el coeficiente de Gini (\(G\)), el cual mide el grado de concentración del ingreso comparando la distribución observada con una situación hipotética de igualdad perfecta. Este coeficiente toma valores entre 0 y 1, donde \(G = 0\) representa igualdad perfecta y \(G = 1\) corresponde al máximo nivel de desigualdad posible. Valores más elevados indican una mayor concentración del ingreso en un grupo reducido de la población.

En el contexto de las encuestas de hogares, la estimación del coeficiente de Gini debe incorporar adecuadamente las características del diseño muestral, incluyendo los factores de expansión, la estratificación y la conglomeración. En muchos casos, los pesos muestrales son previamente normalizados con el fin de simplificar los procedimientos de estimación y facilitar el procesamiento computacional. De acuerdo con \citet{binder1995}, el estimador del coeficiente de Gini puede expresarse como:

\[
\hat{G} = \frac{2\sum_{h}\sum_{i}\sum_{k}w_{hik}^{*}
\ \hat{F}_{hik} \ y_{hik}-1}{\hat{\bar{y}}}
\]

donde \(w_{hik}^{*}=\dfrac{w_{hik}}{\sum_{h}\sum_{i}\sum_{k}w_{hik}}\) corresponde al peso muestral normalizado, \(\hat{F}_{hik}\) representa la función de distribución acumulada (FDA) estimada para la observación \(k\) dentro del conglomerado \(i\) del estrato \(h\), y \(\hat{\bar{y}}\) denota la media ponderada de la variable de ingreso en la población.

Autores como \citet{osier2009variance} y \citet{langel2013variance} profundizan en aspectos técnicos adicionales, especialmente en lo relacionado con la estimación de la varianza del coeficiente de Gini bajo diseños complejos. El paquete \texttt{convey} \citep{convey2024} implementa los procedimientos recomendados para calcular el coeficiente de Gini y su varianza en encuestas de hogares. Primero, se prepara el diseño muestral mediante \texttt{convey\_prep()}. Luego, la función \texttt{svygini()} calcula el índice de Gini utilizando la variable de ingreso y el diseño de encuesta complejo especificado. A continuación se muestra la sintaxis para estimarlo en la base de datos de ejemplo, presentada en la Tabla \ref{tab:tab-gini}:

\begin{Shaded}
\begin{Highlighting}[]
\FunctionTok{library}\NormalTok{(convey)}
\NormalTok{gini\_design }\OtherTok{\textless{}{-}} \FunctionTok{convey\_prep}\NormalTok{(survey\_design)}
\NormalTok{gini }\OtherTok{\textless{}{-}} \FunctionTok{svygini}\NormalTok{(}\SpecialCharTok{\textasciitilde{}}\NormalTok{Income, }\AttributeTok{design =}\NormalTok{ gini\_design)}

\FunctionTok{data.frame}\NormalTok{(}\AttributeTok{Gini =} \FunctionTok{coef}\NormalTok{(gini),}
           \AttributeTok{standard\_error =} \FunctionTok{SE}\NormalTok{(gini),}
           \AttributeTok{ci\_lower =} \FunctionTok{confint}\NormalTok{(gini)[}\DecValTok{1}\NormalTok{],}
           \AttributeTok{ci\_upper =} \FunctionTok{confint}\NormalTok{(gini)[}\DecValTok{2}\NormalTok{])}
\end{Highlighting}
\end{Shaded}

\begin{longtable}[]{@{}llccc@{}}
\caption{\label{tab:tab-gini}Estimación del coeficiente de Gini para el ingreso de los hogares}\tabularnewline
\toprule\noalign{}
& Gini & Income & ci\_lower & ci\_upper \\
\midrule\noalign{}
\endfirsthead
\toprule\noalign{}
& Gini & Income & ci\_lower & ci\_upper \\
\midrule\noalign{}
\endhead
\bottomrule\noalign{}
\endlastfoot
Income & 0.413 & 0.019 & 0.377 & 0.45 \\
\end{longtable}

Si el interés se centra en la estimación de la curva de Lorenz, es importante recordar que, de acuerdo con \citet{kovacevic1997}, esta representa la relación entre el porcentaje acumulado de la población (ordenada desde los hogares con menores ingresos hasta aquellos con mayores ingresos) y la proporción acumulada del ingreso total que concentra dicha población. La línea diagonal de 45 grados corresponde a una situación hipotética de igualdad perfecta, en la cual todos los individuos participan proporcionalmente del ingreso total.

El área comprendida entre la curva de Lorenz y la diagonal de igualdad se conoce como el área de Lorenz, y el coeficiente de Gini puede interpretarse como el doble de esta área relativa. En consecuencia, cuanto más cercana se encuentre la curva de Lorenz a la diagonal, mayor será el nivel de equidad en la distribución del ingreso; por el contrario, curvas más alejadas reflejan mayores niveles de desigualdad económica.

Para realizar la curva de Lorenz en \texttt{R} se utiliza la función \texttt{svylorenz()}, presentada en la Figura \ref{fig:fig-lorenz}:

\begin{Shaded}
\begin{Highlighting}[]
\NormalTok{clorenz }\OtherTok{\textless{}{-}} \FunctionTok{svylorenz}\NormalTok{(}
  \AttributeTok{formula =} \SpecialCharTok{\textasciitilde{}}\NormalTok{Income,}
  \AttributeTok{design =}\NormalTok{ gini\_design,}
  \AttributeTok{quantiles =} \FunctionTok{seq}\NormalTok{(}\DecValTok{0}\NormalTok{, }\DecValTok{1}\NormalTok{, .}\DecValTok{05}\NormalTok{),}
  \AttributeTok{alpha =}\NormalTok{ .}\DecValTok{01}
\NormalTok{)}
\end{Highlighting}
\end{Shaded}

\begin{figure}

{\centering \includegraphics[width=0.65\linewidth]{docs/figures/fig-lorenz-1} 

}

\caption{Curva de Lorenz estimada para el ingreso de los hogares}\label{fig:fig-lorenz}
\end{figure}

El argumento \texttt{formula\ =\ \textasciitilde{}Income} especifica la variable de ingreso sobre la cual se calculará la distribución acumulada. El argumento \texttt{design\ =\ gini\_design} indica el objeto de diseño muestral previamente definido, el cual contiene la información relacionada con pesos, estratos y conglomerados de la encuesta. Por su parte, \texttt{quantiles\ =\ seq(0,\ 1,\ .05)} define los puntos de la distribución en los cuales se evaluará la curva de Lorenz, generando percentiles desde 0 hasta 1 en incrementos de 0.05. Finalmente, el argumento \texttt{alpha\ =\ .01} establece un nivel de significación del 1\%, lo que equivale a construir intervalos de confianza del 99\% para las estimaciones asociadas a la curva.

\subsection{Estimación de la correlación}\label{estimacion-de-la-correlacion}

En el estudio de encuestas de hogares, además de describir variables de forma individual, resulta fundamental analizar cómo se relacionan entre sí. Una de las herramientas más utilizadas para este propósito es el coeficiente de correlación de Pearson, el cual mide la fuerza y la dirección de la relación lineal entre dos variables numéricas. Este coeficiente toma valores entre --1 y 1. Un valor positivo indica que ambas variables tienden a aumentar simultáneamente, mientras que un valor negativo señala que cuando una variable aumenta, la otra tiende a disminuir. Por su parte, valores cercanos a cero sugieren la ausencia de una relación lineal fuerte entre las variables analizadas.

Por ejemplo, en una encuesta de hogares puede ser de interés estudiar si existe una asociación entre el ingreso de los hogares y su nivel de gasto, así como evaluar la magnitud de dicha relación. Este tipo de análisis permite comprender mejor los patrones económicos y sociales observados en la población.

En encuestas complejas, es necesario incorporar los pesos muestrales para que la estimación sea representativa de la población. Este ajuste tiene en cuenta la estratificación, la conglomeración y las probabilidades de selección desiguales. El cálculo ponderado implica evaluar la covarianza entre las dos variables y dividirla entre el producto de sus desviaciones estándar ponderadas, eliminando así la influencia de las unidades de medida. Asumiendo que \(\hat{\bar{x}}\) es una estimación de la media poblacional de la variabnel \(x\), entonces el coeficiente de correlación de Pearson ajustado por pesos se expresa como:

\[
\hat{\rho}_{xy} = \frac{\displaystyle \sum_{h} \sum_{i} \sum_{k} w_{hik} (y_{hik} - \hat{\bar{y}})(x_{hik} - \hat{\bar{x}})} {\sqrt{\displaystyle \sum_{h} \sum_{i} \sum_{k} w_{hik} (y_{hik} - \hat{\bar{y}})^2} \sqrt{\displaystyle \sum_{h} \sum_{i} \sum_{k} w_{hik} (x_{hik} - \hat{\bar{x}})^2}}
\]

El paquete \texttt{survey} cuenta con la función \texttt{svyvar()} que permite obtener matrices de covarianzas ponderadas, a partir de las cuales se puede calcular la correlación. La Tabla \ref{tab:tab-covarianza} presenta un ejemplo de estimación para la matriz de varianzas covarianzas entre el ingreso y el gasto de los hogares.

\begin{Shaded}
\begin{Highlighting}[]
\NormalTok{cov\_matrix }\OtherTok{\textless{}{-}} \FunctionTok{svyvar}\NormalTok{(}\SpecialCharTok{\textasciitilde{}}\NormalTok{Income }\SpecialCharTok{+}\NormalTok{ Expenditure, }\AttributeTok{design =}\NormalTok{ survey\_design)}
\end{Highlighting}
\end{Shaded}

\begin{longtable}[]{@{}llcc@{}}
\caption{\label{tab:tab-covarianza}Matriz de varianzas y covarianzas ponderada entre ingreso y gasto}\tabularnewline
\toprule\noalign{}
& Variable & Income & Expenditure \\
\midrule\noalign{}
\endfirsthead
\toprule\noalign{}
& Variable & Income & Expenditure \\
\midrule\noalign{}
\endhead
\bottomrule\noalign{}
\endlastfoot
Income & Income & 243719 & 98337 \\
Expenditure & Expenditure & 98337 & 77887 \\
\end{longtable}

La función \texttt{svyvar} estima la matriz de varianzas y covarianzas considerando el diseño muestral. A partir de dicha matriz puede obtenerse la correlación aplicando las siguientes instrucciones:

\begin{Shaded}
\begin{Highlighting}[]
\NormalTok{cov\_xy }\OtherTok{\textless{}{-}}\NormalTok{ cov\_matrix[}\DecValTok{1}\NormalTok{, }\DecValTok{2}\NormalTok{]}
\NormalTok{sd\_x }\OtherTok{\textless{}{-}} \FunctionTok{sqrt}\NormalTok{(cov\_matrix[}\DecValTok{1}\NormalTok{, }\DecValTok{1}\NormalTok{])}
\NormalTok{sd\_y }\OtherTok{\textless{}{-}} \FunctionTok{sqrt}\NormalTok{(cov\_matrix[}\DecValTok{2}\NormalTok{, }\DecValTok{2}\NormalTok{])}

\NormalTok{weighted\_cor }\OtherTok{\textless{}{-}}\NormalTok{ cov\_xy }\SpecialCharTok{/}\NormalTok{ (sd\_x }\SpecialCharTok{*}\NormalTok{ sd\_y)}
\NormalTok{weighted\_cor}
\end{Highlighting}
\end{Shaded}

\begin{verbatim}
## [1] 0.714
\end{verbatim}

Si además se desea realizar inferencia estadística sobre la correlación, por ejemplo obtener errores estándar o intervalos de confianza, puede utilizarse métodos de replicación (bootstrap/jackknife) compatibles con el paquete survey. Otra opción es utilizar la función \texttt{svycor} del paquete \texttt{jtools} para estimar directamente la correlación:

\begin{Shaded}
\begin{Highlighting}[]
\FunctionTok{library}\NormalTok{(jtools)}

\FunctionTok{svycor}\NormalTok{(}\SpecialCharTok{\textasciitilde{}}\NormalTok{Income }\SpecialCharTok{+}\NormalTok{ Expenditure, }\AttributeTok{design =}\NormalTok{ survey\_design)}
\end{Highlighting}
\end{Shaded}

\begin{verbatim}
##             Income Expenditure
## Income        1.00        0.71
## Expenditure   0.71        1.00
\end{verbatim}

\section{Pruebas de hipótesis}\label{pruebas-de-hipotesis}

En el análisis de datos provenientes de encuestas de hogares no basta con estudiar cada variable de manera aislada, como estimar el ingreso medio de hombres y mujeres en un país. También es fundamental comparar grupos y evaluar si las diferencias observadas entre ellos reflejan desigualdades reales en la población o si podrían explicarse simplemente por el error de muestreo. Por ejemplo, una pregunta frecuente consiste en determinar si existen diferencias estadísticamente significativas en el ingreso medio entre hogares dirigidos por hombres y hogares dirigidos por mujeres. Este tipo de análisis permite identificar brechas sociales y económicas, además de generar evidencia útil para el diseño y evaluación de políticas públicas.

Para responder este tipo de interrogantes se utilizan pruebas de hipótesis, es decir, procedimientos estadísticos que permiten contrastar afirmaciones sobre parámetros poblacionales utilizando la información observada en la muestra. En el contexto de encuestas de hogares, estas pruebas deben incorporar adecuadamente las características del diseño muestral con el fin de garantizar inferencias válidas y resultados confiables.

Las pruebas de hipótesis constituyen una herramienta fundamental para evaluar afirmaciones sobre parámetros poblacionales a partir de la información observada en una muestra. En términos generales, toda prueba de hipótesis se basa en el contraste entre dos proposiciones opuestas: la hipótesis nula, denotada por \(H_0\), y la hipótesis alternativa, denotada por \(H_1\). En el caso más común de una prueba bilateral, estas hipótesis se expresan como:

\[
\begin{cases}
H_{0}: & \theta = \theta_0 \\
H_{1}: & \theta \neq \theta_0
\end{cases}
\]

donde \(\theta\) representa el parámetro poblacional de interés y \(\theta_0\) un valor de referencia especificado bajo la hipótesis nula. Dependiendo del objetivo del análisis, la hipótesis alternativa también puede formularse de manera unilateral, por ejemplo, cuando se desea evaluar si \(\theta > \theta_0\) o si \(\theta < \theta_0\). El propósito de la prueba es determinar si la evidencia contenida en la muestra resulta suficientemente fuerte para rechazar la hipótesis nula en favor de la hipótesis alternativa.

Supóngase que \(\bar{y}_1\) y \(\bar{y}_2\) representan las medias poblacionales de una variable \(y\) en dos dominios distintos, por ejemplo, el ingreso medio de hogares con jefatura masculina y el ingreso medio de hogares con jefatura femenina. Entonces, el parámetro de interés corresponde entonces a la diferencia:

\[
\Delta = \bar{y}_1 - \bar{y}_2
\]

Su estimador muestral se define como:

\[
\hat{\Delta} = \hat{\bar{y}}_1 - \hat{\bar{y}}_2
\]

Para el cual, su error estándar se estima mediante la siguiente expresión:

\[
\widehat{se}(\hat{\Delta}) =
\sqrt{
\widehat{Var}(\hat{\bar{y}}_1)
+
\widehat{Var}(\hat{\bar{y}}_2)
-
2\widehat{Cov}(\hat{\bar{y}}_1, \hat{\bar{y}}_2)
}
\]

La presencia del término de covarianza es importante porque las estimaciones de ambos dominios suelen provenir de la misma muestra y, por tanto, no necesariamente son independientes. Una vez calculado el estimador y su error estándar, el contraste de hipótesis se realiza utilizando el siguiente estadístico de prueba:

\[
t =
\frac{\hat{\Delta}}
{\widehat{se}(\hat{\Delta})} \sim t_{df}
\]

el cual sigue aproximadamente una distribución \(t\) de Student con \((df)\) grados de libertad. En la práctica, estos grados de libertad suelen aproximarse mediante \(df = n_I - H\), donde \(n_I\) representa el número de unidades primarias de muestreo y \(H\) el número de estratos. Bajo la hipótesis nula, valores absolutos elevados del estadístico \(t\) constituyen evidencia en contra de \(H_0\). De manera complementaria, también puede construirse un intervalo de confianza para el parámetro \(\Delta\). Para un nivel de confianza de \((1-\alpha)100\%\), dicho intervalo está dado por:

\[
\hat{\Delta}
\;\pm\;
t_{(1-\alpha/2,df)}
\,
\hat{se}(\hat{\Delta}).
\]

Este intervalo proporciona un rango de valores plausibles para la diferencia poblacional y constituye una herramienta útil para evaluar tanto la magnitud como la significación estadística de las diferencias observadas entre grupos.

En \texttt{R}, estas pruebas pueden implementarse mediante la función \texttt{svyttest()} del paquete \texttt{survey}, la cual incorpora automáticamente los ajustes asociados al diseño muestral. A partir de los resultados presentados en la Tabla \ref{tab:tab-ttest-sexo} puede concluirse que, con un nivel de confianza del 95\%, no existe evidencia estadística suficiente para afirmar que los ingresos medios difieren según el sexo.

\begin{Shaded}
\begin{Highlighting}[]
\NormalTok{ttest\_sex }\OtherTok{\textless{}{-}} \FunctionTok{svyttest}\NormalTok{(Income }\SpecialCharTok{\textasciitilde{}}\NormalTok{ Sex, }\AttributeTok{design =}\NormalTok{ survey\_design, }\AttributeTok{level =} \FloatTok{0.95}\NormalTok{)}
\end{Highlighting}
\end{Shaded}

\begin{longtable}[]{@{}
  >{\raggedright\arraybackslash}p{(\linewidth - 10\tabcolsep) * \real{0.1515}}
  >{\centering\arraybackslash}p{(\linewidth - 10\tabcolsep) * \real{0.1364}}
  >{\centering\arraybackslash}p{(\linewidth - 10\tabcolsep) * \real{0.0758}}
  >{\centering\arraybackslash}p{(\linewidth - 10\tabcolsep) * \real{0.1515}}
  >{\centering\arraybackslash}p{(\linewidth - 10\tabcolsep) * \real{0.1515}}
  >{\centering\arraybackslash}p{(\linewidth - 10\tabcolsep) * \real{0.3333}}@{}}
\caption{\label{tab:tab-ttest-sexo}Prueba de diferencia de ingresos medios por sexo (población total)}\tabularnewline
\toprule\noalign{}
\begin{minipage}[b]{\linewidth}\raggedright
statistic
\end{minipage} & \begin{minipage}[b]{\linewidth}\centering
p\_value
\end{minipage} & \begin{minipage}[b]{\linewidth}\centering
gl
\end{minipage} & \begin{minipage}[b]{\linewidth}\centering
ci\_lower
\end{minipage} & \begin{minipage}[b]{\linewidth}\centering
ci\_upper
\end{minipage} & \begin{minipage}[b]{\linewidth}\centering
estimated\_difference
\end{minipage} \\
\midrule\noalign{}
\endfirsthead
\toprule\noalign{}
\begin{minipage}[b]{\linewidth}\raggedright
statistic
\end{minipage} & \begin{minipage}[b]{\linewidth}\centering
p\_value
\end{minipage} & \begin{minipage}[b]{\linewidth}\centering
gl
\end{minipage} & \begin{minipage}[b]{\linewidth}\centering
ci\_lower
\end{minipage} & \begin{minipage}[b]{\linewidth}\centering
ci\_upper
\end{minipage} & \begin{minipage}[b]{\linewidth}\centering
estimated\_difference
\end{minipage} \\
\midrule\noalign{}
\endhead
\bottomrule\noalign{}
\endlastfoot
1.36 & 0.176 & 118 & -12.8 & 69.4 & 28.3 \\
\end{longtable}

Otra hipótesis de interés consiste en evaluar si los ingresos medios de los hogares difieren según el sexo de la jefatura dentro del área urbana, como se presenta en la Tabla \ref{tab:tab-ttest-sexo-urbano}. Los resultados indican que no se rechaza la hipótesis nula \(H_0\), por lo que no existe evidencia estadística suficiente para afirmar que los ingresos medios difieren por sexo en la zona urbana.

\begin{Shaded}
\begin{Highlighting}[]
\NormalTok{urban\_subset }\OtherTok{\textless{}{-}}\NormalTok{ survey\_design }\SpecialCharTok{\%\textgreater{}\%}
  \FunctionTok{filter}\NormalTok{(Zone }\SpecialCharTok{==} \StringTok{"Urban"}\NormalTok{)}

\NormalTok{ttest\_urban }\OtherTok{\textless{}{-}} \FunctionTok{svyttest}\NormalTok{(Income }\SpecialCharTok{\textasciitilde{}}\NormalTok{ Sex, }\AttributeTok{design =}\NormalTok{ urban\_subset, }\AttributeTok{level =} \FloatTok{0.95}\NormalTok{)}
\end{Highlighting}
\end{Shaded}

\begin{longtable}[]{@{}lccccc@{}}
\caption{\label{tab:tab-ttest-sexo-urbano}Prueba de diferencia de ingresos medios por sexo en zona urbana}\tabularnewline
\toprule\noalign{}
statistic & p\_value & gl & ci\_lower & ci\_upper & estimated\_difference \\
\midrule\noalign{}
\endfirsthead
\toprule\noalign{}
statistic & p\_value & gl & ci\_lower & ci\_upper & estimated\_difference \\
\midrule\noalign{}
\endhead
\bottomrule\noalign{}
\endlastfoot
1.57 & 0.122 & 63 & -12.3 & 102 & 44.7 \\
\end{longtable}

El procedimiento descrito no se limita a las medias, sino que también puede aplicarse a proporciones, totales, razones o cualquier función diferenciable de totales. En todos los casos, el contraste se fundamenta en la estimación puntual, su varianza (incluyendo covarianzas cuando corresponde) y la comparación con la distribución \(t\) ajustada al diseño muestral.

\section{Estimación de contrastes}\label{estimacion-de-contrastes}

En el análisis de encuestas de hogares es frecuente que el interés no se limite a comparar únicamente dos poblaciones, sino varias de manera simultánea. Por ejemplo, puede ser necesario comparar los ingresos medios de los hogares entre distintas regiones, áreas geográficas o grupos poblacionales, con el propósito de identificar diferencias y establecer patrones de desigualdad entre ellos. En este tipo de situaciones, las pruebas de diferencia de medias presentadas anteriormente resultan limitadas, ya que están diseñadas para comparar únicamente pares de poblaciones.

Para abordar comparaciones más generales se utilizan los contrastes, que constituyen una herramienta flexible para evaluar combinaciones lineales de parámetros poblacionales. En términos generales, un contraste se define como:

\[
f\left(\theta_1,\theta_2,\ldots,\theta_J\right)
=
\sum_{j=1}^{J} a_j \theta_j,
\]

donde los coeficientes \(a_j\) son constantes conocidas y \(\theta_j\) representan los parámetros de interés. Este enfoque permite formular y evaluar una amplia variedad de comparaciones entre grupos, incluyendo diferencias entre medias, comparaciones múltiples y combinaciones más complejas de parámetros.

En \texttt{R}, los procedimientos metodológicos para implementar contrastes en diseños de muestreo complejos se encuentran disponibles mediante la función \texttt{svycontrast()} del paquete \texttt{survey}. Por ejemplo, si se quiere comparar los ingresos medios de dos subpoblaciones particulares (las regiones Norte y Sur) se puede usar el contraste \(\bar{y}_{Norte} - \bar{y}_{Sur}\). Dado que la muestra cubre cinco regiones en total, el contraste debe construirse asignando coeficientes únicamente a las regiones involucradas en la comparación y dejando coeficientes iguales a cero para las demás regiones. De esta manera, el contraste se define como:

\[
1\times\hat{\bar{y}}_{Norte}+\left(-1\right)\times\hat{\bar{y}}_{Sur}+0\times\hat{\bar{y}}_{Centro}+0\times\hat{\bar{y}}_{Occidente}+0\times\hat{\bar{y}}_{Oriente}.
\]

Esta expresión indica que el contraste corresponde a la diferencia entre las medias estimadas de las regiones Norte y Sur, mientras que las demás regiones no participan en la comparación. De forma matricial, el contraste puede escribirse como:

\[
 \left[1,\,-1,\,0,\,0,\,0\right]
 \times
 \left[\begin{array}{c}
\hat{\bar{y}}_{Norte}\\
\hat{\bar{y}}_{Sur}\\
\hat{\bar{y}}_{Centro}\\
\hat{\bar{y}}_{Occidente}\\
\hat{\bar{y}}_{Oriente}
\end{array}\right].
\]

En consecuencia, el vector de contraste asociado a esta comparacion es \(\left[1,\,-1,\,0,\,0,\,0\right]\), donde los coeficientes positivos y negativos indican las poblaciones que se desean comparar y los ceros representan las poblaciones excluidas del contraste. El primer paso consiste en calcular las medias estimadas para cada region mediante \texttt{group\_by()} y \texttt{survey\_mean()}, como se presenta en la Tabla \ref{tab:tab-prom-region}. El objeto visible de resultados se construye con \texttt{srvyr}; adicionalmente se crea un objeto auxiliar con matriz de covarianzas para los contrastes posteriores. Como resultado se obtienen las medias estimadas de los ingresos por region.

\begin{Shaded}
\begin{Highlighting}[]
\NormalTok{region\_mean }\OtherTok{\textless{}{-}}\NormalTok{ survey\_design }\SpecialCharTok{\%\textgreater{}\%}
  \FunctionTok{group\_by}\NormalTok{(Region) }\SpecialCharTok{\%\textgreater{}\%}
  \FunctionTok{summarise}\NormalTok{(}\AttributeTok{mean =} \FunctionTok{survey\_mean}\NormalTok{(Income,}
                               \AttributeTok{na.rm =} \ConstantTok{TRUE}\NormalTok{,}
                               \AttributeTok{vartype =} \FunctionTok{c}\NormalTok{(}\StringTok{"se"}\NormalTok{, }\StringTok{"ci"}\NormalTok{)))}

\NormalTok{region\_mean\_contrast }\OtherTok{\textless{}{-}} \FunctionTok{svyby}\NormalTok{(}\AttributeTok{formula =} \SpecialCharTok{\textasciitilde{}}\NormalTok{Income,}
                              \AttributeTok{by =} \SpecialCharTok{\textasciitilde{}}\NormalTok{Region,}
                              \AttributeTok{design =}\NormalTok{ survey\_design,}
                              \AttributeTok{FUN =}\NormalTok{ svymean,}
                              \AttributeTok{na.rm =} \ConstantTok{TRUE}\NormalTok{,}
                              \AttributeTok{covmat =} \ConstantTok{TRUE}\NormalTok{,}
                              \AttributeTok{vartype =} \FunctionTok{c}\NormalTok{(}\StringTok{"se"}\NormalTok{, }\StringTok{"ci"}\NormalTok{))}
\end{Highlighting}
\end{Shaded}

\begin{longtable}[]{@{}lcccc@{}}
\caption{\label{tab:tab-prom-region}Ingreso medio estimado por región}\tabularnewline
\toprule\noalign{}
Region & mean & mean\_se & mean\_low & mean\_upp \\
\midrule\noalign{}
\endfirsthead
\toprule\noalign{}
Region & mean & mean\_se & mean\_low & mean\_upp \\
\midrule\noalign{}
\endhead
\bottomrule\noalign{}
\endlastfoot
Norte & 552 & 55.4 & 443 & 662 \\
Sur & 626 & 62.4 & 502 & 749 \\
Centro & 651 & 61.5 & 529 & 772 \\
Occidente & 517 & 46.2 & 425 & 609 \\
Oriente & 542 & 71.7 & 400 & 684 \\
\end{longtable}

La función \texttt{svycontrast()} devuelve el contraste estimado y su error estándar. Los argumentos de esta función son los promedios de los ingresos estimados (\texttt{stat}) y las constantes de contraste (\texttt{contrasts}), como se muestra en la Tabla \ref{tab:tab-contraste-ns}.

\begin{Shaded}
\begin{Highlighting}[]
\NormalTok{contrast\_region\_ns }\OtherTok{\textless{}{-}} \FunctionTok{svycontrast}\NormalTok{(}
  \AttributeTok{stat =}\NormalTok{ region\_mean\_contrast,}
  \AttributeTok{contrasts =} \FunctionTok{list}\NormalTok{(}\AttributeTok{north\_south\_difference =} \FunctionTok{c}\NormalTok{(}\DecValTok{1}\NormalTok{, }\SpecialCharTok{{-}}\DecValTok{1}\NormalTok{, }\DecValTok{0}\NormalTok{, }\DecValTok{0}\NormalTok{, }\DecValTok{0}\NormalTok{))}
\NormalTok{)}
\end{Highlighting}
\end{Shaded}

\begin{longtable}[]{@{}lccc@{}}
\caption{\label{tab:tab-contraste-ns}Contraste de ingresos medios: región Norte menos región Sur}\tabularnewline
\toprule\noalign{}
contrast & estimate & variance & standard\_error \\
\midrule\noalign{}
\endfirsthead
\toprule\noalign{}
contrast & estimate & variance & standard\_error \\
\midrule\noalign{}
\endhead
\bottomrule\noalign{}
\endlastfoot
north\_south\_difference & -73.4 & 6959 & 83.4 \\
\end{longtable}

Nótese que estos mismos resultados podrían obtenerse realizando explícitamente el proceso de construcción del estimador del contraste y de su varianza estimada. Si se consideran únicamente los ingresos medios estimados de las regiones Norte y Sur, su diferencia es:

\begin{Shaded}
\begin{Highlighting}[]
\CommentTok{\# Paso 1: diferencia de las estimaciones (Norte {-} Sur)}
\FloatTok{552.4} \SpecialCharTok{{-}} \FloatTok{625.8}
\end{Highlighting}
\end{Shaded}

\begin{verbatim}
## [1] -73.4
\end{verbatim}

El paso siguiente es calcular la matriz de varianzas y covarianzas y de allí extraer las varianzas y covarianzas de las regiones Norte y Sur, presentada en la Tabla \ref{tab:tab-vcov-region}:

\begin{Shaded}
\begin{Highlighting}[]
\CommentTok{\# Paso 2: matriz de varianzas y covarianzas}
\NormalTok{region\_vcov }\OtherTok{\textless{}{-}} \FunctionTok{vcov}\NormalTok{(region\_mean\_contrast)}
\end{Highlighting}
\end{Shaded}

\begin{longtable}[]{@{}lccccc@{}}
\caption{\label{tab:tab-vcov-region}Matriz de varianzas y covarianzas de los ingresos medios por región}\tabularnewline
\toprule\noalign{}
Region & Norte & Sur & Centro & Occidente & Oriente \\
\midrule\noalign{}
\endfirsthead
\toprule\noalign{}
Region & Norte & Sur & Centro & Occidente & Oriente \\
\midrule\noalign{}
\endhead
\bottomrule\noalign{}
\endlastfoot
Norte & 3065 & 0 & 0 & 0 & 0 \\
Sur & 0 & 3894 & 0 & 0 & 0 \\
Centro & 0 & 0 & 3778 & 0 & 0 \\
Occidente & 0 & 0 & 0 & 2136 & 0 \\
Oriente & 0 & 0 & 0 & 0 & 5136 \\
\end{longtable}

Para calcular el error estándar de la diferencia (contraste) se usará las propiedades de la varianza, de la siguiente forma:

\[
\widehat{se}\left(\hat{\bar{y}}_{Norte}-\hat{\bar{y}}_{Sur}\right)=\sqrt{\widehat {Var}\left(\hat{\bar{y}}_{Norte}\right)+
\widehat{Var}\left(\hat{\bar{y}}_{Sur}\right)-
2\, \widehat{Cov}\left(\hat{\bar{y}}_{Norte},\hat{\bar{y}}_{Sur}\right)}
\]

Por tanto:

\begin{Shaded}
\begin{Highlighting}[]
\FunctionTok{sqrt}\NormalTok{(}\DecValTok{2650} \SpecialCharTok{+} \DecValTok{3460} \SpecialCharTok{{-}} \DecValTok{2} \SpecialCharTok{*} \DecValTok{0}\NormalTok{)}
\end{Highlighting}
\end{Shaded}

\begin{verbatim}
## [1] 78.2
\end{verbatim}

El cual corresponde al mismo resultado obtenido anteriormente con la función \texttt{svycontrast()}. Además de comparar únicamente dos poblaciones, los contrastes permiten evaluar simultáneamente varias comparaciones de interés entre diferentes grupos.

Por otro lado, supóngase que se desea comparar los ingresos medios entre distintas regiones geográficas. En particular, podrían considerarse contrastes como \(\bar{y}_{Norte} - \bar{y}_{Centro}\), \(\bar{y}_{Sur} - \bar{y}_{Centro}\) y \(\bar{y}_{Occidente} - \bar{y}_{Oriente}\). Cada una de estas expresiones representa un contraste lineal específico entre medias regionales. De manera conjunta, dichos contrastes pueden organizarse mediante la siguiente matriz de contrastes:

\[
\left[\begin{array}{ccccc}
1 & 0 & -1 & 0 & 0\\
0 & 1 & -1 & 0 & 0\\
0 & 0 & 0 & 1 & -1
\end{array}\right]
\]

En esta matriz, cada fila define un contraste distinto y cada columna corresponde a una de las regiones consideradas en el análisis. Los coeficientes positivos y negativos indican las medias que participan en cada comparación, mientras que los ceros representan las regiones que no intervienen en el contraste correspondiente. A continuación, se muestra la implementación de los contrastes en \texttt{R}, cuyos resultados se presentan en la Tabla \ref{tab:tab-contraste-regiones}. A partir de estos resultados, puede concluirse que las regiones Sur y Centro presentan los ingresos medios de los hogares más similares, dado que la diferencia estimada entre ambas regiones es la menor entre los contrastes evaluados.

\begin{Shaded}
\begin{Highlighting}[]
\NormalTok{contrast\_regions }\OtherTok{\textless{}{-}} \FunctionTok{svycontrast}\NormalTok{(}
  \AttributeTok{stat =}\NormalTok{ region\_mean\_contrast,}
  \AttributeTok{contrasts =} \FunctionTok{list}\NormalTok{(}
    \AttributeTok{north\_south =} \FunctionTok{c}\NormalTok{(}\DecValTok{1}\NormalTok{, }\DecValTok{0}\NormalTok{, }\SpecialCharTok{{-}}\DecValTok{1}\NormalTok{, }\DecValTok{0}\NormalTok{, }\DecValTok{0}\NormalTok{),}
    \AttributeTok{south\_center =} \FunctionTok{c}\NormalTok{(}\DecValTok{0}\NormalTok{, }\DecValTok{1}\NormalTok{, }\SpecialCharTok{{-}}\DecValTok{1}\NormalTok{, }\DecValTok{0}\NormalTok{, }\DecValTok{0}\NormalTok{),}
    \AttributeTok{west\_east =} \FunctionTok{c}\NormalTok{(}\DecValTok{0}\NormalTok{, }\DecValTok{0}\NormalTok{, }\DecValTok{0}\NormalTok{, }\DecValTok{1}\NormalTok{, }\SpecialCharTok{{-}}\DecValTok{1}\NormalTok{)}
\NormalTok{  )}
\NormalTok{)}
\end{Highlighting}
\end{Shaded}

\begin{longtable}[]{@{}lccc@{}}
\caption{\label{tab:tab-contraste-regiones}Contrastes de ingresos medios entre regiones}\tabularnewline
\toprule\noalign{}
contrast & estimate & variance & standard\_error \\
\midrule\noalign{}
\endfirsthead
\toprule\noalign{}
contrast & estimate & variance & standard\_error \\
\midrule\noalign{}
\endhead
\bottomrule\noalign{}
\endlastfoot
north\_south & -98.4 & 6843 & 82.7 \\
south\_center & -25.0 & 7673 & 87.6 \\
west\_east & -24.7 & 7272 & 85.3 \\
\end{longtable}

También es posible construir contrastes entre variables asociadas a diferentes grupos poblacionales, como ocurre al comparar el ingreso medio según el sexo. Siguiendo el mismo procedimiento del ejemplo anterior, el primer paso consiste en estimar el ingreso medio para cada grupo, en este caso hombres y mujeres, como se presenta en la Tabla \ref{tab:tab-prom-sexo}.

\begin{Shaded}
\begin{Highlighting}[]
\NormalTok{sex\_mean }\OtherTok{\textless{}{-}}\NormalTok{ survey\_design }\SpecialCharTok{\%\textgreater{}\%}
  \FunctionTok{group\_by}\NormalTok{(Sex) }\SpecialCharTok{\%\textgreater{}\%}
  \FunctionTok{summarise}\NormalTok{(}\AttributeTok{mean =} \FunctionTok{survey\_mean}\NormalTok{(Income,}
                               \AttributeTok{na.rm =} \ConstantTok{TRUE}\NormalTok{,}
                               \AttributeTok{vartype =} \FunctionTok{c}\NormalTok{(}\StringTok{"se"}\NormalTok{, }\StringTok{"ci"}\NormalTok{)))}

\NormalTok{sex\_mean\_contrast }\OtherTok{\textless{}{-}} \FunctionTok{svyby}\NormalTok{(}\AttributeTok{formula =} \SpecialCharTok{\textasciitilde{}}\NormalTok{Income,}
                           \AttributeTok{by =} \SpecialCharTok{\textasciitilde{}}\NormalTok{Sex,}
                           \AttributeTok{design =}\NormalTok{ survey\_design,}
                           \AttributeTok{FUN =}\NormalTok{ svymean,}
                           \AttributeTok{na.rm =} \ConstantTok{TRUE}\NormalTok{,}
                           \AttributeTok{covmat =} \ConstantTok{TRUE}\NormalTok{,}
                           \AttributeTok{vartype =} \FunctionTok{c}\NormalTok{(}\StringTok{"se"}\NormalTok{, }\StringTok{"ci"}\NormalTok{))}
\end{Highlighting}
\end{Shaded}

\begin{longtable}[]{@{}lcccc@{}}
\caption{\label{tab:tab-prom-sexo}Ingreso medio estimado por sexo}\tabularnewline
\toprule\noalign{}
Sex & mean & mean\_se & mean\_low & mean\_upp \\
\midrule\noalign{}
\endfirsthead
\toprule\noalign{}
Sex & mean & mean\_se & mean\_low & mean\_upp \\
\midrule\noalign{}
\endhead
\bottomrule\noalign{}
\endlastfoot
Female & 558 & 25.8 & 506 & 609 \\
Male & 586 & 34.6 & 517 & 654 \\
\end{longtable}

En este caso, el contraste de interés está dado por \(\bar{y}_{F} - \bar{y}_{M}\), el cual permite cuantificar la diferencia entre el ingreso medio de las mujeres y el correspondiente a los hombres. A continuación, utilizando la función \texttt{svycontrast()} del paquete \texttt{survey}, se obtiene la estimación del contraste, cuyos resultados se presentan en la Tabla \ref{tab:tab-contraste-sexo}. A partir de estos resultados, se concluye que, en promedio, los hombres perciben 28.3 unidades monetarias más que las mujeres, con un error estándar estimado de 19.8 unidades monetarias.

\begin{Shaded}
\begin{Highlighting}[]
\NormalTok{contrast\_sex }\OtherTok{\textless{}{-}} \FunctionTok{svycontrast}\NormalTok{(}
  \AttributeTok{stat =}\NormalTok{ sex\_mean\_contrast,}
  \AttributeTok{contrasts =} \FunctionTok{list}\NormalTok{(}\AttributeTok{sex\_difference =} \FunctionTok{c}\NormalTok{(}\DecValTok{1}\NormalTok{, }\SpecialCharTok{{-}}\DecValTok{1}\NormalTok{))}
\NormalTok{)}
\end{Highlighting}
\end{Shaded}

\begin{longtable}[]{@{}lccc@{}}
\caption{\label{tab:tab-contraste-sexo}Contraste de ingresos medios entre mujeres y hombres}\tabularnewline
\toprule\noalign{}
contrast & estimate & variance & standard\_error \\
\midrule\noalign{}
\endfirsthead
\toprule\noalign{}
contrast & estimate & variance & standard\_error \\
\midrule\noalign{}
\endhead
\bottomrule\noalign{}
\endlastfoot
sex\_difference & -28.3 & 431 & 20.8 \\
\end{longtable}

Dado que los contrastes corresponden a funciones lineales de parámetros, también es posible aplicarlos a estimadores de razón. Un ejemplo de ello consiste en comparar la relación entre gastos e ingresos según el sexo. En este caso, primero se estima la razón gasto-ingreso para cada grupo y posteriormente se construye el contraste entre dichas razones. A continuación, se presentan los códigos computacionales utilizados para realizar este análisis, cuyos resultados se muestran en la Tabla \ref{tab:tab-razon-ingreso-gasto-sexo}:

\begin{Shaded}
\begin{Highlighting}[]
\NormalTok{sex\_ratio }\OtherTok{\textless{}{-}}\NormalTok{ survey\_design }\SpecialCharTok{\%\textgreater{}\%}
  \FunctionTok{group\_by}\NormalTok{(Sex) }\SpecialCharTok{\%\textgreater{}\%}
  \FunctionTok{summarise}\NormalTok{(}\AttributeTok{ratio =} \FunctionTok{survey\_ratio}\NormalTok{(Income,}
\NormalTok{                                 Expenditure,}
                                 \AttributeTok{na.rm =} \ConstantTok{TRUE}\NormalTok{,}
                                 \AttributeTok{vartype =} \FunctionTok{c}\NormalTok{(}\StringTok{"se"}\NormalTok{, }\StringTok{"ci"}\NormalTok{)))}

\NormalTok{sex\_ratio\_contrast }\OtherTok{\textless{}{-}} \FunctionTok{svyby}\NormalTok{(}\AttributeTok{formula =} \SpecialCharTok{\textasciitilde{}}\NormalTok{Income,}
                            \AttributeTok{by =} \SpecialCharTok{\textasciitilde{}}\NormalTok{Sex,}
                            \AttributeTok{denominator =} \SpecialCharTok{\textasciitilde{}}\NormalTok{Expenditure,}
                            \AttributeTok{design =}\NormalTok{ survey\_design,}
                            \AttributeTok{FUN =}\NormalTok{ svyratio,}
                            \AttributeTok{na.rm =} \ConstantTok{TRUE}\NormalTok{,}
                            \AttributeTok{covmat =} \ConstantTok{TRUE}\NormalTok{,}
                            \AttributeTok{vartype =} \FunctionTok{c}\NormalTok{(}\StringTok{"se"}\NormalTok{, }\StringTok{"ci"}\NormalTok{))}
\end{Highlighting}
\end{Shaded}

\begin{longtable}[]{@{}lcccc@{}}
\caption{\label{tab:tab-razon-ingreso-gasto-sexo}Razón ingreso/gasto desagregada por sexo}\tabularnewline
\toprule\noalign{}
Sex & ratio & ratio\_se & ratio\_low & ratio\_upp \\
\midrule\noalign{}
\endfirsthead
\toprule\noalign{}
Sex & ratio & ratio\_se & ratio\_low & ratio\_upp \\
\midrule\noalign{}
\endhead
\bottomrule\noalign{}
\endlastfoot
Female & 1.52 & 0.046 & 1.43 & 1.61 \\
Male & 1.56 & 0.070 & 1.42 & 1.70 \\
\end{longtable}

La función \texttt{svycontrast()} no se limita únicamente a la comparación de medias o totales, sino que también permite construir contrastes entre otros parámetros de interés, como razones y proporciones estimadas a partir de diseños de muestreo complejos. A continuación, se presentan los códigos computacionales utilizados para realizar este análisis, cuyos resultados se muestran en la Tabla \ref{tab:tab-contraste-razon-sexo}. A partir de dichos resultados, puede concluirse que la diferencia entre las razones estimadas es de 0.045 a favor de los hombres.

\begin{Shaded}
\begin{Highlighting}[]
\NormalTok{contrast\_sex\_ratio }\OtherTok{\textless{}{-}} \FunctionTok{svycontrast}\NormalTok{(}
  \AttributeTok{stat =}\NormalTok{ sex\_ratio\_contrast,}
  \AttributeTok{contrasts =} \FunctionTok{list}\NormalTok{(}\AttributeTok{sex\_difference =} \FunctionTok{c}\NormalTok{(}\DecValTok{1}\NormalTok{, }\SpecialCharTok{{-}}\DecValTok{1}\NormalTok{))}
\NormalTok{)}
\end{Highlighting}
\end{Shaded}

\begin{longtable}[]{@{}lccc@{}}
\caption{\label{tab:tab-contraste-razon-sexo}Contraste de la razón ingreso/gasto entre sexos}\tabularnewline
\toprule\noalign{}
contrast & estimate & standard\_error & variance \\
\midrule\noalign{}
\endfirsthead
\toprule\noalign{}
contrast & estimate & standard\_error & variance \\
\midrule\noalign{}
\endhead
\bottomrule\noalign{}
\endlastfoot
sex\_difference & -0.046 & 0.042 & 0.002 \\
\end{longtable}

\section{Visualización de variables continuas}\label{viz-continuas}

La representación gráfica de variables continuas es un componente esencial en el análisis de encuestas de hogares. Los gráficos bien diseñados permiten identificar distribuciones, tendencias y relaciones entre variables, facilitando la interpretación de los hallazgos y su comunicación a distintos públicos \citep{wilkinson2005grammar}. Un aspecto fundamental en este contexto es que los gráficos deben incorporar los pesos muestrales para reflejar la distribución de la población de interés.

En esta sección se utilizan los paquetes \texttt{ggplot2} \citep{ggplot22026} y \texttt{patchwork} \citep{patchwork} para la construcción y composición de gráficos, combinados con \texttt{survey} \citep{Lumley2024} y \texttt{srvyr} \citep{Freedman2024}.

\begin{Shaded}
\begin{Highlighting}[]
\FunctionTok{library}\NormalTok{(ggplot2)}
\FunctionTok{library}\NormalTok{(patchwork)}
\FunctionTok{data}\NormalTok{(BigCity, }\AttributeTok{package =} \StringTok{"TeachingSampling"}\NormalTok{)}
\end{Highlighting}
\end{Shaded}

\subsection{Histogramas}\label{histogramas}

Un histograma es una representación gráfica de la distribución de una variable numérica continua, construida a partir de rectángulos cuya altura es proporcional a la frecuencia de observaciones dentro de determinados intervalos de clase. En el contexto de encuestas de hogares, los histogramas ponderados permiten aproximar la distribución de la variable en la población, incorporando los factores de expansión asociados al diseño muestral.

A continuación, se presenta el histograma ponderado de la variable \texttt{Income}, construido incorporando los pesos muestrales (\texttt{wk}) con el fin de representar la distribución estimada del ingreso en la población. Los resultados obtenidos se muestran en la Figura \ref{fig:fig-hist-ingreso-ponde}:

\begin{Shaded}
\begin{Highlighting}[]
\NormalTok{weighted\_income\_hist }\OtherTok{\textless{}{-}} \FunctionTok{ggplot}\NormalTok{(}
  \AttributeTok{data =}\NormalTok{ survey\_data,}
  \FunctionTok{aes}\NormalTok{(}\AttributeTok{x =}\NormalTok{ Income, }\AttributeTok{weight =}\NormalTok{ wk)) }\SpecialCharTok{+}
  \FunctionTok{geom\_histogram}\NormalTok{(}\FunctionTok{aes}\NormalTok{(}\AttributeTok{y =}\NormalTok{ ..density..)) }\SpecialCharTok{+}
  \FunctionTok{ylab}\NormalTok{(}\StringTok{""}\NormalTok{) }\SpecialCharTok{+}
  \FunctionTok{ggtitle}\NormalTok{(}\StringTok{"Histograma ponderado"}\NormalTok{)}
\NormalTok{weighted\_income\_hist}
\end{Highlighting}
\end{Shaded}

\begin{figure}

{\centering \includegraphics[width=0.65\linewidth]{docs/figures/fig-hist-ingreso-ponde-1} 

}

\caption{Histograma ponderado del ingreso de los hogares}\label{fig:fig-hist-ingreso-ponde}
\end{figure}

En este código, se inicia la construcción del gráfico utilizando la base de datos \texttt{survey\_data}. Dentro de \texttt{aes()}, el argumento \texttt{x\ =\ Income} indica que la variable representada en el eje horizontal corresponde al ingreso, mientras que \texttt{weight\ =\ wk} incorpora los pesos muestrales para construir un histograma ponderado que refleje la distribución estimada en la población. La función \texttt{geom\_histogram()} genera el histograma y, mediante \texttt{aes(y\ =\ ..density..)}, la altura de las barras se expresa en términos de densidad en lugar de frecuencias absolutas. Posteriormente, \texttt{ylab("")} elimina la etiqueta del eje vertical, y \texttt{ggtitle("Histograma\ ponderado")} agrega un título descriptivo al gráfico.

Los histogramas también permiten comparar visualmente la distribución de una variable entre distintos subgrupos de la población mediante el argumento \texttt{fill}, el cual asigna un color diferente a cada categoría. Por ejemplo, es posible comparar la distribución del ingreso entre zonas urbanas y rurales, como se muestra en la Figura \ref{fig:fig-hist-zona}:

\begin{Shaded}
\begin{Highlighting}[]
\NormalTok{zone\_colors }\OtherTok{\textless{}{-}} \FunctionTok{c}\NormalTok{(}\AttributeTok{Urban =} \StringTok{"\#48C9B0"}\NormalTok{, }\AttributeTok{Rural =} \StringTok{"\#117864"}\NormalTok{)}

\NormalTok{weighted\_zone\_hist }\OtherTok{\textless{}{-}} \FunctionTok{ggplot}\NormalTok{(survey\_data, }\FunctionTok{aes}\NormalTok{(}\AttributeTok{x =}\NormalTok{ Income, }\AttributeTok{weight =}\NormalTok{ wk)) }\SpecialCharTok{+}
  \FunctionTok{geom\_histogram}\NormalTok{(}
    \FunctionTok{aes}\NormalTok{(}\AttributeTok{y =}\NormalTok{ ..density.., }\AttributeTok{fill =}\NormalTok{ Zone),}
    \AttributeTok{alpha =} \FloatTok{0.5}\NormalTok{, }\AttributeTok{position =} \StringTok{"identity"}\NormalTok{) }\SpecialCharTok{+}
  \FunctionTok{scale\_fill\_manual}\NormalTok{(}\AttributeTok{values =}\NormalTok{ zone\_colors) }\SpecialCharTok{+}
  \FunctionTok{ylab}\NormalTok{(}\StringTok{""}\NormalTok{) }\SpecialCharTok{+}
  \FunctionTok{ggtitle}\NormalTok{(}\StringTok{"Weighted"}\NormalTok{)}

\NormalTok{weighted\_zone\_hist}
\end{Highlighting}
\end{Shaded}

\begin{figure}

{\centering \includegraphics[width=0.65\linewidth]{docs/figures/fig-hist-zona-1} 

}

\caption{Histograma ponderado del ingreso por zona geográfica}\label{fig:fig-hist-zona}
\end{figure}

En este código, se construye un histograma ponderado utilizando la base de datos \texttt{survey\_data}, donde \texttt{x\ =\ Income} define la variable representada en el eje horizontal y \texttt{weight\ =\ wk} incorpora los pesos muestrales para reflejar la distribución poblacional del ingreso. La función \texttt{geom\_histogram()} genera el histograma y, mediante \texttt{aes(y\ =\ ..density..,\ fill\ =\ Zone)}, representa las barras en términos de densidad y asigna colores diferentes según la variable \texttt{Zone}, permitiendo comparar visualmente la distribución del ingreso entre zonas. El argumento \texttt{alpha\ =\ 0.5} agrega transparencia a las barras para facilitar la superposición de los histogramas, mientras que \texttt{position\ =\ "identity"} permite que las distribuciones se dibujen una encima de otra en lugar de apilarse. Posteriormente, \texttt{scale\_fill\_manual(values\ =\ zone\_colors)} define manualmente los colores asociados a cada zona. Finalmente, \texttt{ylab("")} elimina la etiqueta del eje vertical, \texttt{ggtitle("Weighted")} agrega un título al gráfico.

Adicionalmente, es posible superponer curvas de densidad suavizadas mediante la función \texttt{geom\_density()}, lo que permite visualizar de manera continua la forma de la distribución de la variable analizada. Este tipo de representación facilita la identificación de patrones como asimetrías, concentración de observaciones o la presencia de múltiples modas. Un ejemplo de este tipo de gráfico se presenta en la Figura \ref{fig:fig-densidades}:

\begin{Shaded}
\begin{Highlighting}[]
\NormalTok{weighted\_income\_hist }\SpecialCharTok{+} \FunctionTok{geom\_density}\NormalTok{(}\AttributeTok{fill =} \StringTok{"blue"}\NormalTok{, }\AttributeTok{alpha =} \FloatTok{0.3}\NormalTok{) }\SpecialCharTok{|}
\NormalTok{  weighted\_zone\_hist }\SpecialCharTok{+} \FunctionTok{geom\_density}\NormalTok{(}\FunctionTok{aes}\NormalTok{(}\AttributeTok{fill =}\NormalTok{ Zone), }\AttributeTok{alpha =} \FloatTok{0.3}\NormalTok{) }\SpecialCharTok{+}
  \FunctionTok{theme}\NormalTok{(}\AttributeTok{legend.position =} \StringTok{"top"}\NormalTok{)}
\end{Highlighting}
\end{Shaded}

\begin{figure}

{\centering \includegraphics[width=0.9\linewidth]{docs/figures/fig-densidades-1} 

}

\caption{Histogramas con curvas de densidad superpuestas para ingreso y gasto}\label{fig:fig-densidades}
\end{figure}

\subsection{\texorpdfstring{Diagramas de caja (\emph{boxplots})}{Diagramas de caja (boxplots)}}\label{diagramas-de-caja-boxplots}

Los diagramas de caja (\emph{boxplots}) permiten identificar valores atípicos y evaluar la dispersión de la distribución. Introducido por John Tukey en 1977 son una representación gráfica que resume la distribución de una variable mediante estadísticas basadas en cuartiles, permitiendo identificar la mediana, la dispersión de los datos y la presencia de valores atípicos. En el contexto de encuestas de hogares, este tipo de gráfico resulta útil para comparar distribuciones entre distintos grupos poblacionales de manera compacta y visualmente intuitiva.

Para construir \emph{boxplots} ponderados en \texttt{ggplot2} se utiliza la función \texttt{geom\_boxplot()}, como se muestra en la Figura \ref{fig:fig-boxplot-zona-sexo}. En este caso, la función \texttt{ggplot()} inicia la construcción del gráfico utilizando la base de datos \texttt{survey\_data}, donde \texttt{x\ =\ Income} define la variable numérica analizada y \texttt{weight\ =\ wk} incorpora los pesos muestrales para representar la distribución estimada en la población. Posteriormente, \texttt{geom\_boxplot()} genera los diagramas de caja y, mediante \texttt{aes(fill\ =\ Zone)} o \texttt{aes(fill\ =\ Sex)}, asigna colores diferentes a cada categoría de las variables \texttt{Zone} y \texttt{Sex}, respectivamente, facilitando la comparación visual entre grupos. La función \texttt{coord\_flip()} invierte los ejes para presentar los \emph{boxplots} de forma horizontal, mejorando así su legibilidad. Por su parte, \texttt{scale\_fill\_manual()} permite definir manualmente los colores asociados a cada categoría mediante los vectores \texttt{zone\_colors} y \texttt{sex\_colors}.

\begin{Shaded}
\begin{Highlighting}[]
\NormalTok{weighted\_zone\_boxplot }\OtherTok{\textless{}{-}} \FunctionTok{ggplot}\NormalTok{(survey\_data, }\FunctionTok{aes}\NormalTok{(}\AttributeTok{x =}\NormalTok{ Income, }\AttributeTok{weight =}\NormalTok{ wk)) }\SpecialCharTok{+}
  \FunctionTok{geom\_boxplot}\NormalTok{(}\FunctionTok{aes}\NormalTok{(}\AttributeTok{fill =}\NormalTok{ Zone)) }\SpecialCharTok{+}
  \FunctionTok{ggtitle}\NormalTok{(}\StringTok{"Weighted"}\NormalTok{) }\SpecialCharTok{+}
  \FunctionTok{coord\_flip}\NormalTok{() }\SpecialCharTok{+}
  \FunctionTok{scale\_fill\_manual}\NormalTok{(}\AttributeTok{values =}\NormalTok{ zone\_colors)}

\NormalTok{sex\_colors }\OtherTok{\textless{}{-}} \FunctionTok{c}\NormalTok{(}\AttributeTok{Male =} \StringTok{"\#5DADE2"}\NormalTok{, }\AttributeTok{Female =} \StringTok{"\#2874A6"}\NormalTok{)}

\NormalTok{weighted\_sex\_boxplot }\OtherTok{\textless{}{-}} \FunctionTok{ggplot}\NormalTok{(survey\_data, }\FunctionTok{aes}\NormalTok{(}\AttributeTok{x =}\NormalTok{ Income, }\AttributeTok{weight =}\NormalTok{ wk)) }\SpecialCharTok{+}
  \FunctionTok{geom\_boxplot}\NormalTok{(}\FunctionTok{aes}\NormalTok{(}\AttributeTok{fill =}\NormalTok{ Sex)) }\SpecialCharTok{+}
  \FunctionTok{ggtitle}\NormalTok{(}\StringTok{"Weighted"}\NormalTok{) }\SpecialCharTok{+}
  \FunctionTok{coord\_flip}\NormalTok{() }\SpecialCharTok{+}
  \FunctionTok{scale\_fill\_manual}\NormalTok{(}\AttributeTok{values =}\NormalTok{ sex\_colors)}

\NormalTok{weighted\_zone\_boxplot }\SpecialCharTok{|}\NormalTok{ weighted\_sex\_boxplot}
\end{Highlighting}
\end{Shaded}

\begin{figure}

{\centering \includegraphics[width=0.9\linewidth]{docs/figures/fig-boxplot-zona-sexo-1} 

}

\caption{Boxplots ponderados del ingreso por zona y sexo}\label{fig:fig-boxplot-zona-sexo}
\end{figure}

Además de las herramientas disponibles en \texttt{ggplot2}, el paquete \texttt{survey} ofrece funciones especializadas como \texttt{svyhist()} y \texttt{svyboxplot()}, las cuales incorporan directamente las características del diseño muestral complejo en la construcción de los gráficos. Estas funciones permiten obtener representaciones gráficas coherentes con el diseño de la encuesta, facilitando el análisis exploratorio de las variables numéricas. Un ejemplo de su aplicación se presenta en la Figura \ref{fig:fig-svyhist-svybox}.

\begin{Shaded}
\begin{Highlighting}[]
\FunctionTok{par}\NormalTok{(}\AttributeTok{mfrow =} \FunctionTok{c}\NormalTok{(}\DecValTok{1}\NormalTok{, }\DecValTok{2}\NormalTok{))}
\FunctionTok{svyhist}\NormalTok{(}\SpecialCharTok{\textasciitilde{}}\NormalTok{Income, survey\_design,}
  \AttributeTok{main =} \StringTok{"Household income"}\NormalTok{,}
  \AttributeTok{col  =} \StringTok{"grey80"}\NormalTok{, }\AttributeTok{xlab =} \StringTok{"Income"}\NormalTok{,}
  \AttributeTok{probability =} \ConstantTok{FALSE}\NormalTok{)}

\FunctionTok{svyboxplot}\NormalTok{(Income }\SpecialCharTok{\textasciitilde{}}\NormalTok{ Zone, survey\_design,}
  \AttributeTok{col  =} \StringTok{"grey80"}\NormalTok{,}
  \AttributeTok{ylab =} \StringTok{"Income"}\NormalTok{, }\AttributeTok{xlab =} \StringTok{"Zone"}\NormalTok{)}
\end{Highlighting}
\end{Shaded}

\begin{figure}

{\centering \includegraphics[width=0.8\linewidth]{docs/figures/fig-svyhist-svybox-1} 

}

\caption{Histograma y boxplot del ingreso ajustados al diseño muestral complejo}\label{fig:fig-svyhist-svybox}
\end{figure}

\subsection{Diagramas de dispersión}\label{diagramas-de-dispersion}

Los diagramas de dispersión permiten explorar visualmente la relación entre dos variables continuas, facilitando la identificación de patrones de asociación, tendencias y posibles valores atípicos. En el contexto de encuestas con diseños muestrales complejos, resulta importante incorporar la información de los pesos de muestreo para que la representación gráfica refleje adecuadamente la contribución relativa de cada observación en la población.

Según \citet{Lumley2010}, cuando se trabaja con conjuntos de datos de gran tamaño, representar todos los puntos simultáneamente puede generar gráficos difíciles de interpretar debido a la sobreposición de observaciones. Sin embargo, en muestras de tamaño moderado, una estrategia sencilla y efectiva consiste en representar los pesos mediante el tamaño de los símbolos en el gráfico de dispersión. De esta manera, las observaciones con mayor peso poblacional aparecen resaltadas visualmente. Un ejemplo de este tipo de representación se presenta en la Figura \ref{fig:fig-scatter-basico}:

\begin{Shaded}
\begin{Highlighting}[]
\NormalTok{weighted\_basic\_scatter }\OtherTok{\textless{}{-}} \FunctionTok{ggplot}\NormalTok{(}
  \AttributeTok{data =}\NormalTok{ survey\_data,}
  \FunctionTok{aes}\NormalTok{(}\AttributeTok{y =}\NormalTok{ Income, }\AttributeTok{x =}\NormalTok{ Expenditure)) }\SpecialCharTok{+}
  \FunctionTok{geom\_point}\NormalTok{(}\FunctionTok{aes}\NormalTok{(}\AttributeTok{size =}\NormalTok{ wk), }\AttributeTok{alpha =} \FloatTok{0.3}\NormalTok{)}

\NormalTok{weighted\_basic\_scatter}
\end{Highlighting}
\end{Shaded}

\begin{figure}

{\centering \includegraphics[width=0.65\linewidth]{docs/figures/fig-scatter-basico-1} 

}

\caption{Diagrama de dispersión entre ingreso y gasto de los hogares}\label{fig:fig-scatter-basico}
\end{figure}

En este código, se construye un diagrama de dispersión entre las variables \texttt{x\ =\ Expenditure} y \texttt{y\ =\ Income}. La función \texttt{geom\_point()} agrega los puntos del gráfico y, mediante \texttt{aes(size\ =\ wk)}, utiliza los pesos muestrales (\texttt{wk}) para controlar el tamaño de cada punto, de modo que las observaciones con mayor representación poblacional aparezcan visualmente más destacadas. El argumento \texttt{alpha\ =\ 0.3} incorpora transparencia a los puntos, reduciendo el problema de sobreposición y facilitando la visualización de áreas con alta concentración de observaciones.

Adicionalmente, es posible incorporar variables de agrupamiento en los diagramas de dispersión mediante los argumentos \texttt{shape} y \texttt{color}, los cuales permiten diferenciar visualmente las observaciones según categorías específicas de la población. Esto facilita la identificación de patrones diferenciados entre grupos y mejora la interpretación de la relación entre las variables analizadas. Un ejemplo de este tipo de representación se presenta en la Figura \ref{fig:fig-scatter-zona}.

\begin{Shaded}
\begin{Highlighting}[]
\NormalTok{weighted\_zone\_scatter }\OtherTok{\textless{}{-}} \FunctionTok{ggplot}\NormalTok{(}
  \AttributeTok{data =}\NormalTok{ survey\_data,}
  \FunctionTok{aes}\NormalTok{(}\AttributeTok{y =}\NormalTok{ Income, }\AttributeTok{x =}\NormalTok{ Expenditure, }\AttributeTok{shape =}\NormalTok{ Zone)) }\SpecialCharTok{+}
  \FunctionTok{geom\_point}\NormalTok{(}\FunctionTok{aes}\NormalTok{(}\AttributeTok{size =}\NormalTok{ wk, }\AttributeTok{color =}\NormalTok{ Zone), }\AttributeTok{alpha =} \FloatTok{0.3}\NormalTok{) }\SpecialCharTok{+}
  \FunctionTok{labs}\NormalTok{(}\AttributeTok{size =} \StringTok{"Peso"}\NormalTok{) }\SpecialCharTok{+}
  \FunctionTok{scale\_color\_manual}\NormalTok{(}\AttributeTok{values =}\NormalTok{ zone\_colors)}

\NormalTok{weighted\_zone\_scatter}
\end{Highlighting}
\end{Shaded}

\begin{figure}

{\centering \includegraphics[width=0.65\linewidth]{docs/figures/fig-scatter-zona-1} 

}

\caption{Diagrama de dispersión entre ingreso y gasto por zona geográfica}\label{fig:fig-scatter-zona}
\end{figure}

\chapter{Análisis de variables categóricas}\label{cap4}

En el análisis de encuestas de hogares, uno de los productos más habituales es la estimación de parámetros descriptivos asociados a variables categóricas, los cuales permiten sintetizar las principales características de la población. Estas medidas facilitan una representación clara y comprensible de fenómenos sociales a partir de información recolectada mediante muestras probabilísticas.

Entre los indicadores descriptivos más utilizados para analizar variables categóricas se encuentran las frecuencias y las proporciones. Las frecuencias indican cuántas personas u hogares pertenecen a cada categoría, por ejemplo, el número de personas en condición de pobreza, mientras que las proporciones expresan el peso relativo de cada categoría dentro del conjunto poblacional.

Aunque el análisis descriptivo de variables categóricas suele apoyarse en estos parámetros básicos, también es posible combinarlas con medidas derivadas de variables numéricas, como cuantiles o indicadores de desigualdad, las cuales se abordan en detalle en el Capítulo \ref{cap4}. No obstante, en este capítulo nos concentraremos exclusivamente en la descripción y estimación de parámetros asociados a variables categóricas, empleando diseños muestrales complejos.

El análisis comienza con la definición del diseño de muestreo, siguiendo los lineamientos explicados en capítulos anteriores y utilizando la misma base de datos. Este paso es fundamental, ya que permite incorporar adecuadamente las características del diseño complejo de la encuesta, tales como los factores de expansión, las unidades primarias de muestreo y los estratos de selección. La correcta especificación de estos elementos garantiza que las estimaciones obtenidas sean representativas de la población objetivo y que las medidas de precisión reflejen apropiadamente la variabilidad introducida por el esquema de muestreo.

\begin{Shaded}
\begin{Highlighting}[]
\FunctionTok{library}\NormalTok{(tidyverse)}
\FunctionTok{library}\NormalTok{(survey)}
\FunctionTok{library}\NormalTok{(srvyr)}

\NormalTok{survey\_data }\OtherTok{\textless{}{-}} \FunctionTok{readRDS}\NormalTok{(}\StringTok{"Data/encuesta.rds"}\NormalTok{)}
\FunctionTok{options}\NormalTok{(}\AttributeTok{survey.lonely.psu =} \StringTok{"adjust"}\NormalTok{)}

\NormalTok{survey\_design }\OtherTok{\textless{}{-}}\NormalTok{ survey\_data }\SpecialCharTok{\%\textgreater{}\%}
  \FunctionTok{as\_survey\_design}\NormalTok{(}
    \AttributeTok{strata =}\NormalTok{ Stratum,}
    \AttributeTok{ids =}\NormalTok{ PSU,}
    \AttributeTok{weights =}\NormalTok{ wk,}
    \AttributeTok{nest =} \ConstantTok{TRUE}
\NormalTok{  )}
\end{Highlighting}
\end{Shaded}

Luego, se generan algunas variables categóricas derivadas de la información original de la encuesta, las cuales serán utilizadas en los ejercicios desarrollados a lo largo de este capítulo. Entre ellas, se construye una variable que identifica si la persona se encuentra o no en condición de pobreza. De manera complementaria, también se crean otras variables categóricas de interés analítico que permitirán realizar desagregaciones y comparaciones entre distintos grupos de la población:

\begin{Shaded}
\begin{Highlighting}[]
\NormalTok{survey\_design }\OtherTok{\textless{}{-}}\NormalTok{ survey\_design }\SpecialCharTok{\%\textgreater{}\%}
  \FunctionTok{mutate}\NormalTok{(}
    \AttributeTok{poor =} \FunctionTok{ifelse}\NormalTok{(Poverty }\SpecialCharTok{!=} \StringTok{"NotPoor"}\NormalTok{, }\DecValTok{1}\NormalTok{, }\DecValTok{0}\NormalTok{),}
    \AttributeTok{unemployed =} \FunctionTok{ifelse}\NormalTok{(Employment }\SpecialCharTok{==} \StringTok{"Unemployed"}\NormalTok{, }\DecValTok{1}\NormalTok{, }\DecValTok{0}\NormalTok{),}
    \AttributeTok{age\_18 =} \FunctionTok{case\_when}\NormalTok{(}
\NormalTok{      Age }\SpecialCharTok{\textless{}} \DecValTok{18} \SpecialCharTok{\textasciitilde{}} \StringTok{"\textless{} 18 years"}\NormalTok{,}
      \ConstantTok{TRUE} \SpecialCharTok{\textasciitilde{}} \StringTok{"\textgreater{}= 18 years"}
\NormalTok{    )}
\NormalTok{  )}
\end{Highlighting}
\end{Shaded}

En el código anterior se crean nuevas variables derivadas utilizando la función \texttt{mutate()} del paquete \texttt{dplyr} \citep{dplyr2026}. Las variables \texttt{poor} y \texttt{unemployed} se construyen mediante la función \texttt{ifelse()}, que evalúa una condición lógica y asigna un valor dependiendo de si esta se cumple o no. En este caso, las variables toman el valor \texttt{1} cuando la persona se encuentra en condición de pobreza o desempleo, respectivamente, y \texttt{0} en caso contrario. Por su parte, la variable \texttt{age\_18} se genera mediante la función \texttt{case\_when()}, la cual permite definir categorías a partir de una o varias condiciones de forma más clara y flexible que \texttt{ifelse()}. Esta función resulta especialmente útil cuando se requiere construir variables categóricas con múltiples niveles. En el ejemplo, las personas menores de 18 años se clasifican en la categoría \texttt{"\textless{}\ 18\ years"}, mientras que el resto se agrupa en la categoría \texttt{"\textgreater{}=\ 18\ years"}.

Además de lo anterior, se hace necesario generar subgrupos de la población para realizar estimaciones desagregadas. En este capítulo se emplean cuatro subpoblaciones de ejemplo:

\begin{Shaded}
\begin{Highlighting}[]
\NormalTok{urban\_subset }\OtherTok{\textless{}{-}}\NormalTok{ survey\_design }\SpecialCharTok{\%\textgreater{}\%}
  \FunctionTok{filter}\NormalTok{(Zone }\SpecialCharTok{==} \StringTok{"Urban"}\NormalTok{)}

\NormalTok{rural\_subset }\OtherTok{\textless{}{-}}\NormalTok{ survey\_design }\SpecialCharTok{\%\textgreater{}\%}
  \FunctionTok{filter}\NormalTok{(Zone }\SpecialCharTok{==} \StringTok{"Rural"}\NormalTok{)}

\NormalTok{female\_subset }\OtherTok{\textless{}{-}}\NormalTok{ survey\_design }\SpecialCharTok{\%\textgreater{}\%}
  \FunctionTok{filter}\NormalTok{(Sex }\SpecialCharTok{==} \StringTok{"Female"}\NormalTok{)}

\NormalTok{male\_subset }\OtherTok{\textless{}{-}}\NormalTok{ survey\_design }\SpecialCharTok{\%\textgreater{}\%}
  \FunctionTok{filter}\NormalTok{(Sex }\SpecialCharTok{==} \StringTok{"Male"}\NormalTok{)}
\end{Highlighting}
\end{Shaded}

\section{Estimación del tamaño poblacional}\label{estimacion-del-tamano-poblacional}

En el análisis de encuestas de hogares, es fundamental estimar el tamaño de las subpoblaciones, entendido como el número de personas u hogares que pertenecen a categorías específicas y la proporción que representan dentro de la población. Estas estimaciones, basadas en variables categóricas, permiten caracterizar el perfil demográfico y socioeconómico de la población, información clave para orientar la asignación de recursos, diseñar políticas públicas y formular programas sociales. Por ejemplo, estimar cuántas personas se encuentran por debajo de la línea de pobreza, cuántas están desempleadas o cuántas alcanzaron determinado nivel educativo es esencial para identificar brechas y evaluar condiciones de bienestar.

La estimación del tamaño de una población o subpoblación se realiza a partir de variables categóricas, que segmentan a la población en grupos mutuamente excluyentes (particiones). El tamaño poblacional hace referencia al número total de individuos u hogares que, en la base de datos de la encuesta, pertenecen a una categoría determinada. Estas categorías pueden corresponder, por ejemplo, a estados de ocupación o niveles educativos alcanzados. Para obtener estas estimaciones, se combinan las respuestas de los encuestados con los pesos muestrales, que indican cuántas personas u hogares representa cada unidad de la muestra dentro de la población total.

El estimador del tamaño de la población se define como:

\[
\hat{N} = \sum_{h} \sum_{i} \sum_{k} w_{hik}
\]

En donde, \(w_{hik}\) es el peso o factor de expansión de la unidad \(k\) en la UPM \(i\) del estrato \(h\). Asimismo, la estimación del tamaño de una subpoblación sigue el mismo principio, pero se enmarcado en un subconjunto definido por una característica específica. Para determinar cuántas personas pertenecen a una categoría particular, se identifica dicho grupo en la base de datos y se suman sus pesos muestrales. Esto permite cuantificar grupos de interés específicos y conocer su tamaño dentro de la población:

\[
\hat{N}_d = \sum_{h} \sum_{i} \sum_{k} w_{hik} \ I(y_{hik}=d)
\]

En donde \(I(y_{hik}=d)\) es una variable binaria que toma el valor de 1 si la unidad \(k\) de la UPM \(i\) en el estrato \(h\) pertenece a la categoría \(d\) de la variable de interés \(y\), y 0 en caso contrario. Nótese además que, si \(d\) fue utilizada en la calibración de los pesos, el valor de \(\hat{N}_d\) coincidirá con el control externo aplicado.

En \texttt{R}, y utilizando el diseño muestral definido previamente, es posible generar resúmenes de la información de la encuesta para cada categoría de la variable \texttt{Zone}. Para ello, los datos se agrupan según la zona geográfica, lo que permite obtener resultados independientes para cada uno de estos dominios de estudio. A partir de este procedimiento se producen dos tipos principales de resultados. En primer lugar, se calcula el número de observaciones efectivamente recolectadas en la muestra, sin considerar los factores de expansión, lo que corresponde al tamaño muestral disponible en cada zona. En segundo lugar, se estima el tamaño poblacional asociado a cada dominio incorporando la información del diseño de muestreo.

Junto con las estimaciones poblacionales también se obtienen medidas de precisión, como el error estándar y el intervalo de confianza, las cuales permiten evaluar el nivel de incertidumbre asociado a los resultados producidos a partir de la muestra. Los resultados se presentan en la Tabla \ref{tab:tab-tamano-zona}. En dicha tabla, \texttt{n} representa el número de observaciones de la muestra en cada zona, mientras que \texttt{size} corresponde a la estimación del total poblacional de la subpoblación. Por su parte, la función \texttt{unweighted()} se utiliza para calcular resúmenes no ponderados directamente a partir de los datos observados en la muestra.

\begin{Shaded}
\begin{Highlighting}[]
\NormalTok{survey\_design }\SpecialCharTok{\%\textgreater{}\%}
  \FunctionTok{group\_by}\NormalTok{(Zone) }\SpecialCharTok{\%\textgreater{}\%}
  \FunctionTok{summarise}\NormalTok{(}
    \AttributeTok{n =} \FunctionTok{unweighted}\NormalTok{(}\FunctionTok{n}\NormalTok{()),}
    \AttributeTok{size =} \FunctionTok{survey\_total}\NormalTok{(}\AttributeTok{vartype =} \FunctionTok{c}\NormalTok{(}\StringTok{"se"}\NormalTok{, }\StringTok{"ci"}\NormalTok{))}
\NormalTok{  )}
\end{Highlighting}
\end{Shaded}

\begin{longtable}[]{@{}lccccc@{}}
\caption{\label{tab:tab-tamano-zona}Tamaño poblacional estimado por zona geográfica}\tabularnewline
\toprule\noalign{}
Zone & n & size & size\_se & size\_low & size\_upp \\
\midrule\noalign{}
\endfirsthead
\toprule\noalign{}
Zone & n & size & size\_se & size\_low & size\_upp \\
\midrule\noalign{}
\endhead
\bottomrule\noalign{}
\endlastfoot
Rural & 1297 & 72102 & 3062 & 66039 & 78165 \\
Urban & 1308 & 78164 & 2847 & 72526 & 83802 \\
\end{longtable}

En efecto, el tamaño de muestra fue de 1297 personas en la zona rural y de 1308 en la zona urbana. A partir de estas muestras se estimó una población de 72,102 personas para la zona rural, con un error estándar de 3,062, y de 78,164 personas para la zona urbana, con un error estándar de 2,847. Estos resultados reflejan no solo el tamaño estimado de la población en cada dominio, sino también el nivel de precisión asociado a las estimaciones obtenidas a partir del diseño muestral.

Considerando un nivel de confianza del 95\%, los intervalos de confianza para las estimaciones poblacionales fueron de 66,038.5 a 78,165.4 personas en la zona rural y de 72,526.2 a 83,801.7 personas en la zona urbana. De manera análoga, también es posible estimar el número de personas según los distintos niveles de pobreza, cuyos resultados se presentan en la Tabla \ref{tab:tab-tamano-pobreza}.

\begin{Shaded}
\begin{Highlighting}[]
\NormalTok{survey\_design }\SpecialCharTok{\%\textgreater{}\%}
  \FunctionTok{group\_by}\NormalTok{(Poverty) }\SpecialCharTok{\%\textgreater{}\%}
  \FunctionTok{summarise}\NormalTok{(}\AttributeTok{size =} \FunctionTok{survey\_total}\NormalTok{(}\AttributeTok{vartype =} \FunctionTok{c}\NormalTok{(}\StringTok{"se"}\NormalTok{, }\StringTok{"ci"}\NormalTok{)))}
\end{Highlighting}
\end{Shaded}

\begin{longtable}[]{@{}lcccc@{}}
\caption{\label{tab:tab-tamano-pobreza}Tamaño poblacional estimado por condición de pobreza}\tabularnewline
\toprule\noalign{}
Poverty & size & size\_se & size\_low & size\_upp \\
\midrule\noalign{}
\endfirsthead
\toprule\noalign{}
Poverty & size & size\_se & size\_low & size\_upp \\
\midrule\noalign{}
\endhead
\bottomrule\noalign{}
\endlastfoot
NotPoor & 91398 & 4395 & 82696 & 100101 \\
Extreme & 21519 & 4949 & 11719 & 31319 \\
Relative & 37349 & 3695 & 30032 & 44666 \\
\end{longtable}

Otra variable categórica relevante en las encuestas de hogares es el estado de ocupación. A continuación se presenta el código correspondiente, cuyos resultados se muestran en la Tabla \ref{tab:tab-tamano-empleo}. En este caso, la información se agrupa según las categorías de la variable \texttt{Employment}, lo que permite obtener estimaciones separadas para cada condición de actividad laboral de la población. Mediante el uso de la función \texttt{group\_by()} es posible desagregar las estimaciones en niveles más detallados, facilitando el análisis comparativo entre los distintos grupos ocupacionales.

Antes de realizar las estimaciones, se excluyen los registros con valores faltantes en la variable de empleo, garantizando que los resultados se calculen únicamente con información válida. Posteriormente, para cada categoría de ocupación se estima el tamaño poblacional utilizando el diseño muestral definido previamente. Además de la estimación puntual, también se calculan medidas de precisión, específicamente el error estándar y los intervalos de confianza, lo que permite evaluar el grado de incertidumbre asociado a cada estimación.

\begin{Shaded}
\begin{Highlighting}[]
\NormalTok{survey\_design }\SpecialCharTok{\%\textgreater{}\%}
  \FunctionTok{filter}\NormalTok{(}\SpecialCharTok{!}\FunctionTok{is.na}\NormalTok{(Employment)) }\SpecialCharTok{\%\textgreater{}\%}
  \FunctionTok{group\_by}\NormalTok{(Employment) }\SpecialCharTok{\%\textgreater{}\%}
  \FunctionTok{summarise}\NormalTok{(}\AttributeTok{size =} \FunctionTok{survey\_total}\NormalTok{(}\AttributeTok{vartype =} \FunctionTok{c}\NormalTok{(}\StringTok{"se"}\NormalTok{, }\StringTok{"ci"}\NormalTok{)))}
\end{Highlighting}
\end{Shaded}

\begin{longtable}[]{@{}lcccc@{}}
\caption{\label{tab:tab-tamano-empleo}Tamaño poblacional estimado por estado de ocupación}\tabularnewline
\toprule\noalign{}
Employment & size & size\_se & size\_low & size\_upp \\
\midrule\noalign{}
\endfirsthead
\toprule\noalign{}
Employment & size & size\_se & size\_low & size\_upp \\
\midrule\noalign{}
\endhead
\bottomrule\noalign{}
\endlastfoot
Unemployed & 4635 & 761 & 3129 & 6141 \\
Inactive & 41465 & 2163 & 37183 & 45748 \\
Employed & 61877 & 2540 & 56847 & 66907 \\
\end{longtable}

A partir de los resultados obtenidos se estima que 4,634.8 personas se encuentran desempleadas, con un intervalo de confianza del 95\% entre 3,128.6 y 6,140.9 personas. Asimismo, se estima que 41,465.2 personas pertenecen a la población inactiva, con un intervalo de confianza entre 37,182.6 y 45,747.8. Finalmente, la población ocupada se estimó en 61,877.0 personas, con un intervalo de confianza de 36,784.2 a 47,793.5 personas.

Finalmente, las estimaciones también pueden desagregarse simultáneamente por más de una variable categórica. Por ejemplo, es posible analizar el estado ocupacional según el nivel de pobreza, cuyos resultados se presentan en la Tabla \ref{tab:tab-tamano-empleo-pobreza}, de donde se puede concluir, entre otros aspectos, que aproximadamente 44,600.3 personas ocupadas no se encuentran en condición de pobreza, con un intervalo de confianza del 95\% entre 39,459.6 y 49,741.0 personas. Asimismo, se estima que 6,421.8 personas inactivas se encuentran en situación de pobreza extrema, con un intervalo de confianza al 95\% entre 3,806.6 y 9,037.0 personas.

\begin{Shaded}
\begin{Highlighting}[]
\NormalTok{survey\_design }\SpecialCharTok{\%\textgreater{}\%}
  \FunctionTok{filter}\NormalTok{(}\SpecialCharTok{!}\FunctionTok{is.na}\NormalTok{(Employment)) }\SpecialCharTok{\%\textgreater{}\%}
  \FunctionTok{group\_by}\NormalTok{(Employment, Poverty) }\SpecialCharTok{\%\textgreater{}\%}
  \FunctionTok{cascade}\NormalTok{(}
    \AttributeTok{size =} \FunctionTok{survey\_total}\NormalTok{(}\AttributeTok{vartype =} \FunctionTok{c}\NormalTok{(}\StringTok{"se"}\NormalTok{, }\StringTok{"ci"}\NormalTok{)),}
    \AttributeTok{.fill =} \StringTok{"Total"}
\NormalTok{  )}
\end{Highlighting}
\end{Shaded}

\begin{longtable}[]{@{}lccccc@{}}
\caption{\label{tab:tab-tamano-empleo-pobreza}Tamaño poblacional estimado por estado de ocupación y condición de pobreza}\tabularnewline
\toprule\noalign{}
Employment & Poverty & size & size\_se & size\_low & size\_upp \\
\midrule\noalign{}
\endfirsthead
\toprule\noalign{}
Employment & Poverty & size & size\_se & size\_low & size\_upp \\
\midrule\noalign{}
\endhead
\bottomrule\noalign{}
\endlastfoot
Unemployed & NotPoor & 1768 & 405 & 966 & 2571 \\
Unemployed & Extreme & 1169 & 348 & 480 & 1859 \\
Unemployed & Relative & 1697 & 458 & 791 & 2604 \\
Unemployed & Total & 4635 & 761 & 3129 & 6141 \\
Inactive & NotPoor & 24346 & 1736 & 20908 & 27784 \\
Inactive & Extreme & 6422 & 1321 & 3807 & 9037 \\
Inactive & Relative & 10697 & 1460 & 7806 & 13589 \\
Inactive & Total & 41465 & 2163 & 37183 & 45748 \\
Employed & NotPoor & 44600 & 2596 & 39460 & 49741 \\
Employed & Extreme & 5128 & 1122 & 2907 & 7349 \\
Employed & Relative & 12149 & 1347 & 9483 & 14816 \\
Employed & Total & 61877 & 2540 & 56847 & 66907 \\
Total & Total & 107977 & 3466 & 101115 & 114839 \\
\end{longtable}

\section{Estimación de proporciones}\label{estimacion-de-proporciones}

Las proporciones permiten expresar el peso relativo que tienen determinados grupos dentro de la población. Por ejemplo, conocer el porcentaje de hogares que se encuentran por debajo de la línea de pobreza es esencial para evaluar desigualdades y caracterizar condiciones de bienestar. Para obtener este tipo de indicadores, se calcula el promedio ponderado de la variable dicotómica, lo que asegura que la estimación represente adecuadamente la distribución poblacional. De acuerdo con \citet{heeringa2017applied}, al convertir las categorías originales en variables indicadoras, la proporción estimada se calcula como:

\[
\hat{p}_d = \frac{\hat{N}_d}{\hat{N}} = \frac{\displaystyle\sum_{h} \sum_{i} \sum_{k} w_{hik} \ I(y_{hik}=d)} {\displaystyle\sum_{h} \sum_{i} \sum_{k} w_{hik}}
\]

Si bien este estimador es conceptualmente sencillo, corresponde a una función no lineal de los totales poblacionales. Por esta razón, para derivar sus propiedades estadísticas, particularmente su varianza y error estándar, es necesario recurrir a métodos de aproximación como la linealización de Taylor. En este contexto, se utiliza la función \(z_{hik}=I(y_{hik}=d)-\hat{p}_d\). En la práctica, los programas estadísticos implementan automáticamente estos procedimientos y generan las estimaciones de proporciones junto con sus respectivos errores estándar e intervalos de confianza.

Cuando las proporciones toman valores muy cercanos a 0 o a 1, la precisión de las estimaciones puede verse afectada, por lo que en muchos casos es necesario contar con tamaños muestrales mayores para obtener resultados estables y confiables. Esta situación es especialmente relevante en el análisis de fenómenos poco frecuentes o altamente concentrados, donde pequeñas variaciones en la muestra pueden producir cambios importantes en las estimaciones.

Adicionalmente, los intervalos de confianza construidos mediante aproximaciones normales pueden presentar límites inferiores menores que 0 o superiores mayores que 1, lo que dificulta su interpretación al tratarse de proporciones. Para evitar este problema, una alternativa ampliamente utilizada consiste en aplicar la transformación logit antes de construir los intervalos de confianza. Este procedimiento garantiza que, al regresar a la escala original, los límites obtenidos permanezcan dentro del rango válido de las proporciones.

\[
CI(\hat{p}_d; 1-\alpha) = \frac{exp \left[ln\left(\frac{\hat{p}_d}{1-\hat{p}_d}\right) \pm \frac{\widehat{me}}{\hat{p}_d(1-\hat{p}_d)}\right]}{1 + exp \left[ln\left(\frac{\hat{p}_d}{1-\hat{p}_d}\right) \pm \frac{\widehat{me}}{\hat{p}_d(1-\hat{p}_d)}\right]}
\]

En donde \(\widehat{me}\) corresponde al margen de error estimado y se obtiene a partir del producto entre el valor crítico de la distribución \(t\) de Student y el error estándar estimado de la proporción \(\hat{p}_d\). El término \(t_{1-\alpha/2, df}\) representa el percentil asociado al nivel de confianza seleccionado, considerando \(df\) grados de libertad. En el contexto de encuestas complejas, estos grados de libertad suelen definirse como la diferencia entre el número de unidades primarias de muestreo (UPM) y el número de estratos presentes en el subgrupo o dominio de análisis.

A diferencia de otros programas estadísticos que presentan los resultados en porcentaje, \texttt{R} reporta las proporciones en la escala {[}0,1{]}. Las funciones de estos paquetes permiten calcular directamente las estimaciones, sus errores estándar e intervalos de confianza, lo que facilita el análisis de variables categóricas en encuestas de hogares. A continuación se ilustra cómo obtener la proporción de personas por zona usando el diseño muestral previamente definido, presentada en la Tabla \ref{tab:tab-prop-zona}:

\begin{Shaded}
\begin{Highlighting}[]
\NormalTok{survey\_design }\SpecialCharTok{\%\textgreater{}\%}
  \FunctionTok{group\_by}\NormalTok{(Zone) }\SpecialCharTok{\%\textgreater{}\%}
  \FunctionTok{summarise}\NormalTok{(}\AttributeTok{proportion =} \FunctionTok{survey\_mean}\NormalTok{(}
    \AttributeTok{vartype =} \FunctionTok{c}\NormalTok{(}\StringTok{"se"}\NormalTok{, }\StringTok{"ci"}\NormalTok{),}
    \AttributeTok{proportion =} \ConstantTok{TRUE}
\NormalTok{  ))}
\end{Highlighting}
\end{Shaded}

\begin{longtable}[]{@{}lcccc@{}}
\caption{\label{tab:tab-prop-zona}Proporción de personas por zona geográfica}\tabularnewline
\toprule\noalign{}
Zone & proportion & proportion\_se & proportion\_low & proportion\_upp \\
\midrule\noalign{}
\endfirsthead
\toprule\noalign{}
Zone & proportion & proportion\_se & proportion\_low & proportion\_upp \\
\midrule\noalign{}
\endhead
\bottomrule\noalign{}
\endlastfoot
Rural & 0.48 & 0.014 & 0.452 & 0.508 \\
Urban & 0.52 & 0.014 & 0.492 & 0.548 \\
\end{longtable}

En este ejemplo, la función \texttt{survey\_mean()} se utiliza para calcular la estimación de proporciones poblacionales y, mediante el argumento \texttt{proportion\ =\ TRUE}, se especifica que el interés se centra en estimar la distribución relativa de la población entre las distintas categorías de la variable analizada. Los resultados indican que aproximadamente el 47.9\% de las personas reside en la zona rural, con un intervalo de confianza del 95\% entre 45.2\% y 50.7\%, mientras que el 52.0\% corresponde a población residente en la zona urbana, con un intervalo de confianza entre 49.2\% y 54.7\%.

La librería \texttt{survey} también dispone de la función \texttt{survey\_prop()}, diseñada específicamente para la estimación de proporciones en encuestas complejas. Esta función produce resultados equivalentes a los obtenidos con \texttt{survey\_mean()}, facilitando además una sintaxis más directa e intuitiva cuando el objetivo del análisis es exclusivamente el cálculo de proporciones. Los resultados obtenidos mediante esta alternativa se presentan en la Tabla \ref{tab:tab-prop-zona2}.

A partir de estas estimaciones se observa que, en la zona urbana, el 53.6\% de la población corresponde a mujeres, con un intervalo de confianza del 95\% entre 51.0\% y 56.2\%, mientras que el 46.4\% corresponde a hombres, con un intervalo de confianza entre 43.7\% y 48.9\%. Estos resultados permiten apreciar no solo la distribución porcentual de la población según sexo, sino también el nivel de precisión asociado a las estimaciones.

\begin{Shaded}
\begin{Highlighting}[]
\NormalTok{survey\_design }\SpecialCharTok{\%\textgreater{}\%}
  \FunctionTok{group\_by}\NormalTok{(Zone) }\SpecialCharTok{\%\textgreater{}\%}
  \FunctionTok{summarise}\NormalTok{(}\AttributeTok{proportion =} \FunctionTok{survey\_prop}\NormalTok{(}\AttributeTok{vartype =} \FunctionTok{c}\NormalTok{(}\StringTok{"se"}\NormalTok{, }\StringTok{"ci"}\NormalTok{)))}
\end{Highlighting}
\end{Shaded}

\begin{longtable}[]{@{}lcccc@{}}
\caption{\label{tab:tab-prop-zona2}Proporción de personas por zona geográfica (función survey\_prop)}\tabularnewline
\toprule\noalign{}
Zone & proportion & proportion\_se & proportion\_low & proportion\_upp \\
\midrule\noalign{}
\endfirsthead
\toprule\noalign{}
Zone & proportion & proportion\_se & proportion\_low & proportion\_upp \\
\midrule\noalign{}
\endhead
\bottomrule\noalign{}
\endlastfoot
Rural & 0.48 & 0.014 & 0.452 & 0.508 \\
Urban & 0.52 & 0.014 & 0.492 & 0.548 \\
\end{longtable}

Si el interés se centra ahora en estimar proporciones para subpoblaciones específicas, por ejemplo, la proporción de hombres y mujeres que residen en la zona urbana, es necesario restringir el análisis a dicho dominio de estudio y posteriormente calcular las estimaciones correspondientes para cada categoría de sexo. El código computacional correspondiente se presenta en la Tabla \ref{tab:tab-prop-sexo-urbano}.

\begin{Shaded}
\begin{Highlighting}[]
\NormalTok{urban\_subset }\SpecialCharTok{\%\textgreater{}\%}
  \FunctionTok{group\_by}\NormalTok{(Sex) }\SpecialCharTok{\%\textgreater{}\%}
  \FunctionTok{summarise}\NormalTok{(}\AttributeTok{proportion =} \FunctionTok{survey\_prop}\NormalTok{(}\AttributeTok{vartype =} \FunctionTok{c}\NormalTok{(}\StringTok{"se"}\NormalTok{, }\StringTok{"ci"}\NormalTok{)))}
\end{Highlighting}
\end{Shaded}

\begin{longtable}[]{@{}lcccc@{}}
\caption{\label{tab:tab-prop-sexo-urbano}Proporción de hombres y mujeres en zona urbana}\tabularnewline
\toprule\noalign{}
Sex & proportion & proportion\_se & proportion\_low & proportion\_upp \\
\midrule\noalign{}
\endfirsthead
\toprule\noalign{}
Sex & proportion & proportion\_se & proportion\_low & proportion\_upp \\
\midrule\noalign{}
\endhead
\bottomrule\noalign{}
\endlastfoot
Female & 0.537 & 0.013 & 0.511 & 0.563 \\
Male & 0.463 & 0.013 & 0.437 & 0.489 \\
\end{longtable}

De manera análoga, es posible realizar el mismo ejercicio para la población residente en la zona rural, estimando la proporción de hombres y mujeres. Este tipo de desagregación permite comparar la composición poblacional según sexo entre las áreas urbana y rural, proporcionando una visión más detallada de las características demográficas de la población encuestada. Los resultados correspondientes se presentan en la Tabla \ref{tab:tab-prop-sexo-rural}.

\begin{Shaded}
\begin{Highlighting}[]
\NormalTok{rural\_subset }\SpecialCharTok{\%\textgreater{}\%}
  \FunctionTok{group\_by}\NormalTok{(Sex) }\SpecialCharTok{\%\textgreater{}\%}
  \FunctionTok{summarise}\NormalTok{(}
    \AttributeTok{n =} \FunctionTok{unweighted}\NormalTok{(}\FunctionTok{n}\NormalTok{()),}
    \AttributeTok{proportion =} \FunctionTok{survey\_prop}\NormalTok{(}\AttributeTok{vartype =} \FunctionTok{c}\NormalTok{(}\StringTok{"se"}\NormalTok{, }\StringTok{"ci"}\NormalTok{))}
\NormalTok{  )}
\end{Highlighting}
\end{Shaded}

\begin{longtable}[]{@{}
  >{\raggedright\arraybackslash}p{(\linewidth - 10\tabcolsep) * \real{0.0986}}
  >{\centering\arraybackslash}p{(\linewidth - 10\tabcolsep) * \real{0.0704}}
  >{\centering\arraybackslash}p{(\linewidth - 10\tabcolsep) * \real{0.1690}}
  >{\centering\arraybackslash}p{(\linewidth - 10\tabcolsep) * \real{0.2113}}
  >{\centering\arraybackslash}p{(\linewidth - 10\tabcolsep) * \real{0.2254}}
  >{\centering\arraybackslash}p{(\linewidth - 10\tabcolsep) * \real{0.2254}}@{}}
\caption{\label{tab:tab-prop-sexo-rural}Proporción de hombres y mujeres en zona rural}\tabularnewline
\toprule\noalign{}
\begin{minipage}[b]{\linewidth}\raggedright
Sex
\end{minipage} & \begin{minipage}[b]{\linewidth}\centering
n
\end{minipage} & \begin{minipage}[b]{\linewidth}\centering
proportion
\end{minipage} & \begin{minipage}[b]{\linewidth}\centering
proportion\_se
\end{minipage} & \begin{minipage}[b]{\linewidth}\centering
proportion\_low
\end{minipage} & \begin{minipage}[b]{\linewidth}\centering
proportion\_upp
\end{minipage} \\
\midrule\noalign{}
\endfirsthead
\toprule\noalign{}
\begin{minipage}[b]{\linewidth}\raggedright
Sex
\end{minipage} & \begin{minipage}[b]{\linewidth}\centering
n
\end{minipage} & \begin{minipage}[b]{\linewidth}\centering
proportion
\end{minipage} & \begin{minipage}[b]{\linewidth}\centering
proportion\_se
\end{minipage} & \begin{minipage}[b]{\linewidth}\centering
proportion\_low
\end{minipage} & \begin{minipage}[b]{\linewidth}\centering
proportion\_upp
\end{minipage} \\
\midrule\noalign{}
\endhead
\bottomrule\noalign{}
\endlastfoot
Female & 679 & 0.516 & 0.008 & 0.500 & 0.533 \\
Male & 618 & 0.484 & 0.008 & 0.467 & 0.500 \\
\end{longtable}

Ahora bien, si el análisis se restringe únicamente a la población masculina incluida en la base de datos y el interés consiste en estimar la distribución porcentual de los hombres según la zona de residencia, es posible calcular la proporción de hombres que viven en áreas urbanas y rurales utilizando el diseño muestral previamente definido. El código computacional correspondiente se presenta en la Tabla \ref{tab:tab-prop-zona-hombre}.

\begin{Shaded}
\begin{Highlighting}[]
\NormalTok{male\_subset }\SpecialCharTok{\%\textgreater{}\%}
  \FunctionTok{group\_by}\NormalTok{(Zone) }\SpecialCharTok{\%\textgreater{}\%}
  \FunctionTok{summarise}\NormalTok{(}\AttributeTok{proportion =} \FunctionTok{survey\_prop}\NormalTok{(}\AttributeTok{vartype =} \FunctionTok{c}\NormalTok{(}\StringTok{"se"}\NormalTok{, }\StringTok{"ci"}\NormalTok{)))}
\end{Highlighting}
\end{Shaded}

\begin{longtable}[]{@{}lcccc@{}}
\caption{\label{tab:tab-prop-zona-hombre}Distribución por zona en la subpoblación masculina}\tabularnewline
\toprule\noalign{}
Zone & proportion & proportion\_se & proportion\_low & proportion\_upp \\
\midrule\noalign{}
\endfirsthead
\toprule\noalign{}
Zone & proportion & proportion\_se & proportion\_low & proportion\_upp \\
\midrule\noalign{}
\endhead
\bottomrule\noalign{}
\endlastfoot
Rural & 0.491 & 0.018 & 0.455 & 0.526 \\
Urban & 0.509 & 0.018 & 0.474 & 0.545 \\
\end{longtable}

Si se desea estimar resultados con varios niveles de desagregación, al igual que en los capítulos anteriores, el uso de la función \texttt{group\_by()} permite generar distintos niveles de agregación mediante la combinación de dos o más variables categóricas. Por ejemplo, es posible estimar la proporción de hombres según zona de residencia y condición de pobreza de manera simultánea. Este tipo de análisis permite caracterizar con mayor detalle las condiciones socioeconómicas de subpoblaciones específicas. El procedimiento correspondiente se presenta en la Tabla \ref{tab:tab-prop-zona-pobreza-hombre}.

\begin{Shaded}
\begin{Highlighting}[]
\NormalTok{male\_subset }\SpecialCharTok{\%\textgreater{}\%}
  \FunctionTok{group\_by}\NormalTok{(Zone, Poverty) }\SpecialCharTok{\%\textgreater{}\%}
  \FunctionTok{summarise}\NormalTok{(}\AttributeTok{proportion =} \FunctionTok{survey\_prop}\NormalTok{(}\AttributeTok{vartype =} \FunctionTok{c}\NormalTok{(}\StringTok{"se"}\NormalTok{, }\StringTok{"ci"}\NormalTok{)))}
\end{Highlighting}
\end{Shaded}

\begin{longtable}[]{@{}
  >{\raggedright\arraybackslash}p{(\linewidth - 10\tabcolsep) * \real{0.0800}}
  >{\centering\arraybackslash}p{(\linewidth - 10\tabcolsep) * \real{0.1333}}
  >{\centering\arraybackslash}p{(\linewidth - 10\tabcolsep) * \real{0.1600}}
  >{\centering\arraybackslash}p{(\linewidth - 10\tabcolsep) * \real{0.2000}}
  >{\centering\arraybackslash}p{(\linewidth - 10\tabcolsep) * \real{0.2133}}
  >{\centering\arraybackslash}p{(\linewidth - 10\tabcolsep) * \real{0.2133}}@{}}
\caption{\label{tab:tab-prop-zona-pobreza-hombre}Proporción de hombres por zona y condición de pobreza}\tabularnewline
\toprule\noalign{}
\begin{minipage}[b]{\linewidth}\raggedright
Zone
\end{minipage} & \begin{minipage}[b]{\linewidth}\centering
Poverty
\end{minipage} & \begin{minipage}[b]{\linewidth}\centering
proportion
\end{minipage} & \begin{minipage}[b]{\linewidth}\centering
proportion\_se
\end{minipage} & \begin{minipage}[b]{\linewidth}\centering
proportion\_low
\end{minipage} & \begin{minipage}[b]{\linewidth}\centering
proportion\_upp
\end{minipage} \\
\midrule\noalign{}
\endfirsthead
\toprule\noalign{}
\begin{minipage}[b]{\linewidth}\raggedright
Zone
\end{minipage} & \begin{minipage}[b]{\linewidth}\centering
Poverty
\end{minipage} & \begin{minipage}[b]{\linewidth}\centering
proportion
\end{minipage} & \begin{minipage}[b]{\linewidth}\centering
proportion\_se
\end{minipage} & \begin{minipage}[b]{\linewidth}\centering
proportion\_low
\end{minipage} & \begin{minipage}[b]{\linewidth}\centering
proportion\_upp
\end{minipage} \\
\midrule\noalign{}
\endhead
\bottomrule\noalign{}
\endlastfoot
Rural & NotPoor & 0.549 & 0.063 & 0.424 & 0.668 \\
Rural & Extreme & 0.198 & 0.067 & 0.096 & 0.364 \\
Rural & Relative & 0.254 & 0.037 & 0.187 & 0.334 \\
Urban & NotPoor & 0.660 & 0.037 & 0.584 & 0.728 \\
Urban & Extreme & 0.113 & 0.025 & 0.073 & 0.171 \\
Urban & Relative & 0.227 & 0.026 & 0.180 & 0.283 \\
\end{longtable}

A partir de los resultados obtenidos se observa que, en la zona rural, el 19.7\% de los hombres se encuentra en situación de pobreza extrema, mientras que en la zona urbana esta proporción corresponde al 11.3\%. Aunque las estimaciones puntuales sugieren una mayor incidencia de pobreza extrema en la ruralidad, la comparación de los intervalos de confianza indica que no existe evidencia suficiente para concluir que las diferencias sean estadísticamente significativas al nivel de confianza considerado, debido a que ambos intervalos presentan superposición.

También es posible categorizar variables cuantitativas para facilitar su análisis conjunto con otras variables categóricas. Por ejemplo, la edad puede agruparse en rangos etarios específicos y posteriormente cruzarse con variables relacionadas con la situación en el empleo. Este tipo de procedimiento resulta útil para identificar diferencias en la participación laboral entre distintos grupos de edad.

A continuación, la variable edad se clasifica en dos categorías para la población femenina: mujeres entre 18 y 35 años, y mujeres pertenecientes a otros rangos de edad. Posteriormente, estas categorías se analizan en conjunto con la variable de empleabilidad, permitiendo estimar la distribución de la condición laboral según grupo etario. Los resultados correspondientes se presentan en la Tabla \ref{tab:tab-prop-edad-rango-empleo-mujer}.

\begin{Shaded}
\begin{Highlighting}[]
\NormalTok{female\_subset }\SpecialCharTok{\%\textgreater{}\%}
  \FunctionTok{filter}\NormalTok{(}\SpecialCharTok{!}\FunctionTok{is.na}\NormalTok{(Employment)) }\SpecialCharTok{\%\textgreater{}\%}
  \FunctionTok{mutate}\NormalTok{(}\AttributeTok{age\_range =} \FunctionTok{case\_when}\NormalTok{(}
\NormalTok{    Age }\SpecialCharTok{\textgreater{}=} \DecValTok{18} \SpecialCharTok{\&}\NormalTok{ Age }\SpecialCharTok{\textless{}=} \DecValTok{35} \SpecialCharTok{\textasciitilde{}} \StringTok{"18 {-} 35"}\NormalTok{,}
    \ConstantTok{TRUE} \SpecialCharTok{\textasciitilde{}} \StringTok{"Other"}
\NormalTok{  )) }\SpecialCharTok{\%\textgreater{}\%}
  \FunctionTok{group\_by}\NormalTok{(age\_range, Employment) }\SpecialCharTok{\%\textgreater{}\%}
  \FunctionTok{summarise}\NormalTok{(}\AttributeTok{proportion =} \FunctionTok{survey\_prop}\NormalTok{(}\AttributeTok{vartype =} \FunctionTok{c}\NormalTok{(}\StringTok{"se"}\NormalTok{, }\StringTok{"ci"}\NormalTok{)))}
\end{Highlighting}
\end{Shaded}

\begin{longtable}[]{@{}
  >{\raggedright\arraybackslash}p{(\linewidth - 10\tabcolsep) * \real{0.1235}}
  >{\centering\arraybackslash}p{(\linewidth - 10\tabcolsep) * \real{0.1481}}
  >{\centering\arraybackslash}p{(\linewidth - 10\tabcolsep) * \real{0.1481}}
  >{\centering\arraybackslash}p{(\linewidth - 10\tabcolsep) * \real{0.1852}}
  >{\centering\arraybackslash}p{(\linewidth - 10\tabcolsep) * \real{0.1975}}
  >{\centering\arraybackslash}p{(\linewidth - 10\tabcolsep) * \real{0.1975}}@{}}
\caption{\label{tab:tab-prop-edad-rango-empleo-mujer}Proporción de mujeres por rango de edad y estado de ocupación}\tabularnewline
\toprule\noalign{}
\begin{minipage}[b]{\linewidth}\raggedright
age\_range
\end{minipage} & \begin{minipage}[b]{\linewidth}\centering
Employment
\end{minipage} & \begin{minipage}[b]{\linewidth}\centering
proportion
\end{minipage} & \begin{minipage}[b]{\linewidth}\centering
proportion\_se
\end{minipage} & \begin{minipage}[b]{\linewidth}\centering
proportion\_low
\end{minipage} & \begin{minipage}[b]{\linewidth}\centering
proportion\_upp
\end{minipage} \\
\midrule\noalign{}
\endfirsthead
\toprule\noalign{}
\begin{minipage}[b]{\linewidth}\raggedright
age\_range
\end{minipage} & \begin{minipage}[b]{\linewidth}\centering
Employment
\end{minipage} & \begin{minipage}[b]{\linewidth}\centering
proportion
\end{minipage} & \begin{minipage}[b]{\linewidth}\centering
proportion\_se
\end{minipage} & \begin{minipage}[b]{\linewidth}\centering
proportion\_low
\end{minipage} & \begin{minipage}[b]{\linewidth}\centering
proportion\_upp
\end{minipage} \\
\midrule\noalign{}
\endhead
\bottomrule\noalign{}
\endlastfoot
18 - 35 & Unemployed & 0.029 & 0.009 & 0.015 & 0.054 \\
18 - 35 & Inactive & 0.517 & 0.038 & 0.442 & 0.591 \\
18 - 35 & Employed & 0.455 & 0.036 & 0.385 & 0.526 \\
Other & Unemployed & 0.016 & 0.007 & 0.007 & 0.036 \\
Other & Inactive & 0.571 & 0.030 & 0.511 & 0.629 \\
Other & Employed & 0.413 & 0.029 & 0.356 & 0.472 \\
\end{longtable}

\section{Contrastes y diferencias de proporciones}\label{contrastes-y-diferencias-de-proporciones}

Según \citet{heeringa2017applied}, las proporciones obtenidas en las filas de una tabla de doble entrada pueden interpretarse como estimaciones para distintas subpoblaciones, definidas a partir de los niveles de una variable categórica. En este contexto, no solo resulta relevante estimar las proporciones asociadas a cada grupo, sino también evaluar las diferencias existentes entre ellas mediante contrastes lineales.

Por ejemplo, suponga que el interés está en comparar la proporción de mujeres en condición de pobreza con la proporción de hombres en esta misma situación. Formalmente, este contraste puede expresarse como \(\hat{p}_F - \hat{p}_M\), donde \(\hat{p}_F\) representa la proporción estimada de mujeres en pobreza y \(\hat{p}_M\) la proporción estimada de hombres en dicha condición. Para construir este contraste, primero se estiman las proporciones de hombres y mujeres en estado de pobreza utilizando el mismo procedimiento descrito en capítulos anteriores. Los resultados correspondientes se presentan en la Tabla \ref{tab:tab-svyby-pobreza-sexo2}:

\begin{Shaded}
\begin{Highlighting}[]
\NormalTok{sex\_poverty\_proportion }\OtherTok{\textless{}{-}}\NormalTok{ survey\_design }\SpecialCharTok{\%\textgreater{}\%}
  \FunctionTok{group\_by}\NormalTok{(Sex) }\SpecialCharTok{\%\textgreater{}\%}
  \FunctionTok{summarise}\NormalTok{(}\AttributeTok{proportion =} \FunctionTok{survey\_mean}\NormalTok{(poor,}
                                     \AttributeTok{na.rm =} \ConstantTok{TRUE}\NormalTok{,}
                                     \AttributeTok{vartype =} \FunctionTok{c}\NormalTok{(}\StringTok{"se"}\NormalTok{, }\StringTok{"ci"}\NormalTok{)))}

\NormalTok{sex\_poverty\_contrast }\OtherTok{\textless{}{-}} \FunctionTok{svyby}\NormalTok{(}
  \AttributeTok{formula =} \SpecialCharTok{\textasciitilde{}}\NormalTok{poor,}
  \AttributeTok{by =} \SpecialCharTok{\textasciitilde{}}\NormalTok{Sex,}
  \AttributeTok{design =}\NormalTok{ survey\_design,}
  \AttributeTok{FUN =}\NormalTok{ svymean,}
  \AttributeTok{na.rm =} \ConstantTok{TRUE}\NormalTok{,}
  \AttributeTok{covmat =} \ConstantTok{TRUE}\NormalTok{,}
  \AttributeTok{vartype =} \FunctionTok{c}\NormalTok{(}\StringTok{"se"}\NormalTok{, }\StringTok{"ci"}\NormalTok{)}
\NormalTok{)}
\end{Highlighting}
\end{Shaded}

\begin{longtable}[]{@{}lcccc@{}}
\caption{\label{tab:tab-svyby-pobreza-sexo2}Proporción en condición de pobreza por sexo}\tabularnewline
\toprule\noalign{}
Sex & proportion & proportion\_se & proportion\_low & proportion\_upp \\
\midrule\noalign{}
\endfirsthead
\toprule\noalign{}
Sex & proportion & proportion\_se & proportion\_low & proportion\_upp \\
\midrule\noalign{}
\endhead
\bottomrule\noalign{}
\endlastfoot
Female & 0.389 & 0.032 & 0.327 & 0.452 \\
Male & 0.395 & 0.037 & 0.322 & 0.467 \\
\end{longtable}

Una vez estimadas las proporciones de hombres y mujeres en condición de pobreza, es posible calcular la diferencia entre ambas proporciones junto con su correspondiente error estándar. En este caso, la diferencia estimada se obtiene restando la proporción de hombres en pobreza a la proporción de mujeres en esta misma condición:

\begin{Shaded}
\begin{Highlighting}[]
\FloatTok{0.3892} \SpecialCharTok{{-}} \FloatTok{0.3946}
\end{Highlighting}
\end{Shaded}

\begin{verbatim}
## [1] -0.0054
\end{verbatim}

Para calcular correctamente la variabilidad asociada a esta diferencia, no basta con considerar únicamente las varianzas de cada estimación por separado. También es necesario incorporar la covarianza entre ambas proporciones, dado que las estimaciones provienen de una misma muestra y, por tanto, no son independientes. La matriz de varianzas y covarianzas puede obtenerse mediante la función \texttt{vcov}, cuyos resultados se presentan en la Tabla \ref{tab:tab-vcov-sex-pobreza}:

\begin{Shaded}
\begin{Highlighting}[]
\FunctionTok{vcov}\NormalTok{(sex\_poverty\_contrast) }\SpecialCharTok{\%\textgreater{}\%}
  \FunctionTok{data.frame}\NormalTok{()}
\end{Highlighting}
\end{Shaded}

\begin{longtable}[]{@{}llc@{}}
\caption{\label{tab:tab-vcov-sex-pobreza}Matriz de varianzas y covarianzas de la proporción de pobreza por sexo}\tabularnewline
\toprule\noalign{}
& Female & Male \\
\midrule\noalign{}
\endfirsthead
\toprule\noalign{}
& Female & Male \\
\midrule\noalign{}
\endhead
\bottomrule\noalign{}
\endlastfoot
Female & 0.000998 & 0.000918 \\
Male & 0.000918 & 0.001342 \\
\end{longtable}

A partir de esta matriz, el error estándar de la diferencia de proporciones se estima utilizando la expresión general para la varianza de una diferencia entre estimadores, la cual corresponde a la suma de las varianzas de cada estimación menos dos veces la covarianza entre ellas. En este caso, el cálculo se realiza de la siguiente manera:

\begin{Shaded}
\begin{Highlighting}[]
\FunctionTok{sqrt}\NormalTok{(}\FloatTok{0.0009983} \SpecialCharTok{+} \FloatTok{0.0013416} \SpecialCharTok{{-}} \DecValTok{2} \SpecialCharTok{*} \FloatTok{0.0009183}\NormalTok{)}
\end{Highlighting}
\end{Shaded}

\begin{verbatim}
## [1] 0.0224
\end{verbatim}

No obstante, el paquete \texttt{survey} permite realizar este procedimiento de forma directa mediante la función \texttt{svycontrast}, la cual calcula tanto el contraste lineal como su error estándar asociado. Para obtener la diferencia entre las proporciones de mujeres y hombres en condición de pobreza, se utiliza el siguiente código:

\begin{Shaded}
\begin{Highlighting}[]
\FunctionTok{svycontrast}\NormalTok{(}
  \AttributeTok{stat =}\NormalTok{ sex\_poverty\_contrast,}
  \AttributeTok{contrasts =} \FunctionTok{list}\NormalTok{(}\AttributeTok{sex\_difference =} \FunctionTok{c}\NormalTok{(}\DecValTok{1}\NormalTok{, }\SpecialCharTok{{-}}\DecValTok{1}\NormalTok{))}
\NormalTok{) }\SpecialCharTok{\%\textgreater{}\%}
  \FunctionTok{data.frame}\NormalTok{()}
\end{Highlighting}
\end{Shaded}

\begin{verbatim}
##                contrast sex_difference
## sex_difference -0.00532         0.0224
\end{verbatim}

De lo que se concluye que, la diferencia entre las proporciones de mujeres y hombres en estado de pobreza es -0.005 (-0.5\%) con una error estándar de 0.022, indicando que la proporción estimada de mujeres en condición de pobreza es ligeramente menor que la observada en los hombres.

Otro ejercicio que puede desarrollarse a partir de una encuesta de hogares consiste en estimar la proporción de personas desempleadas según la región de residencia. Para ello, estima las proporciones de desempleo para cada una de las regiones consideradas en el análisis. Los resultados obtenidos se presentan en la Tabla \ref{tab:tab-svyby-desempleo-region}:

\begin{Shaded}
\begin{Highlighting}[]
\NormalTok{region\_unemployment\_proportion }\OtherTok{\textless{}{-}}\NormalTok{ survey\_design }\SpecialCharTok{\%\textgreater{}\%}
  \FunctionTok{filter}\NormalTok{(}\SpecialCharTok{!}\FunctionTok{is.na}\NormalTok{(unemployed)) }\SpecialCharTok{\%\textgreater{}\%}
  \FunctionTok{group\_by}\NormalTok{(Region) }\SpecialCharTok{\%\textgreater{}\%}
  \FunctionTok{summarise}\NormalTok{(}\AttributeTok{proportion =} \FunctionTok{survey\_mean}\NormalTok{(unemployed,}
                                     \AttributeTok{na.rm =} \ConstantTok{TRUE}\NormalTok{,}
                                     \AttributeTok{vartype =} \FunctionTok{c}\NormalTok{(}\StringTok{"se"}\NormalTok{, }\StringTok{"ci"}\NormalTok{)))}

\NormalTok{region\_unemployment\_contrast }\OtherTok{\textless{}{-}} \FunctionTok{svyby}\NormalTok{(}
  \AttributeTok{formula =} \SpecialCharTok{\textasciitilde{}}\NormalTok{unemployed,}
  \AttributeTok{by =} \SpecialCharTok{\textasciitilde{}}\NormalTok{Region,}
  \AttributeTok{design =}\NormalTok{ survey\_design }\SpecialCharTok{\%\textgreater{}\%}
    \FunctionTok{filter}\NormalTok{(}\SpecialCharTok{!}\FunctionTok{is.na}\NormalTok{(unemployed)),}
  \AttributeTok{FUN =}\NormalTok{ svymean,}
  \AttributeTok{na.rm =} \ConstantTok{TRUE}\NormalTok{,}
  \AttributeTok{covmat =} \ConstantTok{TRUE}\NormalTok{,}
  \AttributeTok{vartype =} \FunctionTok{c}\NormalTok{(}\StringTok{"se"}\NormalTok{, }\StringTok{"ci"}\NormalTok{)}
\NormalTok{)}
\end{Highlighting}
\end{Shaded}

\begin{longtable}[]{@{}
  >{\raggedright\arraybackslash}p{(\linewidth - 8\tabcolsep) * \real{0.1449}}
  >{\centering\arraybackslash}p{(\linewidth - 8\tabcolsep) * \real{0.1739}}
  >{\centering\arraybackslash}p{(\linewidth - 8\tabcolsep) * \real{0.2174}}
  >{\centering\arraybackslash}p{(\linewidth - 8\tabcolsep) * \real{0.2319}}
  >{\centering\arraybackslash}p{(\linewidth - 8\tabcolsep) * \real{0.2319}}@{}}
\caption{\label{tab:tab-svyby-desempleo-region}Proporción en condición de desempleo por región}\tabularnewline
\toprule\noalign{}
\begin{minipage}[b]{\linewidth}\raggedright
Region
\end{minipage} & \begin{minipage}[b]{\linewidth}\centering
proportion
\end{minipage} & \begin{minipage}[b]{\linewidth}\centering
proportion\_se
\end{minipage} & \begin{minipage}[b]{\linewidth}\centering
proportion\_low
\end{minipage} & \begin{minipage}[b]{\linewidth}\centering
proportion\_upp
\end{minipage} \\
\midrule\noalign{}
\endfirsthead
\toprule\noalign{}
\begin{minipage}[b]{\linewidth}\raggedright
Region
\end{minipage} & \begin{minipage}[b]{\linewidth}\centering
proportion
\end{minipage} & \begin{minipage}[b]{\linewidth}\centering
proportion\_se
\end{minipage} & \begin{minipage}[b]{\linewidth}\centering
proportion\_low
\end{minipage} & \begin{minipage}[b]{\linewidth}\centering
proportion\_upp
\end{minipage} \\
\midrule\noalign{}
\endhead
\bottomrule\noalign{}
\endlastfoot
Norte & 0.049 & 0.020 & 0.009 & 0.088 \\
Sur & 0.066 & 0.024 & 0.019 & 0.113 \\
Centro & 0.039 & 0.012 & 0.014 & 0.063 \\
Occidente & 0.040 & 0.012 & 0.016 & 0.064 \\
Oriente & 0.030 & 0.013 & 0.005 & 0.054 \\
\end{longtable}

Una vez estimadas las proporciones de desempleo por región, el siguiente paso consiste en evaluar si existen diferencias estadísticamente significativas entre algunas de estas proporciones mediante contrastes lineales. En particular, el interés se centra en comparar las diferencias entre las proporciones estimadas de desempleo de distintas regiones. Los contrastes considerados corresponden a \(\hat{p}_{Norte} - \hat{p}_{Centro} = 0.01004\), \(\hat{p}_{Sur} - \hat{p}_{Centro} = 0.02691\) y \(\hat{p}_{Occidente} - \hat{p}_{Oriente} = 0.01046\). De manera matricial, estos contrastes entre regiones pueden escribirse así:

\[
\left[\begin{array}{ccccc}
1 & 0 & -1 & 0 & 0\\
0 & 1 & -1 & 0 & 0\\
0 & 0 & 0 & 1 & -1
\end{array}\right]
\]

Para calcular los errores estándar asociados a cada uno de los contrastes anteriores, es necesario considerar no solo las varianzas de las estimaciones regionales, sino también las covarianzas entre ellas. Esta información se resume en la matriz de varianzas y covarianzas de las proporciones estimadas de desempleo por región, presentada a continuación:

\begin{Shaded}
\begin{Highlighting}[]
\FunctionTok{vcov}\NormalTok{(region\_unemployment\_contrast) }\SpecialCharTok{\%\textgreater{}\%}
  \FunctionTok{data.frame}\NormalTok{()}
\end{Highlighting}
\end{Shaded}

\begin{longtable}[]{@{}llcccc@{}}
\caption{\label{tab:tab-vcov-desempleo-region}Matriz de varianzas y covarianzas de la proporción de desempleo por región}\tabularnewline
\toprule\noalign{}
& Norte & Sur & Centro & Occidente & Oriente \\
\midrule\noalign{}
\endfirsthead
\toprule\noalign{}
& Norte & Sur & Centro & Occidente & Oriente \\
\midrule\noalign{}
\endhead
\bottomrule\noalign{}
\endlastfoot
Norte & 0 & 0.000 & 0 & 0 & 0 \\
Sur & 0 & 0.001 & 0 & 0 & 0 \\
Centro & 0 & 0.000 & 0 & 0 & 0 \\
Occidente & 0 & 0.000 & 0 & 0 & 0 \\
Oriente & 0 & 0.000 & 0 & 0 & 0 \\
\end{longtable}

Nótese que las covarianzas entre las estimaciones regionales son iguales a cero debido a que las muestras seleccionadas en cada región son independientes entre sí. Por tanto, el error estándar de cada contraste se obtiene aplicando la expresión de la varianza para la diferencia entre dos estimadores. En este caso, dado que, el cálculo se reduce a la raíz cuadrada de la suma de las varianzas correspondientes:

\begin{Shaded}
\begin{Highlighting}[]
\FunctionTok{sqrt}\NormalTok{(}\FloatTok{0.0002981} \SpecialCharTok{+} \FloatTok{0.0002884} \SpecialCharTok{{-}} \DecValTok{2} \SpecialCharTok{*} \DecValTok{0}\NormalTok{)}
\end{Highlighting}
\end{Shaded}

\begin{verbatim}
## [1] 0.0242
\end{verbatim}

\begin{Shaded}
\begin{Highlighting}[]
\FunctionTok{sqrt}\NormalTok{(}\FloatTok{0.0001968} \SpecialCharTok{+} \FloatTok{0.0002884} \SpecialCharTok{{-}} \DecValTok{2} \SpecialCharTok{*} \DecValTok{0}\NormalTok{)}
\end{Highlighting}
\end{Shaded}

\begin{verbatim}
## [1] 0.022
\end{verbatim}

\begin{Shaded}
\begin{Highlighting}[]
\FunctionTok{sqrt}\NormalTok{(}\FloatTok{0.0001267} \SpecialCharTok{+} \FloatTok{0.0004093} \SpecialCharTok{{-}} \DecValTok{2} \SpecialCharTok{*} \DecValTok{0}\NormalTok{)}
\end{Highlighting}
\end{Shaded}

\begin{verbatim}
## [1] 0.0232
\end{verbatim}

Alternativamente, el paquete \texttt{survey} permite calcular estos contrastes de manera directa mediante la función \texttt{svycontrast}, la cual obtiene tanto las diferencias estimadas entre proporciones como sus respectivos errores estándar. En este caso, la estimación de los contrastes definidos anteriormente se realiza de la siguiente manera:

\begin{Shaded}
\begin{Highlighting}[]
\FunctionTok{svycontrast}\NormalTok{(}
  \AttributeTok{stat =}\NormalTok{ region\_unemployment\_contrast,}
  \AttributeTok{contrasts =} \FunctionTok{list}\NormalTok{(}
    \AttributeTok{north\_south =} \FunctionTok{c}\NormalTok{(}\DecValTok{1}\NormalTok{, }\DecValTok{0}\NormalTok{, }\SpecialCharTok{{-}}\DecValTok{1}\NormalTok{, }\DecValTok{0}\NormalTok{, }\DecValTok{0}\NormalTok{),}
    \AttributeTok{south\_center =} \FunctionTok{c}\NormalTok{(}\DecValTok{0}\NormalTok{, }\DecValTok{1}\NormalTok{, }\SpecialCharTok{{-}}\DecValTok{1}\NormalTok{, }\DecValTok{0}\NormalTok{, }\DecValTok{0}\NormalTok{),}
    \AttributeTok{west\_east =} \FunctionTok{c}\NormalTok{(}\DecValTok{0}\NormalTok{, }\DecValTok{0}\NormalTok{, }\DecValTok{0}\NormalTok{, }\DecValTok{1}\NormalTok{, }\SpecialCharTok{{-}}\DecValTok{1}\NormalTok{)}
\NormalTok{  )}
\NormalTok{) }\SpecialCharTok{\%\textgreater{}\%}
  \FunctionTok{data.frame}\NormalTok{()}
\end{Highlighting}
\end{Shaded}

\begin{verbatim}
##              contrast     SE
## north_south    0.0100 0.0236
## south_center   0.0269 0.0268
## west_east      0.0105 0.0176
\end{verbatim}

\section{Tablas cruzadas}\label{tablas-cruzadas}

Además de analizar cada variable de manera individual, resulta especialmente relevante estudiar si existe asociación entre dos variables categóricas. Este tipo de análisis permite identificar patrones y relaciones que aportan información valiosa para la toma de decisiones y la comprensión de fenómenos sociales. Por ejemplo, en el ámbito de las políticas públicas es posible relacionar educación y empleo para diseñar estrategias orientadas al mercado laboral; en la evaluación de programas sociales, analizar diferencias en el acceso a servicios de salud según el nivel de ingresos; y en investigación social, estudiar cómo variables demográficas y condiciones de vida se relacionan entre sí para comprender dinámicas y tendencias de la población.

Formalmente, sean \(x\) y \(y\) dos variables categóricas con \(R\) categorías para las filas y \(C\) categorías para las columnas, respectivamente. En el contexto de encuestas por muestreo, el interés suele centrarse en estimar la distribución conjunta de ambas variables, es decir, la proporción de unidades poblacionales que pertenecen simultáneamente a cada combinación posible de categorías. Para ello, las entradas de una tabla cruzada, también conocida como tabla de contingencia, se estiman mediante frecuencias ponderadas obtenidas a partir de los factores de expansión y del diseño muestral. La estimación del total poblacional asociado a cada celda \((r,c)\) se define como:

\[
\hat{N}_{rc} = \sum_{h} \sum_{i} \sum_{k} w_{hik} \ I(x_{hik}=r,\ y_{hik}=c),
\]

donde \(w_{hik}\) representa el factor de expansión asociado al elemento \(k\) de la UPM \(i\) perteneciente al estrato \(h\), mientras que \(I(x_{hik}=r,\ y_{hik}=c)\) corresponde a una función indicadora que toma el valor de uno cuando la unidad observada pertenece simultáneamente a la categoría \(r\) de la variable \(x\) y a la categoría \(c\) de la variable \(y\), y toma el valor de cero en cualquier otro caso.

Además de las frecuencias conjuntas, también es posible estimar los tamaños marginales por filas y columnas, definidos respectivamente como \(\hat{N}_{r+} = \sum_{c=1}^{C} \hat{N}_{rc}\) y \(\hat{N}_{+c} = \sum_{r=1}^{R} \hat{N}_{rc}\). Mientras que el total general se expresa como \(\hat{N}_{++} = \sum_{r=1}^{R}\sum_{c=1}^{C} \hat{N}_{rc}\). La estructura general de una tabla de contingencia con \(R\) filas y \(C\) columnas puede representarse de la siguiente forma:

{\def\LTcaptype{none} % do not increment counter
\begin{longtable}[]{@{}
  >{\raggedright\arraybackslash}p{(\linewidth - 10\tabcolsep) * \real{0.2135}}
  >{\raggedright\arraybackslash}p{(\linewidth - 10\tabcolsep) * \real{0.1461}}
  >{\raggedright\arraybackslash}p{(\linewidth - 10\tabcolsep) * \real{0.1685}}
  >{\raggedright\arraybackslash}p{(\linewidth - 10\tabcolsep) * \real{0.1236}}
  >{\raggedright\arraybackslash}p{(\linewidth - 10\tabcolsep) * \real{0.1685}}
  >{\raggedright\arraybackslash}p{(\linewidth - 10\tabcolsep) * \real{0.1798}}@{}}
\toprule\noalign{}
\begin{minipage}[b]{\linewidth}\raggedright
Variable 2
\end{minipage} & \begin{minipage}[b]{\linewidth}\raggedright
Variable 1
\end{minipage} & \begin{minipage}[b]{\linewidth}\raggedright
\end{minipage} & \begin{minipage}[b]{\linewidth}\raggedright
\(\cdots\)
\end{minipage} & \begin{minipage}[b]{\linewidth}\raggedright
\end{minipage} & \begin{minipage}[b]{\linewidth}\raggedright
Marginal fila
\end{minipage} \\
\midrule\noalign{}
\endhead
\bottomrule\noalign{}
\endlastfoot
& \(1\) & \(2\) & \(\cdots\) & \(C\) & \\
\(1\) & \(\hat{N}_{11}\) & \(\hat{N}_{12}\) & \(\cdots\) & \(\hat{N}_{1C}\) & \(\hat{N}_{1+}\) \\
\(2\) & \(\hat{N}_{21}\) & \(\hat{N}_{22}\) & \(\cdots\) & \(\hat{N}_{2C}\) & \(\hat{N}_{2+}\) \\
\(\vdots\) & \(\vdots\) & \(\vdots\) & \(\ddots\) & \(\vdots\) & \(\vdots\) \\
\(R\) & \(\hat{N}_{R1}\) & \(\hat{N}_{R2}\) & \(\cdots\) & \(\hat{N}_{RC}\) & \(\hat{N}_{R+}\) \\
Marginal columna & \(\hat{N}_{+1}\) & \(\hat{N}_{+2}\) & \(\cdots\) & \(\hat{N}_{+C}\) & \(\hat{N}_{++}\) \\
\end{longtable}
}

De manera tradicional, las tablas de contingencia suelen representarse como arreglos bidimensionales de dimensión \(R \times C\). Sin embargo, también pueden extenderse para incorporar una o más variables adicionales, generando subconjuntos o subtablas que permiten estudiar relaciones más complejas entre variables categóricas. A partir de los tamaños estimados, también pueden estimarse proporciones para cada celda de la tabla de contingencia, de la siguiente manera:

\[
\hat{p}_{rc}=\frac{\hat{N}_{rc}}{\hat{N}_{++}}
\]

A continuación, y siguiendo con la base de datos de ejemplo, se estima la distribución porcentual de hombres y mujeres según su condición de pobreza, incorporando además sus respectivos errores estándar e intervalos de confianza. Los resultados se presentan en la Tabla \ref{tab:tab-contingencia-pobreza-sexo}. Este tipo de análisis permite describir la composición de la población pobre según sexo y evaluar la precisión de las estimaciones obtenidas a partir del diseño muestral.

\begin{Shaded}
\begin{Highlighting}[]
\NormalTok{poverty\_design }\OtherTok{\textless{}{-}}\NormalTok{ survey\_design }\SpecialCharTok{\%\textgreater{}\%}
  \FunctionTok{mutate}\NormalTok{(}\AttributeTok{poor =} \FunctionTok{factor}\NormalTok{(}
\NormalTok{    poor,}
    \AttributeTok{levels =} \DecValTok{0}\SpecialCharTok{:}\DecValTok{1}\NormalTok{,}
    \AttributeTok{labels =} \FunctionTok{c}\NormalTok{(}\StringTok{"Not poor"}\NormalTok{, }\StringTok{"Poor"}\NormalTok{)}
\NormalTok{  ))}
\end{Highlighting}
\end{Shaded}

\begin{Shaded}
\begin{Highlighting}[]
\NormalTok{poverty\_design }\SpecialCharTok{\%\textgreater{}\%}
  \FunctionTok{group\_by}\NormalTok{(poor, Sex) }\SpecialCharTok{\%\textgreater{}\%}
  \FunctionTok{summarise}\NormalTok{(}\AttributeTok{proportion =} \FunctionTok{survey\_prop}\NormalTok{(}\AttributeTok{vartype =} \FunctionTok{c}\NormalTok{(}\StringTok{"se"}\NormalTok{, }\StringTok{"ci"}\NormalTok{)))}
\end{Highlighting}
\end{Shaded}

\begin{longtable}[]{@{}
  >{\raggedright\arraybackslash}p{(\linewidth - 10\tabcolsep) * \real{0.1184}}
  >{\centering\arraybackslash}p{(\linewidth - 10\tabcolsep) * \real{0.1053}}
  >{\centering\arraybackslash}p{(\linewidth - 10\tabcolsep) * \real{0.1579}}
  >{\centering\arraybackslash}p{(\linewidth - 10\tabcolsep) * \real{0.1974}}
  >{\centering\arraybackslash}p{(\linewidth - 10\tabcolsep) * \real{0.2105}}
  >{\centering\arraybackslash}p{(\linewidth - 10\tabcolsep) * \real{0.2105}}@{}}
\caption{\label{tab:tab-contingencia-pobreza-sexo}Proporción de hombres y mujeres en condición de pobreza y no pobreza}\tabularnewline
\toprule\noalign{}
\begin{minipage}[b]{\linewidth}\raggedright
poor
\end{minipage} & \begin{minipage}[b]{\linewidth}\centering
Sex
\end{minipage} & \begin{minipage}[b]{\linewidth}\centering
proportion
\end{minipage} & \begin{minipage}[b]{\linewidth}\centering
proportion\_se
\end{minipage} & \begin{minipage}[b]{\linewidth}\centering
proportion\_low
\end{minipage} & \begin{minipage}[b]{\linewidth}\centering
proportion\_upp
\end{minipage} \\
\midrule\noalign{}
\endfirsthead
\toprule\noalign{}
\begin{minipage}[b]{\linewidth}\raggedright
poor
\end{minipage} & \begin{minipage}[b]{\linewidth}\centering
Sex
\end{minipage} & \begin{minipage}[b]{\linewidth}\centering
proportion
\end{minipage} & \begin{minipage}[b]{\linewidth}\centering
proportion\_se
\end{minipage} & \begin{minipage}[b]{\linewidth}\centering
proportion\_low
\end{minipage} & \begin{minipage}[b]{\linewidth}\centering
proportion\_upp
\end{minipage} \\
\midrule\noalign{}
\endhead
\bottomrule\noalign{}
\endlastfoot
Not poor & Female & 0.529 & 0.012 & 0.505 & 0.554 \\
Not poor & Male & 0.471 & 0.012 & 0.446 & 0.495 \\
Poor & Female & 0.524 & 0.016 & 0.492 & 0.555 \\
Poor & Male & 0.476 & 0.016 & 0.445 & 0.508 \\
\end{longtable}

Los resultados indican que el 52.3\% de la población en condición de pobreza corresponde a mujeres, mientras que el 47.6\% corresponde a hombres. Para las mujeres, el intervalo de confianza al 95\% se encuentra entre 49.2\% y 55.5\%, mientras que para los hombres el intervalo correspondiente se ubica entre 44.5\% y 50.7\%. Estas estimaciones permiten observar la distribución relativa de la pobreza según sexo, considerando además la incertidumbre inherente al proceso de muestreo.

Con el mismo enfoque de \texttt{srvyr}, tambien puede estimarse la tabla de contingencia mediante \texttt{group\_by()} y \texttt{summarise()}. En el siguiente ejemplo se estima la distribucion de la condicion de pobreza segun sexo, obteniendo ademas medidas de precision como errores estandar e intervalos de confianza. El codigo correspondiente y sus resultados se presentan en la Tabla \ref{tab:tab-svyby-pobreza-sexo}.

\begin{Shaded}
\begin{Highlighting}[]
\NormalTok{sex\_by\_poverty\_proportion }\OtherTok{\textless{}{-}}\NormalTok{ poverty\_design }\SpecialCharTok{\%\textgreater{}\%}
  \FunctionTok{group\_by}\NormalTok{(poor, Sex) }\SpecialCharTok{\%\textgreater{}\%}
  \FunctionTok{summarise}\NormalTok{(}\AttributeTok{proportion =} \FunctionTok{survey\_prop}\NormalTok{(}\AttributeTok{vartype =} \FunctionTok{c}\NormalTok{(}\StringTok{"se"}\NormalTok{, }\StringTok{"ci"}\NormalTok{)))}
\end{Highlighting}
\end{Shaded}

\begin{longtable}[]{@{}
  >{\raggedright\arraybackslash}p{(\linewidth - 10\tabcolsep) * \real{0.1184}}
  >{\centering\arraybackslash}p{(\linewidth - 10\tabcolsep) * \real{0.1053}}
  >{\centering\arraybackslash}p{(\linewidth - 10\tabcolsep) * \real{0.1579}}
  >{\centering\arraybackslash}p{(\linewidth - 10\tabcolsep) * \real{0.1974}}
  >{\centering\arraybackslash}p{(\linewidth - 10\tabcolsep) * \real{0.2105}}
  >{\centering\arraybackslash}p{(\linewidth - 10\tabcolsep) * \real{0.2105}}@{}}
\caption{\label{tab:tab-svyby-pobreza-sexo}Proporción de hombres y mujeres por condición de pobreza}\tabularnewline
\toprule\noalign{}
\begin{minipage}[b]{\linewidth}\raggedright
poor
\end{minipage} & \begin{minipage}[b]{\linewidth}\centering
Sex
\end{minipage} & \begin{minipage}[b]{\linewidth}\centering
proportion
\end{minipage} & \begin{minipage}[b]{\linewidth}\centering
proportion\_se
\end{minipage} & \begin{minipage}[b]{\linewidth}\centering
proportion\_low
\end{minipage} & \begin{minipage}[b]{\linewidth}\centering
proportion\_upp
\end{minipage} \\
\midrule\noalign{}
\endfirsthead
\toprule\noalign{}
\begin{minipage}[b]{\linewidth}\raggedright
poor
\end{minipage} & \begin{minipage}[b]{\linewidth}\centering
Sex
\end{minipage} & \begin{minipage}[b]{\linewidth}\centering
proportion
\end{minipage} & \begin{minipage}[b]{\linewidth}\centering
proportion\_se
\end{minipage} & \begin{minipage}[b]{\linewidth}\centering
proportion\_low
\end{minipage} & \begin{minipage}[b]{\linewidth}\centering
proportion\_upp
\end{minipage} \\
\midrule\noalign{}
\endhead
\bottomrule\noalign{}
\endlastfoot
Not poor & Female & 0.529 & 0.012 & 0.505 & 0.554 \\
Not poor & Male & 0.471 & 0.012 & 0.446 & 0.495 \\
Poor & Female & 0.524 & 0.016 & 0.492 & 0.555 \\
Poor & Male & 0.476 & 0.016 & 0.445 & 0.508 \\
\end{longtable}

Adicionalmente, los intervalos de confianza pueden obtenerse mediante la función \texttt{confint()}, la cual permite calcular automáticamente los límites inferior y superior asociados a las estimaciones realizadas. El procedimiento para la estimación de los intervalos de confianza se presenta en la Tabla \ref{tab:tab-confint-sexo-pobreza}. Nótese que lso intervalos coinciden con los generados anteriormente usando la función \texttt{group\_by}.

\begin{Shaded}
\begin{Highlighting}[]
\NormalTok{sex\_by\_poverty\_proportion }\SpecialCharTok{\%\textgreater{}\%}
\NormalTok{  dplyr}\SpecialCharTok{::}\FunctionTok{select}\NormalTok{(poor, Sex, proportion\_low, proportion\_upp)}
\end{Highlighting}
\end{Shaded}

\begin{longtable}[]{@{}lccc@{}}
\caption{\label{tab:tab-confint-sexo-pobreza}Intervalos de confianza para la proporción por sexo y pobreza}\tabularnewline
\toprule\noalign{}
poor & Sex & proportion\_low & proportion\_upp \\
\midrule\noalign{}
\endfirsthead
\toprule\noalign{}
poor & Sex & proportion\_low & proportion\_upp \\
\midrule\noalign{}
\endhead
\bottomrule\noalign{}
\endlastfoot
Not poor & Female & 0.505 & 0.554 \\
Not poor & Male & 0.446 & 0.495 \\
Poor & Female & 0.492 & 0.555 \\
Poor & Male & 0.445 & 0.508 \\
\end{longtable}

Otro análisis de interés relacionado con tablas de doble entrada en encuestas de hogares consiste en estimar el porcentaje de personas desempleadas según sexo. El procedimiento correspondiente y sus resultados se presentan en la Tabla \ref{tab:tab-svyby-empleo-sexo}.

\begin{Shaded}
\begin{Highlighting}[]
\NormalTok{sex\_by\_employment\_proportion }\OtherTok{\textless{}{-}}\NormalTok{ survey\_design }\SpecialCharTok{\%\textgreater{}\%}
  \FunctionTok{filter}\NormalTok{(}\SpecialCharTok{!}\FunctionTok{is.na}\NormalTok{(Employment)) }\SpecialCharTok{\%\textgreater{}\%}
  \FunctionTok{group\_by}\NormalTok{(Employment, Sex) }\SpecialCharTok{\%\textgreater{}\%}
  \FunctionTok{summarise}\NormalTok{(}\AttributeTok{proportion =} \FunctionTok{survey\_prop}\NormalTok{(}\AttributeTok{vartype =} \FunctionTok{c}\NormalTok{(}\StringTok{"se"}\NormalTok{, }\StringTok{"ci"}\NormalTok{)))}
\end{Highlighting}
\end{Shaded}

\begin{longtable}[]{@{}
  >{\raggedright\arraybackslash}p{(\linewidth - 10\tabcolsep) * \real{0.1410}}
  >{\centering\arraybackslash}p{(\linewidth - 10\tabcolsep) * \real{0.1026}}
  >{\centering\arraybackslash}p{(\linewidth - 10\tabcolsep) * \real{0.1538}}
  >{\centering\arraybackslash}p{(\linewidth - 10\tabcolsep) * \real{0.1923}}
  >{\centering\arraybackslash}p{(\linewidth - 10\tabcolsep) * \real{0.2051}}
  >{\centering\arraybackslash}p{(\linewidth - 10\tabcolsep) * \real{0.2051}}@{}}
\caption{\label{tab:tab-svyby-empleo-sexo}Proporción de hombres y mujeres por estado de ocupación}\tabularnewline
\toprule\noalign{}
\begin{minipage}[b]{\linewidth}\raggedright
Employment
\end{minipage} & \begin{minipage}[b]{\linewidth}\centering
Sex
\end{minipage} & \begin{minipage}[b]{\linewidth}\centering
proportion
\end{minipage} & \begin{minipage}[b]{\linewidth}\centering
proportion\_se
\end{minipage} & \begin{minipage}[b]{\linewidth}\centering
proportion\_low
\end{minipage} & \begin{minipage}[b]{\linewidth}\centering
proportion\_upp
\end{minipage} \\
\midrule\noalign{}
\endfirsthead
\toprule\noalign{}
\begin{minipage}[b]{\linewidth}\raggedright
Employment
\end{minipage} & \begin{minipage}[b]{\linewidth}\centering
Sex
\end{minipage} & \begin{minipage}[b]{\linewidth}\centering
proportion
\end{minipage} & \begin{minipage}[b]{\linewidth}\centering
proportion\_se
\end{minipage} & \begin{minipage}[b]{\linewidth}\centering
proportion\_low
\end{minipage} & \begin{minipage}[b]{\linewidth}\centering
proportion\_upp
\end{minipage} \\
\midrule\noalign{}
\endhead
\bottomrule\noalign{}
\endlastfoot
Unemployed & Female & 0.273 & 0.054 & 0.180 & 0.390 \\
Unemployed & Male & 0.727 & 0.054 & 0.610 & 0.820 \\
Inactive & Female & 0.770 & 0.023 & 0.721 & 0.813 \\
Inactive & Male & 0.230 & 0.023 & 0.187 & 0.279 \\
Employed & Female & 0.405 & 0.019 & 0.369 & 0.442 \\
Employed & Male & 0.595 & 0.019 & 0.558 & 0.631 \\
\end{longtable}

De la anterior salida se puede observar que, el 27.2\% de las mujeres y el 72.7\% de los hombres están desempleados con errores estándares para estas estimaciones de 5.3\% para mujeres y hombres. cuyos intervalos de confianza se calculan a continuación, presentada en la Tabla \ref{tab:tab-confint-empleo-sexo}:

\begin{Shaded}
\begin{Highlighting}[]
\NormalTok{sex\_by\_employment\_proportion }\SpecialCharTok{\%\textgreater{}\%}
\NormalTok{  dplyr}\SpecialCharTok{::}\FunctionTok{select}\NormalTok{(Employment, Sex, proportion\_low, proportion\_upp)}
\end{Highlighting}
\end{Shaded}

\begin{longtable}[]{@{}lccc@{}}
\caption{\label{tab:tab-confint-empleo-sexo}Intervalos de confianza para la proporción por sexo y ocupación}\tabularnewline
\toprule\noalign{}
Employment & Sex & proportion\_low & proportion\_upp \\
\midrule\noalign{}
\endfirsthead
\toprule\noalign{}
Employment & Sex & proportion\_low & proportion\_upp \\
\midrule\noalign{}
\endhead
\bottomrule\noalign{}
\endlastfoot
Unemployed & Female & 0.180 & 0.390 \\
Unemployed & Male & 0.610 & 0.820 \\
Inactive & Female & 0.721 & 0.813 \\
Inactive & Male & 0.187 & 0.279 \\
Employed & Female & 0.369 & 0.442 \\
Employed & Male & 0.558 & 0.631 \\
\end{longtable}

Si ahora el objetivo es estimar la pobreza por region, se utiliza la variable numerica \texttt{poor} dentro de un flujo \texttt{srvyr}, como se muestra en la Tabla \ref{tab:tab-svyby-region-pobreza}.

\begin{Shaded}
\begin{Highlighting}[]
\NormalTok{region\_poverty\_proportion }\OtherTok{\textless{}{-}}\NormalTok{ survey\_design }\SpecialCharTok{\%\textgreater{}\%}
  \FunctionTok{group\_by}\NormalTok{(Region, poor) }\SpecialCharTok{\%\textgreater{}\%}
  \FunctionTok{summarise}\NormalTok{(}\AttributeTok{proportion =} \FunctionTok{survey\_prop}\NormalTok{(}\AttributeTok{vartype =} \FunctionTok{c}\NormalTok{(}\StringTok{"se"}\NormalTok{, }\StringTok{"ci"}\NormalTok{)))}
\end{Highlighting}
\end{Shaded}

\begin{longtable}[]{@{}
  >{\raggedright\arraybackslash}p{(\linewidth - 10\tabcolsep) * \real{0.1333}}
  >{\centering\arraybackslash}p{(\linewidth - 10\tabcolsep) * \real{0.0800}}
  >{\centering\arraybackslash}p{(\linewidth - 10\tabcolsep) * \real{0.1600}}
  >{\centering\arraybackslash}p{(\linewidth - 10\tabcolsep) * \real{0.2000}}
  >{\centering\arraybackslash}p{(\linewidth - 10\tabcolsep) * \real{0.2133}}
  >{\centering\arraybackslash}p{(\linewidth - 10\tabcolsep) * \real{0.2133}}@{}}
\caption{\label{tab:tab-svyby-region-pobreza}Proporción en condición de pobreza por región}\tabularnewline
\toprule\noalign{}
\begin{minipage}[b]{\linewidth}\raggedright
Region
\end{minipage} & \begin{minipage}[b]{\linewidth}\centering
poor
\end{minipage} & \begin{minipage}[b]{\linewidth}\centering
proportion
\end{minipage} & \begin{minipage}[b]{\linewidth}\centering
proportion\_se
\end{minipage} & \begin{minipage}[b]{\linewidth}\centering
proportion\_low
\end{minipage} & \begin{minipage}[b]{\linewidth}\centering
proportion\_upp
\end{minipage} \\
\midrule\noalign{}
\endfirsthead
\toprule\noalign{}
\begin{minipage}[b]{\linewidth}\raggedright
Region
\end{minipage} & \begin{minipage}[b]{\linewidth}\centering
poor
\end{minipage} & \begin{minipage}[b]{\linewidth}\centering
proportion
\end{minipage} & \begin{minipage}[b]{\linewidth}\centering
proportion\_se
\end{minipage} & \begin{minipage}[b]{\linewidth}\centering
proportion\_low
\end{minipage} & \begin{minipage}[b]{\linewidth}\centering
proportion\_upp
\end{minipage} \\
\midrule\noalign{}
\endhead
\bottomrule\noalign{}
\endlastfoot
Norte & 0 & 0.641 & 0.055 & 0.526 & 0.742 \\
Norte & 1 & 0.359 & 0.055 & 0.258 & 0.474 \\
Sur & 0 & 0.656 & 0.043 & 0.566 & 0.737 \\
Sur & 1 & 0.344 & 0.043 & 0.263 & 0.434 \\
Centro & 0 & 0.635 & 0.079 & 0.470 & 0.773 \\
Centro & 1 & 0.365 & 0.079 & 0.227 & 0.530 \\
Occidente & 0 & 0.599 & 0.047 & 0.504 & 0.687 \\
Occidente & 1 & 0.401 & 0.047 & 0.313 & 0.496 \\
Oriente & 0 & 0.548 & 0.088 & 0.374 & 0.711 \\
Oriente & 1 & 0.452 & 0.088 & 0.289 & 0.626 \\
\end{longtable}

De lo anterior se puede concluir que, en la región Norte, el 35\% de las personas están en estado de pobreza mientras que en el sur es el 34\%. La pobreza más alta se tiene en la región oriente con un 45\% de pobres. Los errores estándares de las estimaciones.

\section{\texorpdfstring{La razón de \emph{odds}}{La razón de odds}}\label{la-razon-de-odds}

La razón de \emph{odds} es una medida ampliamente utilizada para estudiar la asociación entre dos variables categóricas, especialmente cuando el interés se centra en comparar la probabilidad de ocurrencia de un evento entre distintos grupos poblacionales. Su interpretación se basa en comparar las ventajas relativas (\emph{odds}) de ocurrencia de un evento entre dos categorías de análisis.

En este sentido, este parámetro permite evaluar cómo cambia dicha ventaja entre diferentes grupos de interés y constituye una herramienta fundamental en estudios sociales, económicos y de salud; además, esta medida también puede utilizarse para cuantificar la relación entre los niveles de una variable y un factor categórico, comparando las ventajas relativas de ocurrencia del evento en cada grupo \citep{heeringa2017applied}.

Por ejemplo, suponga que se desea estudiar la asociación entre el sexo de las personas y la condición de pobreza. En particular, interesa evaluar si la ventaja relativa de pertenecer al grupo de personas no pobres frente al grupo de personas pobres es diferente entre mujeres y hombres. Para ello, puede definirse la siguiente razón de \emph{odds}:

\[
\frac{
P(\text{Sex}=\text{Female} \mid \text{pobreza}=0)
\big/
P(\text{Sex}=\text{Female} \mid \text{pobreza}=1)
}{
P(\text{Sex}=\text{Male} \mid \text{pobreza}=1)
\big/
P(\text{Sex}=\text{Male} \mid \text{pobreza}=0)
}
\]

La expresión del numerador representa la ventaja relativa de observar mujeres en condición de no pobreza respecto a observar mujeres en condición de pobreza. De manera análoga, el denominador representa la ventaja relativa de observar hombres en condición de pobreza respecto a hombres en condición de no pobreza. Por tanto, la razón de \emph{odds} compara ambas ventajas relativas y permite cuantificar si la asociación entre sexo y pobreza favorece más a un grupo que a otro.

Un valor igual a uno indicaría ausencia de asociación entre las variables, es decir, que la relación entre sexo y condición de pobreza es similar para hombres y mujeres. Valores superiores a uno indicarían una mayor ventaja relativa para las mujeres en comparación con los hombres, mientras que valores inferiores a uno sugerirían una mayor ventaja relativa para los hombres. En el contexto de encuestas de hogares, este tipo de medidas resulta especialmente útil para estudiar desigualdades sociodemográficas y brechas en condiciones de bienestar entre distintos grupos poblacionales.

El procedimiento para realizarlo en \texttt{R} sería, primero estimar las proporciones de la tabla cruzada entre las variables sexo y pobreza, presentada en la Tabla \ref{tab:tab-interaccion-sexo-pobreza}:

\begin{Shaded}
\begin{Highlighting}[]
\NormalTok{sex\_poverty\_interaction\_proportion }\OtherTok{\textless{}{-}}\NormalTok{ survey\_design }\SpecialCharTok{\%\textgreater{}\%}
  \FunctionTok{group\_by}\NormalTok{(Sex, poor) }\SpecialCharTok{\%\textgreater{}\%}
  \FunctionTok{summarise}\NormalTok{(}\AttributeTok{proportion =} \FunctionTok{survey\_prop}\NormalTok{(}\AttributeTok{vartype =} \FunctionTok{c}\NormalTok{(}\StringTok{"se"}\NormalTok{, }\StringTok{"ci"}\NormalTok{)))}

\NormalTok{sex\_poverty\_odds\_contrast }\OtherTok{\textless{}{-}} \FunctionTok{svymean}\NormalTok{(}\AttributeTok{x =} \SpecialCharTok{\textasciitilde{}}\FunctionTok{interaction}\NormalTok{(Sex, poor),}
                                     \AttributeTok{design =}\NormalTok{ survey\_design,}
                                     \AttributeTok{se =} \ConstantTok{TRUE}\NormalTok{, }\AttributeTok{na.rm =} \ConstantTok{TRUE}\NormalTok{, }\AttributeTok{ci =} \ConstantTok{TRUE}\NormalTok{,}
                                     \AttributeTok{keep.vars =} \ConstantTok{TRUE}\NormalTok{)}
\end{Highlighting}
\end{Shaded}

\begin{longtable}[]{@{}
  >{\raggedright\arraybackslash}p{(\linewidth - 10\tabcolsep) * \real{0.0972}}
  >{\centering\arraybackslash}p{(\linewidth - 10\tabcolsep) * \real{0.0833}}
  >{\centering\arraybackslash}p{(\linewidth - 10\tabcolsep) * \real{0.1667}}
  >{\centering\arraybackslash}p{(\linewidth - 10\tabcolsep) * \real{0.2083}}
  >{\centering\arraybackslash}p{(\linewidth - 10\tabcolsep) * \real{0.2222}}
  >{\centering\arraybackslash}p{(\linewidth - 10\tabcolsep) * \real{0.2222}}@{}}
\caption{\label{tab:tab-interaccion-sexo-pobreza}Proporciones conjuntas de sexo y condición de pobreza}\tabularnewline
\toprule\noalign{}
\begin{minipage}[b]{\linewidth}\raggedright
Sex
\end{minipage} & \begin{minipage}[b]{\linewidth}\centering
poor
\end{minipage} & \begin{minipage}[b]{\linewidth}\centering
proportion
\end{minipage} & \begin{minipage}[b]{\linewidth}\centering
proportion\_se
\end{minipage} & \begin{minipage}[b]{\linewidth}\centering
proportion\_low
\end{minipage} & \begin{minipage}[b]{\linewidth}\centering
proportion\_upp
\end{minipage} \\
\midrule\noalign{}
\endfirsthead
\toprule\noalign{}
\begin{minipage}[b]{\linewidth}\raggedright
Sex
\end{minipage} & \begin{minipage}[b]{\linewidth}\centering
poor
\end{minipage} & \begin{minipage}[b]{\linewidth}\centering
proportion
\end{minipage} & \begin{minipage}[b]{\linewidth}\centering
proportion\_se
\end{minipage} & \begin{minipage}[b]{\linewidth}\centering
proportion\_low
\end{minipage} & \begin{minipage}[b]{\linewidth}\centering
proportion\_upp
\end{minipage} \\
\midrule\noalign{}
\endhead
\bottomrule\noalign{}
\endlastfoot
Female & 0 & 0.611 & 0.032 & 0.547 & 0.671 \\
Female & 1 & 0.389 & 0.032 & 0.329 & 0.453 \\
Male & 0 & 0.605 & 0.037 & 0.531 & 0.675 \\
Male & 1 & 0.395 & 0.037 & 0.325 & 0.469 \\
\end{longtable}

Luego, se realiza el contraste dividiendo cada uno de los elementos de la expresión mostrada anteriormente:

\begin{Shaded}
\begin{Highlighting}[]
\NormalTok{odds\_ratio }\OtherTok{\textless{}{-}} \FunctionTok{quote}\NormalTok{(}
\NormalTok{ (}\StringTok{\textasciigrave{}}\AttributeTok{interaction(Sex, poor)Female.0}\StringTok{\textasciigrave{}} \SpecialCharTok{/}
  \StringTok{\textasciigrave{}}\AttributeTok{interaction(Sex, poor)Female.1}\StringTok{\textasciigrave{}}\NormalTok{) }\SpecialCharTok{/}
\NormalTok{ (}\StringTok{\textasciigrave{}}\AttributeTok{interaction(Sex, poor)Male.0}\StringTok{\textasciigrave{}} \SpecialCharTok{/}
  \StringTok{\textasciigrave{}}\AttributeTok{interaction(Sex, poor)Male.1}\StringTok{\textasciigrave{}}\NormalTok{)}
\NormalTok{)}

\FunctionTok{svycontrast}\NormalTok{(}
  \AttributeTok{stat =}\NormalTok{ sex\_poverty\_odds\_contrast,}
  \AttributeTok{contrasts =}\NormalTok{ odds\_ratio}
\NormalTok{)}
\end{Highlighting}
\end{Shaded}

\begin{verbatim}
##          nlcon  SE
## contrast  1.02 0.1
\end{verbatim}

Obteniendo que el \emph{odds} estimado de que una mujer no se encuentre en situación de pobreza, en comparación con un hombre, es igual a 1.02. Esto significa que, sin considerar otras variables de la encuesta, las mujeres presentan una razón de probabilidades de no estar en pobreza aproximadamente 2\% mayor que la observada en los hombres. En otras palabras, la probabilidad relativa de no encontrarse en condición de pobreza es ligeramente superior para las mujeres respecto a los hombres.

\section{\texorpdfstring{Prueba de independencia \(\chi^{2}\)}{Prueba de independencia \textbackslash chi\^{}\{2\}}}\label{prueba-de-independencia-chi}

Como se introdujo en los capítulos anteriores, una prueba de hipótesis es un procedimiento estadístico utilizado para evaluar si la evidencia observada en una muestra es compatible con una afirmación acerca de la población. En el contexto de tablas de contingencia, las pruebas de independencia permiten analizar si existe asociación entre dos variables categóricas.

En este caso, la hipótesis nula (\(H_{0}\)) establece que ambas variables son independientes; es decir, que la distribución de una variable no depende de las categorías de la otra. Bajo este supuesto, las diferencias observadas en la muestra se atribuyen únicamente a la variabilidad muestral y no a una relación real entre las variables en la población. Matemáticamente, esta hipótesis puede expresarse como:

\[
H_0: p_{rc}^0 = p_{r+} \times p_{+c}, \quad \text{para todo } r = 1,\ldots,R \text{ y } c = 1,\ldots,C
\]

En donde \(p_{rc}^0\) representa la proporción esperada en la celda \((r,c)\) bajo el supuesto de independencia, mientras que \(P_{r+}\) y \(P_{+c}\) corresponden a las proporciones marginales de las filas y columnas, respectivamente. Bajo esta formulación, la hipótesis de independencia implica que la probabilidad conjunta de cada celda puede expresarse como el producto de sus probabilidades marginales.

En consecuencia, la prueba de independencia consiste en comparar las proporciones observadas o estimadas \(\hat{p}_{rc}\) con las proporciones esperadas \(p_{rc}^0\) bajo la hipótesis nula. Cuando las diferencias entre ambas son pequeñas, la evidencia empírica resulta consistente con el supuesto de independencia. Por el contrario, discrepancias suficientemente grandes sugieren la existencia de asociación entre las variables y conducen al rechazo de \(H_0\).

Como se ha venido mencionando, en las encuestas de hogares, características propias del diseño muestral complejo, como la estratificación, la selección por conglomerados y el uso de factores de expansión, afectan la variabilidad de los estimadores en comparación con la que se obtendría bajo un muestreo aleatorio simple. Como consecuencia, la prueba clásica de independencia \(\chi^2\) de Pearson no resulta adecuada para analizar tablas de contingencia provenientes de este tipo de diseños, ya que puede subestimar o sobreestimar la variabilidad real de las estimaciones.

Con el propósito de corregir este problema, \citet{fay1979contingency} y \citet{fellegi1980approximate} propusieron los primeros ajustes al estadístico chi-cuadrado de Pearson basados en el efecto de diseño generalizado (\(GDEFF\)). Posteriormente, \citet{rao1984chisquare} y \citet{thomas1987chisquare} ampliaron y formalizaron el marco teórico de estas correcciones, dando origen a lo que actualmente se conoce como la prueba de Rao-Scott, considerada el procedimiento de referencia para el análisis de datos categóricos obtenidos mediante encuestas complejas.

La idea central de este enfoque consiste en adaptar la prueba clásica de independencia mediante la incorporación del efecto de diseño generalizado, obteniendo así un estadístico robusto frente a las complejidades del diseño muestral:

\[
\chi_{RS}^2 = \frac{n_{++}}{GDEFF} \sum_r \sum_c \frac{(\hat{p}_{rc} - p_{rc}^0)^2}{p_{rc}^0}
\]

donde \(n_{++}\) representa el tamaño total de la muestra, \(\hat{p}_{rc}\) corresponde a la proporción estimada en la celda \((r,c)\) de la tabla de contingencia y \(p_{rc}^0\) denota la proporción esperada bajo la hipótesis nula de independencia, calculada a partir del producto de las proporciones marginales de filas y columnas.

Por su parte, el término \(GDEFF\) \citep{heeringa2017applied} representa el efecto de diseño generalizado y mide cuánto se incrementa o disminuye la variabilidad de las estimaciones como consecuencia del diseño complejo de la encuesta respecto a la variabilidad que se observaría bajo un muestreo aleatorio simple. De esta manera, el ajuste de Rao--Scott permite que las inferencias derivadas de la prueba de independencia reflejen adecuadamente la estructura del diseño muestral.

Bajo la hipótesis nula \(H_0\), el estadístico \(\chi_{RS}^2\) sigue aproximadamente una distribución \(\chi^2\) con \((R-1)(C-1)\) grados de libertad, donde \(R\) y \(C\) corresponden al número de categorías de las filas y columnas, respectivamente. En situaciones de muestras pequeñas o con pocos grados de libertad, es frecuente utilizar ajustes basados en la distribución \(F\), los cuales mejoran la precisión de la inferencia estadística. Actualmente, las pruebas de Rao-Scott constituyen el estándar para el análisis de asociación entre variables categóricas en encuestas complejas \citep{heeringa2017applied}.

El paquete \texttt{survey} en \texttt{R} implementa este tipo de pruebas mediante la función \texttt{svychisq()}, la cual incorpora automáticamente los ajustes de Rao-Scott para tener en cuenta las características del diseño muestral complejo. Esta función permite evaluar la existencia de asociación entre dos variables categóricas utilizando tablas de contingencia ponderadas y corrigiendo adecuadamente la variabilidad de los estimadores. A modo de ejemplo, para evaluar si la condición de pobreza es independiente del sexo, se ejecuta el siguiente código:

\begin{Shaded}
\begin{Highlighting}[]
\FunctionTok{svychisq}\NormalTok{(}\AttributeTok{formula =} \SpecialCharTok{\textasciitilde{}}\NormalTok{Sex }\SpecialCharTok{+}\NormalTok{ poor, }\AttributeTok{design =}\NormalTok{ survey\_design, }\AttributeTok{statistic =} \StringTok{"F"}\NormalTok{)}
\end{Highlighting}
\end{Shaded}

\begin{verbatim}
## 
##  Pearson's X^2: Rao & Scott adjustment
## 
## data:  NextMethod()
## F = 0.06, ndf = 1, ddf = 119, p-value = 0.8
\end{verbatim}

En este caso, el argumento \texttt{formula} especifica las dos variables categóricas cuya independencia se desea evaluar; \texttt{design} corresponde al diseño muestral previamente definido; y \texttt{statistic\ =\ "F"} indica que la prueba se realizará utilizando la versión ajustada basada en la distribución \(F\). Si el valor-\emph{p} es superior al 5\%, no se rechaza la hipótesis nula \(H_0\), concluyendo que no existe evidencia estadística suficiente para afirmar que pobreza y sexo están asociadas. Por el contrario, si el valor-\emph{p} es inferior al 5\%, se rechaza \(H_0\), lo que sugiere la existencia de dependencia o asociación entre ambas variables categóricas. De manera similar, se pueden evaluar otras relaciones, como desempleo y sexo o pobreza y región.

\begin{Shaded}
\begin{Highlighting}[]
\FunctionTok{svychisq}\NormalTok{(}\AttributeTok{formula =} \SpecialCharTok{\textasciitilde{}}\NormalTok{Sex }\SpecialCharTok{+}\NormalTok{ Employment, }\AttributeTok{design =}\NormalTok{ survey\_design, }\AttributeTok{statistic =} \StringTok{"F"}\NormalTok{)}
\end{Highlighting}
\end{Shaded}

\begin{verbatim}
## 
##  Pearson's X^2: Rao & Scott adjustment
## 
## data:  NextMethod()
## F = 62, ndf = 2, ddf = 201, p-value <0.0000000000000002
\end{verbatim}

En este primer caso, para el contraste de independencia entre las variables \texttt{Sex} y \texttt{Employment}, el estadístico obtenido fue \(F = 62.251\), con un valor-\(p\) menor que \(2.2\times10^{-16}\), lo que proporciona evidencia estadísticamente significativa para rechazar la hipótesis nula de independencia entre ambas variables. En consecuencia, los resultados sugieren que existe una asociación significativa entre el sexo y la condición de empleo en la población de estudio.

\begin{Shaded}
\begin{Highlighting}[]
\FunctionTok{svychisq}\NormalTok{(}\AttributeTok{formula =} \SpecialCharTok{\textasciitilde{}}\NormalTok{Region }\SpecialCharTok{+}\NormalTok{ poor, }\AttributeTok{design =}\NormalTok{ survey\_design, }\AttributeTok{statistic =} \StringTok{"F"}\NormalTok{)}
\end{Highlighting}
\end{Shaded}

\begin{verbatim}
## 
##  Pearson's X^2: Rao & Scott adjustment
## 
## data:  NextMethod()
## F = 0.5, ndf = 3, ddf = 358, p-value = 0.7
\end{verbatim}

Por otro lado, esta prueba de independencia entre las variables \texttt{Region} y \texttt{poor}, da como resultado \(F = 0.48794\), con un valor-\(p = 0.6914\). Dado que este valor-\(p\) es considerablemente mayor que niveles de significancia convencionales, no existe evidencia suficiente para rechazar la hipótesis nula de independencia. Por tanto, los resultados sugieren que, bajo el diseño muestral considerado, no se observa una asociación estadísticamente significativa entre la región y la condición de pobreza.

\section{Visualización de variables categóricas}\label{viz-categoricas}

La visualización de variables categóricas es fundamental en el análisis de encuestas de hogares, ya que permite comunicar de manera clara la distribución de grupos poblacionales y las comparaciones entre ellos. Al igual que en el análisis de variables continuas, los gráficos deben incorporar los pesos muestrales para que la representación visual corresponda a estimaciones válidas a nivel poblacional. Cuando la precisión estadística es relevante, es recomendable acompañar las barras con intervalos de confianza, que comunican tanto el valor central como la incertidumbre asociada al proceso de estimación.

En esta sección se utiliza el paquete \texttt{ggplot2} \citep{ggplot22026}.

\begin{Shaded}
\begin{Highlighting}[]
\FunctionTok{library}\NormalTok{(ggplot2)}
\end{Highlighting}
\end{Shaded}

\subsection{Gráficos de barras}\label{graficos-de-barras}

La construcción de gráficos de barras en \texttt{R} requiere, como paso previo, obtener las estimaciones que se desean representar visualmente. En el contexto de encuestas complejas, estas estimaciones pueden calcularse mediante las funciones \texttt{survey\_total()} o \texttt{survey\_prop()} del paquete \texttt{srvyr}, dependiendo de si el interés se centra en totales poblacionales o proporciones.

A continuación, se ilustra este procedimiento utilizando el tamaño poblacional estimado por zona de residencia. Los resultados gráficos correspondientes se presentan en la Figura \ref{fig:fig-barras-zona}:

\begin{Shaded}
\begin{Highlighting}[]
\NormalTok{zone\_size }\OtherTok{\textless{}{-}}\NormalTok{ survey\_design }\SpecialCharTok{\%\textgreater{}\%}
  \FunctionTok{group\_by}\NormalTok{(Zone) }\SpecialCharTok{\%\textgreater{}\%}
  \FunctionTok{summarise}\NormalTok{(}\AttributeTok{size =} \FunctionTok{survey\_total}\NormalTok{(}\AttributeTok{vartype =} \FunctionTok{c}\NormalTok{(}\StringTok{"se"}\NormalTok{, }\StringTok{"ci"}\NormalTok{)))}

\NormalTok{zone\_colors }\OtherTok{\textless{}{-}} \FunctionTok{c}\NormalTok{(}\AttributeTok{Urban =} \StringTok{"\#48C9B0"}\NormalTok{, }\AttributeTok{Rural =} \StringTok{"\#117864"}\NormalTok{)}

\FunctionTok{ggplot}\NormalTok{(}
  \AttributeTok{data =}\NormalTok{ zone\_size,}
  \FunctionTok{aes}\NormalTok{(}
    \AttributeTok{x =}\NormalTok{ Zone, }\AttributeTok{y =}\NormalTok{ size,}
    \AttributeTok{ymax =}\NormalTok{ size\_upp, }\AttributeTok{ymin =}\NormalTok{ size\_low,}
    \AttributeTok{fill =}\NormalTok{ Zone}
\NormalTok{  )}
\NormalTok{) }\SpecialCharTok{+}
  \FunctionTok{geom\_bar}\NormalTok{(}\AttributeTok{stat =} \StringTok{"identity"}\NormalTok{, }\AttributeTok{position =} \StringTok{"dodge"}\NormalTok{) }\SpecialCharTok{+}
  \FunctionTok{geom\_errorbar}\NormalTok{(}\AttributeTok{position =} \FunctionTok{position\_dodge}\NormalTok{(}\AttributeTok{width =} \FloatTok{0.9}\NormalTok{), }\AttributeTok{width =} \FloatTok{0.3}\NormalTok{) }\SpecialCharTok{+}
  \FunctionTok{scale\_fill\_manual}\NormalTok{(}\AttributeTok{values =}\NormalTok{ zone\_colors) }\SpecialCharTok{+}
  \FunctionTok{theme\_bw}\NormalTok{() }\SpecialCharTok{+}
  \FunctionTok{theme}\NormalTok{(}\AttributeTok{legend.position =} \StringTok{"top"}\NormalTok{)}
\end{Highlighting}
\end{Shaded}

\begin{figure}

{\centering \includegraphics[width=0.7\linewidth]{docs/figures/fig-barras-zona-1} 

}

\caption{Gráfico de barras del tamaño poblacional por zona geográfica}\label{fig:fig-barras-zona}
\end{figure}

El procedimiento anterior también puede aplicarse a variables con más de dos categorías. En estos casos, el gráfico de barras permite comparar de manera sencilla las estimaciones asociadas a cada categoría de la variable de interés. Por ejemplo, a continuación se presenta la distribución estimada de la población según la condición de pobreza, cuyos resultados gráficos se muestran en la Figura \ref{fig:fig-barras-pobreza}.

\begin{Shaded}
\begin{Highlighting}[]
\NormalTok{poverty\_size }\OtherTok{\textless{}{-}}\NormalTok{ survey\_design }\SpecialCharTok{\%\textgreater{}\%}
  \FunctionTok{group\_by}\NormalTok{(Poverty) }\SpecialCharTok{\%\textgreater{}\%}
  \FunctionTok{summarise}\NormalTok{(}\AttributeTok{size =} \FunctionTok{survey\_total}\NormalTok{(}\AttributeTok{vartype =} \FunctionTok{c}\NormalTok{(}\StringTok{"se"}\NormalTok{, }\StringTok{"ci"}\NormalTok{)))}

\FunctionTok{ggplot}\NormalTok{(}
  \AttributeTok{data =}\NormalTok{ poverty\_size,}
  \FunctionTok{aes}\NormalTok{(}
    \AttributeTok{x =}\NormalTok{ Poverty, }\AttributeTok{y =}\NormalTok{ size,}
    \AttributeTok{ymax =}\NormalTok{ size\_upp, }\AttributeTok{ymin =}\NormalTok{ size\_low,}
    \AttributeTok{fill =}\NormalTok{ Poverty}
\NormalTok{  )}
\NormalTok{) }\SpecialCharTok{+}
  \FunctionTok{geom\_bar}\NormalTok{(}\AttributeTok{stat =} \StringTok{"identity"}\NormalTok{, }\AttributeTok{position =} \StringTok{"dodge"}\NormalTok{) }\SpecialCharTok{+}
  \FunctionTok{geom\_errorbar}\NormalTok{(}\AttributeTok{position =} \FunctionTok{position\_dodge}\NormalTok{(}\AttributeTok{width =} \FloatTok{0.9}\NormalTok{), }\AttributeTok{width =} \FloatTok{0.3}\NormalTok{) }\SpecialCharTok{+}
  \FunctionTok{theme\_bw}\NormalTok{() }\SpecialCharTok{+}
  \FunctionTok{theme}\NormalTok{(}\AttributeTok{legend.position =} \StringTok{"top"}\NormalTok{)}
\end{Highlighting}
\end{Shaded}

\begin{figure}

{\centering \includegraphics[width=0.7\linewidth]{docs/figures/fig-barras-pobreza-1} 

}

\caption{Gráfico de barras del tamaño poblacional por condición de pobreza}\label{fig:fig-barras-pobreza}
\end{figure}

Una de las ventajas de los gráficos de barras es que pueden extenderse fácilmente a comparaciones entre dos variables categóricas. Por ejemplo, se puede analizar el tamaño poblacional por nivel de pobreza cruzado con el estado de ocupación, presentada en la Figura \ref{fig:fig-barras-empleo-pobreza}:

\begin{Shaded}
\begin{Highlighting}[]
\NormalTok{employment\_poverty\_size }\OtherTok{\textless{}{-}}\NormalTok{ survey\_design }\SpecialCharTok{\%\textgreater{}\%}
  \FunctionTok{group\_by}\NormalTok{(unemployed, Poverty) }\SpecialCharTok{\%\textgreater{}\%}
  \FunctionTok{summarise}\NormalTok{(}\AttributeTok{size =} \FunctionTok{survey\_total}\NormalTok{(}\AttributeTok{vartype =} \FunctionTok{c}\NormalTok{(}\StringTok{"se"}\NormalTok{, }\StringTok{"ci"}\NormalTok{))) }\SpecialCharTok{\%\textgreater{}\%}
  \FunctionTok{as.data.frame}\NormalTok{() }\SpecialCharTok{\%\textgreater{}\%}
  \FunctionTok{mutate}\NormalTok{(}
    \AttributeTok{unemployed =} \FunctionTok{case\_when}\NormalTok{(}
\NormalTok{      unemployed }\SpecialCharTok{==} \DecValTok{0} \SpecialCharTok{\textasciitilde{}} \StringTok{"Occupied"}\NormalTok{,}
\NormalTok{      unemployed }\SpecialCharTok{==} \DecValTok{1} \SpecialCharTok{\textasciitilde{}} \StringTok{"Unemployed"}\NormalTok{,}
      \FunctionTok{is.na}\NormalTok{(unemployed) }\SpecialCharTok{\textasciitilde{}} \StringTok{"Not in the Labour Force"}
\NormalTok{    )}
\NormalTok{  )}

\FunctionTok{ggplot}\NormalTok{(}
  \AttributeTok{data =}\NormalTok{ employment\_poverty\_size,}
  \FunctionTok{aes}\NormalTok{(}
    \AttributeTok{x =}\NormalTok{ Poverty, }\AttributeTok{y =}\NormalTok{ size,}
    \AttributeTok{ymax =}\NormalTok{ size\_upp, }\AttributeTok{ymin =}\NormalTok{ size\_low,}
    \AttributeTok{fill =} \FunctionTok{as.factor}\NormalTok{(unemployed)}
\NormalTok{  )}
\NormalTok{) }\SpecialCharTok{+}
  \FunctionTok{geom\_bar}\NormalTok{(}\AttributeTok{stat =} \StringTok{"identity"}\NormalTok{, }\AttributeTok{position =} \StringTok{"dodge"}\NormalTok{) }\SpecialCharTok{+}
  \FunctionTok{geom\_errorbar}\NormalTok{(}\AttributeTok{position =} \FunctionTok{position\_dodge}\NormalTok{(}\AttributeTok{width =} \FloatTok{0.9}\NormalTok{), }\AttributeTok{width =} \FloatTok{0.3}\NormalTok{) }\SpecialCharTok{+}
  \FunctionTok{theme\_bw}\NormalTok{() }\SpecialCharTok{+}
  \FunctionTok{theme}\NormalTok{(}\AttributeTok{legend.position =} \StringTok{"top"}\NormalTok{)}
\end{Highlighting}
\end{Shaded}

\begin{figure}

{\centering \includegraphics[width=0.8\linewidth]{docs/figures/fig-barras-empleo-pobreza-1} 

}

\caption{Gráfico de barras del tamaño poblacional por estado de ocupación y pobreza}\label{fig:fig-barras-empleo-pobreza}
\end{figure}

Este tipo de análisis gráfico resulta especialmente útil para identificar intersecciones de vulnerabilidad, al mostrar de manera simultánea dos o más características poblacionales. También es posible presentar proporciones en lugar de conteos. Por ejemplo, la proporción de hombres en condición de pobreza por zona, presentada en la Figura \ref{fig:fig-barras-prop-zona-hombre}:

\begin{Shaded}
\begin{Highlighting}[]
\NormalTok{male\_zone\_poverty\_proportion }\OtherTok{\textless{}{-}}\NormalTok{ male\_subset }\SpecialCharTok{\%\textgreater{}\%}
  \FunctionTok{group\_by}\NormalTok{(Zone, Poverty) }\SpecialCharTok{\%\textgreater{}\%}
  \FunctionTok{summarise}\NormalTok{(}\AttributeTok{proportion =} \FunctionTok{survey\_prop}\NormalTok{(}\AttributeTok{vartype =} \FunctionTok{c}\NormalTok{(}\StringTok{"se"}\NormalTok{, }\StringTok{"ci"}\NormalTok{)))}

\FunctionTok{ggplot}\NormalTok{(}
  \AttributeTok{data =}\NormalTok{ male\_zone\_poverty\_proportion,}
  \FunctionTok{aes}\NormalTok{(}
    \AttributeTok{x =}\NormalTok{ Poverty, }\AttributeTok{y =}\NormalTok{ proportion,}
    \AttributeTok{ymax =}\NormalTok{ proportion\_upp, }\AttributeTok{ymin =}\NormalTok{ proportion\_low,}
    \AttributeTok{fill =}\NormalTok{ Zone}
\NormalTok{  )}
\NormalTok{) }\SpecialCharTok{+}
  \FunctionTok{geom\_bar}\NormalTok{(}\AttributeTok{stat =} \StringTok{"identity"}\NormalTok{, }\AttributeTok{position =} \StringTok{"dodge"}\NormalTok{) }\SpecialCharTok{+}
  \FunctionTok{geom\_errorbar}\NormalTok{(}\AttributeTok{position =} \FunctionTok{position\_dodge}\NormalTok{(}\AttributeTok{width =} \FloatTok{0.9}\NormalTok{), }\AttributeTok{width =} \FloatTok{0.3}\NormalTok{) }\SpecialCharTok{+}
  \FunctionTok{scale\_fill\_manual}\NormalTok{(}\AttributeTok{values =}\NormalTok{ zone\_colors) }\SpecialCharTok{+}
  \FunctionTok{theme\_bw}\NormalTok{() }\SpecialCharTok{+}
  \FunctionTok{theme}\NormalTok{(}\AttributeTok{legend.position =} \StringTok{"top"}\NormalTok{)}
\end{Highlighting}
\end{Shaded}

\begin{figure}

{\centering \includegraphics[width=0.7\linewidth]{docs/figures/fig-barras-prop-zona-hombre-1} 

}

\caption{Proporción de hombres en condición de pobreza por zona geográfica}\label{fig:fig-barras-prop-zona-hombre}
\end{figure}

De forma análoga, se puede graficar la proporción de hombres en condición de pobreza desagregada por región, presentada en la Figura \ref{fig:fig-barras-prop-region-hombre}:

\begin{Shaded}
\begin{Highlighting}[]
\NormalTok{male\_region\_poverty\_proportion }\OtherTok{\textless{}{-}}\NormalTok{ male\_subset }\SpecialCharTok{\%\textgreater{}\%}
  \FunctionTok{group\_by}\NormalTok{(Region, poor) }\SpecialCharTok{\%\textgreater{}\%}
  \FunctionTok{summarise}\NormalTok{(}\AttributeTok{proportion =} \FunctionTok{survey\_prop}\NormalTok{(}\AttributeTok{vartype =} \FunctionTok{c}\NormalTok{(}\StringTok{"se"}\NormalTok{, }\StringTok{"ci"}\NormalTok{))) }\SpecialCharTok{\%\textgreater{}\%}
  \FunctionTok{data.frame}\NormalTok{()}

\FunctionTok{ggplot}\NormalTok{(}
  \AttributeTok{data =}\NormalTok{ male\_region\_poverty\_proportion,}
  \FunctionTok{aes}\NormalTok{(}
    \AttributeTok{x =}\NormalTok{ Region, }\AttributeTok{y =}\NormalTok{ proportion,}
    \AttributeTok{ymax =}\NormalTok{ proportion\_upp, }\AttributeTok{ymin =}\NormalTok{ proportion\_low,}
    \AttributeTok{fill =} \FunctionTok{as.factor}\NormalTok{(poor)}
\NormalTok{  )}
\NormalTok{) }\SpecialCharTok{+}
  \FunctionTok{geom\_bar}\NormalTok{(}\AttributeTok{stat =} \StringTok{"identity"}\NormalTok{, }\AttributeTok{position =} \StringTok{"dodge"}\NormalTok{) }\SpecialCharTok{+}
  \FunctionTok{geom\_errorbar}\NormalTok{(}\AttributeTok{position =} \FunctionTok{position\_dodge}\NormalTok{(}\AttributeTok{width =} \FloatTok{0.9}\NormalTok{), }\AttributeTok{width =} \FloatTok{0.3}\NormalTok{) }\SpecialCharTok{+}
  \FunctionTok{theme\_bw}\NormalTok{() }\SpecialCharTok{+}
  \FunctionTok{theme}\NormalTok{(}\AttributeTok{legend.position =} \StringTok{"top"}\NormalTok{)}
\end{Highlighting}
\end{Shaded}

\begin{figure}

{\centering \includegraphics[width=0.75\linewidth]{docs/figures/fig-barras-prop-region-hombre-1} 

}

\caption{Proporción de hombres en condición de pobreza por región}\label{fig:fig-barras-prop-region-hombre}
\end{figure}

En síntesis, los gráficos de barras junto con sus correspondientes errores estánar permiten una descripción clara de los patrones de pobreza, empleo u otras variables categóricas en la población, y cuando se aplican correctamente, constituyen una herramienta poderosa para comunicar hallazgos de encuestas de hogares manteniendo el rigor estadístico en la representación de los resultados.

\subsection{Mapas}\label{viz-mapas}

Los mapas temáticos constituyen una herramienta de visualización especialmente útil en el análisis de encuestas de hogares, ya que permiten representar la distribución territorial de indicadores sociales y demográficos. En este contexto, las estimaciones de proporciones calculadas en secciones anteriores ---como el porcentaje de pobreza o desempleo por región--- pueden proyectarse directamente sobre unidades geográficas, facilitando así la identificación de patrones espaciales y desigualdades regionales.

Para construir este tipo de mapas en \texttt{R} se requiere información geoespacial en formato \emph{shapefile}, la cual contiene los polígonos asociados a cada unidad geográfica de análisis. En este capítulo se utiliza el archivo \texttt{bigcity\_shape/BigCity.shp}, cuyas regiones corresponden a las definidas en la base de datos \texttt{BigCity}. Los paquetes \texttt{sf} \citep{sf2026} y \texttt{tmap} \citep{tmap2026} permiten cargar, manipular y visualizar esta información geográfica de manera integrada con los resultados obtenidos a partir de la encuesta.

El siguiente código carga los paquetes necesarios para trabajar con información geoespacial y visualización cartográfica en \texttt{R}. En primer lugar, la librería \texttt{sf} permite leer y manipular objetos espaciales modernos, mientras que \texttt{tmap} proporciona herramientas para la elaboración de mapas temáticos. La instrucción \texttt{tmap\_mode("plot")} establece el modo de visualización estática de los mapas, apropiado para documentos y reportes. Finalmente, la función \texttt{read\_sf()} importa el archivo geográfico \texttt{BigCity.shp}, almacenándolo en el objeto \texttt{shapeBigCity}, el cual contiene los polígonos correspondientes a las regiones de análisis.

\begin{Shaded}
\begin{Highlighting}[]
\FunctionTok{library}\NormalTok{(sf)}
\FunctionTok{library}\NormalTok{(tmap)}
\FunctionTok{tmap\_mode}\NormalTok{(}\StringTok{"plot"}\NormalTok{)}
\NormalTok{bigcity\_shape }\OtherTok{\textless{}{-}} \FunctionTok{read\_sf}\NormalTok{(}\StringTok{"Data/shapeBigCity/BigCity.shp"}\NormalTok{)}
\end{Highlighting}
\end{Shaded}

Con el shapefile cargado, se puede generar un mapa de las regiones usando \texttt{tm\_shape()} y \texttt{tm\_polygons()}, como se muestra en la Figura \ref{fig:fig-mapa-regiones}:

\begin{Shaded}
\begin{Highlighting}[]
\FunctionTok{tm\_shape}\NormalTok{(bigcity\_shape) }\SpecialCharTok{+}
  \FunctionTok{tm\_polygons}\NormalTok{(}\AttributeTok{col =} \StringTok{"Region"}\NormalTok{)}
\end{Highlighting}
\end{Shaded}

\begin{figure}

{\centering \includegraphics[width=0.7\linewidth]{docs/figures/fig-mapa-regiones-1} 

}

\caption{Mapa de regiones geográficas de BigCity}\label{fig:fig-mapa-regiones}
\end{figure}

Por ejemplo, si se desea representar geográficamente la proporción de hombres en condición de pobreza por región, estimada previamente en la Figura \ref{fig:fig-barras-prop-region-hombre}, es necesario integrar las estimaciones obtenidas con la información espacial del shapefile. Para ello, el objeto \texttt{male\_region\_poverty\_proportion}, que contiene las proporciones estimadas para cada región, se vincula con el archivo geográfico mediante la función \texttt{left\_join()}. Posteriormente, se filtran únicamente las observaciones correspondientes a la categoría de pobreza (\texttt{poor\ ==\ 1}), de manera que el mapa represente exclusivamente la distribución espacial de la población masculina en condición de pobreza.

\begin{Shaded}
\begin{Highlighting}[]
\NormalTok{shape\_region\_map }\OtherTok{\textless{}{-}} \FunctionTok{tm\_shape}\NormalTok{(}
\NormalTok{  bigcity\_shape }\SpecialCharTok{\%\textgreater{}\%}
    \FunctionTok{left\_join}\NormalTok{(}
\NormalTok{      male\_region\_poverty\_proportion }\SpecialCharTok{\%\textgreater{}\%}
        \FunctionTok{filter}\NormalTok{(poor }\SpecialCharTok{==} \DecValTok{1}\NormalTok{),}
      \AttributeTok{by =} \StringTok{"Region"}
\NormalTok{    )}
\NormalTok{)}
\end{Highlighting}
\end{Shaded}

Una vez construido el objeto espacial con las estimaciones de interés, el argumento \texttt{breaks} ayuda a definir los puntos de corte que determinarán los intervalos de la escala de color utilizada en el mapa. Posteriormente, se genera el mapa temático, permitiendo visualizar la distribución regional de la proporción de hombres en condición de pobreza, como se presenta en la Figura \ref{fig:fig-mapa-pobreza-hombre}.

\begin{Shaded}
\begin{Highlighting}[]
\NormalTok{poverty\_breaks }\OtherTok{\textless{}{-}} \FunctionTok{c}\NormalTok{(}\DecValTok{0}\NormalTok{, }\FloatTok{0.2}\NormalTok{, }\FloatTok{0.4}\NormalTok{, }\FloatTok{0.6}\NormalTok{, }\FloatTok{0.8}\NormalTok{, }\DecValTok{1}\NormalTok{)}
\NormalTok{shape\_region\_map }\SpecialCharTok{+}
  \FunctionTok{tm\_polygons}\NormalTok{(}
    \AttributeTok{col =} \StringTok{"proportion"}\NormalTok{,}
    \AttributeTok{breaks =}\NormalTok{ poverty\_breaks,}
    \AttributeTok{title =} \StringTok{"Poverty proportion"}\NormalTok{,}
    \AttributeTok{palette =} \StringTok{"YlOrRd"}
\NormalTok{  )}
\end{Highlighting}
\end{Shaded}

\begin{figure}

{\centering \includegraphics[width=0.7\linewidth]{docs/figures/fig-mapa-pobreza-hombre-1} 

}

\caption{Proporción de hombres en condición de pobreza por región}\label{fig:fig-mapa-pobreza-hombre}
\end{figure}

Otro uso de los mapas en encuestas es la visualización de la precisión de las estimaciones. Por ejemplo, usando el coeficiente de variación del ingreso medio por región es posible identificar las áreas donde las estimaciones son menos confiables, como se muestra en la Figura \ref{fig:fig-mapa-cv-ingreso}.

El siguiente código calcula y representa geográficamente el coeficiente de variación del ingreso medio por región. En primer lugar, la función se estima el ingreso promedio para cada región. El argumento \texttt{vartype\ =\ "cv"} solicita el cálculo del coeficiente de variación de cada estimación. Posteriormente, se definen los intervalos de clasificación de la escala de color mediante el objeto \texttt{cv\_breaks}, y las estimaciones se integran con la información geográfica del shapefile usando \texttt{left\_join()}. Finalmente, mediante las funciones \texttt{tm\_shape()} y \texttt{tm\_polygons()} del paquete \texttt{tmap}, se construye el mapa que representa espacialmente los coeficientes de variación utilizando una escala de colores en blanco y negro y ajustando la proporción del mapa con \texttt{tm\_layout(asp\ =\ 0)}.

\begin{Shaded}
\begin{Highlighting}[]
\NormalTok{region\_mean }\OtherTok{\textless{}{-}}\NormalTok{ survey\_design }\SpecialCharTok{\%\textgreater{}\%}
  \FunctionTok{group\_by}\NormalTok{(Region) }\SpecialCharTok{\%\textgreater{}\%}
  \FunctionTok{summarise}\NormalTok{(}\AttributeTok{mean =} \FunctionTok{survey\_mean}\NormalTok{(Income,}
                               \AttributeTok{na.rm =} \ConstantTok{TRUE}\NormalTok{,}
                               \AttributeTok{vartype =} \FunctionTok{c}\NormalTok{(}\StringTok{"cv"}\NormalTok{)))}

\NormalTok{cv\_breaks }\OtherTok{\textless{}{-}} \FunctionTok{c}\NormalTok{(}\DecValTok{0}\NormalTok{, }\FloatTok{0.1}\NormalTok{, }\FloatTok{0.2}\NormalTok{, }\DecValTok{1}\NormalTok{)}
\NormalTok{shape\_cv\_map }\OtherTok{\textless{}{-}} \FunctionTok{tm\_shape}\NormalTok{(}
\NormalTok{  bigcity\_shape }\SpecialCharTok{\%\textgreater{}\%}
    \FunctionTok{left\_join}\NormalTok{(region\_mean, }\AttributeTok{by =} \StringTok{"Region"}\NormalTok{)}
\NormalTok{)}

\NormalTok{shape\_cv\_map }\SpecialCharTok{+}
  \FunctionTok{tm\_polygons}\NormalTok{(}
    \StringTok{"mean\_cv"}\NormalTok{,}
    \AttributeTok{breaks =}\NormalTok{ cv\_breaks,}
    \AttributeTok{title =} \StringTok{"Mean income CV"}\NormalTok{,}
    \AttributeTok{palette =} \FunctionTok{c}\NormalTok{(}\StringTok{"\#FFFFFF"}\NormalTok{, }\StringTok{"\#000000"}\NormalTok{)}
\NormalTok{  ) }\SpecialCharTok{+}
  \FunctionTok{tm\_layout}\NormalTok{(}\AttributeTok{asp =} \DecValTok{0}\NormalTok{)}
\end{Highlighting}
\end{Shaded}

\begin{figure}

{\centering \includegraphics[width=0.7\linewidth]{docs/figures/fig-mapa-cv-ingreso-1} 

}

\caption{Coeficiente de variación del ingreso medio por región}\label{fig:fig-mapa-cv-ingreso}
\end{figure}

Los polígonos con tonos más oscuros corresponden a regiones donde la estimación del ingreso presenta mayor variabilidad relativa, lo que puede indicar tamaños muestrales insuficientes para obtener estimaciones precisas en esas áreas.

Cuando se desea representar simultáneamente dos variables ---por ejemplo, la proporción de pobreza y su coeficiente de variación--- es posible construir un mapa bivariante utilizando \texttt{ggplot2} junto con los paquetes \texttt{biscale} \citep{biscale2025} y \texttt{cowplot} \citep{cowplot2025}. Este tipo de visualización permite identificar regiones con alta pobreza y alta incertidumbre en la estimación de forma conjunta, como se presenta en la Figura \ref{fig:fig-mapa-bivariante}. El siguiente código estima la proporción de mujeres en condición de pobreza en zonas rurales para cada región, incorporando además el cálculo del coeficiente de variación de las estimaciones, generando así una base de datos lista para la construcción de un mapa bivariado.

\begin{Shaded}
\begin{Highlighting}[]
\NormalTok{region\_sex\_poverty\_proportion }\OtherTok{\textless{}{-}}\NormalTok{ survey\_design }\SpecialCharTok{\%\textgreater{}\%}
  \FunctionTok{group\_by}\NormalTok{(Region, Zone, Sex, poor) }\SpecialCharTok{\%\textgreater{}\%}
  \FunctionTok{summarise}\NormalTok{(}\AttributeTok{proportion =} \FunctionTok{survey\_mean}\NormalTok{(}\AttributeTok{vartype =} \StringTok{"cv"}\NormalTok{)) }\SpecialCharTok{\%\textgreater{}\%}
  \FunctionTok{filter}\NormalTok{(poor }\SpecialCharTok{==} \DecValTok{1}\NormalTok{, Zone }\SpecialCharTok{==} \StringTok{"Rural"}\NormalTok{, Sex }\SpecialCharTok{==} \StringTok{"Female"}\NormalTok{)}
\end{Highlighting}
\end{Shaded}

El siguiente código construye un mapa bivariado que permite visualizar simultáneamente la proporción de pobreza femenina rural y el coeficiente de variación asociado a dicha estimación por región. En primer lugar, se integran las estimaciones contenidas en \texttt{region\_sex\_poverty\_proportion} con la información geográfica del shapefile mediante \texttt{left\_join()}. Posteriormente, utilizando la función \texttt{bi\_class()} del paquete \texttt{biscale}, se clasifican conjuntamente las variables \texttt{proportion} y \texttt{proportion\_cv} en una cuadrícula de \(3 \times 3\) categorías (\texttt{dim\ =\ 3}), empleando el método de clasificación de Fisher (\texttt{style\ =\ "fisher"}). Luego, con \texttt{ggplot2} y \texttt{geom\_sf()}, se genera el mapa temático utilizando colores bivariados definidos por \texttt{bi\_scale\_fill()}, mientras que \texttt{bi\_theme()} aplica un estilo minimalista y se elimina la leyenda predeterminada.

A continuación, la función \texttt{bi\_legend()} crea una leyenda especializada que permite interpretar simultáneamente ambas dimensiones del mapa: el coeficiente de variación y la pobreza femenina rural. Finalmente, mediante las funciones \texttt{ggdraw()} y \texttt{draw\_plot()} del paquete \texttt{cowplot}, se combinan el mapa y la leyenda en una única visualización.

\begin{Shaded}
\begin{Highlighting}[]
\FunctionTok{library}\NormalTok{(biscale)}
\FunctionTok{library}\NormalTok{(cowplot)}

\NormalTok{bivariate\_shape }\OtherTok{\textless{}{-}}\NormalTok{ bigcity\_shape }\SpecialCharTok{\%\textgreater{}\%}
  \FunctionTok{left\_join}\NormalTok{(region\_sex\_poverty\_proportion, }\AttributeTok{by =} \StringTok{"Region"}\NormalTok{)}

\NormalTok{k }\OtherTok{\textless{}{-}} \DecValTok{3}
\NormalTok{bivariate\_data }\OtherTok{\textless{}{-}} \FunctionTok{bi\_class}\NormalTok{(}
\NormalTok{  bivariate\_shape,}
  \AttributeTok{y =}\NormalTok{ proportion,}
  \AttributeTok{x =}\NormalTok{ proportion\_cv,}
  \AttributeTok{dim =}\NormalTok{ k,}
  \AttributeTok{style =} \StringTok{"fisher"}
\NormalTok{)}

\NormalTok{bivariate\_map }\OtherTok{\textless{}{-}} \FunctionTok{ggplot}\NormalTok{() }\SpecialCharTok{+}
  \FunctionTok{geom\_sf}\NormalTok{(}
    \AttributeTok{data =}\NormalTok{ bivariate\_data,}
    \FunctionTok{aes}\NormalTok{(}\AttributeTok{fill =}\NormalTok{ bi\_class, }\AttributeTok{geometry =}\NormalTok{ geometry),}
    \AttributeTok{colour =} \StringTok{"white"}\NormalTok{,}
    \AttributeTok{size =} \FloatTok{0.1}
\NormalTok{  ) }\SpecialCharTok{+}
  \FunctionTok{bi\_scale\_fill}\NormalTok{(}\AttributeTok{pal =} \StringTok{"GrPink"}\NormalTok{, }\AttributeTok{dim =}\NormalTok{ k) }\SpecialCharTok{+}
  \FunctionTok{bi\_theme}\NormalTok{() }\SpecialCharTok{+}
  \FunctionTok{theme}\NormalTok{(}\AttributeTok{legend.position =} \StringTok{"none"}\NormalTok{)}

\NormalTok{bivariate\_legend }\OtherTok{\textless{}{-}} \FunctionTok{bi\_legend}\NormalTok{(}
  \AttributeTok{pal =} \StringTok{"GrPink"}\NormalTok{,}
  \AttributeTok{dim =}\NormalTok{ k,}
  \AttributeTok{xlab =} \StringTok{"Coefficient of variation"}\NormalTok{,}
  \AttributeTok{ylab =} \StringTok{"Rural female poverty"}\NormalTok{,}
  \AttributeTok{size =} \DecValTok{8}
\NormalTok{)}

\FunctionTok{ggdraw}\NormalTok{() }\SpecialCharTok{+}
  \FunctionTok{draw\_plot}\NormalTok{(bivariate\_map, }\DecValTok{0}\NormalTok{, }\DecValTok{0}\NormalTok{, }\DecValTok{1}\NormalTok{, }\AttributeTok{scale =} \FloatTok{0.7}\NormalTok{) }\SpecialCharTok{+}
  \FunctionTok{draw\_plot}\NormalTok{(bivariate\_legend, }\FloatTok{0.75}\NormalTok{, }\FloatTok{0.4}\NormalTok{, }\FloatTok{0.2}\NormalTok{, }\FloatTok{0.2}\NormalTok{)}
\end{Highlighting}
\end{Shaded}

\begin{figure}

{\centering \includegraphics[width=0.8\linewidth]{docs/figures/fig-mapa-bivariante-1} 

}

\caption{Mapa bivariante: proporción de pobreza femenina rural y su coeficiente de variación por región}\label{fig:fig-mapa-bivariante}
\end{figure}

En este mapa, cada celda de la leyenda combina un tono de la proporción de pobreza (eje vertical) con un tono del coeficiente de variación (eje horizontal), permitiendo identificar simultáneamente regiones con alta pobreza y alta incertidumbre en la estimación. Este tipo de representación es especialmente valioso para orientar decisiones sobre la necesidad de ampliar el tamaño muestral en zonas geográficas específicas.

\chapter{Modelos de regresión}\label{cap5}

Los modelos de regresión constituyen una de las herramientas más utilizadas en el análisis de datos provenientes de encuestas, ya que permiten estudiar la relación entre una variable de interés y un conjunto de variables explicativas. A través de estos modelos, es posible evaluar cómo cambian determinadas características de la población según factores demográficos, sociales o económicos observados en la muestra. No obstante, la validez de los resultados obtenidos depende de una adecuada especificación del modelo y de la correcta consideración de las características del diseño muestral de la encuesta.

En términos generales, los modelos de regresión buscan describir y cuantificar la asociación entre una variable respuesta (dependiente) y una o más variables explicativas (independientes), proporcionando elementos para la interpretación estadística y la formulación de inferencias sobre la población de estudio. No obstante, la validez de los resultados obtenidos depende, en gran medida, de una adecuada especificación del modelo y de la correcta consideración de las características del diseño muestral de la encuesta \citep{heeringa2017applied}.

Por ejemplo, es posible analizar el comportamiento del ingreso de los hogares como variable de respuesta, en función de características como el nivel educativo alcanzado y la situación laboral de sus integrantes, consideradas variables explicativas. Utilizando información proveniente de encuestas de hogares, este tipo de modelos permite identificar patrones de asociación, cuantificar el efecto de distintos factores socioeconómicos sobre el ingreso y generar evidencia empírica útil para el diseño, seguimiento y evaluación de políticas públicas.

No obstante, dado que las encuestas de hogares están sustentadas en diseños de muestreo complejos, los métodos clásicos de regresión, desarrollados bajo supuestos de muestreo aleatorio simple, pueden resultar inapropiados. Ignorar las características del diseño de muestreo puede generar estimaciones sesgadas de los coeficientes de regresión, así como una subestimación de sus errores estándar y varianzas, comprometiendo la validez de las inferencias estadísticas.

En consecuencia, el análisis de datos provenientes de encuestas exige una atención detallada al diseño de muestreo. La incorporación de los pesos de la encuesta y los ajustes correspondientes a la estratificación y la conglomeración permite obtener inferencias válidas y precisas. Adicionalmente, en algunos casos se han propuesto alternativas simplificadas, como el uso de pesos normalizados o enfoques de ponderación aproximada, que buscan equilibrar la complejidad metodológica con la factibilidad práctica del análisis.

El estudio de la regresión bajo diseños de muestreo complejos cuenta con una trayectoria bien documentada. \citet{kish1974inference} fueron de los primeros en discutir el impacto de estos diseños sobre las inferencias derivadas de modelos de regresión. Posteriormente, \citet{fuller1975regression} desarrolló un estimador de varianza apoyado en técnicas de linealización para modelos de regresión lineal múltiple con ponderación desigual, e introdujo métodos específicos para diseños estratificados y de dos etapas.

Más adelante, \citet{shah1977inference} abordaron el problema de las violaciones a los supuestos clásicos al trabajar con datos de encuestas, proponiendo alternativas de inferencia robusta para los parámetros. En paralelo, \citet{binder1983variances} se enfocó en las distribuciones muestrales de los estimadores de regresión en poblaciones finitas, estableciendo procedimientos para estimar varianzas bajo esquemas complejos.

En los años siguientes, \citet{skinner1989analysis} ampliaron estos aportes mediante estimadores de varianza para los coeficientes de regresión que incorporaban la estratificación y la conglomeración, recomendando explícitamente el uso de métodos de linealización o técnicas alternativas para la estimación de la varianza. Posteriormente, \citet{fuller2002regression} realizó un compendio de los métodos de estimación aplicables a modelos de regresión en encuestas complejas, mientras que \citet{pfeffermann2011modelling} discutió desarrollos más recientes, como los métodos de ponderación \emph{q-weighted}, mostrando evidencia empírica de su utilidad.

\section{Formulación del modelo}\label{formulacion-del-modelo}

Los modelos de regresión bajo diseños de muestreo complejos permiten trascender las estadísticas descriptivas y aproximarse a explicaciones causales o predictivas, siempre que se reconozcan y ajusten las particularidades del diseño muestral. Su correcta aplicación abre la posibilidad de examinar cómo las características sociodemográficas y económicas se asocian con distintos resultados de interés, aportando evidencia clave para la formulación de políticas públicas.

Como punto de partida, resulta conveniente presentar la estructura general de los modelos de regresión. La formulación más básica corresponde al modelo de regresión lineal simple, el cual describe la relación entre una variable respuesta y una única variable explicativa mediante la siguiente expresión:

\[
y = \beta_{0} + \beta_{1}x + \varepsilon
\]

En donde \(y\) representa la variable de respuesta (dependiente), \(x\) la variable explicativa (independiente), \(\beta_{0}\) el intercepto del modelo, \(\beta_{1}\) el coeficiente asociado a la variable independiente y \(\varepsilon\) el término de error, que recoge la variabilidad no explicada por el modelo y puede interpretarse como la diferencia entre el valor observado y el valor estimado por el modelo, denotado por \(\hat{y}_k\).

En muchas aplicaciones empíricas, y especialmente en el análisis de encuestas de hogares, el fenómeno de interés depende simultáneamente de múltiples factores. En estos casos se emplean modelos de regresión lineal múltiple, los cuales incorporan varias variables explicativas:

\[
y = \beta_{0} + \beta_{1}x_{1} + \cdots + \beta_{p}x_{p} + \varepsilon,
\]

donde cada coeficiente \(\beta_{j}\) cuantifica la asociación entre la variable respuesta y la correspondiente covariable \(x_{j}\), manteniendo constantes las demás variables incluidas en el modelo. De manera más compacta, el modelo puede expresarse utilizando notación matricial como:

\[
y_{k} = \mathbf{x}_{k}\boldsymbol{\beta} + \varepsilon_{k}, \quad k=1,\ldots,n,
\]

donde \(\mathbf{x}_{k} = [1, x_{1i}, \ldots, x_{pi}]\) representa el vector de covariables asociado a la unidad \(k\), mientras que \(\boldsymbol{\beta} = [\beta_{0}, \beta_{1}, \ldots, \beta_{p}]'\) corresponde al vector de parámetros desconocidos del modelo. En este contexto, el valor esperado de la variable respuesta condicionado al conjunto de covariables puede expresarse como:

\[
E(y \mid \mathbf{x}) = {\beta}_{0} + {\beta}_{1}x_{1} + \cdots + {\beta}_{p}x_{p}
\]

Los modelos de regresión lineal se sustentan en una serie de supuestos teóricos que garantizan la validez de las estimaciones y de las inferencias derivadas del modelo. En primer lugar, se asume que el valor esperado de los residuos condicionado a las covariables es igual a cero, es decir, \(E(\varepsilon_{k} \mid \mathbf{x}_{k}) = 0\), lo que implica la ausencia de sesgos sistemáticos en la estimación de la variable respuesta. Asimismo, se supone homogeneidad de la varianza de los errores, de modo que la variabilidad residual permanece constante para todos los valores de las covariables, esto es, \(Var(\varepsilon_{k} \mid \mathbf{x}_{k}) = \sigma^2\). Adicionalmente, se considera que los errores siguen una distribución normal con media cero y varianza constante, expresada como \(\varepsilon_{k} \mid \mathbf{x}_{k} \sim N(0,\sigma^2)\), y que los residuos asociados a distintas observaciones son independientes entre sí, de manera que \(Cov(\varepsilon_{k},\varepsilon_{j}\mid \mathbf{x}_{k},\mathbf{x}_{j})=0\).

La validez de los modelos de regresión lineal depende del cumplimiento de algunos supuestos clásicos ampliamente discutidos en la literatura especializada. En conjunto, estos supuestos permiten que los estimadores obtenidos mediante el modelo posean propiedades estadísticas deseables, tales como insesgamiento, eficiencia y consistencia.

\section{Uso de los pesos de muestreo}\label{uso-de-los-pesos-de-muestreo}

Cuando se trabaja con datos provenientes de encuestas basadas en diseños muestrales complejos, los supuestos clásicos de los modelos de regresión rara vez se cumplen de manera estricta, puesto que las observaciones no provienen de muestras aleatorias simples e independientes, sino de esquemas de selección que incorporan elementos como estratificación, conglomeración y probabilidades desiguales de selección. Por ejemplo, la presencia de conglomeración puede violar el supuesto de independencia de los errores, ya que individuos pertenecientes a un mismo hogar, segmento o área geográfica tienden a compartir características similares. Del mismo modo, las probabilidades desiguales de selección y los ajustes de ponderación pueden generar heterogeneidad en la varianza de las observaciones, incumpliendo el supuesto de homocedasticidad.

En este contexto, surge una pregunta fundamental: ¿de qué manera deben incorporarse los pesos muestrales en los modelos de regresión? La respuesta no es sencilla, ya que los pesos de la encuesta reflejan no solo las probabilidades de selección, sino también diversos ajustes asociados a la no respuesta, calibración y correcciones de cobertura. Aunque su utilización permite producir estimaciones representativas de la población, la incorporación directa de ponderaciones en los modelos también puede reducir la eficiencia estadística e incrementar la variabilidad de los estimadores.

En términos generales, la literatura distingue dos enfoques principales para abordar este problema \citep{heeringa2017applied}. El primero corresponde al enfoque basado en el diseño muestral, cuyo objetivo es obtener inferencias válidas para la población objetivo respetando las características del proceso de selección de la muestra. Desde esta perspectiva, los pesos muestrales son esenciales para corregir las probabilidades desiguales de inclusión y producir estimaciones insesgadas de los coeficientes de regresión. Sin embargo, este enfoque no protege frente a posibles errores de especificación del modelo; es decir, aun cuando las estimaciones sean válidas desde el punto de vista del diseño, el modelo puede no representar adecuadamente las relaciones existentes en la población.

El segundo corresponde al enfoque orientado al modelo, según el cual los pesos no son necesariamente requeridos siempre que el modelo esté correctamente especificado y el mecanismo de muestreo sea no informativo. Bajo este supuesto, las relaciones entre las variables observadas en la muestra coinciden con las de la población, de manera que el diseño muestral no introduce sesgos relevantes en la estimación de los parámetros. Desde esta perspectiva, incorporar ponderaciones podría incluso resultar contraproducente, al aumentar innecesariamente la varianza de los estimadores y, en consecuencia, los errores estándar.

La discusión sobre la conveniencia de utilizar ponderaciones en modelos de regresión ha sido ampliamente desarrollada en la literatura especializada, particularmente en trabajos como los de \citet{skinner1989analysis} y \citet{pfeffermann2011modelling}. En la práctica, una recomendación metodológica frecuente consiste en estimar los modelos tanto con ponderaciones como sin ellas y comparar posteriormente los resultados obtenidos. Si la inclusión de los pesos produce cambios importantes en los coeficientes estimados o modifica las conclusiones sustantivas del análisis, ello sugiere que el diseño de muestreo es informativo o que el modelo presenta problemas de especificación, por lo que el uso de ponderaciones resulta aconsejable \citep{UNhandbook}. Por el contrario, si los coeficientes permanecen relativamente estables y las ponderaciones únicamente incrementan los errores estándar, puede considerarse que el modelo captura adecuadamente la estructura de los datos y que el uso de pesos no es estrictamente necesario.

En términos aplicados, esta decisión suele depender del propósito analítico del estudio. Cuando el objetivo es realizar inferencia descriptiva y producir estimaciones representativas de la población, la utilización de ponderaciones es indispensable. En cambio, en contextos de inferencia analítica orientados al estudio de asociaciones, relaciones causales o contrastes de hipótesis, pueden emplearse tanto modelos ponderados como no ponderados, especialmente cuando el modelo incorpora variables relacionadas con el diseño muestral, como estratos o conglomerados. No obstante, el uso de modelos no ponderados debe justificarse cuidadosamente, dado que implica asumir condiciones más restrictivas sobre el mecanismo de selección de la muestra y la correcta especificación del modelo estadístico.

En resumen, cuando se opta por ponderar, las ponderaciones corrigen posibles sesgos de sobre o subrepresentación de determinados grupos y contribuyen a obtener estimaciones de varianza más exactas. Dentro del enfoque basado en el diseño, esto permite que los resultados se aproximen a valores insesgados comparables a los que se obtendrían en un censo completo, incluso cuando el modelo no esté formulado de manera óptima. Sin embargo, cuando los pesos presentan gran dispersión, pueden aumentar la varianza de los parámetros estimados y volver inestables las estimaciones, razón por la que en contextos explicativos o analíticos los modelos sin ponderar pueden, en ocasiones, arrojar resultados más consistentes y eficientes.

En cualquier caso, si el modelo se encuentra mal especificado, ignorar las ponderaciones muestrales puede conducir a estimaciones sesgadas, poco informativas o incluso carentes de validez inferencial. Por esta razón, la selección adecuada de las variables incluidas en el modelo constituye un aspecto fundamental del análisis. Con el fin de abordar estas problemáticas, se han propuesto diversas estrategias de ajuste orientadas a lograr un equilibrio entre ambos objetivos. Entre los procedimientos más utilizados se destacan los siguientes:

\begin{enumerate}
\def\labelenumi{\arabic{enumi}.}
\item
  Pesos tipo Senado: este procedimiento ajusta los pesos de manera que su suma coincida con el tamaño de la muestra, en lugar del tamaño de la población. El objetivo es mantener la representatividad relativa de las unidades, pero reduciendo la dispersión de los pesos originales, lo que resulta particularmente ventajoso en encuestas donde existe alta variabilidad entre los factores de expansión. De esta forma, las nuevas ponderaciones son \(w_k^{Senate} = w_k \times \frac{n}{\sum w_k}\).
\item
  Pesos normalizados: en este enfoque los pesos originales se reescalan para que su suma sea igual a uno, lo que evita un incremento innecesario de la varianza en los modelos. Esta técnica es especialmente útil cuando se trabaja con diferentes subconjuntos de datos (por ejemplo, modelos estimados en subpoblaciones) o cuando se busca minimizar la inflación de la varianza. De esta forma, las nuevas ponderaciones son \(w_k^{Normalized} = \frac{w_k}{\sum w_k}\).
\end{enumerate}

Es importante destacar que, en ambos métodos, los pesos ajustados se obtienen mediante transformaciones multiplicativas directas de los pesos muestrales originales. Por ello, no deben emplearse para calcular tamaños o totales poblacionales. Asimismo, estos procedimientos no alteran las estimaciones de razones, como medias o proporciones, ya que en tales casos los pesos se cancelan en el cociente. En la práctica, el uso de estos ajustes puede considerarse una solución pragmática en contextos donde no se dispone de software especializado para encuestas. No obstante, cuando se cuenta con herramientas como los paquetes \texttt{survey} \citep{Lumley2024} y \texttt{srvyr} \citep{Freedman2024} de R, descritos en los capítulos anteriores, el reescalamiento deja de ser necesario, pues dichos paquetes realizan el tratamiento adecuado para preservar tanto la representatividad como las propiedades inferenciales del modelo.

\section{Estimación de los parámetros del modelo}\label{estimacion-de-los-parametros-del-modelo}

En términos del ajuste de modelos de regresión con encuestas de hogares, se han desarrollado metodologías inferenciales avanzadas que integran las fuentes de incertidumbre en un mismo marco analítico, buscando reflejar tanto la estructura del diseño como los supuestos y limitaciones del modelo. Entre las aproximaciones más relevantes se encuentran la pseudo-verosimilitud \citep{molina1992pseudo} y la inferencia combinada \citep{binder2011estimating}.

El método de pseudo-verosimilitud extiende las técnicas tradicionales de máxima verosimilitud para ajustarlas a las particularidades de los diseños muestrales complejos. En este enfoque, la distribución de muestreo definida por el diseño ocupa un rol central, mientras que la distribución del modelo pasa a un segundo plano. Si bien en contextos de modelos bien especificados los estimadores basados en pseudo-verosimilitud tienden a ser insesgados o consistentes, su principal virtud radica en evitar los sesgos que se originarían al ignorar el diseño muestral. En términos prácticos, este método traduce el modelo tradicional en uno que respeta la forma en que los datos fueron obtenidos, garantizando inferencias más sólidas.

En contraste, la inferencia combinada propone un marco unificado en el que se integran simultáneamente la variabilidad del muestreo y la incertidumbre del modelo. Al considerar ambas fuentes, este enfoque ofrece una visión más completa de la variabilidad total y permite obtener estimaciones más precisas y confiables. Su principal aporte radica en evitar los sesgos que pueden surgir cuando el análisis se realiza exclusivamente desde la perspectiva del diseño o del modelo. De este modo, la inferencia combinada resulta especialmente útil en aplicaciones donde se requiere equilibrar representatividad poblacional y solidez estadística de los modelos ajustados.

Cuando se desea estimar los parámetros de un modelo de regresión lineal a partir de una muestra con un diseño complejo, el enfoque estándar cambia. En consecuencia, los estimadores tradicionales, como los obtenidos por máxima verosimilitud o mínimos cuadrados, pueden resultar sesgados y producir errores estándar poco confiables. Frente a esta situación, \citet{wolter2007} propone el uso de métodos no paramétricos robustos, como la linealización de Taylor o Bootstrap, para estimar varianzas y obtener inferencias válidas.

En el caso de un modelo de regresión lineal simple, la estimación de \(\beta_1\) bajo un esquema de muestreo complejo se realiza mediante un estimador ponderado, que puede expresarse de la siguiente forma:

\[
\hat{\beta_{1}}  =
\frac{\sum_{h} \sum_{i} \sum_{k} w_{hik}\,(y_{hik}-\hat{\bar{y}})(x_{hik}-\hat{\bar{x}})}
{\sum_{h} \sum_{i} \sum_{k} w_{hik}\,(x_{hik}-\hat{\bar{x}})^{2}}
=
\frac{\hat{t}_{xy}}{\hat{t}_{xx}}
\]

En donde \(\hat{\bar{y}}=\hat{t}_y/\hat{N}\) y \(\hat{\bar{x}}=\hat{t}_x/\hat{N}\) corresponden a las medias estimadas de las variables de estudio y auxiliar, respectivamente, mientras que \(w_{hik}\) representa los pesos de muestreo asociados a cada unidad de observación. Esta formulación incorpora explícitamente dichos pesos, lo que permite considerar las probabilidades desiguales de selección y las características propias del diseño muestral. Asimismo, la varianza estimada de \(\hat{\beta}_1\) puede aproximarse como:

\[
\widehat{Var}\left(\hat{\beta_{1}}\right)
\approx  \frac{\widehat{Var}(\hat{t}_{xy})+\hat{\beta}_{1}^{2}\widehat{Var}(\hat{t}_{xx})-2\hat{\beta}_{1}\widehat{Cov}(\hat{t}_{xy},\hat{t}_{xx})}{(\hat{t}_{xx})^{2}}
\]

Si \(\hat{\beta}_{1}\) corresponde al estimador de la pendiente en un modelo de regresión lineal simple ponderado, entonces el estimador del intercepto \(\hat{\beta}_{0}\) está dado por

\[
\hat{\beta}_{0} = \hat{\bar{y}} - \hat{\beta}_{1}\hat{\bar{x}}
\]

Nótese que la varianza de \(\hat{\beta}_{0}\) depende simultáneamente de los estimadores \(\hat{\bar{y}}\), \(\hat{\bar{x}}\) y \(\hat{\beta}_{1}\). Por esta razón, su aproximación puede obtenerse mediante las propiedades de la varianza de combinaciones lineales de estimadores, en conjunto con técnicas de linearización de Taylor. A partir de este procedimiento se deriva la matriz de varianzas y covarianzas de los coeficientes del modelo, de la cual se obtienen tanto los errores estándar como las covarianzas asociadas a las estimaciones de los parámetros.

En el caso de la regresión múltiple, si \(\hat{\boldsymbol{\beta}}\) representa el estimador de \(\boldsymbol{\beta}\), entonces la varianza de cada coeficiente se estima considerando su interdependencia con los demás parámetros, lo que se refleja en la construcción de una matriz de varianzas-covarianzas que recoge tanto la variabilidad individual como las covarianzas entre todos los estimadores. De acuerdo con \citet{kish1974inference}, este cálculo requiere utilizar totales ponderados de cuadrados y productos cruzados de todas las combinaciones entre la variable dependiente \(y\) y el conjunto de predictores \(\mathbf{x} = (1, x_1, \ldots, x_p)\). En términos generales:

\[
\widehat{Var}(\hat{\boldsymbol{\beta}}) = \hat{\boldsymbol\Sigma} =
\begin{bmatrix}
\widehat{Var}(\hat{\beta}_{0}) & \widehat{Cov}(\hat{\beta}_{0},\hat{\beta}_{1}) & \cdots & \widehat{Cov}(\hat{\beta}_{0},\hat{\beta}_{p}) \\
\widehat{Cov}(\hat{\beta}_{0},\hat{\beta}_{1}) & \widehat{Var}(\hat{\beta}_{1}) & \cdots & \widehat{Cov}(\hat{\beta}_{1},\hat{\beta}_{p}) \\
\vdots & \vdots & \ddots & \vdots \\
\widehat{Cov}(\hat{\beta}_{0},\hat{\beta}_{p}) & \widehat{Cov}(\hat{\beta}_{1},\hat{\beta}_{p}) & \cdots & \widehat{Var}(\hat{\beta}_{p})
\end{bmatrix}
\]

En este capítulo se mostrará cómo implementar estos modelos en \texttt{R} mediante la librería \texttt{survey} \citep{Lumley2024}, junto con \texttt{srvyr} \citep{Freedman2024} y \texttt{tidyverse} \citep{tidyverse2026} para organizar el flujo de trabajo. Se abordará la especificación del diseño muestral, la estimación de modelos lineales y logísticos, así como la obtención de errores estándar y pruebas de hipótesis ajustadas al diseño, integrando los fundamentos teóricos con ejemplos prácticos reproducibles. Para ejemplificar los conceptos desarrollados hasta este punto, se empleará la misma base de datos utilizada a lo largo del libro. El proceso inicia con la carga de las librerías, los datos y la definición del diseño de muestreo:

\begin{Shaded}
\begin{Highlighting}[]
\FunctionTok{library}\NormalTok{(survey)}
\FunctionTok{library}\NormalTok{(srvyr)}
\FunctionTok{library}\NormalTok{(tidyverse)}

\FunctionTok{data}\NormalTok{(BigCity, }\AttributeTok{package =} \StringTok{"TeachingSampling"}\NormalTok{)}
\NormalTok{survey\_data }\OtherTok{\textless{}{-}} \FunctionTok{readRDS}\NormalTok{(}\StringTok{"Data/encuesta.rds"}\NormalTok{)}

\NormalTok{survey\_design }\OtherTok{\textless{}{-}}\NormalTok{ survey\_data }\SpecialCharTok{\%\textgreater{}\%}
  \FunctionTok{as\_survey\_design}\NormalTok{(}
    \AttributeTok{strata =}\NormalTok{ Stratum,}
    \AttributeTok{ids =}\NormalTok{ PSU,}
    \AttributeTok{weights =}\NormalTok{ wk,}
    \AttributeTok{nest =} \ConstantTok{TRUE}
\NormalTok{  )}
\end{Highlighting}
\end{Shaded}

En los ejemplos desarrollados a lo largo de este capítulo se emplerán los mismos subgrupos considerados en los capítulos anteriores.

\begin{Shaded}
\begin{Highlighting}[]
\NormalTok{urban\_subset }\OtherTok{\textless{}{-}}\NormalTok{ survey\_design }\SpecialCharTok{\%\textgreater{}\%}
  \FunctionTok{filter}\NormalTok{(Zone }\SpecialCharTok{==} \StringTok{"Urban"}\NormalTok{)}

\NormalTok{rural\_subset }\OtherTok{\textless{}{-}}\NormalTok{ survey\_design }\SpecialCharTok{\%\textgreater{}\%}
  \FunctionTok{filter}\NormalTok{(Zone }\SpecialCharTok{==} \StringTok{"Rural"}\NormalTok{)}

\NormalTok{female\_subset }\OtherTok{\textless{}{-}}\NormalTok{ survey\_design }\SpecialCharTok{\%\textgreater{}\%}
  \FunctionTok{filter}\NormalTok{(Sex }\SpecialCharTok{==} \StringTok{"Female"}\NormalTok{)}

\NormalTok{male\_subset }\OtherTok{\textless{}{-}}\NormalTok{ survey\_design }\SpecialCharTok{\%\textgreater{}\%}
  \FunctionTok{filter}\NormalTok{(Sex }\SpecialCharTok{==} \StringTok{"Male"}\NormalTok{)}
\end{Highlighting}
\end{Shaded}

Con el objetivo de ajustar un modelo de regresión entre las variables ingreso y gasto, resulta conveniente explorar previamente la relación existente entre ambas variables. Para ello, se construye un diagrama de dispersión mediante \texttt{ggplot2}, el cual permite evaluar visualmente la forma y la intensidad de la asociación. El resultado se presenta en la Figura \ref{fig:fig-scatter-muestra-sinP}.

\begin{Shaded}
\begin{Highlighting}[]
\NormalTok{unweighted\_plot }\OtherTok{\textless{}{-}} \FunctionTok{ggplot}\NormalTok{(}
  \AttributeTok{data =}\NormalTok{ survey\_data,}
  \FunctionTok{aes}\NormalTok{(}\AttributeTok{x =}\NormalTok{ Expenditure, }\AttributeTok{y =}\NormalTok{ Income)}
\NormalTok{) }\SpecialCharTok{+}
  \FunctionTok{geom\_point}\NormalTok{() }\SpecialCharTok{+}
  \FunctionTok{geom\_smooth}\NormalTok{(}\AttributeTok{method =} \StringTok{"lm"}\NormalTok{, }\AttributeTok{se =} \ConstantTok{FALSE}\NormalTok{)}

\NormalTok{unweighted\_plot}
\end{Highlighting}
\end{Shaded}

\begin{figure}

{\centering \includegraphics[width=0.65\linewidth]{docs/figures/fig-scatter-muestra-sinP-1} 

}

\caption{Diagrama de dispersión entre gasto e ingreso en la muestra (sin ponderar)}\label{fig:fig-scatter-muestra-sinP}
\end{figure}

Los datos muestrales conservan una relación aproximadamente lineal entre ingreso y gasto, aunque se observa una mayor dispersión en los valores correspondientes a hogares con niveles de gasto más elevados. Una vez completada la exploración gráfica de los datos, se procede al ajuste de los modelos de regresión lineal.

Para ajustar modelos incorporando los factores de expansión se usa la función \texttt{svyglm} de la librería \texttt{survey}, que permite estimar modelos lineales incorporando explícitamente las características del diseño muestral. En este caso, la expresión \texttt{Income\ \textasciitilde{}\ Expenditure} especifica que la variable \texttt{Income} se modela en función de la variable \texttt{Expenditure}; es decir, se ajusta una regresión lineal simple en la que el gasto actúa como variable explicativa del ingreso. El argumento \texttt{design\ =\ survey\_design} indica que el modelo debe utilizar el objeto de diseño muestral previamente definido, mientras que \texttt{family\ =\ stats::gaussian()} especifica que se ajusta un modelo lineal con distribución normal, equivalente a una regresión lineal clásica.

\begin{Shaded}
\begin{Highlighting}[]
\NormalTok{fit\_svy }\OtherTok{\textless{}{-}} \FunctionTok{svyglm}\NormalTok{(}
\NormalTok{  Income }\SpecialCharTok{\textasciitilde{}}\NormalTok{ Expenditure }\SpecialCharTok{+}\NormalTok{ Zone }\SpecialCharTok{+}\NormalTok{ Sex,}
  \AttributeTok{design =}\NormalTok{ survey\_design,}
  \AttributeTok{family =}\NormalTok{ stats}\SpecialCharTok{::}\FunctionTok{gaussian}\NormalTok{()}
\NormalTok{)}

\NormalTok{fit\_svy}
\end{Highlighting}
\end{Shaded}

\begin{verbatim}
## Stratified 1 - level Cluster Sampling design (with replacement)
## With (238) clusters.
## Called via srvyr
## Sampling variables:
##   - ids: PSU 
##   - strata: Stratum 
##   - weights: wk 
## 
## Call:  svyglm(formula = Income ~ Expenditure + Zone + Sex, design = survey_design, 
##     family = stats::gaussian())
## 
## Coefficients:
## (Intercept)  Expenditure    ZoneUrban      SexMale  
##       73.58         1.22        66.65        20.64  
## 
## Degrees of Freedom: 2604 Total (i.e. Null);  116 Residual
## Null Deviance:       635000000 
## Residual Deviance: 309000000     AIC: 38300
\end{verbatim}

La salida anterior corresponde al ajuste de un modelo de regresión lineal en el que la variable \texttt{Income} se explica a partir de las variables \texttt{Expenditure}, \texttt{Zone} y \texttt{Sex}. Los resultados indican la existencia de una relación positiva entre ingreso y gasto, de modo que, manteniendo constantes las demás variables, el ingreso esperado aumenta en aproximadamente 1.22 unidades por cada unidad adicional de gasto.

Asimismo, se evidencian diferencias sistemáticas asociadas a las características sociodemográficas incluidas en el modelo: residir en zona urbana se asocia con un incremento promedio de 66.65 unidades en el ingreso respecto de la categoría de referencia, mientras que ser hombre se asocia con un aumento de 20.64 unidades en comparación con la categoría base del sexo. El intercepto estimado, igual a 73.58, representa el nivel esperado de ingreso para la categoría de referencia del modelo (zona rural y sexo femenino) cuando el gasto es igual a cero.

\section{Diagnóstico del modelo}\label{diagnostico-del-modelo}

En el análisis de encuestas de hogares, la evaluación de un modelo estadístico constituye un aspecto importante que garantiza la validez de las inferencias. La literatura señala que un modelo adecuadamente especificado no depende únicamente de la selección de las covariables, sino también del cumplimiento de los supuestos que sustentan la consistencia y confiabilidad de los resultados \citep{tellez2015estadistica}.

En el caso de los modelos de regresión lineal aplicados a encuestas complejas, resulta necesario examinar distintos aspectos relacionados con la calidad del ajuste y el comportamiento de los errores. Como se indicó anteriormente, entre ellos se encuentran la capacidad explicativa del modelo, la normalidad y homogeneidad de la varianza de los residuos, la independencia entre errores y la posible presencia de observaciones influyentes o valores atípicos que puedan afectar las estimaciones.

Estas verificaciones adquieren particular importancia en el contexto de diseños muestrales complejos, donde características como la estratificación, la conglomeración y el uso de factores de expansión pueden intensificar problemas asociados con la heterocedasticidad, la dependencia entre observaciones o la sensibilidad a valores extremos. Por esta razón, el diagnóstico del modelo no debe limitarse a los supuestos clásicos de la regresión, sino incorporar también las particularidades derivadas del diseño de la encuesta. La aplicación sistemática de estos procedimientos permite evaluar la solidez del modelo, fortalecer la confianza en las inferencias y asegurar que los resultados obtenidos sean representativos y útiles tanto para la investigación social aplicada como para el análisis de políticas públicas.

\subsection{Coeficiente de determinación}\label{coeficiente-de-determinacion}

Una de las medidas más utilizadas para evaluar el ajuste de un modelo de regresión es el coeficiente de determinación, denotado por \(R^{2}\) y conocido también como coeficiente de correlación múltiple. Este indicador cuantifica la proporción de la variabilidad de la variable dependiente que es explicada por el modelo. Sus valores oscilan entre cero y uno: cuanto más cercano se encuentre a la unidad, mayor será la capacidad explicativa del modelo; en contraste, valores próximos a cero indican un bajo nivel de explicación de la variabilidad observada. La magnitud del \(R^{2}\) no debe evaluarse en términos absolutos, sino considerando el fenómeno estudiado y el tipo de información disponible.

El coeficiente se calcula a partir de las sumas de cuadrados totales y de error, como \(R^{2} = 1 - \frac{SSE}{SST}\); en donde \(SST\) representa la suma de cuadrados totales y \(SSE\) la suma de cuadrados del error. En encuestas con diseños de muestreo complejos, esta medida debe ajustarse para incorporar la estructura del diseño y los pesos muestrales. El estimador ponderado del coeficiente de determinación se define como:

\[
\hat{R}^2 = 1 - \frac{\widehat{SSE}}{\widehat{SST}}
\]

donde \(\widehat{SSE}\) es la suma ponderada de errores al cuadrado, calculada como \(\widehat{SSE} = \sum_{h} \sum_{i} \sum_{k} w_{hik} \,(y_{hik} - x_{hik}\hat{\boldsymbol{\beta}})^2\); mientras que \(\widehat{SST}\) representa la suma total ponderada de cuadrados, definida por \(\widehat{SST} = \sum_{h} \sum_{i} \sum_{k} w_{hik}\,(y_{hik} - \hat{\bar{y}})^2\).

Dado que \(R^{2}\) tiende a incrementarse conforme se incorporan más variables al modelo, se recomienda complementar su interpretación con el coeficiente de determinación ajustado (\(R_{adj}^{2}\)), el cual introduce una penalización en función del número de covariables incluidas y del tamaño muestral \citep{pfeffermann2011modelling}. Este coeficiente se define como \(\hat{R}^2_{adj} = 1 - \frac{(n-1)}{(n-p)}(1 - \hat{R}^2)\), donde \(n\) corresponde al número de observaciones, \(p\) representa el número de parámetros estimados y \(\hat{R}^{2}\) denota el coeficiente de determinación ponderado definido previamente.

De esta manera, el ajuste permite evaluar la capacidad explicativa del modelo controlando el efecto derivado de la incorporación de nuevas variables. Este ajuste facilita una comparación más equitativa entre modelos con distinto número de predictores y resulta particularmente útil en el análisis de encuestas, donde la complejidad del diseño y el uso de ponderaciones pueden influir notablemente en la magnitud del \(R^{2}\).

En \texttt{R}, el coeficiente de determinación \(R^{2}\) puede estimarse a partir de los modelos ajustados previamente. Para ello, se ajusta inicialmente un modelo nulo (incluyendo únicamente el intercepto), el cual permite calcular la suma total ponderada de cuadrados (\(\widehat{SST}\)).

\begin{Shaded}
\begin{Highlighting}[]
\NormalTok{null\_model }\OtherTok{\textless{}{-}} \FunctionTok{svyglm}\NormalTok{(Income }\SpecialCharTok{\textasciitilde{}} \DecValTok{1}\NormalTok{, }\AttributeTok{design =}\NormalTok{ survey\_design)}

\NormalTok{s1 }\OtherTok{\textless{}{-}} \FunctionTok{summary}\NormalTok{(fit\_svy)}
\NormalTok{s0 }\OtherTok{\textless{}{-}} \FunctionTok{summary}\NormalTok{(null\_model)}

\NormalTok{wsst }\OtherTok{\textless{}{-}}\NormalTok{ s0}\SpecialCharTok{$}\NormalTok{dispersion[}\DecValTok{1}\NormalTok{]}
\NormalTok{wsse }\OtherTok{\textless{}{-}}\NormalTok{ s1}\SpecialCharTok{$}\NormalTok{dispersion[}\DecValTok{1}\NormalTok{]}
\end{Highlighting}
\end{Shaded}

La estimación del coeficiente de determinación y su contraparte ajustada se obtiene a partir del siguiente cálculo:

\begin{Shaded}
\begin{Highlighting}[]
\NormalTok{n }\OtherTok{\textless{}{-}} \FunctionTok{nrow}\NormalTok{(survey\_design)}
\NormalTok{p }\OtherTok{\textless{}{-}} \DecValTok{4}
\NormalTok{r2 }\OtherTok{\textless{}{-}} \DecValTok{1} \SpecialCharTok{{-}}\NormalTok{ wsse }\SpecialCharTok{/}\NormalTok{ wsst}
\NormalTok{r2\_adj }\OtherTok{\textless{}{-}} \DecValTok{1} \SpecialCharTok{{-}}\NormalTok{ ((n }\SpecialCharTok{{-}} \DecValTok{1}\NormalTok{) }\SpecialCharTok{/}\NormalTok{ (n }\SpecialCharTok{{-}}\NormalTok{ p)) }\SpecialCharTok{*}\NormalTok{ (}\DecValTok{1} \SpecialCharTok{{-}}\NormalTok{ r2)}

\NormalTok{r2}
\end{Highlighting}
\end{Shaded}

\begin{verbatim}
## [1] 0.5138
\end{verbatim}

\begin{Shaded}
\begin{Highlighting}[]
\NormalTok{r2\_adj}
\end{Highlighting}
\end{Shaded}

\begin{verbatim}
## [1] 0.5133
\end{verbatim}

\subsection{Residuales}\label{residuales}

En el diagnóstico de modelos, el análisis de los residuales es una de las herramientas más relevantes. Si el modelo está correctamente especificado, los residuales actúan como una aproximación de los errores no observados y permiten evaluar, de forma indirecta, si los supuestos del modelo se cumplen razonablemente en los datos. Su revisión sistemática permite identificar posibles desviaciones en la especificación, lo que ayuda a determinar si el ajuste es adecuado o si, por el contrario, es necesario reconsiderar la forma funcional del modelo o el método de estimación utilizado.

En encuestas con diseños de muestreo complejos, los residuales de Pearson constituyen una herramienta habitual para evaluar las discrepancias entre los valores observados y los valores esperados bajo el modelo ajustado. Siguiendo a \citet{heeringa2017applied}, estos se definen como:

\[
r_{k} = \frac{y_k - \mu_k}{\sqrt{\widehat{Var}(\mu_k)/w_k}},
\]

donde \(\hat{\mu}_{k} = \mathbf{x}_{k}^{\top}\hat{\boldsymbol{\beta}}\) representa el valor esperado de \(y_{k}\) bajo el modelo ajustado, \(w_{k}\) corresponde al peso muestral asociado a la unidad \(k\), y \(\widehat{Var}(\hat{\mu}_{k})\) denota la varianza estimada bajo la familia del modelo.

Los residuales de Pearson también se utilizan para evaluar aspectos clave del ajuste del modelo, en particular la normalidad y la homogeneidad de la varianza de los errores. De forma complementaria, el gráfico de residuos frente a valores ajustados constituye una herramienta diagnóstica especialmente informativa, ya que su inspección visual permite detectar posibles patrones sistemáticos que sugieran incumplimientos de los supuestos de independencia o heterocedasticidad.

El supuesto de normalidad de los errores también debe verificarse de manera explícita. Para ello, una herramienta estándar es el gráfico cuantil-cuantil (QQ-plot), el cual compara los cuantiles de los residuos observados con los cuantiles teóricos de una distribución normal con la misma media y varianza. Una alineación aproximada de los puntos sobre la recta de 45° indica que el supuesto de normalidad resulta razonable dentro del contexto del modelo.

A continuación, estos diagnósticos se aplican a los modelos previamente ajustados. Para ello, se utiliza el paquete \texttt{svydiags} \citep{Valliant2024}, desarrollado como extensión del paquete \texttt{survey} para el diagnóstico de modelos estimados bajo diseños muestrales complejos. Este paquete permite obtener los residuales estandarizados de manera directa, facilitando la evaluación de los supuestos del modelo en este tipo de contextos:

\begin{Shaded}
\begin{Highlighting}[]
\FunctionTok{library}\NormalTok{(svydiags)}
\NormalTok{stdresids }\OtherTok{\textless{}{-}} \FunctionTok{as.numeric}\NormalTok{(}\FunctionTok{svystdres}\NormalTok{(fit\_svy)}\SpecialCharTok{$}\NormalTok{stdresids)}
\NormalTok{survey\_design}\SpecialCharTok{$}\NormalTok{variables }\OtherTok{\textless{}{-}}\NormalTok{ survey\_design}\SpecialCharTok{$}\NormalTok{variables }\SpecialCharTok{\%\textgreater{}\%}
  \FunctionTok{mutate}\NormalTok{(}\AttributeTok{stdresids =}\NormalTok{ stdresids)}
\end{Highlighting}
\end{Shaded}

El análisis de normalidad puede complementarse con el histograma de los residuales estandarizados, presentado en la Figura \ref{fig:fig-hist-residuales}. El histograma de los residuos estandarizados revela una falta de ajuste significativa respecto a la curva de distribución normal teórica (en rojo), lo cual corrobora una violación del supuesto de normalidad. La distribución de los residuos (representada en azul) muestra un pico mucho más agudo y estrecho que la curva normal, además de una asimetría positiva marcada con una cola extendida hacia la derecha. Esta discrepancia indica que el modelo no logra capturar la variabilidad real de los datos, lo que sugiere la necesidad de aplicar una transformación de la variable dependiente o emplear un modelo con una familia de distribución más flexible y robusta.

\begin{Shaded}
\begin{Highlighting}[]
\FunctionTok{ggplot}\NormalTok{(}
  \AttributeTok{data =}\NormalTok{ survey\_design}\SpecialCharTok{$}\NormalTok{variables,}
  \FunctionTok{aes}\NormalTok{(}\AttributeTok{x =}\NormalTok{ stdresids)}
\NormalTok{) }\SpecialCharTok{+}
  \FunctionTok{geom\_histogram}\NormalTok{(}
    \FunctionTok{aes}\NormalTok{(}\AttributeTok{y =}\NormalTok{ ..density..),}
    \AttributeTok{colour =} \StringTok{"black"}\NormalTok{,}
    \AttributeTok{fill =} \StringTok{"blue"}\NormalTok{,}
    \AttributeTok{alpha =} \FloatTok{0.3}
\NormalTok{  ) }\SpecialCharTok{+}
  \FunctionTok{geom\_density}\NormalTok{(}\AttributeTok{size =} \DecValTok{1}\NormalTok{, }\AttributeTok{colour =} \StringTok{"blue"}\NormalTok{) }\SpecialCharTok{+}
  \FunctionTok{geom\_function}\NormalTok{(}\AttributeTok{fun =}\NormalTok{ dnorm, }\AttributeTok{colour =} \StringTok{"red"}\NormalTok{, }\AttributeTok{size =} \DecValTok{1}\NormalTok{) }\SpecialCharTok{+}
  \FunctionTok{labs}\NormalTok{(}\AttributeTok{y =} \StringTok{""}\NormalTok{)}
\end{Highlighting}
\end{Shaded}

\begin{figure}

{\centering \includegraphics[width=0.65\linewidth]{docs/figures/fig-hist-residuales-1} 

}

\caption{Histograma de residuales estandarizados con curva de densidad estimada y distribución normal teórica}\label{fig:fig-hist-residuales}
\end{figure}

La homocedasticidad de los errores (varianza constante), constituye uno de los supuestos centrales en los modelos de regresión. Su incumplimiento afecta la eficiencia de los estimadores del modelo. Dado que los residuos representan las discrepancias entre los valores observados y los valores ajustados por el modelo, su análisis en función de los valores predichos o de las covariables constituye una herramienta diagnóstica fundamental. Bajo un modelo adecuadamente especificado, se espera que los residuos se distribuyan de forma aleatoria alrededor de cero, sin patrones estructurados; en cambio, la aparición de formas sistemáticas (como estructuras en embudo o curvaturas) sugiere la presencia de heterocedasticidad (varianza no constante) o posibles relaciones funcionales no captadas por la especificación del modelo.

Para evaluar este supuesto, se recomienda analizar gráficamente los residuos en función de los valores estimados \(\hat{\mu}_{k}\) o de covariables específicas del modelo. La aparición de patrones sistemáticos en estos gráficos constituye evidencia de heterocedasticidad o de posibles problemas de especificación funcional. En particular, para evaluar la homocedasticidad respecto de las covariables incluidas en el modelo, se construyen gráficos de residuos frente a cada una de ellas y se combinan con \texttt{patchwork} \citep{patchwork}.

\begin{Shaded}
\begin{Highlighting}[]
\FunctionTok{library}\NormalTok{(patchwork)}

\NormalTok{g1 }\OtherTok{\textless{}{-}} \FunctionTok{ggplot}\NormalTok{(}
  \AttributeTok{data =}\NormalTok{ survey\_design}\SpecialCharTok{$}\NormalTok{variables,}
  \FunctionTok{aes}\NormalTok{(}\AttributeTok{x =}\NormalTok{ Expenditure, }\AttributeTok{y =}\NormalTok{ stdresids)}
\NormalTok{) }\SpecialCharTok{+}
  \FunctionTok{geom\_point}\NormalTok{() }\SpecialCharTok{+}
  \FunctionTok{geom\_hline}\NormalTok{(}\AttributeTok{yintercept =} \DecValTok{0}\NormalTok{)}

\NormalTok{g2 }\OtherTok{\textless{}{-}} \FunctionTok{ggplot}\NormalTok{(}
  \AttributeTok{data =}\NormalTok{ survey\_design}\SpecialCharTok{$}\NormalTok{variables,}
  \FunctionTok{aes}\NormalTok{(}\AttributeTok{x =}\NormalTok{ Zone, }\AttributeTok{y =}\NormalTok{ stdresids)}
\NormalTok{) }\SpecialCharTok{+}
  \FunctionTok{geom\_point}\NormalTok{() }\SpecialCharTok{+}
  \FunctionTok{geom\_hline}\NormalTok{(}\AttributeTok{yintercept =} \DecValTok{0}\NormalTok{)}

\NormalTok{g3 }\OtherTok{\textless{}{-}} \FunctionTok{ggplot}\NormalTok{(}
  \AttributeTok{data =}\NormalTok{ survey\_design}\SpecialCharTok{$}\NormalTok{variables,}
  \FunctionTok{aes}\NormalTok{(}\AttributeTok{x =}\NormalTok{ Sex, }\AttributeTok{y =}\NormalTok{ stdresids)}
\NormalTok{) }\SpecialCharTok{+}
  \FunctionTok{geom\_point}\NormalTok{() }\SpecialCharTok{+}
  \FunctionTok{geom\_hline}\NormalTok{(}\AttributeTok{yintercept =} \DecValTok{0}\NormalTok{)}

\NormalTok{(g1 }\SpecialCharTok{|}\NormalTok{ g2 }\SpecialCharTok{|}\NormalTok{ g3)}
\end{Highlighting}
\end{Shaded}

\begin{figure}

{\centering \includegraphics[width=0.8\linewidth]{docs/figures/fig-residuales-covariables-1} 

}

\caption{Gráficos de residuales estandarizados frente a cada covariable del modelo}\label{fig:fig-residuales-covariables}
\end{figure}

La Figura \ref{fig:fig-residuales-covariables} presenta los gráficos de dispersión de los residuos estandarizados frente a las covariables del modelo. En conjunto, estos resultados sugieren posibles limitaciones en la especificación funcional. En particular, para \texttt{Expenditure} se observa una ligera curvatura en la tendencia de los residuos, lo que podría indicar la presencia de una relación no lineal no completamente capturada por el modelo. Por su parte, en los gráficos asociados a \texttt{Zone} y \texttt{Sex} se aprecia una dispersión desigual de los residuos entre categorías, lo que sugiere diferencias en la variabilidad condicional o en el ajuste de la media entre grupos.

\subsection{Observaciones influyentes}\label{observaciones-influyentes}

En el análisis diagnóstico, la identificación de observaciones influyentes es una técnica de especial importancia. Estas son unidades muestrales cuyo impacto sobre el ajuste del modelo resulta desproporcionado respecto al resto de la muestra. Conviene subrayar que una observación influyente no necesariamente corresponde a un valor atípico: mientras un atípico puede estar alejado del patrón general de los datos sin afectar sustancialmente el ajuste, una observación influyente puede alterar de manera significativa las estimaciones de los parámetros, incluso si no luce inusual.

La distancia de Cook mide el impacto que tendría la eliminación de la observación \(k\) sobre el ajuste global del modelo, al combinar simultáneamente la magnitud del residuo, la varianza estimada del modelo y su apalancamiento (\emph{leverage}). En el contexto de encuestas con diseños muestrales complejos, la distancia de Cook se denota como \(c_k\) y su expresión incorpora explícitamente los pesos muestrales \citep[página 213]{heeringa2017applied}. Para evaluar la magnitud de este estadístico, se utiliza una aproximación basada en su comparación con una distribución F, mediante la expresión

\[
\frac{(df-p+1)\,c_{k}}{df} \sim F_{(p,\,df-p)},
\]

donde \(p\) denota los parámetros del modelo y \(df\) representa los grados de libertad asociados al diseño muestral. En la práctica, se considera que una observación puede ser influyente si su distancia de Cook supera valores de referencia aproximados entre 2 y 3, según lo propuesto en la literatura \citep{heeringa2017applied}.

En \texttt{R}, se evalúa la presencia de observaciones influyentes mediante la distancia de Cook, estimada con la función \texttt{svyCooksD} del paquete \texttt{svydiags}. Los resultados obtenidos se presentan en la Figura \ref{fig:fig-cook}, en donde se observa que ninguna de las distancias de Cook supera el umbral de referencia de 3, por lo que, bajo este criterio, no se identifican observaciones con influencia significativa sobre el ajuste del modelo.

\begin{Shaded}
\begin{Highlighting}[]
\FunctionTok{svyCooksD}\NormalTok{(fit\_svy, }\AttributeTok{doplot =} \ConstantTok{TRUE}\NormalTok{) }\SpecialCharTok{\%\textgreater{}\%}
  \FunctionTok{head}\NormalTok{(}\DecValTok{10}\NormalTok{)}
\end{Highlighting}
\end{Shaded}

\begin{figure}

{\centering \includegraphics[width=0.65\linewidth]{docs/figures/fig-cook-1} 

}

\caption{Distancia de Cook para cada observación de la muestra}\label{fig:fig-cook}
\end{figure}

\begin{verbatim}
##         1         2         3         4         5         6         7         8 
## 0.0008667 0.0013961 0.0013979 0.0010179 0.0016359 0.0016423 0.0018452 0.0013093 
##         9        10 
## 0.0021197 0.0012425
\end{verbatim}

Otro estadístico que mide el grado de influencia de una observación es el DFBETA, denotado como \(D_f\text{Beta}_{(k)j}\), que mide el cambio que experimenta el coeficiente de regresión \(\hat{\beta}_{j}\) cuando la observación \(k\) es eliminada del proceso de estimación. En consecuencia, esta medida evalúa la influencia individual de cada observación sobre los parámetros del modelo, identificando aquellas unidades cuya presencia altera de manera importante las estimaciones obtenidas. Valores elevados de \(D_f\text{Beta}_{(k)j}\) indican que la observación ejerce una influencia considerable sobre el coeficiente correspondiente, lo que puede afectar la estabilidad e interpretación del modelo.

Segun \citet{Li_Valliant_2015}, como criterio práctico de diagnóstico, una observación suele considerarse influyente sobre el coeficiente \(\hat{\beta}_{j}\) cuando

\[
\left|D_f\text{Beta}_{(k)j}\right|
\geq
\frac{z}{\sqrt{n_I\times DEFF}},
\]

donde \(z\) toma comúnmente valores de 2 o 3, \(n_I\) representa el número de UPMs seleccionadas en la muestra de la primera etapa y \(DEFF\) corresponde al efecto de diseño. A continuación, se evalúa la influencia de las observaciones sobre los coeficientes del modelo de ejemplo mediante la función \texttt{svydfbetas}. Los resultados correspondientes a las primeras diez observaciones se presentan en la Tabla \ref{tab:tab-dfbetas-primeros}.

\begin{Shaded}
\begin{Highlighting}[]
\NormalTok{d\_dfbetas }\OtherTok{\textless{}{-}} \FunctionTok{data.frame}\NormalTok{(}\FunctionTok{t}\NormalTok{(}\FunctionTok{svydfbetas}\NormalTok{(fit\_svy)}\SpecialCharTok{$}\NormalTok{Dfbetas))}
\FunctionTok{colnames}\NormalTok{(d\_dfbetas) }\OtherTok{\textless{}{-}} \FunctionTok{paste0}\NormalTok{(}\StringTok{"Beta\_"}\NormalTok{, }\DecValTok{1}\SpecialCharTok{:}\DecValTok{4}\NormalTok{)}
\NormalTok{d\_dfbetas }\SpecialCharTok{\%\textgreater{}\%}
  \FunctionTok{slice}\NormalTok{(}\DecValTok{1}\SpecialCharTok{:}\DecValTok{10}\DataTypeTok{L}\NormalTok{)}
\end{Highlighting}
\end{Shaded}

\begin{longtable}[]{@{}lccc@{}}
\caption{\label{tab:tab-dfbetas-primeros}Valores DfBetas para las primeras diez observaciones de la muestra}\tabularnewline
\toprule\noalign{}
Beta\_1 & Beta\_2 & Beta\_3 & Beta\_4 \\
\midrule\noalign{}
\endfirsthead
\toprule\noalign{}
Beta\_1 & Beta\_2 & Beta\_3 & Beta\_4 \\
\midrule\noalign{}
\endhead
\bottomrule\noalign{}
\endlastfoot
0.000 & 0 & 0.001 & -0.003 \\
-0.001 & 0 & 0.001 & 0.002 \\
-0.001 & 0 & 0.001 & 0.002 \\
0.000 & 0 & 0.001 & -0.003 \\
-0.001 & 0 & 0.001 & 0.002 \\
0.003 & 0 & -0.003 & -0.008 \\
0.003 & 0 & -0.003 & -0.008 \\
0.000 & 0 & -0.003 & 0.008 \\
0.003 & 0 & -0.003 & -0.008 \\
-0.001 & 0 & 0.001 & -0.003 \\
\end{longtable}

Una vez obtenidos los valores de la medida DFBETA para cada observación y coeficiente del modelo, se calcula el correspondiente umbral de influencia a partir de los resultados generados por la función \texttt{svydfbetas}. Con base en este criterio, se construye una variable dicotómica que permite clasificar las observaciones según su nivel de influencia sobre las estimaciones del modelo. Los resultados de esta clasificación se presentan en la Tabla \ref{tab:tab-dfbetas-influyentes}.

\begin{Shaded}
\begin{Highlighting}[]
\NormalTok{d\_dfbetas}\SpecialCharTok{$}\NormalTok{id }\OtherTok{\textless{}{-}} \DecValTok{1}\SpecialCharTok{:}\FunctionTok{nrow}\NormalTok{(d\_dfbetas)}
\NormalTok{d\_dfbetas }\OtherTok{\textless{}{-}}\NormalTok{ reshape2}\SpecialCharTok{::}\FunctionTok{melt}\NormalTok{(d\_dfbetas, }\AttributeTok{id.vars =} \StringTok{"id"}\NormalTok{)}
\NormalTok{cutoff }\OtherTok{\textless{}{-}} \FunctionTok{svydfbetas}\NormalTok{(fit\_svy)}\SpecialCharTok{$}\NormalTok{cutoff}
\NormalTok{d\_dfbetas }\OtherTok{\textless{}{-}}\NormalTok{ d\_dfbetas }\SpecialCharTok{\%\textgreater{}\%}
  \FunctionTok{mutate}\NormalTok{(}\AttributeTok{criterion =} \FunctionTok{ifelse}\NormalTok{(}\FunctionTok{abs}\NormalTok{(value) }\SpecialCharTok{\textgreater{}}\NormalTok{ cutoff, }\StringTok{"Yes"}\NormalTok{, }\StringTok{"No"}\NormalTok{))}

\NormalTok{label\_text }\OtherTok{\textless{}{-}}\NormalTok{ d\_dfbetas }\SpecialCharTok{\%\textgreater{}\%}
  \FunctionTok{filter}\NormalTok{(criterion }\SpecialCharTok{==} \StringTok{"Yes"}\NormalTok{) }\SpecialCharTok{\%\textgreater{}\%}
  \FunctionTok{arrange}\NormalTok{(}\FunctionTok{desc}\NormalTok{(}\FunctionTok{abs}\NormalTok{(value))) }\SpecialCharTok{\%\textgreater{}\%}
  \FunctionTok{slice}\NormalTok{(}\DecValTok{1}\SpecialCharTok{:}\DecValTok{10}\DataTypeTok{L}\NormalTok{)}

\NormalTok{label\_text}
\end{Highlighting}
\end{Shaded}

\begin{longtable}[]{@{}lccc@{}}
\caption{\label{tab:tab-dfbetas-influyentes}Las diez observaciones más influyentes según el criterio DfBetas}\tabularnewline
\toprule\noalign{}
id & variable & value & criterion \\
\midrule\noalign{}
\endfirsthead
\toprule\noalign{}
id & variable & value & criterion \\
\midrule\noalign{}
\endhead
\bottomrule\noalign{}
\endlastfoot
889 & Beta\_1 & 0.284 & Yes \\
891 & Beta\_1 & 0.284 & Yes \\
890 & Beta\_2 & -0.280 & Yes \\
889 & Beta\_2 & -0.275 & Yes \\
891 & Beta\_2 & -0.275 & Yes \\
890 & Beta\_1 & 0.250 & Yes \\
890 & Beta\_3 & 0.161 & Yes \\
890 & Beta\_4 & 0.158 & Yes \\
889 & Beta\_3 & 0.156 & Yes \\
891 & Beta\_3 & 0.156 & Yes \\
\end{longtable}

Los resultados muestran que varias observaciones superan el umbral establecido por el criterio \(D_f\text{Beta}_{(k)j}\), lo que indica que dichas unidades ejercen una influencia importante sobre la estimación de algunos coeficientes del modelo, lo que sugiere la presencia de unidades con alto impacto sobre la estabilidad del ajuste. La Figura \ref{fig:fig-dfbetas} presenta la representación gráfica de los valores de \(D_f\)Beta junto con el umbral de referencia utilizado para identificar observaciones influyentes. En la figura, los puntos resaltados en rojo corresponden a las observaciones que exceden dicho umbral y que, por tanto, requieren una revisión más detallada dentro del proceso de diagnóstico del modelo.

\begin{Shaded}
\begin{Highlighting}[]
\FunctionTok{ggplot}\NormalTok{(d\_dfbetas, }\FunctionTok{aes}\NormalTok{(}\AttributeTok{y =} \FunctionTok{abs}\NormalTok{(value), }\AttributeTok{x =}\NormalTok{ id)) }\SpecialCharTok{+}
  \FunctionTok{geom\_point}\NormalTok{(}\FunctionTok{aes}\NormalTok{(}\AttributeTok{col =}\NormalTok{ criterion)) }\SpecialCharTok{+}
  \FunctionTok{geom\_text}\NormalTok{(}
    \AttributeTok{data =}\NormalTok{ label\_text,}
    \AttributeTok{angle =} \DecValTok{45}\NormalTok{,}
    \AttributeTok{vjust =} \SpecialCharTok{{-}}\DecValTok{1}\NormalTok{,}
    \FunctionTok{aes}\NormalTok{(}\AttributeTok{label =}\NormalTok{ id)}
\NormalTok{  ) }\SpecialCharTok{+}
  \FunctionTok{geom\_hline}\NormalTok{(}\FunctionTok{aes}\NormalTok{(}\AttributeTok{yintercept =}\NormalTok{ cutoff)) }\SpecialCharTok{+}
  \FunctionTok{facet\_wrap}\NormalTok{(. }\SpecialCharTok{\textasciitilde{}}\NormalTok{ variable, }\AttributeTok{nrow =} \DecValTok{2}\NormalTok{) }\SpecialCharTok{+}
  \FunctionTok{scale\_color\_manual}\NormalTok{(}
    \AttributeTok{values =} \FunctionTok{c}\NormalTok{(}\StringTok{"Yes"} \OtherTok{=} \StringTok{"red"}\NormalTok{, }\StringTok{"No"} \OtherTok{=} \StringTok{"black"}\NormalTok{)}
\NormalTok{  )}
\end{Highlighting}
\end{Shaded}

\begin{figure}

{\centering \includegraphics[width=0.8\linewidth]{docs/figures/fig-dfbetas-1} 

}

\caption{Valores DfBetas por parámetro: observaciones influyentes señaladas en rojo}\label{fig:fig-dfbetas}
\end{figure}

Por último, el estadístico DFFITS, denotado por \(D_f\text{Fits}_{(k)}\), mide la influencia que ejerce la observación \(k\) sobre los valores ajustados del modelo. En particular, este estadístico evalúa cuánto cambian las predicciones del modelo cuando una observación es eliminada del proceso de estimación. En consecuencia, valores elevados de \(D_f\text{Fits}_{(k)}\) indican observaciones con capacidad de alterar de manera importante el ajuste global del modelo y, por tanto, potencialmente influyentes en el proceso de inferencia. Segun \citet{Li_Valliant_2015}, una observación se considera influyente si:

\[
|D_f\text{Fits}_{(k)}| \geq z\sqrt{\frac{p}{n\times DEFF}}
\]
donde \(z\) toma comúnmente valores de 2 o 3, mientras que \(n\) representa el tamaño muestral total de elementos incluidos en la encuesta. En el contexto de encuestas complejas, el uso de umbrales ajustados por efecto de diseño resulta especialmente relevante, ya que incorpora la pérdida de precisión derivada de la estratificación, la conglomeración y la utilización de pesos desiguales.

En \texttt{R}, el cálculo de este estadístico puede realizarse mediante la función \texttt{svydffits} del paquete \texttt{svydiags}, la cual obtiene las medidas de influencia considerando explícitamente la estructura del diseño muestral. Los resultados gráficos correspondientes se presentan en la Figura \ref{fig:fig-dffits}, en donde se confirma la existencia de algunas observaciones influyentes (señaladas en rojo).

\begin{Shaded}
\begin{Highlighting}[]
\NormalTok{d\_dffits }\OtherTok{\textless{}{-}} \FunctionTok{data.frame}\NormalTok{(}
  \AttributeTok{dffits =} \FunctionTok{svydffits}\NormalTok{(fit\_svy)}\SpecialCharTok{$}\NormalTok{Dffits,}
  \AttributeTok{id =} \DecValTok{1}\SpecialCharTok{:}\FunctionTok{length}\NormalTok{(}\FunctionTok{svydffits}\NormalTok{(fit\_svy)}\SpecialCharTok{$}\NormalTok{Dffits)}
\NormalTok{)}
\NormalTok{cutoff }\OtherTok{\textless{}{-}} \FunctionTok{svydffits}\NormalTok{(fit\_svy)}\SpecialCharTok{$}\NormalTok{cutoff}
\NormalTok{d\_dffits }\OtherTok{\textless{}{-}}\NormalTok{ d\_dffits }\SpecialCharTok{\%\textgreater{}\%}
  \FunctionTok{mutate}\NormalTok{(}\AttributeTok{cutoff\_criterion =} \FunctionTok{ifelse}\NormalTok{(}\FunctionTok{abs}\NormalTok{(dffits) }\SpecialCharTok{\textgreater{}}\NormalTok{ cutoff, }\StringTok{"Yes"}\NormalTok{, }\StringTok{"No"}\NormalTok{))}

\FunctionTok{ggplot}\NormalTok{(d\_dffits, }\FunctionTok{aes}\NormalTok{(}\AttributeTok{y =} \FunctionTok{abs}\NormalTok{(dffits), }\AttributeTok{x =}\NormalTok{ id)) }\SpecialCharTok{+}
  \FunctionTok{geom\_point}\NormalTok{(}\FunctionTok{aes}\NormalTok{(}\AttributeTok{col =}\NormalTok{ cutoff\_criterion)) }\SpecialCharTok{+}
  \FunctionTok{geom\_hline}\NormalTok{(}\AttributeTok{yintercept =}\NormalTok{ cutoff) }\SpecialCharTok{+}
  \FunctionTok{scale\_color\_manual}\NormalTok{(}
    \AttributeTok{values =} \FunctionTok{c}\NormalTok{(}\StringTok{"Yes"} \OtherTok{=} \StringTok{"red"}\NormalTok{, }\StringTok{"No"} \OtherTok{=} \StringTok{"black"}\NormalTok{)}
\NormalTok{  )}
\end{Highlighting}
\end{Shaded}

\begin{figure}

{\centering \includegraphics[width=0.65\linewidth]{docs/figures/fig-dffits-1} 

}

\caption{Valores DfFits: observaciones influyentes sobre el ajuste global señaladas en rojo}\label{fig:fig-dffits}
\end{figure}

\section{Inferencia sobre los parámetros del modelo}\label{inferencia-sobre-los-parametros-del-modelo}

Una vez evaluado el ajuste del modelo y verificado el cumplimiento de los supuestos asociados a los errores, el siguiente paso consiste en determinar la significación estadística de los parámetros estimados. Este análisis permite establecer si las covariables incluidas en el modelo aportan información relevante para explicar la variabilidad de la variable respuesta.

En modelos ajustados con datos provenientes de encuestas complejas, las pruebas de hipótesis sobre los coeficientes de regresión se construyen utilizando estimadores de varianza consistentes con el diseño muestral. Para evaluar el efecto de una covariable asociada al parámetro \(\beta_j\), se considera usualmente el contraste

\[
H_0:\beta_j=0
\qquad \text{frente a} \qquad
H_1:\beta_j\neq 0.
\]

El estadístico de prueba correspondiente se define como

\[
t_j=
\frac{\hat{\beta}_{j}-\beta_{j}}
{\sqrt{\widehat{Var}(\hat{\beta}_{j})}}
\]

el cual, bajo la hipótesis nula sigue aproximadamente una distribución \(t\) de Student con \(df\) grados de libertad asociados al diseño muestral (usualmente definidos como el número de unidades primarias de muestreo menos el número de estratos). La decisión inferencial se basa en comparar el valor absoluto del estadístico \(t_j\) con el valor crítico de la distribución de referencia. Cuando el valor observado del estadístico resulta suficientemente grande en magnitud, se rechaza la hipótesis nula y se concluye que la covariable correspondiente presenta una asociación estadísticamente significativa con la variable respuesta.

Las mismas propiedades distribucionales permiten construir intervalos de confianza para los coeficientes de regresión. Así, un intervalo de confianza al nivel \((1-\alpha)\times100\%\) para el parámetro \(\beta_j\) está dado por

\[
\hat{\beta}_{j}
\pm
t_{1-\frac{\alpha}{2},\,df}
\sqrt{\widehat{Var}(\hat{\beta}_{j})},
\]

donde \(t_{1-\frac{\alpha}{2},\,df}\) representa el cuantil de orden \(1-\alpha/2\) de la distribución \(t\) de Student con \(df\) grados de libertad.

Para implementar estos procedimientos en \texttt{R}, se utilizan las funciones \texttt{summary.svyglm} y \texttt{confint.svyglm}, las cuales permiten obtener, respectivamente, las pruebas de significación basadas en el estadístico \(t\) y los intervalos de confianza de los coeficientes estimados bajo el diseño muestral complejo. Los resultados correspondientes se presentan en la Tabla \ref{tab:tab-inferencia-modelo}.

\begin{Shaded}
\begin{Highlighting}[]
\NormalTok{summary\_table }\OtherTok{\textless{}{-}} \FunctionTok{summary}\NormalTok{(fit\_svy)}\SpecialCharTok{$}\NormalTok{coefficients }\SpecialCharTok{\%\textgreater{}\%}
  \FunctionTok{as.data.frame}\NormalTok{() }\SpecialCharTok{\%\textgreater{}\%}
\NormalTok{  tibble}\SpecialCharTok{::}\FunctionTok{rownames\_to\_column}\NormalTok{(}\StringTok{"Parameter"}\NormalTok{)}
\end{Highlighting}
\end{Shaded}

\begin{longtable}[]{@{}lcccc@{}}
\caption{\label{tab:tab-inferencia-modelo}Pruebas t e intervalos de confianza al 95\% para los parámetros del modelo}\tabularnewline
\toprule\noalign{}
Parameter & Estimate & Std. Error & t value & Pr(\textgreater\textbar t\textbar) \\
\midrule\noalign{}
\endfirsthead
\toprule\noalign{}
Parameter & Estimate & Std. Error & t value & Pr(\textgreater\textbar t\textbar) \\
\midrule\noalign{}
\endhead
\bottomrule\noalign{}
\endlastfoot
(Intercept) & 73.580 & 59.470 & 1.237 & 0.218 \\
Expenditure & 1.222 & 0.197 & 6.212 & 0.000 \\
ZoneUrban & 66.652 & 39.666 & 1.680 & 0.096 \\
SexMale & 20.644 & 15.611 & 1.322 & 0.189 \\
\end{longtable}

Los resultados del modelo indican que la variable \texttt{Expenditure} presenta una asociación estadísticamente significativa con la variable de interés (\texttt{Income}). Este efecto resulta altamente significativo desde el punto de vista estadístico (\(p < 0.001\)), lo que proporciona evidencia sólida de una relación lineal positiva entre ambas variables. Por su parte, las variables categóricas \texttt{ZoneUrban} y \texttt{SexMale}, su nivel de significación estadística es marginal (\(p = 0.096\) y \(p = 0.189\), respectivamente), por lo que la evidencia no resulta concluyente al nivel convencional del 5\%.

\section{Estimación y predicción}\label{estimacion-y-prediccion}

Segun \citet{heeringa2017applied}, los modelos de regresión cumplen dos objetivos fundamentales. El primero consiste en explicar la variabilidad de la variable respuesta a partir de las covariables disponibles, propósito que ha sido desarrollado a lo largo de este capítulo. El segundo corresponde a la predicción del valor esperado de la variable de interés para nuevas observaciones, tanto dentro como fuera del dominio muestral. La siguiente expresión representa el valor esperado de la variable respuesta \(y_k\) condicionada al vector de covariables observadas \(\mathbf{x}_{k}\). En otras palabras, el modelo utiliza la información disponible de las variables explicativas para estimar el valor promedio esperado de la variable de interés para la unidad \(i\).

\[
\hat \mu_k = \hat{E}(y_{k}\mid\mathbf{x}_{k})=\mathbf{x}_{k}\hat{\boldsymbol{\beta}}
\]

Asimismo, la varianza asociada a la estimación del valor esperado, definida como \({Var}(\hat{\mu}_k) = {Var}(\mathbf{x}_{k}\hat{\boldsymbol{\beta}})\), puede estimarse mediante la siguiente expresión:

\[
\widehat{Var}(\hat \mu_k) =
\mathbf{x}'_{k}\ \hat{\boldsymbol{\Sigma}}\ \mathbf{x}_{k}
\]

En donde \(\hat{\boldsymbol{\Sigma}}=\widehat{Var}(\hat{\boldsymbol{\beta}})\) corresponde a la matriz de varianzas y covarianzas estimada de los parámetros del modelo de regresión. Esta expresión cuantifica la incertidumbre asociada a la predicción obtenida para la unidad \(k\), considerando simultáneamente la variabilidad de cada coeficiente estimado y la correlación existente entre ellos. En consecuencia, aquellas unidades cuyos valores de las covariables estén asociados a regiones con mayor incertidumbre en la estimación de los coeficientes presentarán varianzas más elevadas. Por tanto, la precisión de las estimaciones depende directamente de la precisión con la que se estimen los coeficientes del modelo de regresión. Continuando con el modelo del ejemplo, recordemos que los coeficientes estimados del modelo se presentan en la Tabla \ref{tab:tab-coeficientes-modelo}:

Los coeficientes del modelo se convierten a una tabla ordenada mediante \texttt{broom::tidy()} \citep{broom2026}, lo que facilita presentar estimaciones, errores estándar y pruebas en una estructura rectangular.

\begin{longtable}[]{@{}lcccc@{}}
\caption{\label{tab:tab-coeficientes-modelo}Coeficientes estimados del modelo de regresión con diseño muestral complejo}\tabularnewline
\toprule\noalign{}
term & estimate & std.error & statistic & p.value \\
\midrule\noalign{}
\endfirsthead
\toprule\noalign{}
term & estimate & std.error & statistic & p.value \\
\midrule\noalign{}
\endhead
\bottomrule\noalign{}
\endlastfoot
(Intercept) & 73.580 & 59.470 & 1.237 & 0.218 \\
Expenditure & 1.222 & 0.197 & 6.212 & 0.000 \\
ZoneUrban & 66.652 & 39.666 & 1.680 & 0.096 \\
SexMale & 20.644 & 15.611 & 1.322 & 0.189 \\
\end{longtable}

A partir de los coeficientes estimados del modelo, la estimación del valor esperado de la variable respuesta para la unidad \(k\) puede expresarse como:

\[
\hat{E}(y_{k}\mid\mathbf{x}_{k})=73.58+1.22 \times x_{1i}+66.65 \times x_{2i}+20.64  \times  x_{3i}
\]

El siguiente código calcula las predicciones del modelo de regresión para las observaciones incluidas en la muestra, utilizando la matriz de diseño y el vector de coeficientes estimados:

\begin{Shaded}
\begin{Highlighting}[]
\NormalTok{x\_obs }\OtherTok{\textless{}{-}} \FunctionTok{model.matrix}\NormalTok{(fit\_svy) }\SpecialCharTok{\%\textgreater{}\%}
  \FunctionTok{data.frame}\NormalTok{() }\SpecialCharTok{\%\textgreater{}\%}
  \FunctionTok{as.matrix}\NormalTok{()}
\NormalTok{hat\_beta }\OtherTok{\textless{}{-}}\NormalTok{ fit\_svy}\SpecialCharTok{$}\NormalTok{coe}
\NormalTok{hat\_mu }\OtherTok{\textless{}{-}} \FunctionTok{as.numeric}\NormalTok{(x\_obs }\SpecialCharTok{\%*\%}\NormalTok{ hat\_beta)}
\NormalTok{hat\_mu }\SpecialCharTok{\%\textgreater{}\%}
  \FunctionTok{head}\NormalTok{(}\DecValTok{10}\NormalTok{)}
\end{Highlighting}
\end{Shaded}

\begin{verbatim}
##  [1] 517.6 496.9 496.9 517.6 496.9 553.1 553.1 573.7 553.1 184.8
\end{verbatim}

En primer lugar, la instrucción \texttt{model.matrix(fit\_svy)} construye la matriz de diseño asociada al modelo ajustado, incorporando automáticamente las variables explicativas. A continuación, el objeto \texttt{hat\_beta} almacena el vector de coeficientes estimados del modelo de regresión. Posteriormente, se obtienen las estimaciones mediante el producto matricial \texttt{x\_obs\ \%*\%\ hat\_beta}. Finalmente, la función \texttt{head(10)} permite visualizar las primeras diez estimaciones obtenidas.

Para calcular la varianza de las estimaciones se requiere la matriz de varianzas y covarianzas de los parámetros del modelo. Los elementos de la diagonal principal corresponden a las varianzas de los parámetros estimados, mientras que los elementos fuera de la diagonal representan las covarianzas entre pares de coeficientes. La matriz estimada de varianzas y covarianzas se presenta en la Tabla \ref{tab:tab-vcov-modelo}:

\begin{Shaded}
\begin{Highlighting}[]
\FunctionTok{vcov}\NormalTok{(fit\_svy) }\SpecialCharTok{\%\textgreater{}\%}
  \FunctionTok{as.data.frame}\NormalTok{() }\SpecialCharTok{\%\textgreater{}\%}
\NormalTok{  tibble}\SpecialCharTok{::}\FunctionTok{rownames\_to\_column}\NormalTok{(}\StringTok{"Parameter"}\NormalTok{)}
\end{Highlighting}
\end{Shaded}

\begin{longtable}[]{@{}lcccc@{}}
\caption{\label{tab:tab-vcov-modelo}Matriz de varianzas y covarianzas de los coeficientes del modelo}\tabularnewline
\toprule\noalign{}
Parameter & (Intercept) & Expenditure & ZoneUrban & SexMale \\
\midrule\noalign{}
\endfirsthead
\toprule\noalign{}
Parameter & (Intercept) & Expenditure & ZoneUrban & SexMale \\
\midrule\noalign{}
\endhead
\bottomrule\noalign{}
\endlastfoot
(Intercept) & 3536.68 & -10.551 & 123.689 & 326.13 \\
Expenditure & -10.55 & 0.039 & -3.295 & -0.57 \\
ZoneUrban & 123.69 & -3.295 & 1573.353 & -121.45 \\
SexMale & 326.13 & -0.570 & -121.446 & 243.72 \\
\end{longtable}

En términos computacionales, esto implica combinar el vector de covariables de cada observación con la matriz \(\hat{\boldsymbol{\Sigma}}\) para obtener la precisión de las predicciones individuales. El siguiente código calcula la varianza asociada a las predicciones obtenidas a partir del modelo de regresión:

\begin{Shaded}
\begin{Highlighting}[]
\NormalTok{cov\_beta }\OtherTok{\textless{}{-}} \FunctionTok{vcov}\NormalTok{(fit\_svy) }\SpecialCharTok{\%\textgreater{}\%}
  \FunctionTok{as.matrix}\NormalTok{()}
\NormalTok{mu\_variance }\OtherTok{\textless{}{-}} \FunctionTok{as.numeric}\NormalTok{(x\_obs }\SpecialCharTok{\%*\%}\NormalTok{ cov\_beta }\SpecialCharTok{\%*\%} \FunctionTok{t}\NormalTok{(x\_obs))}
\NormalTok{mu\_variance }\SpecialCharTok{\%\textgreater{}\%}
  \FunctionTok{head}\NormalTok{(}\DecValTok{10}\NormalTok{)}
\end{Highlighting}
\end{Shaded}

\begin{verbatim}
##  [1] 1374.9 1002.6 1002.6 1374.9 1002.6 1107.8 1107.8 1480.1 1107.8  750.8
\end{verbatim}

El intervalo de confianza, que proporciona un rango de valores plausibles para el valor esperado de la variable respuesta condicionado a las covariables observadas, se obtiene mediante:

\[
\mathbf{x}_{k}\hat{\beta}\pm t_{\left(1-\frac{\alpha}{2},df\right)}\sqrt{\mathbf{x}'_{k} \
\hat{\boldsymbol{\Sigma}} \
\mathbf{x}_{k}}
\]

En R, la función \texttt{predict(fit\_svy,\ type\ =\ "link")} genera las predicciones lineales del modelo ajustado. El argumento \texttt{type\ =\ "link"} indica que las estimaciones se obtienen en la escala del predictor lineal, es decir, utilizando directamente la combinación lineal de las covariables y los coeficientes estimados. Posteriormente, la función \texttt{confint()} calcula los intervalos de confianza asociados a dichas predicciones.

\begin{Shaded}
\begin{Highlighting}[]
\NormalTok{hat\_mu }\OtherTok{\textless{}{-}} \FunctionTok{data.frame}\NormalTok{(}\FunctionTok{predict}\NormalTok{(fit\_svy, }\AttributeTok{type =} \StringTok{"link"}\NormalTok{))}
\NormalTok{mu\_ci }\OtherTok{\textless{}{-}} \FunctionTok{data.frame}\NormalTok{(}\FunctionTok{confint}\NormalTok{(}\FunctionTok{predict}\NormalTok{(fit\_svy, }\AttributeTok{type =} \StringTok{"link"}\NormalTok{)))}
\FunctionTok{colnames}\NormalTok{(mu\_ci) }\OtherTok{\textless{}{-}} \FunctionTok{c}\NormalTok{(}\StringTok{"lower\_limit"}\NormalTok{, }\StringTok{"upper\_limit"}\NormalTok{)}
\end{Highlighting}
\end{Shaded}

Las estimaciones obtenidas a partir del modelo, junto con sus respectivos intervalos de confianza, pueden visualizarse gráficamente en la Figura \ref{fig:fig-pred-intervalos}.

\begin{Shaded}
\begin{Highlighting}[]
\NormalTok{pred }\OtherTok{\textless{}{-}} \FunctionTok{cbind}\NormalTok{(hat\_mu, mu\_ci)}
\NormalTok{pred}\SpecialCharTok{$}\NormalTok{Expenditure }\OtherTok{\textless{}{-}}\NormalTok{ survey\_data}\SpecialCharTok{$}\NormalTok{Expenditure}
\NormalTok{pd }\OtherTok{\textless{}{-}} \FunctionTok{position\_dodge}\NormalTok{(}\AttributeTok{width =} \FloatTok{0.2}\NormalTok{)}
\FunctionTok{ggplot}\NormalTok{(}
\NormalTok{  pred }\SpecialCharTok{\%\textgreater{}\%}
    \FunctionTok{slice}\NormalTok{(}\DecValTok{1}\SpecialCharTok{:}\DecValTok{100}\DataTypeTok{L}\NormalTok{),}
  \FunctionTok{aes}\NormalTok{(}\AttributeTok{x =}\NormalTok{ Expenditure, }\AttributeTok{y =}\NormalTok{ link)}
\NormalTok{) }\SpecialCharTok{+}
  \FunctionTok{geom\_errorbar}\NormalTok{(}
    \FunctionTok{aes}\NormalTok{(}\AttributeTok{ymin =}\NormalTok{ lower\_limit, }\AttributeTok{ymax =}\NormalTok{ upper\_limit),}
    \AttributeTok{width =}\NormalTok{ .}\DecValTok{1}\NormalTok{,}
    \AttributeTok{linetype =} \DecValTok{1}
\NormalTok{  ) }\SpecialCharTok{+}
  \FunctionTok{geom\_point}\NormalTok{(}\AttributeTok{size =} \DecValTok{2}\NormalTok{, }\AttributeTok{position =}\NormalTok{ pd) }\SpecialCharTok{+}
  \FunctionTok{theme\_bw}\NormalTok{()}
\end{Highlighting}
\end{Shaded}

\begin{figure}

{\centering \includegraphics[width=0.65\linewidth]{docs/figures/fig-pred-intervalos-1} 

}

\caption{Predicciones puntuales e intervalos de confianza para las primeras 100 observaciones de la muestra}\label{fig:fig-pred-intervalos}
\end{figure}

Cuando el interés está en realizar predicciones para valores de las covariables que se encuentran fuera del rango observado en la muestra, la incertidumbre asociada a la estimación aumenta. En estos casos, no basta con considerar únicamente la variabilidad derivada de la estimación de los parámetros del modelo, sino que también debe incorporarse la variabilidad inherente al término aleatorio de la regresión. Por esta razón, la varianza de predicción incluye un componente adicional asociado al error residual del modelo:

\[
\widehat{Var}\left[\hat{E}\left(y_{k}\mid\mathbf{x}_{k}\right)\right]
=
\mathbf{x}'_{k} \
\hat{\boldsymbol{\Sigma}} \
\mathbf{x}_{k}
+
\hat{\sigma}^2_{yx}
\]

En esta expresión, el término \(\hat{\sigma}^2_{yx}\) corresponde a la varianza residual estimada, es decir, la parte de la variabilidad de la variable respuesta que no es explicada por las covariables incluidas en el modelo. La incorporación de este segundo componente es fundamental cuando se desea predecir sobre unidades no observadas, ya que reconoce explícitamente la existencia de variabilidad individual alrededor de la media estimada.

En consecuencia, los intervalos de predicción suelen ser más amplios que los intervalos de confianza para la media, reflejando una mayor incertidumbre asociada a la predicción de valores individuales. De esta manera, el intervalo de predicción para una nueva observación puede expresarse como:

\[
\mathbf{x}_{k}\hat{\boldsymbol{\beta}}
\pm
t_{\left(1-\frac{\alpha}{2},df\right)}
\sqrt{
\mathbf{x}_{k}^{t}
\hat{\boldsymbol{\Sigma}}_{\hat{\boldsymbol{\beta}}}
\mathbf{x}_{k}
+
\hat{\sigma}^2_{yx}
}
\]

Supóngase, por ejemplo, que se desea obtener la predicción del modelo para una nueva observación correspondiente a un individuo de sexo masculino (\texttt{Male}), residente en zona urbana (\texttt{Urban}) y con un nivel de gasto (\texttt{Expenditure}) igual a 1600. Este perfil puede definirse en \texttt{R} mediante el siguiente conjunto de datos:

\begin{Shaded}
\begin{Highlighting}[]
\NormalTok{new\_data }\OtherTok{\textless{}{-}} \FunctionTok{data.frame}\NormalTok{(}
  \AttributeTok{Expenditure =} \DecValTok{1600}\NormalTok{,}
  \AttributeTok{Sex =} \StringTok{"Male"}\NormalTok{,}
  \AttributeTok{Zone =} \StringTok{"Urban"}
\NormalTok{)}

\FunctionTok{predict}\NormalTok{(fit\_svy, }\AttributeTok{newdata =}\NormalTok{ new\_data, }\AttributeTok{type =} \StringTok{"link"}\NormalTok{)}
\end{Highlighting}
\end{Shaded}

\begin{verbatim}
##   link  SE
## 1 2117 243
\end{verbatim}

\begin{Shaded}
\begin{Highlighting}[]
\FunctionTok{confint}\NormalTok{(}\FunctionTok{predict}\NormalTok{(fit\_svy, }\AttributeTok{newdata =}\NormalTok{ new\_data))}
\end{Highlighting}
\end{Shaded}

\begin{verbatim}
##   2.5 % 97.5 %
## 1  1641   2593
\end{verbatim}

En primer lugar, el objeto \texttt{new\_data} especifica los valores de las covariables incluidas en el modelo. Posteriormente, la función \texttt{predict()} calcula la predicción puntual asociada a dicha observación. El argumento \texttt{newdata\ =\ new\_data} indica que la predicción debe realizarse utilizando la nueva información suministrada, mientras que \texttt{type\ =\ "link"} solicita el resultado en la escala del predictor lineal. Finalmente, la función \texttt{confint()} permite obtener el intervalo de confianza asociado a la predicción realizada. Este intervalo cuantifica la incertidumbre de la estimación y proporciona un rango de valores plausibles para el valor esperado de la variable respuesta correspondiente al perfil especificado.

\chapter{Modelos lineales generalizados}\label{cap6}

Cuando la variable de interés no es continua o los supuestos del modelo lineal clásico no se satisfacen, los modelos lineales generalizados constituyen una alternativa adecuada para modelar las relaciones de interés. Este enfoque fue introducido por \citet{nelder1972generalized}, quienes demostraron que diversos métodos estadísticos ampliamente utilizados, como la regresión logística, los modelos log-lineales para datos de conteo y ciertos modelos avanzados de regresión para variables continuas, podían integrarse dentro de un mismo marco teórico unificado. Esta formulación permitió extender las herramientas de regresión más allá de los escenarios en los que la variable respuesta sigue una distribución normal.

La necesidad de esta generalización surge porque los modelos lineales clásicos asumen, entre otras cosas, normalidad y homogeneidad de varianza en la variable respuesta, condiciones que rara vez se cumplen cuando se analizan variables categóricas, proporciones o conteos. En este tipo de situaciones, los modelos lineales generalizados ofrecen un marco flexible para modelar adecuadamente la relación entre las variables mediante la selección de una distribución de probabilidad apropiada para la respuesta y una función de enlace que conecte la media de dicha distribución con el predictor lineal.

Para ilustrar los modelos presentados en este capítulo se utilizan la base de datos \texttt{survey\_data} y la población \texttt{BigCity}, siguiendo la misma estructura utilizada en capítulos anteriores. Después de cargar \texttt{tidyverse} \citep{tidyverse2026}, \texttt{survey} \citep{Lumley2024}, \texttt{srvyr} \citep{Freedman2024}, \texttt{broom} \citep{broom2026} y \texttt{jtools} \citep{jtools2026}, el siguiente bloque de código crea dos nuevas variables categóricas (\texttt{poor} y \texttt{unemployed}) a partir de las variables originales \texttt{Poverty} y \texttt{Employment}, respectivamente.

\begin{Shaded}
\begin{Highlighting}[]
\FunctionTok{options}\NormalTok{(}\AttributeTok{digits =} \DecValTok{4}\NormalTok{)}
\FunctionTok{library}\NormalTok{(tidyverse)}
\FunctionTok{library}\NormalTok{(survey)}
\FunctionTok{library}\NormalTok{(srvyr)}
\FunctionTok{library}\NormalTok{(broom)}
\FunctionTok{library}\NormalTok{(jtools)}

\NormalTok{survey\_data }\OtherTok{\textless{}{-}} \FunctionTok{readRDS}\NormalTok{(}\StringTok{"Data/encuesta.rds"}\NormalTok{)}
\FunctionTok{data}\NormalTok{(}\StringTok{"BigCity"}\NormalTok{, }\AttributeTok{package =} \StringTok{"TeachingSampling"}\NormalTok{)}

\NormalTok{survey\_design }\OtherTok{\textless{}{-}}\NormalTok{ survey\_data }\SpecialCharTok{\%\textgreater{}\%}
  \FunctionTok{as\_survey\_design}\NormalTok{(}
    \AttributeTok{strata =}\NormalTok{ Stratum,}
    \AttributeTok{ids =}\NormalTok{ PSU,}
    \AttributeTok{weights =}\NormalTok{ wk,}
    \AttributeTok{nest =} \ConstantTok{TRUE}
\NormalTok{  )}

\NormalTok{survey\_design }\OtherTok{\textless{}{-}}\NormalTok{ survey\_design }\SpecialCharTok{\%\textgreater{}\%}
  \FunctionTok{mutate}\NormalTok{(}
    \AttributeTok{poor =} \FunctionTok{factor}\NormalTok{(}
      \FunctionTok{ifelse}\NormalTok{(Poverty }\SpecialCharTok{!=} \StringTok{"NotPoor"}\NormalTok{, }\DecValTok{1}\NormalTok{, }\DecValTok{0}\NormalTok{),}
      \AttributeTok{levels =} \FunctionTok{c}\NormalTok{(}\DecValTok{0}\NormalTok{, }\DecValTok{1}\NormalTok{),}
      \AttributeTok{labels =} \FunctionTok{c}\NormalTok{(}\StringTok{"=Poor"}\NormalTok{, }\StringTok{"=Not poor"}\NormalTok{)}
\NormalTok{    ),}
    \AttributeTok{unemployed =} \FunctionTok{factor}\NormalTok{(}\FunctionTok{ifelse}\NormalTok{(Employment }\SpecialCharTok{==} \StringTok{"Unemployed"}\NormalTok{, }\DecValTok{1}\NormalTok{, }\DecValTok{0}\NormalTok{),}
      \AttributeTok{levels =} \FunctionTok{c}\NormalTok{(}\DecValTok{0}\NormalTok{, }\DecValTok{1}\NormalTok{),}
      \AttributeTok{labels =} \FunctionTok{c}\NormalTok{(}\StringTok{"=Unemployed"}\NormalTok{, }\StringTok{"=Not unemployed"}\NormalTok{)}
\NormalTok{    )}
\NormalTok{  )}
\end{Highlighting}
\end{Shaded}

\section{Modelo de regresión logística}\label{sec-logistica}

Cuando la variable de interés es dicotómica, tomando valores de cero o uno, el modelo lineal clásico no resulta adecuado debido a que puede generar predicciones fuera del intervalo \((0,1)\) y porque la varianza de la respuesta depende de su media. Para abordar estas limitaciones, se emplea el modelo de regresión logística, uno de los modelos lineales generalizados más utilizados para analizar variables binarias. En el contexto de encuestas de hogares, la regresión logística puede emplearse para analizar la asociación entre una condición de interés, como la pobreza, y características sociodemográficas de los individuos o los hogares.

Considere una variable dicotómica \(Y_k \in \{0,1\}\), cuya esperanza es \(\theta_k=Pr(Y_k=1)\). El modelo de regresión logística relaciona dicha probabilidad con un conjunto de variables explicativas mediante la función de enlace logit, definida como

\[
\log\left(\frac{\theta_k}{1-\theta_k}\right)
=
\beta_0+\beta_1x_{k1}+\cdots+\beta_px_{kp}
=\mathbf{x}_{k}\boldsymbol{\beta}
\]

En donde \(\boldsymbol{\beta}\) es el vector de coeficientes de regresión asociado a las covariables. Esta transformación logit convierte el intervalo restringido de probabilidades \((0,1)\) en toda la recta real, permitiendo modelar la relación entre la respuesta y las variables auxiliares mediante un predictor lineal. De manera equivalente, la probabilidad de éxito puede expresarse mediante la siguiente relación:

\[
\theta_k
=
Pr(Y_k = 1 \mid \mathbf{x}_k)
=
\frac{\exp(\mathbf{x}_k\boldsymbol{\beta})}
     {1+\exp(\mathbf{x}_k\boldsymbol{\beta})}
\]

De manera equivalente, la probabilidad de que la unidad no presente la característica de interés es \(Pr(Y_k = 0 \mid \mathbf{x}_k) = 1-\theta_k\). Bajo esta parametrización, cada coeficiente de \(\boldsymbol{\beta}\) representa el efecto de una covariable sobre el logaritmo de la razón de probabilidades de que ocurra el evento de interés, manteniendo constantes las demás variables del modelo.
La estimación de los parámetros puede realizarse mediante el enfoque de pseudo-máxima verosimilitud. Este método incorpora los factores de expansión en la función de verosimilitud convencional con el fin de obtener estimadores consistentes respecto de la población objetivo. Siguiendo a \citet{heeringa2017applied}, la función de pseudo-verosimilitud para el modelo logístico se define como

\[
PL(\boldsymbol{\beta})
=
\prod_{k}
\left[
\theta_{k}^{y_{k}}
\left\{
1-\theta_{k}
\right\}^{1-y_{k}}
\right]^{w_{k}}
\]

En donde el subíndice \(k\) representa una unidad observada de la muestra. Asimismo, \(\theta_{k}\) denota la probabilidad de que la unidad \(k\) de la UPM \(i\) en el estrato \(h\) presente la característica de interés, \(y_{k}\) corresponde al valor observado de la variable dicotómica para dicha unidad y \(w_{k}\) representa el factor de expansión asociado a la observación.

La contribución de cada observación al proceso de estimación está determinada por su respectivo factor de expansión, permitiendo que la función refleje adecuadamente las características del diseño complejo de la encuesta. Dado que la maximización de la pseudo-verosimilitud resulta equivalente a maximizar su logaritmo, entonces el problema se reduce a optimizar la siguiente función:

\[
\ell(\boldsymbol{\beta})
=
\sum_{k}
w_{k}
\left[
y_{k}\log(\theta_{k})
+
(1-y_{k})
\log(1-\theta_{k})
\right]
\]

Esta expresión constituye la función objetivo utilizada por los algoritmos de optimización numérica para obtener el estimador de los parámetros del modelo. Es importante señalar que la función anterior no corresponde a una verosimilitud en sentido estricto bajo el diseño complejo de la encuesta. Por esta razón, se utiliza el término \emph{pseudo-verosimilitud} y los estimadores obtenidos mediante su maximización son conocidos como estimadores de pseudo-máxima verosimilitud (\emph{Pseudo Maximum Likelihood Estimators}, PMLE).

Según \citet{molina1992pseudo}, bajo condiciones de regularidad, estos estimadores son consistentes y asintóticamente normales, constituyendo la base de los procedimientos de inferencia implementados en la mayoría de los paquetes estadísticos para el análisis de datos de encuestas complejas.

Una vez obtenidos los estimadores de pseudo-máxima verosimilitud, sus varianzas se estiman mediante técnicas de linealización de Taylor \citep{binder1983variances}, las cuales incorporan explícitamente las características del diseño complejo de muestreo. La matriz de varianzas y covarianzas de los parámetros estimados, denotada por \(\text{Var}(\hat{\boldsymbol{\beta}})=\boldsymbol{\Sigma}\), se obtiene a partir de una aproximación lineal de las ecuaciones de estimación alrededor del verdadero valor de los parámetros.

Esta matriz describe la precisión de los estimadores, ya que sus elementos diagonales corresponden a las varianzas de cada coeficiente, mientras que los elementos fuera de la diagonal representan las covarianzas entre pares de coeficientes. Su estimación constituye la base para el cálculo de errores estándar, la construcción de intervalos de confianza y la realización de pruebas de hipótesis sobre los parámetros del modelo.

Como los datos provienen de un diseño muestral complejo, es necesario incorporar adecuadamente las ponderaciones, la estratificación y la conglomeración para obtener estimadores consistentes de los parámetros y de sus errores estándar. En R, esta tarea puede realizarse mediante la función \texttt{svyglm()} del paquete \texttt{survey}, la cual extiende los modelos lineales generalizados al contexto de encuestas complejas.

Siguiendo con la encuesta de ejemplo, y con el fin de analizar los factores asociados a la condición de pobreza, se ajusta un modelo de regresión logística utilizando como variables explicativas el gasto, la situación en el empleo y el sexo. Dado que los datos provienen de una encuesta con diseño muestral complejo, la estimación se realiza mediante la función \texttt{svyglm()} del paquete \texttt{survey}, especificando \texttt{family\ =\ binomial} para indicar que la variable respuesta es dicotómica y que se utilizará la función de enlace logit. Los resultados del ajuste se presentan en la Tabla \ref{tab:tab-ic-logistica}.

\begin{Shaded}
\begin{Highlighting}[]
\NormalTok{logit\_model }\OtherTok{\textless{}{-}} \FunctionTok{svyglm}\NormalTok{(}
  \AttributeTok{formula =}\NormalTok{ poor }\SpecialCharTok{\textasciitilde{}}\NormalTok{ Expenditure }\SpecialCharTok{+}\NormalTok{ Employment }\SpecialCharTok{+}\NormalTok{ Sex,}
  \AttributeTok{family =}\NormalTok{ binomial,}
  \AttributeTok{design =}\NormalTok{ survey\_design}
\NormalTok{)}
\end{Highlighting}
\end{Shaded}

En la especificación del modelo, el argumento \texttt{formula\ =\ poor\ \textasciitilde{}\ Expenditure\ +\ Employment\ +\ Sex} define la relación entre la variable respuesta \texttt{poor} y las covariables incluidas en el predictor lineal. Para las variables categóricas, la función genera automáticamente las variables indicadoras correspondientes y estima sus efectos con respecto a una categoría de referencia. Por su parte, el argumento \texttt{design\ =\ survey\_design} incorpora la información del diseño muestral. El objeto resultante, \texttt{logit\_model}, contiene los parámetros estimados, sus medidas de precisión y las estadísticas necesarias para realizar inferencia sobre la relación entre la condición de pobreza y las variables explicativas consideradas en el modelo.

Los resultados de la Tabla \ref{tab:tab-ic-logistica} muestran las sestimaciones de los coeficientes de regresión, junto con su error estándar, los intervalos de confirza y el valor-p.~Del modelo se deduce que un mayor nivel de gasto se asocia con una menor probabilidad de encontrarse en condición de pobreza. Asimismo, en comparación con la categoría de referencia (desocupados), tanto las personas inactivas como las ocupadas presentan una menor propensión a ser pobres, siendo este efecto más marcado entre estas últimas. En contraste, el sexo tiene una influencia muy reducida sobre la probabilidad de pobreza una vez controlado el efecto de las demás covariables incluidas en el modelo.

\begin{verbatim}
##                        2.5 %    97.5 %
## (Intercept)         1.448618  2.926867
## Expenditure        -0.007044 -0.003752
## EmploymentInactive -1.618341 -0.108547
## EmploymentEmployed -2.076705 -0.737132
## SexMale            -0.309026  0.314060
\end{verbatim}

\begin{longtable}[]{@{}
  >{\raggedright\arraybackslash}p{(\linewidth - 10\tabcolsep) * \real{0.1408}}
  >{\centering\arraybackslash}p{(\linewidth - 10\tabcolsep) * \real{0.2817}}
  >{\centering\arraybackslash}p{(\linewidth - 10\tabcolsep) * \real{0.1408}}
  >{\centering\arraybackslash}p{(\linewidth - 10\tabcolsep) * \real{0.1549}}
  >{\centering\arraybackslash}p{(\linewidth - 10\tabcolsep) * \real{0.1549}}
  >{\centering\arraybackslash}p{(\linewidth - 10\tabcolsep) * \real{0.1268}}@{}}
\caption{\label{tab:tab-ic-logistica}Intervalos de confianza al 95\% para los parámetros del modelo logístico}\tabularnewline
\toprule\noalign{}
\begin{minipage}[b]{\linewidth}\raggedright
Parameter
\end{minipage} & \begin{minipage}[b]{\linewidth}\centering
term
\end{minipage} & \begin{minipage}[b]{\linewidth}\centering
estimate
\end{minipage} & \begin{minipage}[b]{\linewidth}\centering
std.error
\end{minipage} & \begin{minipage}[b]{\linewidth}\centering
statistic
\end{minipage} & \begin{minipage}[b]{\linewidth}\centering
p.value
\end{minipage} \\
\midrule\noalign{}
\endfirsthead
\toprule\noalign{}
\begin{minipage}[b]{\linewidth}\raggedright
Parameter
\end{minipage} & \begin{minipage}[b]{\linewidth}\centering
term
\end{minipage} & \begin{minipage}[b]{\linewidth}\centering
estimate
\end{minipage} & \begin{minipage}[b]{\linewidth}\centering
std.error
\end{minipage} & \begin{minipage}[b]{\linewidth}\centering
statistic
\end{minipage} & \begin{minipage}[b]{\linewidth}\centering
p.value
\end{minipage} \\
\midrule\noalign{}
\endhead
\bottomrule\noalign{}
\endlastfoot
1 & (Intercept) & 2.188 & 0.373 & 5.863 & 0.000 \\
2 & Expenditure & -0.005 & 0.001 & -6.497 & 0.000 \\
3 & EmploymentInactive & -0.863 & 0.381 & -2.266 & 0.025 \\
4 & EmploymentEmployed & -1.407 & 0.338 & -4.161 & 0.000 \\
5 & SexMale & 0.003 & 0.157 & 0.016 & 0.987 \\
\end{longtable}

La Figura \ref{fig:fig-logistica-base} presenta la distribución de los estimadores de los coeficientes de regresión. Se observa que, para todas las covariables excepto el sexo, los intervalos de confianza no incluyen el valor cero, lo que proporciona evidencia de una asociación estadísticamente significativa. En contraste, el intervalo de confianza asociado a la variable sexo contiene el valor cero, por lo que no se encuentra evidencia suficiente para concluir que esta variable tenga un efecto significativo una vez controladas las demás covariables incluidas en el modelo.

Para visualizar la distribución de los coeficientes se emplea \texttt{plot\_summs()} del paquete \texttt{jtools}, junto con \texttt{ggstance} \citep{ggstance2024} para la geometría horizontal de los intervalos.

\begin{Shaded}
\begin{Highlighting}[]
\FunctionTok{library}\NormalTok{(ggstance)}
\FunctionTok{plot\_summs}\NormalTok{(logit\_model,}
  \AttributeTok{scale =} \ConstantTok{TRUE}\NormalTok{,}
  \AttributeTok{plot.distributions =} \ConstantTok{TRUE}
\NormalTok{)}
\end{Highlighting}
\end{Shaded}

\begin{figure}

{\centering \includegraphics[width=0.8\linewidth]{docs/figures/fig-logistica-base-1} 

}

\caption{Distribución de los parámetros del modelo logístico base}\label{fig:fig-logistica-base}
\end{figure}

El modelo puede ampliarse incorporando términos de interacción para evaluar si el efecto de una variable explicativa depende de los valores que toma otra. En este caso, se incluye una interacción entre el sexo y la situación laboral mediante el término \texttt{Sex:Employment}, con el objetivo de analizar si la relación entre el estado ocupacional y la probabilidad de pobreza es la misma para hombres y mujeres. Sin este término, el modelo asume que el efecto de la situación laboral sobre la pobreza es idéntico para ambos sexos. En cambio, al incorporar la interacción, se permite que dicho efecto varíe entre hombres y mujeres, capturando posibles diferencias en la forma en que la condición laboral se asocia con la pobreza en cada grupo. Los resultados del modelo ajustado se presentan en la Tabla \ref{tab:tab-logistica-interaccion}.

\begin{Shaded}
\begin{Highlighting}[]
\NormalTok{interaction\_logit\_model }\OtherTok{\textless{}{-}} \FunctionTok{svyglm}\NormalTok{(}
  \AttributeTok{formula =}\NormalTok{ poor }\SpecialCharTok{\textasciitilde{}}\NormalTok{ Expenditure }\SpecialCharTok{+}\NormalTok{ Employment }\SpecialCharTok{+}\NormalTok{ Sex }\SpecialCharTok{+}\NormalTok{ Sex}\SpecialCharTok{:}\NormalTok{Employment,}
  \AttributeTok{family =}\NormalTok{ binomial,}
  \AttributeTok{design =}\NormalTok{ survey\_design}
\NormalTok{)}
\end{Highlighting}
\end{Shaded}

\begin{longtable}[]{@{}
  >{\raggedright\arraybackslash}p{(\linewidth - 8\tabcolsep) * \real{0.3971}}
  >{\centering\arraybackslash}p{(\linewidth - 8\tabcolsep) * \real{0.1471}}
  >{\centering\arraybackslash}p{(\linewidth - 8\tabcolsep) * \real{0.1618}}
  >{\centering\arraybackslash}p{(\linewidth - 8\tabcolsep) * \real{0.1618}}
  >{\centering\arraybackslash}p{(\linewidth - 8\tabcolsep) * \real{0.1324}}@{}}
\caption{\label{tab:tab-logistica-interaccion}Coeficientes del modelo logístico con interacciones}\tabularnewline
\toprule\noalign{}
\begin{minipage}[b]{\linewidth}\raggedright
term
\end{minipage} & \begin{minipage}[b]{\linewidth}\centering
estimate
\end{minipage} & \begin{minipage}[b]{\linewidth}\centering
std.error
\end{minipage} & \begin{minipage}[b]{\linewidth}\centering
statistic
\end{minipage} & \begin{minipage}[b]{\linewidth}\centering
p.value
\end{minipage} \\
\midrule\noalign{}
\endfirsthead
\toprule\noalign{}
\begin{minipage}[b]{\linewidth}\raggedright
term
\end{minipage} & \begin{minipage}[b]{\linewidth}\centering
estimate
\end{minipage} & \begin{minipage}[b]{\linewidth}\centering
std.error
\end{minipage} & \begin{minipage}[b]{\linewidth}\centering
statistic
\end{minipage} & \begin{minipage}[b]{\linewidth}\centering
p.value
\end{minipage} \\
\midrule\noalign{}
\endhead
\bottomrule\noalign{}
\endlastfoot
(Intercept) & 1.770 & 0.610 & 2.900 & 0.004 \\
Expenditure & -0.005 & 0.001 & -6.473 & 0.000 \\
EmploymentInactive & -0.395 & 0.594 & -0.665 & 0.507 \\
EmploymentEmployed & -1.059 & 0.570 & -1.860 & 0.066 \\
SexMale & 0.587 & 0.748 & 0.785 & 0.434 \\
EmploymentInactive:SexMale & -0.846 & 0.860 & -0.984 & 0.327 \\
EmploymentEmployed:SexMale & -0.481 & 0.760 & -0.633 & 0.528 \\
\end{longtable}

Al incorporar la interacción entre sexo y situación laboral, el gasto continúa mostrando una asociación negativa y estadísticamente significativa con la probabilidad de pobreza, manteniéndose como la covariable con mayor evidencia explicativa dentro del modelo. En contraste, los efectos principales de la situación laboral pierden importancia estadística una vez que se incluyen los términos de interacción, lo que sugiere una mayor incertidumbre en la estimación de sus efectos. Por su parte, ni el efecto principal del sexo ni los coeficientes asociados a las interacciones entre sexo y situación laboral resultan estadísticamente significativos al nivel del 5\%. En consecuencia, la inclusión de los términos de interacción no parece aportar información adicional relevante respecto al modelo sin interacción.

La Figura \ref{fig:fig-comparacion-logistica} presenta una comparación de las estimaciones de los coeficientes y sus respectivos intervalos de confianza para los modelos con y sin interacción. Se observa que la incorporación de los términos de interacción incrementa la incertidumbre asociada a varios de los parámetros, lo que se refleja en intervalos de confianza más amplios. Asimismo, los coeficientes de interacción presentan estimaciones cercanas a cero y amplios intervalos de confianza que incluyen dicho valor, lo que coincide con la falta de evidencia estadística para concluir que el efecto de la situación laboral sobre la probabilidad de pobreza difiere entre hombres y mujeres.

\begin{Shaded}
\begin{Highlighting}[]
\FunctionTok{plot\_summs}\NormalTok{(interaction\_logit\_model, logit\_model,}
  \AttributeTok{scale =} \ConstantTok{TRUE}\NormalTok{,}
  \AttributeTok{plot.distributions =} \ConstantTok{TRUE}
\NormalTok{)}
\end{Highlighting}
\end{Shaded}

\begin{figure}

{\centering \includegraphics[width=0.7\linewidth]{docs/figures/fig-comparacion-logistica-1} 

}

\caption{Comparación de los parámetros del modelo logístico con y sin interacciones}\label{fig:fig-comparacion-logistica}
\end{figure}

\section{Modelo de regresión multinomial}\label{sec-multinomial}

En las encuestas de hogares es frecuente encontrar variables de respuesta con más de dos categorías, como la situación laboral, cuyos posibles estados incluyen empleado, desempleado e inactivo. Cuando estas categorías son mutuamente excluyentes y no poseen un orden natural, una herramienta adecuada para modelar su relación con un conjunto de covariables es la regresión logística multinomial, la cual constituye una extensión del modelo de regresión logística binaria.

Este modelo permite estimar la probabilidad de pertenecer a cada categoría de la variable respuesta en función de variables explicativas de naturaleza continua o categórica. Su aplicación requiere que las categorías de respuesta sean exhaustivas y mutuamente excluyentes y que no exista multicolinealidad severa entre las covariables. Además, para las covariables continuas, se asume una relación lineal entre estas y el logaritmo de las razones de probabilidades (\emph{log-odds}) de cada categoría respecto a una categoría de referencia.

El modelo logístico multinomial especifica la probabilidad \(\theta_{k(j)}\) de que la unidad \(k\) pertenezca a la categoría \(j\), con \(j=1,\ldots,J\), dado el vector de covariables \(\mathbf{x}_k\), mediante la siguiente expresión:

\[
Pr(Y_{k} = j \mid \mathbf{x}_{k})
=
\theta_{k(j)}
=
\frac{\exp(\mathbf{x}_{k}\boldsymbol{\beta}_j)}
     {\sum_{l=1}^{J} \exp(\mathbf{x}_{k}\boldsymbol{\beta}_l)}
\]

En donde \(\boldsymbol{\beta}_j\) es el vector de coeficientes de regresión para la \(j\)-ésima categoría. Esta formulación expresa directamente las probabilidades de pertenencia a cada categoría, en la práctica resulta más conveniente reescribir el modelo utilizando una categoría de referencia. Hay parametrizaciones alternativas que facilitan la interpretación de los parámetros, ya que cada conjunto de coeficientes describe el efecto de las covariables sobre el logaritmo de las razones de probabilidades (\emph{log-odds}) de pertenecer a una determinada categoría en comparación con la categoría de referencia.

La estimación de los parámetros del modelo se realiza mediante el enfoque de pseudo-máxima verosimilitud. Siguiendo a \citet{heeringa2017applied}, la función de pseudo-verosimilitud multinomial se define como

\[
PL(\boldsymbol{\beta} \mid \mathbf{X}) =
\prod_{k}
\left\{
\prod_{j=1}^{J}
\theta_{k(j)}^{y_{k(j)}}
\right\}^{w_{k}}
\]

En donde \(k\) representa las unidades de observación pertenecientes a la muestra compleja. Asimismo, \(y_{k(j)}\) es una variable indicadora que toma el valor de uno cuando dicha unidad pertenece a la categoría \(j\) y cero en caso contrario, y \(w_{k}\) corresponde al factor de expansión asociado a la unidad observada. Por conveniencia matemática y computacional, la estimación de los parámetros se realiza a partir de la pseudo-logverosimilitud, cuya maximización conduce a los mismos estimadores que la pseudo-verosimilitud original.

\[
\ell(\boldsymbol{\beta})
=
\sum_{k} w_k
\sum_{j=1}^{J}
y_{k(j)}
\log(\theta_{k(j)})
\]

Como se mencionó anteriormente, los estimadores obtenidos mediante este procedimiento son consistentes y asintóticamente normales bajo condiciones regulares \citep{molina1992pseudo}. La inferencia estadística para los parámetros de interés se basa en estimaciones de varianza obtenidas mediante técnicas de linealización de Taylor que incorporan explícitamente las características del diseño muestral complejo, permitiendo calcular errores estándar, intervalos de confianza y pruebas de hipótesis apropiados para datos de encuestas \citep{binder1983variances}.

Para el ajuste del modelo usando \texttt{R}, como referencia inicial, se estiman las proporciones de personas por estado de empleo, restringiendo las unidades de interés a mayores de 15 años, como se presenta en la Tabla \ref{tab:tab-prop-empleo-15}.

\begin{Shaded}
\begin{Highlighting}[]
\NormalTok{survey\_design }\SpecialCharTok{\%\textgreater{}\%}
  \FunctionTok{filter}\NormalTok{(Age }\SpecialCharTok{\textgreater{}=} \DecValTok{15}\NormalTok{) }\SpecialCharTok{\%\textgreater{}\%}
  \FunctionTok{group\_by}\NormalTok{(Employment) }\SpecialCharTok{\%\textgreater{}\%}
  \FunctionTok{summarise}\NormalTok{(}\AttributeTok{proportion =} \FunctionTok{survey\_mean}\NormalTok{(}\AttributeTok{vartype =} \FunctionTok{c}\NormalTok{(}\StringTok{"se"}\NormalTok{, }\StringTok{"ci"}\NormalTok{)))}
\end{Highlighting}
\end{Shaded}

\begin{longtable}[]{@{}
  >{\raggedright\arraybackslash}p{(\linewidth - 8\tabcolsep) * \real{0.1571}}
  >{\centering\arraybackslash}p{(\linewidth - 8\tabcolsep) * \real{0.1714}}
  >{\centering\arraybackslash}p{(\linewidth - 8\tabcolsep) * \real{0.2143}}
  >{\centering\arraybackslash}p{(\linewidth - 8\tabcolsep) * \real{0.2286}}
  >{\centering\arraybackslash}p{(\linewidth - 8\tabcolsep) * \real{0.2286}}@{}}
\caption{\label{tab:tab-prop-empleo-15}Proporción estimada de personas por estado de empleo (mayores de 15 años)}\tabularnewline
\toprule\noalign{}
\begin{minipage}[b]{\linewidth}\raggedright
Employment
\end{minipage} & \begin{minipage}[b]{\linewidth}\centering
proportion
\end{minipage} & \begin{minipage}[b]{\linewidth}\centering
proportion\_se
\end{minipage} & \begin{minipage}[b]{\linewidth}\centering
proportion\_low
\end{minipage} & \begin{minipage}[b]{\linewidth}\centering
proportion\_upp
\end{minipage} \\
\midrule\noalign{}
\endfirsthead
\toprule\noalign{}
\begin{minipage}[b]{\linewidth}\raggedright
Employment
\end{minipage} & \begin{minipage}[b]{\linewidth}\centering
proportion
\end{minipage} & \begin{minipage}[b]{\linewidth}\centering
proportion\_se
\end{minipage} & \begin{minipage}[b]{\linewidth}\centering
proportion\_low
\end{minipage} & \begin{minipage}[b]{\linewidth}\centering
proportion\_upp
\end{minipage} \\
\midrule\noalign{}
\endhead
\bottomrule\noalign{}
\endlastfoot
Unemployed & 0.043 & 0.007 & 0.029 & 0.057 \\
Inactive & 0.384 & 0.015 & 0.354 & 0.414 \\
Employed & 0.573 & 0.014 & 0.545 & 0.601 \\
\end{longtable}

Para modelar una variable de respuesta multinomial se emplea la función \texttt{svy\_vglm()} del paquete \texttt{svyVGAM} \citep{svyVGAM2026}, que permite ajustar modelos de regresión multinomial incorporando explícitamente las características del diseño complejo de la encuesta. En este ejemplo, se especifica un modelo donde la condición de actividad se explica en función de la edad, el sexo y la zona de residencia. A diferencia de los modelos multinomiales convencionales, \texttt{svy\_vglm()} calcula las varianzas de los estimadores considerando la complejidad del diseño muestral, incluyendo la estratificación, la conglomeración y los pesos de expansión. De esta forma, tanto las estimaciones puntuales como sus errores estándar e inferencias asociadas son consistentes con el diseño de la encuesta.

\begin{Shaded}
\begin{Highlighting}[]
\FunctionTok{library}\NormalTok{(svyVGAM)}
\NormalTok{survey\_design\_15 }\OtherTok{\textless{}{-}}\NormalTok{ survey\_design }\SpecialCharTok{\%\textgreater{}\%}
  \FunctionTok{filter}\NormalTok{(Age }\SpecialCharTok{\textgreater{}=} \DecValTok{15}\NormalTok{)}
\NormalTok{multinomial\_model }\OtherTok{\textless{}{-}} \FunctionTok{svy\_vglm}\NormalTok{(}
  \AttributeTok{formula =}\NormalTok{ Employment }\SpecialCharTok{\textasciitilde{}}\NormalTok{ Age }\SpecialCharTok{+}\NormalTok{ Sex }\SpecialCharTok{+}\NormalTok{ Zone,}
  \AttributeTok{design =}\NormalTok{ survey\_design\_15,}
  \AttributeTok{crit =} \StringTok{"coef"}\NormalTok{,}
  \AttributeTok{family =} \FunctionTok{multinomial}\NormalTok{(}\AttributeTok{refLevel =} \StringTok{"Unemployed"}\NormalTok{)}
\NormalTok{)}
\end{Highlighting}
\end{Shaded}

El argumento \texttt{family\ =\ multinomial(refLevel\ =\ "Unemployed")} define un modelo logístico multinomial tomando la categoría \texttt{Unemployed} como referencia. Bajo esta especificación, se estiman de manera simultánea los logaritmos de las razones de probabilidades (\emph{log-odds}) entre cada una de las categorías de \texttt{Employment} y la categoría de referencia. Los parámetros se obtienen mediante pseudo-máxima verosimilitud, incorporando los factores de expansión de la encuesta en la función objetivo.

Dado que \texttt{broom::tidy()} no puede aplicarse directamente a objetos de clase \texttt{svyVGAM}, se define la siguiente función auxiliar para estructurar los resultados del modelo en un formato tabular estandarizado\footnote{Adaptada de \url{https://tech.popdata.org/pma-data-hub/posts/2021-08-15-covid-analysis/}}:

\begin{Shaded}
\begin{Highlighting}[]
\NormalTok{tidy.svyVGAM }\OtherTok{\textless{}{-}} \ControlFlowTok{function}\NormalTok{(x, }\AttributeTok{conf.int =} \ConstantTok{FALSE}\NormalTok{, }\AttributeTok{conf.level =} \FloatTok{0.95}\NormalTok{,}
                          \AttributeTok{exponentiate =} \ConstantTok{FALSE}\NormalTok{, ...) \{}
\NormalTok{  ret }\OtherTok{\textless{}{-}} \FunctionTok{as\_tibble}\NormalTok{(}\FunctionTok{summary}\NormalTok{(x)}\SpecialCharTok{$}\NormalTok{coeftable, }\AttributeTok{rownames =} \StringTok{"term"}\NormalTok{)}
  \FunctionTok{colnames}\NormalTok{(ret) }\OtherTok{\textless{}{-}} \FunctionTok{c}\NormalTok{(}\StringTok{"term"}\NormalTok{, }\StringTok{"estimate"}\NormalTok{, }\StringTok{"std.error"}\NormalTok{, }\StringTok{"statistic"}\NormalTok{, }\StringTok{"p.value"}\NormalTok{)}
\NormalTok{  coefs }\OtherTok{\textless{}{-}}\NormalTok{ tibble}\SpecialCharTok{::}\FunctionTok{enframe}\NormalTok{(stats}\SpecialCharTok{::}\FunctionTok{coef}\NormalTok{(x), }\AttributeTok{name =} \StringTok{"term"}\NormalTok{, }\AttributeTok{value =} \StringTok{"estimate"}\NormalTok{)}
\NormalTok{  ret }\OtherTok{\textless{}{-}} \FunctionTok{left\_join}\NormalTok{(coefs, ret, }\AttributeTok{by =} \FunctionTok{c}\NormalTok{(}\StringTok{"term"}\NormalTok{, }\StringTok{"estimate"}\NormalTok{))}
  \ControlFlowTok{if}\NormalTok{ (conf.int) \{}
\NormalTok{    ci }\OtherTok{\textless{}{-}}\NormalTok{ broom}\SpecialCharTok{:::}\FunctionTok{broom\_confint\_terms}\NormalTok{(x, }\AttributeTok{level =}\NormalTok{ conf.level, ...)}
\NormalTok{    ret }\OtherTok{\textless{}{-}}\NormalTok{ dplyr}\SpecialCharTok{::}\FunctionTok{left\_join}\NormalTok{(ret, ci, }\AttributeTok{by =} \StringTok{"term"}\NormalTok{)}
\NormalTok{  \}}
  \ControlFlowTok{if}\NormalTok{ (exponentiate) ret }\OtherTok{\textless{}{-}}\NormalTok{ broom}\SpecialCharTok{:::}\FunctionTok{exponentiate}\NormalTok{(ret)}
\NormalTok{  ret }\SpecialCharTok{\%\textgreater{}\%}
\NormalTok{    tidyr}\SpecialCharTok{::}\FunctionTok{separate}\NormalTok{(term, }\AttributeTok{into =} \FunctionTok{c}\NormalTok{(}\StringTok{"term"}\NormalTok{, }\StringTok{"y.level"}\NormalTok{), }\AttributeTok{sep =} \StringTok{":"}\NormalTok{) }\SpecialCharTok{\%\textgreater{}\%}
    \FunctionTok{arrange}\NormalTok{(y.level) }\SpecialCharTok{\%\textgreater{}\%}
    \FunctionTok{relocate}\NormalTok{(y.level, }\AttributeTok{.before =}\NormalTok{ term)}
\NormalTok{\}}
\end{Highlighting}
\end{Shaded}

La Tabla \ref{tab:tab-multinomial-coef} presenta los coeficientes estimados del modelo logístico multinomial, tomando la categoría \emph{Unemployed} como referencia. Los resultados muestran que la edad tiene un efecto positivo y estadísticamente significativo sobre las probabilidades relativas de pertenecer a las categorías 1 y 2 frente al desempleo. Asimismo, la variable \texttt{SexMale} presenta coeficientes negativos y significativos, indicando que los hombres tienen menores probabilidades relativas que las mujeres de ubicarse en dichas categorías. Por el contrario, la variable de zona de residencia no muestra efectos estadísticamente significativos al nivel convencional del 5\%.

\begin{Shaded}
\begin{Highlighting}[]
\FunctionTok{tidy.svyVGAM}\NormalTok{(multinomial\_model, }\AttributeTok{exponentiate =} \ConstantTok{FALSE}\NormalTok{, }\AttributeTok{conf.int =} \ConstantTok{FALSE}\NormalTok{)}
\end{Highlighting}
\end{Shaded}

\begin{longtable}[]{@{}lccccc@{}}
\caption{\label{tab:tab-multinomial-coef}Coeficientes estimados del modelo de regresión multinomial}\tabularnewline
\toprule\noalign{}
y.level & term & estimate & std.error & statistic & p.value \\
\midrule\noalign{}
\endfirsthead
\toprule\noalign{}
y.level & term & estimate & std.error & statistic & p.value \\
\midrule\noalign{}
\endhead
\bottomrule\noalign{}
\endlastfoot
1 & (Intercept) & 2.556 & 0.550 & 4.648 & 0.000 \\
1 & Age & 0.022 & 0.010 & 2.117 & 0.034 \\
1 & SexMale & -2.185 & 0.317 & -6.901 & 0.000 \\
1 & ZoneUrban & -0.254 & 0.404 & -0.627 & 0.531 \\
2 & (Intercept) & 2.306 & 0.461 & 4.998 & 0.000 \\
2 & Age & 0.018 & 0.009 & 2.056 & 0.040 \\
2 & SexMale & -0.578 & 0.277 & -2.091 & 0.037 \\
2 & ZoneUrban & 0.039 & 0.361 & 0.109 & 0.913 \\
\end{longtable}

La Figura \ref{fig:fig-multinomial-ic} muestra los intervalos de confianza de los coeficientes estimados para cada ecuación del modelo multinomial. La visualización se construye con \texttt{dotwhisker::dwplot()} \citep{dotwhisker2025}, que permite comparar coeficientes e intervalos en una escala común, y se usa \texttt{gtools::stars.pval()} \citep{gtools2025} para codificar la significancia estadística. Las covariables cuyos intervalos no incluyen el valor cero presentan evidencia de asociación estadística significativa con la variable respuesta, mientras que aquellos cuyos intervalos contienen el cero no muestran efectos significativos al nivel de confianza considerado.

\begin{Shaded}
\begin{Highlighting}[]
\NormalTok{multinomial\_results }\SpecialCharTok{\%\textgreater{}\%}
  \FunctionTok{mutate}\NormalTok{(}
    \AttributeTok{model =} \FunctionTok{if\_else}\NormalTok{(y.level }\SpecialCharTok{==} \DecValTok{1}\NormalTok{, }\StringTok{"Inactive"}\NormalTok{, }\StringTok{"Employed"}\NormalTok{),}
    \AttributeTok{sig =}\NormalTok{ gtools}\SpecialCharTok{::}\FunctionTok{stars.pval}\NormalTok{(p.value)}
\NormalTok{  ) }\SpecialCharTok{\%\textgreater{}\%}
\NormalTok{  dotwhisker}\SpecialCharTok{::}\FunctionTok{dwplot}\NormalTok{(}
    \AttributeTok{dodge\_size =} \FloatTok{0.3}\NormalTok{,}
    \AttributeTok{vline =} \FunctionTok{geom\_vline}\NormalTok{(}\AttributeTok{xintercept =} \DecValTok{0}\NormalTok{, }\AttributeTok{colour =} \StringTok{"grey60"}\NormalTok{, }\AttributeTok{linetype =} \DecValTok{2}\NormalTok{)}
\NormalTok{  ) }\SpecialCharTok{+}
  \FunctionTok{guides}\NormalTok{(}\AttributeTok{color =} \FunctionTok{guide\_legend}\NormalTok{(}\AttributeTok{reverse =} \ConstantTok{TRUE}\NormalTok{)) }\SpecialCharTok{+}
  \FunctionTok{theme\_bw}\NormalTok{() }\SpecialCharTok{+}
  \FunctionTok{theme}\NormalTok{(}\AttributeTok{legend.position =} \StringTok{"top"}\NormalTok{)}
\end{Highlighting}
\end{Shaded}

\begin{figure}

{\centering \includegraphics[width=0.75\linewidth]{docs/figures/fig-multinomial-ic-1} 

}

\caption{Intervalos de confianza para los coeficientes del modelo multinomial}\label{fig:fig-multinomial-ic}
\end{figure}

\section{Modelo de regresión Gamma}\label{sec-gamma}

La regresión Gamma constituye una extensión de los modelos lineales generalizados adecuada para variables respuesta continuas, estrictamente positivas y cuya variabilidad aumenta con la media. El modelo Gamma constituye una alternativa adecuada para analizar variables continuas estrictamente positivas cuya variabilidad aumenta a medida que crece su valor esperado. Este comportamiento es frecuente en variables económicas y sociales, como los ingresos de los hogares o los gastos en consumo.

Si \(Y_{k}\), la variable de interés observada para la unidad \(k\), sigue una distribución Gamma con media \(\mu_{k}\) y parámetro de dispersión \(\phi\), el modelo relaciona el valor esperado de la variable respuesta con un conjunto de covariables mediante una función de enlace. La función de enlace más utilizada es el enlace logarítmico, el cual garantiza que las predicciones sean siempre positivas. En este caso, el modelo se expresa como

\[
\log(\mu_{k})
=
\mathbf{x}_{k}\boldsymbol{\beta}
\]

En donde \(\mathbf{x}_{k}\) corresponde al vector de variables explicativas asociado a la unidad observada y \(\boldsymbol{\beta}\) es el vector de parámetros desconocidos del modelo. De forma equivalente, se tiene que:

\[
\mu_{k}
= E(Y_{k}\mid \mathbf{x}_{k})
=
\exp(\mathbf{x}_{k}\boldsymbol{\beta}),
\]

Con el enlace logarítmico, cada parámetro puede interpretarse como el efecto de una covariable sobre el logaritmo de la media esperada, de manera que \(\exp(\beta_j)\) representa el cambio multiplicativo asociado a un incremento unitario en la covariable correspondiente, manteniendo constantes las demás variables del modelo. Bajo la distribución Gamma, la función de densidad condicional de \(Y_{k}\) puede escribirse como

\[
f(y_{k}\mid \mu_{k},\phi)
=
\frac{1}
{\Gamma(1/\phi)
(\phi\mu_{k})^{1/\phi}}
y_{k}^{\frac{1}{\phi}-1}
\exp\left(
-\frac{y_{k}}
{\phi\mu_{k}}
\right),
\qquad y_{k}>0,
\]

donde \(\Gamma(\cdot)\) denota la función Gamma y \(\phi\) representa el parámetro de dispersión. Bajo el enfoque de inferencia basado en pseudo-máxima verosimilitud, la función objetivo que se optimiza para obtener los estimadores de los parámetros del modelo está dada por

\[
\ell(\boldsymbol{\beta},\phi)
=
\sum_{k}
w_{k}
\log
\left[
f(y_{hik}\mid\mu_{k},\phi)
\right]
\]

En donde \(k\) representa las unidades de observación de la muestra y \(w_k\) corresponde al factor de expansión asociado a cada una de ellas. Esta función define el criterio de optimización utilizado para estimar los parámetros del modelo bajo el enfoque de pseudo-máxima verosimilitud. En \texttt{R}, se hace necesario primero definir el diseño de muestreo.

\begin{Shaded}
\begin{Highlighting}[]
\NormalTok{gamma\_design }\OtherTok{\textless{}{-}}\NormalTok{ survey\_data }\SpecialCharTok{\%\textgreater{}\%}
  \FunctionTok{as\_survey\_design}\NormalTok{(}
    \AttributeTok{strata =}\NormalTok{ Stratum,}
    \AttributeTok{ids =}\NormalTok{ PSU,}
    \AttributeTok{weights =}\NormalTok{ wk,}
    \AttributeTok{nest =} \ConstantTok{TRUE}
\NormalTok{  )}
\end{Highlighting}
\end{Shaded}

El siguiente código ajusta un modelo de regresión Gamma para la variable de ingresos del hogar (\texttt{Income}), incorporando el diseño complejo de la encuesta definido en el objeto \texttt{gamma\_design}. El modelo utiliza como variables explicativas el gasto del hogar (\texttt{Expenditure}), la edad (\texttt{Age}), el sexo (\texttt{Sex}) y la zona de residencia (\texttt{Zone}). Asimismo, se especifica una función de enlace logarítmica, garantizando predicciones positivas y permitiendo interpretar los efectos de las variables explicativas en términos multiplicativos sobre el ingreso esperado. Los coeficientes estimados se presentan en la Tabla \ref{tab:tab-gamma-coeficientes}:

\begin{Shaded}
\begin{Highlighting}[]
\NormalTok{gamma\_model }\OtherTok{\textless{}{-}} \FunctionTok{svyglm}\NormalTok{(}
  \AttributeTok{formula =}\NormalTok{ Income }\SpecialCharTok{\textasciitilde{}}\NormalTok{ Expenditure }\SpecialCharTok{+}\NormalTok{ Age }\SpecialCharTok{+}\NormalTok{ Sex }\SpecialCharTok{+}\NormalTok{ Zone,}
  \AttributeTok{design =}\NormalTok{ gamma\_design,}
  \AttributeTok{family =} \FunctionTok{Gamma}\NormalTok{(}\AttributeTok{link =} \StringTok{"log"}\NormalTok{)}
\NormalTok{)}
\end{Highlighting}
\end{Shaded}

\begin{longtable}[]{@{}lcccc@{}}
\caption{\label{tab:tab-gamma-coeficientes}Coeficientes del modelo Gamma para el ingreso de los hogares}\tabularnewline
\toprule\noalign{}
term & estimate & std.error & statistic & p.value \\
\midrule\noalign{}
\endfirsthead
\toprule\noalign{}
term & estimate & std.error & statistic & p.value \\
\midrule\noalign{}
\endhead
\bottomrule\noalign{}
\endlastfoot
(Intercept) & 5.409 & 0.098 & 55.184 & 0.000 \\
Expenditure & 0.002 & 0.000 & 16.891 & 0.000 \\
Age & 0.002 & 0.001 & 2.430 & 0.017 \\
SexMale & 0.048 & 0.038 & 1.260 & 0.210 \\
ZoneUrban & 0.151 & 0.089 & 1.701 & 0.092 \\
\end{longtable}

Dado que el modelo utiliza un enlace logarítmico, los coeficientes pueden interpretarse como efectos multiplicativos sobre el ingreso esperado. El coeficiente asociado al gasto del hogar es positivo y estadísticamente significativo (\(p<0.001\)), indicando que, manteniendo constantes las demás variables del modelo, un incremento de una unidad en el gasto se asocia con un aumento aproximado del \(0.2\%\) en el ingreso esperado, ya que \(\exp(0.002)\approx 1.002\). La edad también presenta un efecto positivo y significativo (\(p=0.017\)). En promedio, cada año adicional de edad se asocia con un incremento cercano al \(0.2\%\) en el ingreso esperado, manteniendo constantes las demás covariables. Por el contrario, las variables de sexo y zona de residencia no muestran evidencia estadística suficiente para concluir que influyen sobre el ingreso.

\chapter{Modelos multinivel}\label{cap7}

Los modelos multinivel, también conocidos como modelos jerárquicos, son una técnica estadística diseñada para el análisis de datos con estructura jerárquica, como los que provienen de encuestas de hogares, en donde las unidades de observación no son independientes entre sí. Los individuos pertenecen a hogares, los hogares se ubican en áreas geográficas específicas (UPMs) y estas, a su vez, forman parte de unidades territoriales más amplias (estratos o dominios). Como consecuencia, las observaciones que comparten un mismo contexto suelen presentar características más similares entre sí que aquellas pertenecientes a contextos diferentes. Esta estructura de dependencia viola uno de los supuestos fundamentales de los modelos de regresión convencionales y puede conducir a estimaciones ineficientes e inferencias incorrectas.

A diferencia de los modelos de regresión convencionales, los modelos multinivel reconocen que las observaciones pueden estar agrupadas en distintos niveles y que cada uno de ellos puede contribuir a explicar la variabilidad observada. En consecuencia, permiten representar simultáneamente los efectos asociados a las características de las unidades de análisis y aquellos derivados del entorno o contexto en el que estas se encuentran. Esta formulación resulta especialmente útil cuando se busca estudiar fenómenos determinados tanto por factores individuales como por características de los hogares, comunidades o áreas geográficas.

Otra característica relevante de los modelos multinivel es que permiten cuantificar la proporción de la variabilidad asociada a cada nivel de agrupación presente en los datos. Mientras los efectos fijos resumen la relación promedio entre la variable respuesta y las covariables incluidas en el modelo, los efectos aleatorios representan las diferencias sistemáticas entre los grupos que no son explicadas por dichas covariables. Gracias a esta estructura, es posible modelar explícitamente la dependencia entre observaciones pertenecientes a un mismo grupo y obtener inferencias más adecuadas cuando los datos presentan organización jerárquica.

El desarrollo teórico y metodológico de los modelos multinivel ha sido ampliamente documentado en la literatura de muestreo. Entre las referencias más influyentes se encuentran los trabajos de \citet{goldstein2011multilevel}, \citet{gelman2006data} y \citet{rabehesketh2012multilevel}, que presentan los fundamentos conceptuales, las estrategias de estimación y diversas aplicaciones de estos modelos en contextos sociales y demográficos. Por otra parte, \citet{browne2006comparison} comparan distintos enfoques de estimación para modelos jerárquicos, evaluando las diferencias entre los métodos basados en máxima verosimilitud y los enfoques bayesianos. En el ámbito de la epidemiología social, \citet{merlo2006brief} destacan la utilidad de los modelos multinivel para estudiar fenómenos contextuales, mostrando cómo factores asociados al entorno pueden contribuir a explicar desigualdades en salud y otros resultados de interés poblacional.

\section{Motivación}\label{sec-motivacionMult}

Para iniciar este capítulo, se cargan \texttt{tidyverse} \citep{tidyverse2026} para la manipulación de datos y la generación de gráficos, y \texttt{lme4} \citep{lme42026} para la estimación de modelos multinivel. Asimismo, se importa la base de datos que será utilizada a lo largo de los ejemplos. Los títulos matemáticos de algunas figuras se formatean con \texttt{latex2exp} \citep{latex2exp2026}, llamado de forma explícita en el código.

\begin{Shaded}
\begin{Highlighting}[]
\FunctionTok{library}\NormalTok{(tidyverse)}
\FunctionTok{library}\NormalTok{(lme4)}
\end{Highlighting}
\end{Shaded}

\begin{Shaded}
\begin{Highlighting}[]
\NormalTok{survey\_data }\OtherTok{\textless{}{-}} \FunctionTok{readRDS}\NormalTok{(}\StringTok{"Data/encuesta.rds"}\NormalTok{) }\SpecialCharTok{\%\textgreater{}\%}
  \FunctionTok{mutate}\NormalTok{(}\AttributeTok{poor =} \FunctionTok{ifelse}\NormalTok{(Poverty }\SpecialCharTok{!=} \StringTok{"NotPoor"}\NormalTok{, }\DecValTok{1}\NormalTok{, }\DecValTok{0}\NormalTok{))}
\end{Highlighting}
\end{Shaded}

Con fines ilustrativos, los ejemplos de esta sección se desarrollan a partir de una submuestra de la encuesta. El siguiente bloque de código construye la base de datos utilizada en los ejemplos del capítulo. Para ello, identifica los tres estratos con menor número de hogares y conserva únicamente las variables requeridas para el análisis: ingreso, gasto, estrato, sexo, región, zona de residencia y condición de pobreza.

\begin{Shaded}
\begin{Highlighting}[]
\NormalTok{plot\_data }\OtherTok{\textless{}{-}}\NormalTok{ survey\_data }\SpecialCharTok{\%\textgreater{}\%}
  \FunctionTok{select}\NormalTok{(HHID, Stratum) }\SpecialCharTok{\%\textgreater{}\%}
  \FunctionTok{distinct}\NormalTok{() }\SpecialCharTok{\%\textgreater{}\%}
  \FunctionTok{group\_by}\NormalTok{(Stratum) }\SpecialCharTok{\%\textgreater{}\%}
  \FunctionTok{tally}\NormalTok{() }\SpecialCharTok{\%\textgreater{}\%}
  \FunctionTok{arrange}\NormalTok{(n) }\SpecialCharTok{\%\textgreater{}\%}
  \FunctionTok{select}\NormalTok{(}\SpecialCharTok{{-}}\NormalTok{n) }\SpecialCharTok{\%\textgreater{}\%}
  \FunctionTok{slice}\NormalTok{(}\DecValTok{1}\SpecialCharTok{:}\DecValTok{3}\DataTypeTok{L}\NormalTok{) }\SpecialCharTok{\%\textgreater{}\%}
  \FunctionTok{inner\_join}\NormalTok{(survey\_data, }\AttributeTok{by =} \StringTok{"Stratum"}\NormalTok{) }\SpecialCharTok{\%\textgreater{}\%}
  \FunctionTok{select}\NormalTok{(Income, Expenditure, Stratum, Sex, Region, Zone, poor)}
\end{Highlighting}
\end{Shaded}

Como punto de partida, se ajusta un modelo de regresión lineal convencional que relaciona el ingreso de los hogares con su nivel de gasto, ignorando temporalmente la estructura jerárquica de los datos. Bajo esta especificación, se asume que todas las observaciones son independientes y que la relación entre ingreso y gasto es homogénea para todos los hogares, independientemente del estrato al que pertenezcan. La Figura \ref{fig:fig-regresion-simple} presenta la recta de regresión estimada junto con los datos observados.

\begin{Shaded}
\begin{Highlighting}[]
\FunctionTok{ggplot}\NormalTok{(}\AttributeTok{data =}\NormalTok{ plot\_data,}
       \FunctionTok{aes}\NormalTok{(}\AttributeTok{y =}\NormalTok{ Income, }\AttributeTok{x =}\NormalTok{ Expenditure)) }\SpecialCharTok{+}
  \FunctionTok{geom\_jitter}\NormalTok{() }\SpecialCharTok{+}
  \FunctionTok{geom\_smooth}\NormalTok{(}\AttributeTok{formula =}\NormalTok{ y }\SpecialCharTok{\textasciitilde{}}\NormalTok{ x, }\AttributeTok{method =} \StringTok{"lm"}\NormalTok{, }\AttributeTok{se =} \ConstantTok{FALSE}\NormalTok{) }\SpecialCharTok{+}
  \FunctionTok{ggtitle}\NormalTok{(latex2exp}\SpecialCharTok{::}\FunctionTok{TeX}\NormalTok{(}
    \StringTok{"$Income\_\{i\} }\SpecialCharTok{\textbackslash{}\textbackslash{}}\StringTok{sim }\SpecialCharTok{\textbackslash{}\textbackslash{}}\StringTok{hat\{}\SpecialCharTok{\textbackslash{}\textbackslash{}}\StringTok{beta\}\_\{0\} + }\SpecialCharTok{\textbackslash{}\textbackslash{}}\StringTok{hat\{}\SpecialCharTok{\textbackslash{}\textbackslash{}}\StringTok{beta\}\_\{1\}Expenditure\_\{i\} + }\SpecialCharTok{\textbackslash{}\textbackslash{}}\StringTok{epsilon\_\{i\}$"}
\NormalTok{  )) }\SpecialCharTok{+}
  \FunctionTok{theme}\NormalTok{(}
    \AttributeTok{legend.position =} \StringTok{"none"}\NormalTok{,}
    \AttributeTok{plot.title =} \FunctionTok{element\_text}\NormalTok{(}\AttributeTok{hjust =} \FloatTok{0.5}\NormalTok{)}
\NormalTok{  )}
\end{Highlighting}
\end{Shaded}

\begin{figure}

{\centering \includegraphics[width=0.7\linewidth]{docs/figures/fig-regresion-simple-1} 

}

\caption{Modelo de regresión lineal simple para el ingreso en función del gasto (sin considerar estratos)}\label{fig:fig-regresion-simple}
\end{figure}

Si bien este modelo resulta útil como referencia inicial, sus supuestos suelen ser poco realistas en el contexto de encuestas de hogares. En particular, la independencia entre observaciones puede verse comprometida cuando los hogares pertenecen a estratos que comparten características socioeconómicas, demográficas o geográficas similares. Asimismo, es posible que los niveles promedio de ingreso difieran sistemáticamente entre estratos, aun cuando la relación entre ingreso y gasto sea semejante. Como señalan \citet{finch2019multilevel}, ignorar esta estructura jerárquica puede conducir a estimaciones incorrectas de los errores estándar y, en consecuencia, a inferencias estadísticas poco confiables.

Con fines exclusivamente ilustrativos, se incorpora gradualmente la estructura jerárquica de los datos al proceso de modelación. Para ello, se ajusta inicialmente un modelo que permite que cada estrato tenga su propio intercepto, mientras mantiene una pendiente común para todas las observaciones. Esta especificación resulta útil para visualizar cómo los niveles promedio de ingreso pueden diferir entre estratos, aun cuando la relación entre ingreso y gasto se considere constante.

\begin{Shaded}
\begin{Highlighting}[]
\NormalTok{beta\_1 }\OtherTok{\textless{}{-}} \FunctionTok{coef}\NormalTok{(}\FunctionTok{lm}\NormalTok{(Income }\SpecialCharTok{\textasciitilde{}}\NormalTok{ Expenditure, }\AttributeTok{data =}\NormalTok{ plot\_data))[}\DecValTok{2}\NormalTok{]}
\NormalTok{model\_coef }\OtherTok{\textless{}{-}}\NormalTok{ plot\_data }\SpecialCharTok{\%\textgreater{}\%}
  \FunctionTok{group\_by}\NormalTok{(Stratum) }\SpecialCharTok{\%\textgreater{}\%}
  \FunctionTok{summarise}\NormalTok{(}\AttributeTok{beta\_0 =} \FunctionTok{coef}\NormalTok{(}\FunctionTok{lm}\NormalTok{(Income }\SpecialCharTok{\textasciitilde{}}\NormalTok{ Expenditure))[}\DecValTok{1}\NormalTok{]) }\SpecialCharTok{\%\textgreater{}\%}
  \FunctionTok{mutate}\NormalTok{(}\AttributeTok{beta\_1 =}\NormalTok{ beta\_1)}
\end{Highlighting}
\end{Shaded}

De esta manera, el modelo introduce una primera fuente de heterogeneidad entre grupos y permite apreciar las limitaciones del modelo de regresión simple presentado anteriormente, el cual asumía una única recta de regresión para toda la población. La Figura \ref{fig:fig-intercepto-por-estrato} presenta las rectas de regresión con interceptos diferenciados por estrato:

\begin{Shaded}
\begin{Highlighting}[]
\FunctionTok{ggplot}\NormalTok{(}\AttributeTok{data =}\NormalTok{ plot\_data,}
       \FunctionTok{aes}\NormalTok{(}\AttributeTok{y =}\NormalTok{ Income, }\AttributeTok{x =}\NormalTok{ Expenditure, }\AttributeTok{colour =}\NormalTok{ Stratum)) }\SpecialCharTok{+}
  \FunctionTok{geom\_jitter}\NormalTok{() }\SpecialCharTok{+}
  \FunctionTok{geom\_abline}\NormalTok{(}\AttributeTok{data =}\NormalTok{ model\_coef,}
              \AttributeTok{mapping =} \FunctionTok{aes}\NormalTok{(}\AttributeTok{slope =}\NormalTok{ beta\_1, }\AttributeTok{intercept =}\NormalTok{ beta\_0, }\AttributeTok{colour =}\NormalTok{ Stratum)) }\SpecialCharTok{+}
  \FunctionTok{ggtitle}\NormalTok{(latex2exp}\SpecialCharTok{::}\FunctionTok{TeX}\NormalTok{(}
    \StringTok{"$Income\_\{kj\} }\SpecialCharTok{\textbackslash{}\textbackslash{}}\StringTok{sim }\SpecialCharTok{\textbackslash{}\textbackslash{}}\StringTok{hat\{}\SpecialCharTok{\textbackslash{}\textbackslash{}}\StringTok{beta\}\_\{0j\} + }\SpecialCharTok{\textbackslash{}\textbackslash{}}\StringTok{hat\{}\SpecialCharTok{\textbackslash{}\textbackslash{}}\StringTok{beta\}\_\{1\}Expenditure\_\{kj\} + }\SpecialCharTok{\textbackslash{}\textbackslash{}}\StringTok{epsilon\_\{kj\}$"}
\NormalTok{  )) }\SpecialCharTok{+}
  \FunctionTok{theme}\NormalTok{(}
    \AttributeTok{legend.position =} \StringTok{"none"}\NormalTok{,}
    \AttributeTok{plot.title =} \FunctionTok{element\_text}\NormalTok{(}\AttributeTok{hjust =} \FloatTok{0.5}\NormalTok{)}
\NormalTok{  )}
\end{Highlighting}
\end{Shaded}

\begin{figure}

{\centering \includegraphics[width=0.7\linewidth]{docs/figures/fig-intercepto-por-estrato-1} 

}

\caption{Modelo con intercepto específico por estrato y pendiente común}\label{fig:fig-intercepto-por-estrato}
\end{figure}

Como siguiente paso, se considera una especificación alternativa en la que todos los estratos comparten un mismo nivel promedio de ingreso, representado por un intercepto común, pero se permite que la relación entre ingreso y gasto varíe entre estratos mediante pendientes específicas para cada uno de ellos.

\begin{Shaded}
\begin{Highlighting}[]
\NormalTok{beta\_0 }\OtherTok{\textless{}{-}} \FunctionTok{coef}\NormalTok{(}\FunctionTok{lm}\NormalTok{(Income }\SpecialCharTok{\textasciitilde{}}\NormalTok{ Expenditure, }\AttributeTok{data =}\NormalTok{ plot\_data))[}\DecValTok{1}\NormalTok{]}
\NormalTok{model\_coef }\OtherTok{\textless{}{-}}\NormalTok{ plot\_data }\SpecialCharTok{\%\textgreater{}\%}
  \FunctionTok{group\_by}\NormalTok{(Stratum) }\SpecialCharTok{\%\textgreater{}\%}
  \FunctionTok{summarise}\NormalTok{(}\AttributeTok{beta\_1 =} \FunctionTok{coef}\NormalTok{(}\FunctionTok{lm}\NormalTok{(Income }\SpecialCharTok{\textasciitilde{}}\NormalTok{ Expenditure))[}\DecValTok{2}\NormalTok{]) }\SpecialCharTok{\%\textgreater{}\%}
  \FunctionTok{mutate}\NormalTok{(}\AttributeTok{beta\_0 =}\NormalTok{ beta\_0)}
\end{Highlighting}
\end{Shaded}

Bajo esta especificación, las diferencias entre estratos se manifiestan en la intensidad de la asociación entre ingreso y gasto, más que en sus niveles promedio. La Figura \ref{fig:fig-pendiente-por-estrato} permite visualizar estas diferencias en las trayectorias de regresión estimadas para cada estrato.

\begin{Shaded}
\begin{Highlighting}[]
\FunctionTok{ggplot}\NormalTok{(}\AttributeTok{data =}\NormalTok{ plot\_data,}
       \FunctionTok{aes}\NormalTok{(}\AttributeTok{y =}\NormalTok{ Income, }\AttributeTok{x =}\NormalTok{ Expenditure, }\AttributeTok{colour =}\NormalTok{ Stratum)) }\SpecialCharTok{+}
  \FunctionTok{geom\_jitter}\NormalTok{() }\SpecialCharTok{+}
  \FunctionTok{geom\_abline}\NormalTok{(}\AttributeTok{data =}\NormalTok{ model\_coef,}
              \AttributeTok{mapping =} \FunctionTok{aes}\NormalTok{(}\AttributeTok{slope =}\NormalTok{ beta\_1, }\AttributeTok{intercept =}\NormalTok{ beta\_0, }\AttributeTok{colour =}\NormalTok{ Stratum)) }\SpecialCharTok{+}
  \FunctionTok{ggtitle}\NormalTok{(latex2exp}\SpecialCharTok{::}\FunctionTok{TeX}\NormalTok{(}
    \StringTok{"$Income\_\{kj\} }\SpecialCharTok{\textbackslash{}\textbackslash{}}\StringTok{sim }\SpecialCharTok{\textbackslash{}\textbackslash{}}\StringTok{hat\{}\SpecialCharTok{\textbackslash{}\textbackslash{}}\StringTok{beta\}\_\{0\} + }\SpecialCharTok{\textbackslash{}\textbackslash{}}\StringTok{hat\{}\SpecialCharTok{\textbackslash{}\textbackslash{}}\StringTok{beta\}\_\{1j\}Expenditure\_\{kj\} + }\SpecialCharTok{\textbackslash{}\textbackslash{}}\StringTok{epsilon\_\{kj\}$"}
\NormalTok{  )) }\SpecialCharTok{+}
  \FunctionTok{theme}\NormalTok{(}
    \AttributeTok{legend.position =} \StringTok{"none"}\NormalTok{,}
    \AttributeTok{plot.title =} \FunctionTok{element\_text}\NormalTok{(}\AttributeTok{hjust =} \FloatTok{0.5}\NormalTok{)}
\NormalTok{  )}
\end{Highlighting}
\end{Shaded}

\begin{figure}

{\centering \includegraphics[width=0.7\linewidth]{docs/figures/fig-pendiente-por-estrato-1} 

}

\caption{Modelo con pendiente específica por estrato e intercepto común}\label{fig:fig-pendiente-por-estrato}
\end{figure}

Finalmente, se considera la especificación más flexible de las presentadas hasta el momento, permitiendo que tanto el nivel promedio de ingreso como la relación entre ingreso y gasto varíen entre estratos. Bajo este enfoque, cada estrato posee su propia recta de regresión, con interceptos y pendientes estimados de manera independiente.

\begin{Shaded}
\begin{Highlighting}[]
\NormalTok{model\_coef }\OtherTok{\textless{}{-}}\NormalTok{ plot\_data }\SpecialCharTok{\%\textgreater{}\%}
  \FunctionTok{group\_by}\NormalTok{(Stratum) }\SpecialCharTok{\%\textgreater{}\%}
  \FunctionTok{summarise}\NormalTok{(}
    \AttributeTok{beta\_0 =} \FunctionTok{coef}\NormalTok{(}\FunctionTok{lm}\NormalTok{(Income }\SpecialCharTok{\textasciitilde{}}\NormalTok{ Expenditure))[}\DecValTok{1}\NormalTok{],}
    \AttributeTok{beta\_1 =} \FunctionTok{coef}\NormalTok{(}\FunctionTok{lm}\NormalTok{(Income }\SpecialCharTok{\textasciitilde{}}\NormalTok{ Expenditure))[}\DecValTok{2}\NormalTok{]}
\NormalTok{  )}
\end{Highlighting}
\end{Shaded}

Esta formulación permite capturar simultáneamente diferencias en los niveles de ingreso y en la intensidad de su asociación con el gasto. La Figura \ref{fig:fig-intercepto-pendiente-libres} muestra las rectas ajustadas para cada estrato bajo esta especificación.

\begin{Shaded}
\begin{Highlighting}[]
\FunctionTok{ggplot}\NormalTok{(}\AttributeTok{data =}\NormalTok{ plot\_data,}
       \FunctionTok{aes}\NormalTok{(}\AttributeTok{y =}\NormalTok{ Income, }\AttributeTok{x =}\NormalTok{ Expenditure, }\AttributeTok{colour =}\NormalTok{ Stratum)) }\SpecialCharTok{+}
  \FunctionTok{geom\_jitter}\NormalTok{() }\SpecialCharTok{+}
  \FunctionTok{geom\_abline}\NormalTok{(}\AttributeTok{data =}\NormalTok{ model\_coef,}
              \AttributeTok{mapping =} \FunctionTok{aes}\NormalTok{(}\AttributeTok{slope =}\NormalTok{ beta\_1, }\AttributeTok{intercept =}\NormalTok{ beta\_0, }\AttributeTok{colour =}\NormalTok{ Stratum)) }\SpecialCharTok{+}
  \FunctionTok{ggtitle}\NormalTok{(latex2exp}\SpecialCharTok{::}\FunctionTok{TeX}\NormalTok{(}
    \StringTok{"$Income\_\{kj\} }\SpecialCharTok{\textbackslash{}\textbackslash{}}\StringTok{sim }\SpecialCharTok{\textbackslash{}\textbackslash{}}\StringTok{hat\{}\SpecialCharTok{\textbackslash{}\textbackslash{}}\StringTok{beta\}\_\{0j\} + }\SpecialCharTok{\textbackslash{}\textbackslash{}}\StringTok{hat\{}\SpecialCharTok{\textbackslash{}\textbackslash{}}\StringTok{beta\}\_\{1j\}Expenditure\_\{kj\} + }\SpecialCharTok{\textbackslash{}\textbackslash{}}\StringTok{epsilon\_\{kj\}$"}
\NormalTok{  )) }\SpecialCharTok{+}
  \FunctionTok{theme}\NormalTok{(}
    \AttributeTok{legend.position =} \StringTok{"none"}\NormalTok{,}
    \AttributeTok{plot.title =} \FunctionTok{element\_text}\NormalTok{(}\AttributeTok{hjust =} \FloatTok{0.5}\NormalTok{)}
\NormalTok{  )}
\end{Highlighting}
\end{Shaded}

\begin{figure}

{\centering \includegraphics[width=0.7\linewidth]{docs/figures/fig-intercepto-pendiente-libres-1} 

}

\caption{Modelo con intercepto y pendiente específicos por estrato (ajuste independiente)}\label{fig:fig-intercepto-pendiente-libres}
\end{figure}

Aunque esta última especificación ofrece una representación más flexible de los datos y permite capturar las particularidades de cada estrato, presenta una limitación importante: los parámetros se estiman de manera completamente independiente para cada grupo. Como consecuencia, no se aprovecha la información compartida entre estratos que podrían presentar patrones similares, lo que puede generar estimaciones inestables, especialmente en aquellos grupos con un número reducido de observaciones.

Los modelos multinivel superan esta dificultad mediante un mecanismo conocido como encogimiento (\emph{shrinkage}), en el cual las estimaciones específicas de cada grupo se obtienen combinando la información propia del grupo con la información proveniente del conjunto de la población. De esta manera, se logra un equilibrio entre la representación de las diferencias entre estratos y la eficiencia estadística de las estimaciones.

\section{\texorpdfstring{Pesos \emph{q-weighted}}{Pesos q-weighted}}\label{pesos-q-weighted}

En relación con la incorporación de pesos de muestreo en modelos multinivel, aunque existe un consenso general sobre la necesidad de utilizarlos para garantizar inferencias representativas de la población objetivo \citep{cai2013investigation}, no existe una única estrategia metodológica para su implementación. En este contexto, \citet{pfeffermann1998weighting} y \citet{asparouhov2006general} proponen enfoques basados en pseudo-máxima verosimilitud, en particular mediante procedimientos de mínimos cuadrados generalizados ponderados. Alternativamente, \citet{rabehesketh2006multilevel} proponen métodos basados en algoritmos EM para la estimación de modelos jerárquicos con ponderaciones.

Desde una perspectiva general, el enfoque de máxima verosimilitud consiste en estimar los parámetros del modelo identificando aquellos valores que maximizan la probabilidad de observar los datos disponibles bajo la especificación asumida. Una extensión importante de este enfoque es la máxima verosimilitud restringida (REML), la cual mejora la estimación de los componentes de varianza al corregir el sesgo asociado a la estimación de los grados de libertad, produciendo en muchos casos estimaciones más estables que la máxima verosimilitud convencional \citep{kreft1998introducing}.

En el contexto de los modelos multinivel con datos provenientes de encuestas, una dificultad adicional radica en la incorporación coherente de los pesos de muestreo a través de los distintos niveles de la jerarquía. Como señalan \citet{pfeffermann1998weighting}, la estructura agrupada de los datos implica que las observaciones no son independientes, por lo que la función de log-verosimilitud no puede descomponerse como una suma simple de contribuciones individuales. En su lugar, es necesario considerar explícitamente la dependencia entre los distintos niveles del diseño muestral para lograr inferencias adecuadas.

Para ajustar un modelo multinivel con datos provenientes de encuestas complejas, se recomienda redefinir los factores de expansión mediante el enfoque de pesos \emph{q-weighted} propuesto por \citet{pfeffermann2011modelling}. Este procedimiento busca eliminar la parte sistemática de los pesos asociada a las covariables incluidas en el modelo, conservando únicamente la componente residual de los mecanismos de selección. Los pasos son los siguientes:

\begin{enumerate}
\def\labelenumi{\arabic{enumi}.}
\item
  Ajustar un modelo de regresión para los factores de expansión finales de la encuesta, utilizando como variables explicativas el mismo conjunto de covariables que será empleado en el modelo multinivel. Siendo \(w_k\) el factor de expansión original de la unidad \(k\) y \(\mathbf{x}_k\) el vector de covariables. Entonces se ajusta el modelo

  \[
  w_k = f(\mathbf{x}_k, \boldsymbol{\beta}) + \varepsilon_k,
  \]

  donde \(f(\cdot)\) representa la relación funcional entre los pesos \(w_k\) y las covariables \(\mathbf{x}_k\), a través de los coeficientes \(\boldsymbol{\beta}\).
\item
  Obtener los valores predichos del modelo para cada unidad de observación, de tal forma que:

  \[
  \hat{w}_k = f(\mathbf{x}_k, \hat{\boldsymbol{\beta}}).
  \]

  Estos valores representan la componente de los factores de expansión explicada por las covariables incluidas en el modelo.
\item
  Construir los pesos \emph{q-weighted} dividiendo los factores de expansión originales por sus correspondientes valores predichos,

  \[
  q_k = \frac{w_k}{\hat{w}_k}.
  \]

  De esta forma, los nuevos pesos reflejan únicamente la variación residual de los factores de expansión una vez controlado el efecto de las covariables.
\item
  Definir el nuevo diseño muestral utilizando los pesos ajustados \(q_k\) en lugar de los pesos originales \(w_k\), y emplear este diseño para la estimación de los parámetros del modelo multinivel.
\end{enumerate}

El siguiente código implementa el procedimiento de construcción de pesos \emph{q-weighted}. En primer lugar, ajusta un modelo de regresión lineal que explica los factores de expansión originales a partir de la variable explicativa gasto. De esta forma, los pesos resultantes conservan únicamente la componente de variación no explicada por las covariables incluidas en el modelo multinivel.

\begin{Shaded}
\begin{Highlighting}[]
\NormalTok{mod\_qw }\OtherTok{\textless{}{-}} \FunctionTok{lm}\NormalTok{(wk }\SpecialCharTok{\textasciitilde{}}\NormalTok{ Expenditure, }\AttributeTok{data =}\NormalTok{ survey\_data)}
\NormalTok{survey\_data}\SpecialCharTok{$}\NormalTok{qk }\OtherTok{\textless{}{-}}\NormalTok{ survey\_data}\SpecialCharTok{$}\NormalTok{wk }\SpecialCharTok{/} \FunctionTok{predict}\NormalTok{(mod\_qw)}
\end{Highlighting}
\end{Shaded}

\section{Modelo lineal multinivel}\label{modelo-lineal-multinivel}

Los modelos multinivel más utilizados en la práctica corresponden a extensiones del modelo lineal clásico para datos con estructura jerárquica. En su formulación básica, estos modelos asumen que la variable respuesta es continua y sigue una distribución normal condicional a los efectos fijos y aleatorios incluidos en el modelo. Asimismo, se supone que los efectos aleatorios y los errores residuales son independientes entre sí y se distribuyen normalmente con media cero y varianzas constantes.

\subsection{Modelo nulo}\label{sec-multinivel-muestras-complejas}

En el análisis de regresión multinivel se distinguen dos tipos de parámetros: los coeficientes de regresión, conocidos como parámetros fijos, y los componentes de varianza asociados a los efectos aleatorios. Una etapa inicial fundamental en este tipo de análisis consiste en descomponer la variabilidad total de la variable dependiente en los distintos niveles de la estructura jerárquica. En el contexto del ejemplo anterior, esta descomposición permite separar la variación del ingreso en una componente atribuible a diferencias dentro de los estratos y otra asociada a diferencias entre estratos.

El punto de partida para esta descomposición es el denominado modelo nulo, el cual no incorpora variables explicativas y se expresa como

\[
y_{kj} = \beta_{0j} + \epsilon_{kj}
\]

En donde \(y_{kj}\) denota el valor observado del ingreso para la unidad \(k\) en el estrato \(j\). El término \(\beta_{0j}\) representa el intercepto específico de cada estrato, capturando las diferencias en los niveles promedio de la variable entre grupos. Finalmente, \(\epsilon_{kj}\) corresponde al error a nivel de unidad, el cual recoge la variabilidad no explicada dentro de cada estrato y se asume con media cero y varianza constante condicional al grupo.

A su vez, el intercepto por estrato puede descomponerse en una parte común a todos los estratos y una desviación específica de cada uno de ellos, de la siguiente forma:

\[
\beta_{0j} = \gamma_{00} + \tau_{0j}
\]

donde \(\gamma_{00}\) representa el intercepto promedio global y \(\tau_{0j}\) es el efecto aleatorio para el intercepto que captura la desviación del estrato \(j\), permitiendo modelar explícitamente la heterogeneidad entre estratos. Además, los componentes aleatorios del modelo se asumen distribuidos normalmente con media cero y varianzas específicas para cada nivel de la jerarquía. En particular, el efecto aleatorio asociado al intercepto por estrato cumple

\[
\tau_{0j} \sim N\left(0, \sigma_{\tau}^{2}\right),
\]

lo que implica que las desviaciones de cada estrato respecto al intercepto global se distribuyen alrededor de cero, con una variabilidad determinada por \(\sigma_{\tau}^{2}\). De manera análoga, el término de error a nivel de unidad sigue la distribución

\[
\epsilon_{kj} \sim N\left(0, \sigma_{\epsilon}^{2}\right),
\]

donde \(\sigma_{\epsilon}^{2}\) representa la variabilidad residual dentro de los estratos. A partir de este modelo se define el coeficiente de correlación intraclase (ICC, por sus siglas en inglés), que cuantifica la proporción de la varianza total explicada por las diferencias entre estratos:

\[
\rho = \frac{\sigma_{\tau}^{2}}{\sigma_{\tau}^{2} + \sigma_{\epsilon}^{2}}
\]

Un valor elevado del ICC indica que una proporción importante de la variabilidad total de la variable de interés se debe a diferencias entre estratos, lo que sugiere la presencia de heterogeneidad relevante entre grupos y la necesidad de incorporar explícitamente dicha estructura en el modelo. En contraste, un ICC bajo implica que la variabilidad se concentra principalmente dentro de los estratos, los cuales presentan una mayor homogeneidad relativa entre sí.

Siguiendo con el análisis de la encuesta de ejemplo, cuando la variable de interés es el ingreso, el modelo nulo constituye el punto de partida para cuantificar qué proporción de la variabilidad total se debe a diferencias entre estratos y qué proporción corresponde a diferencias entre hogares dentro de un mismo estrato.

La estimación de modelos multinivel en \texttt{R} se realiza mediante el paquete \texttt{lme4} \citep{lme42026}. En particular, la función \texttt{lmer()} permite especificar los efectos aleatorios mediante expresiones de la forma \texttt{(efecto\ \textbar{}\ grupo)}. En este contexto, la expresión \texttt{(1\ \textbar{}\ Stratum)} indica que el modelo incluye un intercepto aleatorio para cada estrato. El valor \texttt{1} representa el intercepto del modelo, mientras que \texttt{Stratum} identifica la variable de agrupación.

\begin{Shaded}
\begin{Highlighting}[]
\NormalTok{mod\_null }\OtherTok{\textless{}{-}} \FunctionTok{lmer}\NormalTok{(}
\NormalTok{  Income }\SpecialCharTok{\textasciitilde{}}\NormalTok{ (}\DecValTok{1} \SpecialCharTok{|}\NormalTok{ Stratum),}
  \AttributeTok{data =}\NormalTok{ survey\_data,}
  \AttributeTok{weights =}\NormalTok{ qk}
\NormalTok{)}
\end{Highlighting}
\end{Shaded}

La Tabla \ref{tab:tab-interceptos-nulos} presenta los interceptos estimados para cada estrato a partir del modelo nulo. Dado que este modelo no incorpora variables explicativas, cada intercepto puede interpretarse como el ingreso promedio esperado en el correspondiente estrato. Los resultados evidencian diferencias importantes entre grupos: mientras que el estrato \texttt{idStrt004} presenta el mayor ingreso promedio estimado (959.6), el estrato \texttt{idStrt009} registra el menor valor (207.6). Estas diferencias sugieren la existencia de una variabilidad sustancial entre estratos, justificando la incorporación de efectos aleatorios para modelar explícitamente la estructura jerárquica de los datos.

\begin{Shaded}
\begin{Highlighting}[]
\NormalTok{coef\_mod\_null }\OtherTok{\textless{}{-}} \FunctionTok{coef}\NormalTok{(mod\_null)}\SpecialCharTok{$}\NormalTok{Stratum}
\NormalTok{coef\_mod\_null }\SpecialCharTok{\%\textgreater{}\%}
  \FunctionTok{slice}\NormalTok{(}\DecValTok{1}\SpecialCharTok{:}\DecValTok{12}\DataTypeTok{L}\NormalTok{)}
\end{Highlighting}
\end{Shaded}

\begin{longtable}[]{@{}ll@{}}
\caption{\label{tab:tab-interceptos-nulos}Interceptos estimados por estrato bajo pesos q-weighted y senate-weighted (modelo nulo, primeros 12 estratos)}\tabularnewline
\toprule\noalign{}
& (Intercept) \\
\midrule\noalign{}
\endfirsthead
\toprule\noalign{}
& (Intercept) \\
\midrule\noalign{}
\endhead
\bottomrule\noalign{}
\endlastfoot
idStrt001 & 635 \\
idStrt002 & 507 \\
idStrt003 & 486 \\
idStrt004 & 960 \\
idStrt005 & 518 \\
idStrt006 & 439 \\
idStrt007 & 477 \\
idStrt008 & 377 \\
idStrt009 & 218 \\
idStrt010 & 594 \\
idStrt011 & 590 \\
idStrt012 & 362 \\
\end{longtable}

A partir de los componentes de varianza estimados en el modelo nulo, la correlación intraclase se calcula mediante la función \texttt{icc()} del paquete \texttt{performance} \citep{performance2026}. Este indicador cuantifica la proporción de la variabilidad total del ingreso que puede atribuirse a diferencias entre estratos, constituyendo una medida de la dependencia existente entre observaciones pertenecientes a un mismo grupo. Los resultados obtenidos se presentan en la Tabla \ref{tab:tab-icc-nulo}.

\begin{Shaded}
\begin{Highlighting}[]
\NormalTok{performance}\SpecialCharTok{::}\FunctionTok{icc}\NormalTok{(mod\_null) }\SpecialCharTok{\%\textgreater{}\%}
  \FunctionTok{as.data.frame}\NormalTok{()}
\end{Highlighting}
\end{Shaded}

\begin{longtable}[]{@{}lcc@{}}
\caption{\label{tab:tab-icc-nulo}Correlación intraclase del modelo nulo (ingreso por estrato)}\tabularnewline
\toprule\noalign{}
ICC\_adjusted & ICC\_unadjusted & optional \\
\midrule\noalign{}
\endfirsthead
\toprule\noalign{}
ICC\_adjusted & ICC\_unadjusted & optional \\
\midrule\noalign{}
\endhead
\bottomrule\noalign{}
\endlastfoot
0.329 & 0.329 & FALSE \\
\end{longtable}

La correlación intraclase estimada de 32\% indica que aproximadamente un tercio de la variabilidad total observada en el ingreso se debe a diferencias entre estratos, mientras que el porcentaje restante corresponde a diferencias entre hogares dentro de los mismos estratos. Por último, como el modelo nulo no incluye ningún predictor, la estimación del ingreso dentro de cada estrato es constante e igual al intercepto estimado para ese estrato.

\subsection{Modelo con intercepto aleatorio}\label{sec-intercepto-aleatorio}

El modelo multinivel más simple que incorpora covariables corresponde al modelo de intercepto aleatorio. En esta especificación, se asume que los efectos de las variables explicativas son comunes a todos los estratos, mientras que el nivel promedio de la variable respuesta puede variar entre ellos. El modelo se expresa como

\[
y_{kj} = \beta_{0j} + \mathbf{x}_{kj}\boldsymbol{\beta} + \epsilon_{kj}
\]

En donde \(y_{kj}\) representa el valor observado de la variable respuesta para la unidad \(k\) perteneciente al estrato \(j\), \(\mathbf{x}_{kj}\) es el vector de covariables asociado a dicha unidad, \(\boldsymbol{\beta}\) es el vector de coeficientes de regresión comunes a todos los estratos y \(\epsilon_{kj}\) corresponde al error aleatorio a nivel de unidad.

Al igual que en el modelo nulo, el intercepto se modela como \(\beta_{0j} = \gamma_{00} + \tau_{0j}\); en donde \(\gamma_{00}\) representa el intercepto promedio global y \(\tau_{0j}\) es el efecto aleatorio asociado al estrato \(j\), que mide la desviación de dicho estrato respecto al promedio general.

Los componentes aleatorios del modelo se asumen independientes y normalmente distribuidos, de manera que \(\tau_{0j} \sim N(0,\sigma_{\tau}^{2})\) y \(\epsilon_{kj} \sim N(0,\sigma_{\epsilon}^{2})\). Bajo estas condiciones, la variabilidad total de la variable respuesta puede descomponerse en una componente entre estratos, cuantificada por \(\sigma_{\tau}^{2}\), y una componente dentro de los estratos, representada por \(\sigma_{\epsilon}^{2}\).

El siguiente código ajusta un modelo multinivel de intercepto aleatorio, en donde el ingreso (\texttt{Income}) se explica a partir del gasto del hogar (\texttt{Expenditure}), incorporando además un efecto aleatorio asociado al estrato (\texttt{Stratum}) y utilizando los pesos \emph{q-weighted} almacenados en la variable \texttt{qk}. Una vez ajustado el modelo, se estima la correlación intraclase (ICC) a partir de los componentes de varianza estimados.

\begin{Shaded}
\begin{Highlighting}[]
\NormalTok{random\_intercept\_model }\OtherTok{\textless{}{-}} \FunctionTok{lmer}\NormalTok{(}
\NormalTok{  Income }\SpecialCharTok{\textasciitilde{}}\NormalTok{ Expenditure }\SpecialCharTok{+}\NormalTok{ (}\DecValTok{1} \SpecialCharTok{|}\NormalTok{ Stratum),}
  \AttributeTok{data =}\NormalTok{ survey\_data,}
  \AttributeTok{weights =}\NormalTok{ qk}
\NormalTok{)}
\NormalTok{performance}\SpecialCharTok{::}\FunctionTok{icc}\NormalTok{(random\_intercept\_model)}
\end{Highlighting}
\end{Shaded}

\begin{verbatim}
## # Intraclass Correlation Coefficient
## 
##     Adjusted ICC: 0.203
##   Unadjusted ICC: 0.109
\end{verbatim}

La correlación intraclase ajustada de 0.196 indica que, después de controlar por el gasto del hogar (\texttt{Expenditure}), aproximadamente el 20\% de la variabilidad residual del ingreso sigue siendo atribuible a diferencias entre estratos. Aunque este valor es inferior al obtenido con el modelo nulo, la magnitud del ICC continúa justificando el uso de un modelo multinivel para representar adecuadamente la estructura jerárquica de los datos. Los coeficientes estimados por estrato se presentan en la Tabla \ref{tab:tab-coef-intercepto-aleatorio}:

\begin{Shaded}
\begin{Highlighting}[]
\FunctionTok{coef}\NormalTok{(random\_intercept\_model)}\SpecialCharTok{$}\NormalTok{Stratum }\SpecialCharTok{\%\textgreater{}\%}
  \FunctionTok{slice}\NormalTok{(}\DecValTok{1}\SpecialCharTok{:}\DecValTok{8}\DataTypeTok{L}\NormalTok{)}
\end{Highlighting}
\end{Shaded}

\begin{longtable}[]{@{}llc@{}}
\caption{\label{tab:tab-coef-intercepto-aleatorio}Coeficientes del modelo con intercepto aleatorio por estrato (primeros 8 estratos)}\tabularnewline
\toprule\noalign{}
& (Intercept) & Expenditure \\
\midrule\noalign{}
\endfirsthead
\toprule\noalign{}
& (Intercept) & Expenditure \\
\midrule\noalign{}
\endhead
\bottomrule\noalign{}
\endlastfoot
idStrt001 & 250.4 & 1.19 \\
idStrt002 & 156.8 & 1.19 \\
idStrt003 & 142.9 & 1.19 \\
idStrt004 & 296.8 & 1.19 \\
idStrt005 & -32.4 & 1.19 \\
idStrt006 & 53.8 & 1.19 \\
idStrt007 & 12.5 & 1.19 \\
idStrt008 & 107.1 & 1.19 \\
\end{longtable}

Con el fin de facilitar la visualización de los resultados del modelo, se utiliza la submuestra reducida presentada anteriormente, la cual contiene únicamente tres estratos seleccionados.

\begin{Shaded}
\begin{Highlighting}[]
\NormalTok{estimated\_coef }\OtherTok{\textless{}{-}} \FunctionTok{inner\_join}\NormalTok{(}
  \FunctionTok{coef}\NormalTok{(random\_intercept\_model)}\SpecialCharTok{$}\NormalTok{Stratum }\SpecialCharTok{\%\textgreater{}\%}
\NormalTok{    tibble}\SpecialCharTok{::}\FunctionTok{rownames\_to\_column}\NormalTok{(}\AttributeTok{var =} \StringTok{"Stratum"}\NormalTok{),}
\NormalTok{  plot\_data }\SpecialCharTok{\%\textgreater{}\%}
    \FunctionTok{select}\NormalTok{(Stratum) }\SpecialCharTok{\%\textgreater{}\%}
    \FunctionTok{distinct}\NormalTok{()}
\NormalTok{)}
\end{Highlighting}
\end{Shaded}

La Figura \ref{fig:fig-intercepto-aleatorio} muestra las rectas de regresión estimadas para cada estrato bajo el modelo de intercepto aleatorio. Se observa que las rectas comparten una misma pendiente, reflejando el efecto común del gasto sobre el ingreso, pero difieren en su posición vertical debido a los interceptos específicos estimados para cada estrato.

\begin{Shaded}
\begin{Highlighting}[]
\FunctionTok{ggplot}\NormalTok{(}\AttributeTok{data =}\NormalTok{ plot\_data,}
       \FunctionTok{aes}\NormalTok{(}\AttributeTok{y =}\NormalTok{ Income, }\AttributeTok{x =}\NormalTok{ Expenditure, }\AttributeTok{colour =}\NormalTok{ Stratum)) }\SpecialCharTok{+}
  \FunctionTok{geom\_jitter}\NormalTok{() }\SpecialCharTok{+}
  \FunctionTok{geom\_abline}\NormalTok{(}\AttributeTok{data =}\NormalTok{ estimated\_coef,}
              \AttributeTok{mapping =} \FunctionTok{aes}\NormalTok{(}\AttributeTok{slope =}\NormalTok{ Expenditure,}
                            \AttributeTok{intercept =} \StringTok{\textasciigrave{}}\AttributeTok{(Intercept)}\StringTok{\textasciigrave{}}\NormalTok{,}
                            \AttributeTok{colour =}\NormalTok{ Stratum)) }\SpecialCharTok{+}
  \FunctionTok{theme}\NormalTok{(}
    \AttributeTok{legend.position =} \StringTok{"none"}\NormalTok{,}
    \AttributeTok{plot.title =} \FunctionTok{element\_text}\NormalTok{(}\AttributeTok{hjust =} \FloatTok{0.5}\NormalTok{)}
\NormalTok{  )}
\end{Highlighting}
\end{Shaded}

\begin{figure}

{\centering \includegraphics[width=0.7\linewidth]{docs/figures/fig-intercepto-aleatorio-1} 

}

\caption{Rectas de regresión estimadas por estrato: modelo con intercepto aleatorio}\label{fig:fig-intercepto-aleatorio}
\end{figure}

\subsection{Modelo con intercepto y pendiente aleatoria}\label{sec-intercepto-pendiente-aleatorios}

Una extensión natural del modelo anterior es el modelo de intercepto y pendiente aleatorios. En esta especificación, no sólo se permite que el nivel promedio de la variable respuesta varíe entre estratos, sino también que el efecto de una o más covariables difiera entre grupos. De esta manera, cada estrato puede tener su propia recta de regresión, tanto en términos de intercepto como de pendiente. El modelo puede expresarse como

\[
y_{kj} = \beta_{0j} + x_{kj}\beta_{1j} + \mathbf{z}_{kj}\boldsymbol{\beta} + \epsilon_{kj},
\]

donde \(y_{kj}\) representa el valor observado de la variable respuesta para la unidad \(k\) perteneciente al estrato \(j\), \(x_{kj}\) corresponde a la covariable cuya pendiente se permite variar entre estratos, \(\mathbf{z}_{kj}\) contiene las covariables con efectos fijos comunes a todos los grupos, y \(\epsilon_{kj}\) es el término de error a nivel de unidad. En este modelo, tanto el intercepto como la pendiente asociada a \(x_{kj}\) se consideran aleatorios y se descomponen de la siguiente forma:

\[
\beta_{0j} = \gamma_{00} + \tau_{0j}, \qquad
\beta_{1j} = \gamma_{10} + \tau_{1j}
\]

donde \(\gamma_{00}\) y \(\gamma_{10}\) representan, respectivamente, el intercepto y la pendiente promedio en la población, mientras que \(\tau_{0j}\) y \(\tau_{1j}\) capturan las desviaciones específicas del estrato \(j\) respecto a dichos promedios. Estos efectos aleatorios se asumen normalmente distribuidos como se indica a continuación:

\[
\begin{pmatrix}
\tau_{0j} \\
\tau_{1j}
\end{pmatrix}
\sim
N\left(
\begin{pmatrix}
0 \\
0
\end{pmatrix},
\begin{pmatrix}
\sigma_{\tau_0}^2 & \sigma_{\tau_{01}} \\
\sigma_{\tau_{01}} & \sigma_{\tau_1}^2
\end{pmatrix}
\right),
\]

mientras que el error a nivel de unidad satisface \(\epsilon_{kj} \sim N(0,\sigma_{\epsilon}^{2})\). En consecuencia, el modelo permite cuantificar no sólo la variabilidad entre estratos en los niveles promedio de la variable respuesta, sino también la heterogeneidad existente en el efecto de la covariable cuyo coeficiente se modela como aleatorio.

El siguiente código ajusta un modelo multinivel con intercepto y pendiente aleatorios, en el que el ingreso (\texttt{Income}) se explica a partir del gasto del hogar (\texttt{Expenditure}) y de la zona de residencia (\texttt{Zone}). En este caso, tanto el intercepto como el efecto del gasto pueden variar entre estratos mediante efectos aleatorios asociados a \texttt{Stratum}, empleando además los pesos \emph{q-weighted} almacenados en la variable \texttt{qk}.

\begin{Shaded}
\begin{Highlighting}[]
\NormalTok{random\_slope\_model }\OtherTok{\textless{}{-}} \FunctionTok{lmer}\NormalTok{(}
\NormalTok{  Income }\SpecialCharTok{\textasciitilde{}} \DecValTok{1} \SpecialCharTok{+}\NormalTok{ Expenditure }\SpecialCharTok{+}\NormalTok{ (}\DecValTok{1} \SpecialCharTok{+}\NormalTok{ Expenditure }\SpecialCharTok{|}\NormalTok{ Stratum),}
  \AttributeTok{data =}\NormalTok{ survey\_data,}
  \AttributeTok{weights =}\NormalTok{ qk}
\NormalTok{)}

\NormalTok{performance}\SpecialCharTok{::}\FunctionTok{icc}\NormalTok{(random\_slope\_model)}
\end{Highlighting}
\end{Shaded}

\begin{verbatim}
## # Intraclass Correlation Coefficient
## 
##     Adjusted ICC: 0.690
##   Unadjusted ICC: 0.457
\end{verbatim}

Nótese que en la especificación del código se incluye simultáneamente efectos fijos y efectos aleatorios para el intercepto y la variable \texttt{Expenditure}. Los términos ubicados fuera del paréntesis, \texttt{1\ +\ Expenditure}, corresponden a los efectos fijos del modelo y permiten estimar el intercepto promedio global y la pendiente promedio global asociada al gasto. Por su parte, la expresión \texttt{(1\ +\ Expenditure\ \textbar{}\ Stratum)} incorpora las desviaciones específicas de cada estrato respecto a dichos promedios. La inclusión explícita de los efectos fijos es fundamental, ya que los efectos aleatorios se interpretan como desviaciones alrededor de una media poblacional. Los coeficientes del modelo por estrato se presentan en la Tabla \ref{tab:tab-coef-pendiente-aleatoria}:

\begin{Shaded}
\begin{Highlighting}[]
\FunctionTok{coef}\NormalTok{(random\_slope\_model)}\SpecialCharTok{$}\NormalTok{Stratum }\SpecialCharTok{\%\textgreater{}\%}
  \FunctionTok{slice}\NormalTok{(}\DecValTok{1}\SpecialCharTok{:}\DecValTok{10}\DataTypeTok{L}\NormalTok{)}
\end{Highlighting}
\end{Shaded}

\begin{longtable}[]{@{}llc@{}}
\caption{\label{tab:tab-coef-pendiente-aleatoria}Coeficientes del modelo con intercepto y pendiente aleatoria por estrato (primeros 10 estratos)}\tabularnewline
\toprule\noalign{}
& (Intercept) & Expenditure \\
\midrule\noalign{}
\endfirsthead
\toprule\noalign{}
& (Intercept) & Expenditure \\
\midrule\noalign{}
\endhead
\bottomrule\noalign{}
\endlastfoot
idStrt001 & -222.8 & 2.730 \\
idStrt002 & 35.5 & 1.607 \\
idStrt003 & 151.9 & 1.164 \\
idStrt004 & 224.2 & 1.353 \\
idStrt005 & -89.7 & 1.286 \\
idStrt006 & 28.6 & 1.217 \\
idStrt007 & 41.0 & 1.079 \\
idStrt008 & 163.8 & 0.928 \\
idStrt009 & 16.5 & 0.838 \\
idStrt010 & 89.9 & 1.830 \\
\end{longtable}

Una vez más, con el propósito de facilitar la visualización de los resultados, se utiliza la submuestra reducida definida previamente, la cual considera únicamente tres estratos seleccionados.

\begin{Shaded}
\begin{Highlighting}[]
\NormalTok{estimated\_coef }\OtherTok{\textless{}{-}} \FunctionTok{inner\_join}\NormalTok{(}
  \FunctionTok{coef}\NormalTok{(random\_slope\_model)}\SpecialCharTok{$}\NormalTok{Stratum }\SpecialCharTok{\%\textgreater{}\%}
\NormalTok{    tibble}\SpecialCharTok{::}\FunctionTok{rownames\_to\_column}\NormalTok{(}\AttributeTok{var =} \StringTok{"Stratum"}\NormalTok{),}
\NormalTok{  plot\_data }\SpecialCharTok{\%\textgreater{}\%}
    \FunctionTok{select}\NormalTok{(Stratum) }\SpecialCharTok{\%\textgreater{}\%}
    \FunctionTok{distinct}\NormalTok{()}
\NormalTok{)}
\end{Highlighting}
\end{Shaded}

La Figura \ref{fig:fig-pendiente-aleatoria} presenta las rectas de regresión estimadas para cada estrato bajo el modelo de intercepto y pendiente aleatorios. A diferencia del modelo de intercepto aleatorio, en este caso las rectas difieren tanto en su posición vertical como en su inclinación, reflejando que los estratos no sólo presentan niveles promedio de ingreso distintos, sino también diferentes intensidades en la relación entre ingreso y gasto.

\begin{Shaded}
\begin{Highlighting}[]
\FunctionTok{ggplot}\NormalTok{(}\AttributeTok{data =}\NormalTok{ plot\_data,}
       \FunctionTok{aes}\NormalTok{(}\AttributeTok{y =}\NormalTok{ Income, }\AttributeTok{x =}\NormalTok{ Expenditure, }\AttributeTok{colour =}\NormalTok{ Stratum)) }\SpecialCharTok{+}
  \FunctionTok{geom\_jitter}\NormalTok{() }\SpecialCharTok{+}
  \FunctionTok{geom\_abline}\NormalTok{(}\AttributeTok{data =}\NormalTok{ estimated\_coef,}
              \AttributeTok{mapping =} \FunctionTok{aes}\NormalTok{(}\AttributeTok{slope =}\NormalTok{ Expenditure,}
                            \AttributeTok{intercept =} \StringTok{\textasciigrave{}}\AttributeTok{(Intercept)}\StringTok{\textasciigrave{}}\NormalTok{,}
                            \AttributeTok{colour =}\NormalTok{ Stratum)) }\SpecialCharTok{+}
  \FunctionTok{theme}\NormalTok{(}
    \AttributeTok{legend.position =} \StringTok{"none"}\NormalTok{,}
    \AttributeTok{plot.title =} \FunctionTok{element\_text}\NormalTok{(}\AttributeTok{hjust =} \FloatTok{0.5}\NormalTok{)}
\NormalTok{  )}
\end{Highlighting}
\end{Shaded}

\begin{figure}

{\centering \includegraphics[width=0.7\linewidth]{docs/figures/fig-pendiente-aleatoria-1} 

}

\caption{Rectas de regresión estimadas por estrato: modelo con intercepto y pendiente aleatoria}\label{fig:fig-pendiente-aleatoria}
\end{figure}

\section{Modelo logístico multinivel}\label{sec-logistica-multinivel}

Los modelos logísticos multinivel constituyen una extensión de los modelos multinivel para variables continuas al caso en que la variable respuesta es dicotómica. En lugar de modelar directamente el valor esperado de una variable continua, estos modelos estiman la probabilidad de ocurrencia de un evento, incorporando simultáneamente covariables individuales y efectos aleatorios asociados a los grupos de pertenencia. En el contexto de encuestas de hogares, esta formulación es especialmente útil cuando se analizan resultados binarios, como encontrarse o no en condición de pobreza, acceder o no a un servicio, participar o no en el mercado laboral, o presentar determinada característica sociodemográfica.

Al igual que en los modelos multinivel para variables continuas, el punto de partida consiste en reconocer que las unidades de observación no son independientes cuando pertenecen a un mismo estrato, conglomerado o dominio. Sin embargo, en el caso logístico la respuesta no se representa mediante una distribución normal con un error residual aditivo, sino mediante una distribución Bernoulli condicional a la probabilidad de ocurrencia del evento. Esta probabilidad se vincula con los predictores a través de la función logit, lo que permite expresar el modelo en una escala lineal.

\[
y_{kj} \mid \pi_{kj} \sim \text{Bernoulli}(\pi_{kj})
\]

donde \(y_{kj}\) toma el valor uno si la unidad \(k\) del estrato \(j\) presenta el evento de interés y cero en caso contrario. La probabilidad condicional del evento se denota por \(\pi_{kj}=Pr(y_{kj}=1)\) y se relaciona con el predictor lineal mediante

\[
\text{logit}(\pi_{kj}) =
\log\left(\frac{\pi_{kj}}{1-\pi_{kj}}\right)
= \eta_{kj}
\]

En esta formulación, \(\eta_{kj}\) cumple un papel análogo al valor esperado lineal de los modelos para variables continuas. La diferencia central es que los efectos fijos y aleatorios actúan sobre la escala logit y no directamente sobre la probabilidad. Por tanto, las diferencias entre estratos se interpretan como desplazamientos en los \emph{log-odds} del evento, los cuales luego se transforman a probabilidades mediante la función logística. La Figura \ref{fig:fig-logit-pobreza-gasto} presenta la relación entre gasto y probabilidad de pobreza, con una curva logística ajustada.

\begin{Shaded}
\begin{Highlighting}[]
\FunctionTok{ggplot}\NormalTok{(}\AttributeTok{data =}\NormalTok{ survey\_data,}
       \FunctionTok{aes}\NormalTok{(}\AttributeTok{y =}\NormalTok{ poor, }\AttributeTok{x =}\NormalTok{ Expenditure)) }\SpecialCharTok{+}
  \FunctionTok{geom\_point}\NormalTok{(}\AttributeTok{alpha =} \FloatTok{0.5}\NormalTok{) }\SpecialCharTok{+}
  \FunctionTok{geom\_smooth}\NormalTok{(}
    \AttributeTok{formula =}\NormalTok{ y }\SpecialCharTok{\textasciitilde{}}\NormalTok{ x,}
    \AttributeTok{method =} \StringTok{"glm"}\NormalTok{,}
    \AttributeTok{se =} \ConstantTok{FALSE}\NormalTok{,}
    \AttributeTok{method.args =} \FunctionTok{list}\NormalTok{(}\AttributeTok{family =} \FunctionTok{binomial}\NormalTok{(}\AttributeTok{link =} \StringTok{"logit"}\NormalTok{))}
\NormalTok{  ) }\SpecialCharTok{+}
  \FunctionTok{labs}\NormalTok{(}\AttributeTok{y =} \StringTok{"Poverty (1 = poor)"}\NormalTok{, }\AttributeTok{x =} \StringTok{"Expenditure"}\NormalTok{)}
\end{Highlighting}
\end{Shaded}

\begin{figure}

{\centering \includegraphics[width=0.7\linewidth]{docs/figures/fig-logit-pobreza-gasto-1} 

}

\caption{Relación entre gasto y condición de pobreza con curva logística ajustada}\label{fig:fig-logit-pobreza-gasto}
\end{figure}

La curva ajustada resume la asociación marginal entre el gasto y la condición de pobreza, sin incorporar todavía la estructura jerárquica de los estratos. En términos generales, la forma descendente de la curva indica que, a medida que aumenta el gasto del hogar, disminuye la probabilidad estimada de encontrarse en condición de pobreza. Sin embargo, esta relación promedio puede ocultar diferencias relevantes entre estratos, especialmente cuando los hogares pertenecen a contextos socioeconómicos distintos.

\subsection{Modelo nulo logístico}\label{modelo-nulo-logistico}

Como en el caso de variables continuas, el modelo nulo logístico constituye el punto de partida para estudiar la estructura jerárquica de la variable respuesta. Este modelo no incorpora covariables y permite evaluar si la probabilidad promedio del evento varía entre estratos. En el ejemplo considerado, el evento de interés corresponde a la condición de pobreza, de modo que el modelo permite separar la heterogeneidad atribuible a diferencias entre estratos de la variación individual propia de una respuesta binaria.

El modelo nulo, la variable de interés sigue la siguiente distribución Bernoulli:

\[
y_{kj} \mid \pi_{kj} \sim \text{Bernoulli}(\pi_{kj})
\]

En donde la probabilidad de éxito para la unidad \(k\) perteneciente al grupo \(j\) se modela mediante

\[
\text{logit}(\pi_{kj}) = \beta_{0j},
\]

donde \(\beta_{0j}\) representa el intercepto específico del estrato \(j\) en escala logit. A su vez, dicho intercepto se descompone como

\[
\beta_{0j} = \gamma_{00} + \tau_{0j}
\]

siendo \(\gamma_{00}\) el intercepto promedio global de la población y \(\tau_{0j}\) el efecto aleatorio asociado al estrato \(j\), el cual captura las desviaciones de cada grupo respecto al promedio general. Al igual que en el modelo nulo para variables continuas, este efecto aleatorio permite modelar explícitamente la heterogeneidad entre grupos:

\[
\tau_{0j} \sim N(0,\sigma_{\tau}^{2})
\]

A diferencia del modelo lineal multinivel, aquí no se incorpora un término residual aditivo \(\epsilon_{kj}\) en la ecuación del nivel individual. La variabilidad interna en el grupo está determinada por la distribución Bernoulli de la respuesta. Para cuantificar la dependencia entre unidades pertenecientes al mismo estrato se suele utilizar el enfoque de variable latente, bajo el cual esta varianza residual se aproxima por la varianza de la distribución logística estándar, es decir, \(\pi^2/3\) \citep{snijders2011}. Así, la correlación intraclase del modelo logístico se expresa como

\[
\rho =
\frac{\sigma_\tau^2}
{\sigma_\tau^2 + \frac{\pi^2}{3}}
\]

Un valor elevado de \(\rho\) indica que la probabilidad del evento presenta una estructura de agrupamiento importante, es decir, que dos unidades pertenecientes al mismo estrato tienden a parecerse más entre sí que dos unidades tomadas de estratos distintos. En contraste, un valor bajo sugiere que la mayor parte de la variación se concentra a nivel individual. El ajuste en \texttt{R} se realiza de la siguiente manera:

\begin{Shaded}
\begin{Highlighting}[]
\NormalTok{logistic\_null\_model }\OtherTok{\textless{}{-}} \FunctionTok{glmer}\NormalTok{(}
\NormalTok{  poor }\SpecialCharTok{\textasciitilde{}}\NormalTok{ (}\DecValTok{1} \SpecialCharTok{|}\NormalTok{ Stratum),}
  \AttributeTok{data =}\NormalTok{ survey\_data,}
  \AttributeTok{weights =}\NormalTok{ qk,}
  \AttributeTok{family =} \FunctionTok{binomial}\NormalTok{(}\AttributeTok{link =} \StringTok{"logit"}\NormalTok{)}
\NormalTok{)}
\end{Highlighting}
\end{Shaded}

Los coeficientes del modelo nulo logístico por estrato se presentan en la Tabla \ref{tab:tab-coef-logit-nulo}. Los interceptos estimados en el modelo nulo muestran una heterogeneidad considerable entre estratos. En los primeros estratos reportados, algunos interceptos son fuertemente negativos, como \texttt{idStrt004} e \texttt{idStrt003}, lo que corresponde a probabilidades base muy bajas de pobreza. En contraste, estratos como \texttt{idStrt009} e \texttt{idStrt006} presentan interceptos positivos y elevados, asociados con probabilidades base mucho mayores.

\begin{Shaded}
\begin{Highlighting}[]
\FunctionTok{coef}\NormalTok{(logistic\_null\_model)}\SpecialCharTok{$}\NormalTok{Stratum }\SpecialCharTok{\%\textgreater{}\%}
  \FunctionTok{slice}\NormalTok{(}\DecValTok{1}\SpecialCharTok{:}\DecValTok{12}\DataTypeTok{L}\NormalTok{)}
\end{Highlighting}
\end{Shaded}

\begin{longtable}[]{@{}ll@{}}
\caption{\label{tab:tab-coef-logit-nulo}Interceptos estimados por estrato en el modelo logístico nulo (primeros 12 estratos)}\tabularnewline
\toprule\noalign{}
& (Intercept) \\
\midrule\noalign{}
\endfirsthead
\toprule\noalign{}
& (Intercept) \\
\midrule\noalign{}
\endhead
\bottomrule\noalign{}
\endlastfoot
idStrt001 & -0.852 \\
idStrt002 & -0.038 \\
idStrt003 & -2.378 \\
idStrt004 & -2.618 \\
idStrt005 & -1.037 \\
idStrt006 & 0.902 \\
idStrt007 & -1.018 \\
idStrt008 & 0.156 \\
idStrt009 & 1.913 \\
idStrt010 & -0.609 \\
idStrt011 & -1.275 \\
idStrt012 & 0.203 \\
\end{longtable}

La correlación intraclase asciende a 0.315. Dado que el modelo no incluye covariables, estas diferencias reflejan exclusivamente la variación entre estratos en la prevalencia del evento.

\begin{Shaded}
\begin{Highlighting}[]
\NormalTok{performance}\SpecialCharTok{::}\FunctionTok{icc}\NormalTok{(logistic\_null\_model)}
\end{Highlighting}
\end{Shaded}

\begin{verbatim}
## # Intraclass Correlation Coefficient
## 
##     Adjusted ICC: 0.315
##   Unadjusted ICC: 0.315
\end{verbatim}

\subsection{Modelo logístico con intercepto aleatorio}\label{modelo-logistico-con-intercepto-aleatorio}

El modelo logístico con intercepto aleatorio incorpora covariables individuales o del hogar, permitiendo que el nivel base del evento varíe entre estratos. Su lógica es paralela a la del modelo con intercepto aleatorio para variables continuas: los efectos de las covariables se consideran comunes a todos los grupos, mientras que cada estrato puede tener su propio intercepto. En este caso, las diferencias entre interceptos se interpretan como diferencias en los log-odds base del evento.

El modelo se expresa como \(y_{kj} \mid \pi_{kj} \sim \text{Bernoulli}(\pi_{kj})\). En donde \(\text{logit}(\pi_{kj}) = \beta_{0j} + \mathbf{x}_{kj}\boldsymbol{\beta}\) y \(\beta_{0j} = \gamma_{00} + \tau_{0j}\). En este caso, \(\mathbf{x}_{kj}\) representa el vector de covariables de la unidad \(k\) en el estrato \(j\), \(\boldsymbol{\beta}\) es el vector de efectos fijos comunes a todos los estratos y \(\tau_{0j}\) captura la desviación específica del estrato respecto al intercepto promedio global. Como antes, se asume que \(\tau_{0j} \sim N(0,\sigma_{\tau}^{2})\).

Bajo esta especificación, dos estratos con los mismos valores de las covariables pueden presentar probabilidades base distintas del evento debido a sus interceptos aleatorios. Sin embargo, el efecto de cada covariable sobre la escala logit se mantiene constante entre estratos. En el ejemplo, esto equivale a permitir que la probabilidad base de pobreza cambie entre estratos, mientras que la asociación entre gasto y pobreza se resume mediante una pendiente común.

\begin{Shaded}
\begin{Highlighting}[]
\NormalTok{random\_intercept\_logit\_model }\OtherTok{\textless{}{-}} \FunctionTok{glmer}\NormalTok{(}
\NormalTok{  poor }\SpecialCharTok{\textasciitilde{}}\NormalTok{ Expenditure }\SpecialCharTok{+}\NormalTok{ (}\DecValTok{1} \SpecialCharTok{|}\NormalTok{ Stratum),}
  \AttributeTok{data =}\NormalTok{ survey\_data,}
  \AttributeTok{family =} \FunctionTok{binomial}\NormalTok{(}\AttributeTok{link =} \StringTok{"logit"}\NormalTok{),}
  \AttributeTok{weights =}\NormalTok{ qk}
\NormalTok{)}
\NormalTok{performance}\SpecialCharTok{::}\FunctionTok{icc}\NormalTok{(random\_intercept\_logit\_model)}
\end{Highlighting}
\end{Shaded}

\begin{verbatim}
## # Intraclass Correlation Coefficient
## 
##     Adjusted ICC: 0.298
##   Unadjusted ICC: 0.171
\end{verbatim}

Los coeficientes estimados por estrato se recogen en la Tabla \ref{tab:tab-coef-logit-int-aleatorio}. Al incorporar el gasto del hogar como covariable, el coeficiente fijo asociado a \texttt{Expenditure} es negativo, lo que indica que mayores niveles de gasto se asocian con menores log-odds de pobreza. Esto implica una reducción progresiva de la probabilidad estimada de pobreza a medida que aumenta el gasto, manteniendo constante el efecto del estrato. En este modelo, el efecto del gasto es común, pero cada estrato tiene un nivel base propio de pobreza. Así, los estratos con interceptos positivos elevados presentan una mayor probabilidad base de pobreza para un mismo nivel de gasto, mientras que los estratos con interceptos negativos muestran una probabilidad base menor.

Aun después de controlar por el gasto, la correlación intraclase ajustada permanece alta, alrededor de 0.298, lo que muestra que las diferencias entre estratos continúan siendo relevantes.

\begin{Shaded}
\begin{Highlighting}[]
\FunctionTok{coef}\NormalTok{(random\_intercept\_logit\_model)}\SpecialCharTok{$}\NormalTok{Stratum }\SpecialCharTok{\%\textgreater{}\%}
  \FunctionTok{slice}\NormalTok{(}\DecValTok{1}\SpecialCharTok{:}\DecValTok{10}\DataTypeTok{L}\NormalTok{)}
\end{Highlighting}
\end{Shaded}

\begin{longtable}[]{@{}llc@{}}
\caption{\label{tab:tab-coef-logit-int-aleatorio}Coeficientes del modelo logístico con intercepto aleatorio por estrato (primeros 10 estratos)}\tabularnewline
\toprule\noalign{}
& (Intercept) & Expenditure \\
\midrule\noalign{}
\endfirsthead
\toprule\noalign{}
& (Intercept) & Expenditure \\
\midrule\noalign{}
\endhead
\bottomrule\noalign{}
\endlastfoot
idStrt001 & 1.044 & -0.007 \\
idStrt002 & 1.918 & -0.007 \\
idStrt003 & -0.454 & -0.007 \\
idStrt004 & 0.062 & -0.007 \\
idStrt005 & 1.760 & -0.007 \\
idStrt006 & 3.194 & -0.007 \\
idStrt007 & 0.658 & -0.007 \\
idStrt008 & 1.711 & -0.007 \\
idStrt009 & 3.721 & -0.007 \\
idStrt010 & 1.175 & -0.007 \\
\end{longtable}

La Figura \ref{fig:fig-logit-int-aleatorio} presenta las curvas de probabilidad predichas mediante un modelo logístico con intercepto aleatorio por estrato. Se observa una relación inversa entre el gasto y la probabilidad de pobreza, de modo que hogares con mayores niveles de gasto presentan una menor probabilidad de ser clasificados como pobres. Asimismo, las diferencias entre las curvas reflejan la heterogeneidad existente entre estratos, capturada mediante los interceptos aleatorios del modelo.

\begin{Shaded}
\begin{Highlighting}[]
\NormalTok{prediction\_data }\OtherTok{\textless{}{-}}\NormalTok{ plot\_data }\SpecialCharTok{\%\textgreater{}\%}
  \FunctionTok{group\_by}\NormalTok{(Stratum) }\SpecialCharTok{\%\textgreater{}\%}
  \FunctionTok{summarise}\NormalTok{(}
    \AttributeTok{Expenditure =} \FunctionTok{list}\NormalTok{(}\FunctionTok{seq}\NormalTok{(}\FunctionTok{min}\NormalTok{(Expenditure), }\FunctionTok{max}\NormalTok{(Expenditure), }\AttributeTok{len =} \DecValTok{100}\NormalTok{))}
\NormalTok{  ) }\SpecialCharTok{\%\textgreater{}\%}
\NormalTok{  tidyr}\SpecialCharTok{::}\FunctionTok{unnest\_legacy}\NormalTok{()}

\NormalTok{prediction\_data }\OtherTok{\textless{}{-}}\NormalTok{ prediction\_data }\SpecialCharTok{\%\textgreater{}\%}
  \FunctionTok{mutate}\NormalTok{(}
    \AttributeTok{probability =} \FunctionTok{predict}\NormalTok{(}
\NormalTok{      random\_intercept\_logit\_model,}
      \AttributeTok{newdata =}\NormalTok{ prediction\_data,}
      \AttributeTok{type =} \StringTok{"response"}
\NormalTok{    )}
\NormalTok{  )}
\end{Highlighting}
\end{Shaded}

\begin{Shaded}
\begin{Highlighting}[]
\FunctionTok{ggplot}\NormalTok{(}\AttributeTok{data =}\NormalTok{ prediction\_data,}
       \FunctionTok{aes}\NormalTok{(}\AttributeTok{y =}\NormalTok{ probability, }\AttributeTok{x =}\NormalTok{ Expenditure, }\AttributeTok{colour =}\NormalTok{ Stratum)) }\SpecialCharTok{+}
  \FunctionTok{geom\_line}\NormalTok{() }\SpecialCharTok{+}
  \FunctionTok{geom\_point}\NormalTok{(}\AttributeTok{data =}\NormalTok{ plot\_data,}
             \FunctionTok{aes}\NormalTok{(}\AttributeTok{y =}\NormalTok{ poor, }\AttributeTok{x =}\NormalTok{ Expenditure)) }\SpecialCharTok{+}
  \FunctionTok{labs}\NormalTok{(}\AttributeTok{y =} \StringTok{"Poverty probability"}\NormalTok{, }\AttributeTok{x =} \StringTok{"Expenditure"}\NormalTok{) }\SpecialCharTok{+}
  \FunctionTok{theme}\NormalTok{(}\AttributeTok{legend.position =} \StringTok{"none"}\NormalTok{)}
\end{Highlighting}
\end{Shaded}

\begin{figure}

{\centering \includegraphics[width=0.7\linewidth]{docs/figures/fig-logit-int-aleatorio-1} 

}

\caption{Curvas de probabilidad predichas por estrato: modelo logístico con intercepto aleatorio}\label{fig:fig-logit-int-aleatorio}
\end{figure}

\subsection{Modelo logístico con intercepto y pendiente aleatoria}\label{modelo-logistico-con-intercepto-y-pendiente-aleatoria}

El modelo logístico con intercepto y pendiente aleatoria extiende la especificación anterior permitiendo que no sólo varíe el nivel base del evento entre estratos, sino también el efecto de una covariable específica. Esta formulación es análoga al modelo de intercepto y pendiente aleatorios para variables continuas, con la salvedad de que las diferencias entre estratos se expresan en la escala logit y se traducen en curvas de probabilidad no lineales.

Sea \(x_{kj}\) la covariable cuyo efecto se permite variar entre estratos, y sea \(\mathbf{z}_{kj}\) el conjunto de covariables con efectos fijos comunes. En este modelo, la probabilidad de éxito se asume de tal forma que \(\text{logit}(\pi_{kj}) = \beta_{0j} + x_{kj}\beta_{1j} + \mathbf{z}_{kj}\boldsymbol{\beta}\). En donde, \(\beta_{0j}\) es el intercepto específico del estrato \(j\) y \(\beta_{1j}\) es la pendiente específica asociada a la covariable \(x_{kj}\).

Ambos parámetros se descomponen en una parte promedio poblacional y una desviación específica del estrato, de tal manera que \(\beta_{0j} = \gamma_{00} + \tau_{0j}\) y \(\beta_{1j} = \gamma_{10} + \tau_{1j}\). En donde los términos \(\gamma_{00}\) y \(\gamma_{10}\) representan, respectivamente, el intercepto promedio global y la pendiente promedio global en escala logit. Por su parte, \(\tau_{0j}\) y \(\tau_{1j}\) capturan cuánto se aparta el estrato \(j\) de esos promedios. La implementación en R es la siguiente:

\begin{Shaded}
\begin{Highlighting}[]
\NormalTok{random\_slope\_logit\_model }\OtherTok{\textless{}{-}} \FunctionTok{glmer}\NormalTok{(}
\NormalTok{  poor }\SpecialCharTok{\textasciitilde{}} \DecValTok{1} \SpecialCharTok{+}\NormalTok{ Expenditure }\SpecialCharTok{+}\NormalTok{ (}\DecValTok{1} \SpecialCharTok{+}\NormalTok{ Expenditure }\SpecialCharTok{|}\NormalTok{ Stratum),}
  \AttributeTok{data =}\NormalTok{ survey\_data,}
  \AttributeTok{weights =}\NormalTok{ qk,}
  \AttributeTok{family =} \FunctionTok{binomial}\NormalTok{(}\AttributeTok{link =} \StringTok{"logit"}\NormalTok{)}
\NormalTok{)}
\NormalTok{performance}\SpecialCharTok{::}\FunctionTok{icc}\NormalTok{(random\_slope\_logit\_model)}
\end{Highlighting}
\end{Shaded}

\begin{verbatim}
## # Intraclass Correlation Coefficient
## 
##     Adjusted ICC: 0.875
##   Unadjusted ICC: 0.631
\end{verbatim}

La estimación del modelo con intercepto y pendiente aleatoria indica una heterogeneidad mucho más marcada entre estratos. En este caso, tanto los interceptos como las pendientes asociadas al gasto pueden cambiar entre grupos. La correlación intraclase ajustada es muy alta, lo que sugiere que el componente entre estratos domina la variación latente del modelo. Los coeficientes del modelo por estrato se muestran en la Tabla \ref{tab:tab-coef-logit-pendiente}.

\begin{Shaded}
\begin{Highlighting}[]
\FunctionTok{coef}\NormalTok{(random\_slope\_logit\_model)}\SpecialCharTok{$}\NormalTok{Stratum }\SpecialCharTok{\%\textgreater{}\%}
  \FunctionTok{slice}\NormalTok{(}\DecValTok{1}\SpecialCharTok{:}\DecValTok{10}\DataTypeTok{L}\NormalTok{)}
\end{Highlighting}
\end{Shaded}

\begin{longtable}[]{@{}llc@{}}
\caption{\label{tab:tab-coef-logit-pendiente}Coeficientes del modelo logístico con intercepto y pendiente aleatoria por estrato (primeros 10 estratos)}\tabularnewline
\toprule\noalign{}
& (Intercept) & Expenditure \\
\midrule\noalign{}
\endfirsthead
\toprule\noalign{}
& (Intercept) & Expenditure \\
\midrule\noalign{}
\endhead
\bottomrule\noalign{}
\endlastfoot
idStrt001 & 4.864 & -0.025 \\
idStrt002 & 9.860 & -0.035 \\
idStrt003 & -1.133 & -0.007 \\
idStrt004 & 1.897 & -0.015 \\
idStrt005 & 8.029 & -0.026 \\
idStrt006 & -1.153 & 0.009 \\
idStrt007 & 0.973 & -0.012 \\
idStrt008 & 1.488 & -0.006 \\
idStrt009 & 3.660 & -0.005 \\
idStrt010 & 4.132 & -0.020 \\
\end{longtable}

Los coeficientes por estrato evidencian que la relación entre gasto y pobreza no sólo cambia en el nivel base, sino también en la intensidad y dirección de la pendiente. En varios estratos la pendiente del gasto es negativa, lo que mantiene la interpretación esperada de que mayores niveles de gasto reducen la probabilidad de pobreza. Sin embargo, también aparecen pendientes cercanas a cero e incluso positivas en algunos estratos, lo que sugiere patrones locales distintos. Esta variabilidad es precisamente la que el modelo busca capturar mediante el efecto aleatorio de la pendiente. La Figura \ref{fig:fig-logit-pendiente-aleatoria} muestra que las curvas predichas son más heterogéneas entre estratos que en el modelo anterior:

\begin{Shaded}
\begin{Highlighting}[]
\NormalTok{prediction\_data }\OtherTok{\textless{}{-}}\NormalTok{ plot\_data }\SpecialCharTok{\%\textgreater{}\%}
  \FunctionTok{group\_by}\NormalTok{(Stratum) }\SpecialCharTok{\%\textgreater{}\%}
  \FunctionTok{summarise}\NormalTok{(}
    \AttributeTok{Expenditure =} \FunctionTok{list}\NormalTok{(}\FunctionTok{seq}\NormalTok{(}\FunctionTok{min}\NormalTok{(Expenditure), }\FunctionTok{max}\NormalTok{(Expenditure), }\AttributeTok{len =} \DecValTok{100}\NormalTok{))}
\NormalTok{  ) }\SpecialCharTok{\%\textgreater{}\%}
\NormalTok{  tidyr}\SpecialCharTok{::}\FunctionTok{unnest\_legacy}\NormalTok{()}

\NormalTok{prediction\_data }\OtherTok{\textless{}{-}}\NormalTok{ prediction\_data }\SpecialCharTok{\%\textgreater{}\%}
  \FunctionTok{mutate}\NormalTok{(}
    \AttributeTok{probability =} \FunctionTok{predict}\NormalTok{(}
\NormalTok{      random\_slope\_logit\_model,}
      \AttributeTok{newdata =}\NormalTok{ prediction\_data,}
      \AttributeTok{type =} \StringTok{"response"}
\NormalTok{    )}
\NormalTok{  )}
\end{Highlighting}
\end{Shaded}

\begin{Shaded}
\begin{Highlighting}[]
\FunctionTok{ggplot}\NormalTok{(}\AttributeTok{data =}\NormalTok{ prediction\_data,}
       \FunctionTok{aes}\NormalTok{(}\AttributeTok{y =}\NormalTok{ probability, }\AttributeTok{x =}\NormalTok{ Expenditure, }\AttributeTok{colour =}\NormalTok{ Stratum)) }\SpecialCharTok{+}
  \FunctionTok{geom\_line}\NormalTok{() }\SpecialCharTok{+}
  \FunctionTok{geom\_point}\NormalTok{(}\AttributeTok{data =}\NormalTok{ plot\_data,}
             \FunctionTok{aes}\NormalTok{(}\AttributeTok{y =}\NormalTok{ poor, }\AttributeTok{x =}\NormalTok{ Expenditure)) }\SpecialCharTok{+}
  \FunctionTok{labs}\NormalTok{(}\AttributeTok{y =} \StringTok{"Poverty probability"}\NormalTok{, }\AttributeTok{x =} \StringTok{"Expenditure"}\NormalTok{) }\SpecialCharTok{+}
  \FunctionTok{theme}\NormalTok{(}\AttributeTok{legend.position =} \StringTok{"none"}\NormalTok{)}
\end{Highlighting}
\end{Shaded}

\begin{figure}

{\centering \includegraphics[width=0.7\linewidth]{docs/figures/fig-logit-pendiente-aleatoria-1} 

}

\caption{Curvas de probabilidad predichas por estrato: modelo logístico con intercepto y pendiente aleatoria}\label{fig:fig-logit-pendiente-aleatoria}
\end{figure}

\appendix


\chapter{\texorpdfstring{Manejo de bases de datos con \texttt{R}}{Manejo de bases de datos con R}}\label{ap1}

\texttt{R} es un software estadístico ampliamente utilizado para el análisis de datos y la generación de gráficos y permite un ambiente integrado para el análisis estadístico, en particular de bases de datos. Esto significa que permite realizar procesos de análisis, programación, visualización y manejo de datos dentro de un mismo entorno. En este anexo se utiliza el enfoque del \texttt{tidyverse}, y en particular el paquete \texttt{dplyr}, para organizar una secuencia lógica de trabajo con bases de datos: cargar la información, inspeccionarla, seleccionar variables, filtrar registros, ordenar filas, crear variables y producir resúmenes.

\section{Preparación del entorno de trabajo}\label{preparacion-del-entorno-de-trabajo}

\texttt{R} es un lenguaje de programación de carácter colaborativo, lo que permite que su comunidad de usuarios y desarrolladores contribuya continuamente mediante la creación de funciones, paquetes y librerías especializadas. Gracias a este ecosistema abierto, es posible ampliar sus capacidades para distintos tipos de análisis estadístico y procesamiento de datos. Entre las librerías más utilizadas para el análisis de bases de datos se encuentran las siguientes:

\begin{itemize}
\item
  La librería \texttt{tidyverse} \citep{tidyverse2026} corresponde a un conjunto de paquetes de \texttt{R} diseñados para apoyar distintas etapas del trabajo con datos, desde su importación y transformación hasta su exploración, análisis y visualización. Su uso se ha extendido ampliamente en ciencia de datos porque ofrece una forma de trabajo consistente y ordenada.
\item
  El paquete \texttt{dplyr} \citep{dplyr2026} ofrece herramientas para organizar y transformar bases de datos rectangulares. Su utilidad radica en que reúne, en una sintaxis simple y consistente, las operaciones más frecuentes del trabajo analítico con datos: seleccionar variables, filtrar registros, crear nuevas columnas, ordenar observaciones y resumir información.
\item
  El paquete \texttt{TeachingSampling} \citep{GutierrezT2020} ofrece funciones para seleccionar muestras probabilísticas y realizar inferencia sobre poblaciones finitas bajo diversos diseños de muestreo.
\end{itemize}

Antes de poder utilizar las diferentes funciones que cada librería tiene, es necesario descargarlas de antemano de la web. El comando \texttt{install.packages} permite realizar esta tarea. Note que algunas librerías pueden depender de otras, así que para poder utilizarlas es necesario instalar también las dependencias adicionales.

\begin{Shaded}
\begin{Highlighting}[]
\FunctionTok{install.packages}\NormalTok{(}\StringTok{"dplyr"}\NormalTok{)}
\FunctionTok{install.packages}\NormalTok{(}\StringTok{"tidyverse"}\NormalTok{)}
\FunctionTok{install.packages}\NormalTok{(}\StringTok{"TeachingSampling"}\NormalTok{)}
\end{Highlighting}
\end{Shaded}

Una vez instalados los paquetes, se cargan con \texttt{library()}. El primer comando limpia los objetos existentes en la sesión de trabajo, lo que ayuda a que los resultados del capítulo dependan solo del código que se ejecuta aquí. Debe tenerse en cuenta que un paquete solo puede cargarse si previamente ha sido instalado en el sistema.

\begin{Shaded}
\begin{Highlighting}[]
\FunctionTok{rm}\NormalTok{(}\AttributeTok{list =} \FunctionTok{ls}\NormalTok{())}

\FunctionTok{library}\NormalTok{(tidyverse)}
\FunctionTok{library}\NormalTok{(dplyr)}
\FunctionTok{library}\NormalTok{(TeachingSampling)}
\end{Highlighting}
\end{Shaded}

Una vez se descargan e instalan las librerías o paquetes en \texttt{R}, el paso recomendado es que todos los procesamientos se realicen mediante la creación de proyectos. Un proyecto de \texttt{R} se define como un archivo que contiene los archivos de origen y contenido asociados con el trabajo que se está realizando. Adicionalmente, contiene información que permite la compilación de cada archivo de \texttt{R} que se va a utilizar, mantiene la información para integrarse con sistemas de control de código fuente y ayuda a organizar la aplicación de los procesamientos en componentes lógicos.

\section{Carga e inspección inicial de la base}\label{carga-e-inspeccion-inicial-de-la-base}

Es muy usual que al trabajar proyectos en \texttt{R} sea necesario importar bases de datos con información relevante para un estudio en particular. Los formatos de bases de datos que \texttt{R} permite importar son diversos, entre ellos se tienen \texttt{xlsx}, \texttt{csv}, \texttt{txt}, \texttt{STATA}, etc. Una vez leída la base de datos en el formato pertinente es recomendable transformarla al formato nativo de \texttt{R}, es decir \texttt{.RDS}. Este es un formato más eficiente y propio de \texttt{R}.

Una vez se carga la base de datos se procede a utilizar las funciones en \texttt{R} para poder obtener resultados de los procesamientos agregados y gráficos de interés. Para ejemplificar el uso de funciones que permitan obtener resultados agregados, utilizaremos la base de datos \texttt{BigCity} del paquete \texttt{TeachingSampling}. Esta base corresponde a un conjunto de variables socioeconómicas de 150266 personas en un año en particular.

\begin{Shaded}
\begin{Highlighting}[]
\FunctionTok{data}\NormalTok{(BigCity)}
\NormalTok{bigcity\_data }\OtherTok{\textless{}{-}}\NormalTok{ BigCity}
\end{Highlighting}
\end{Shaded}

El objeto \texttt{bigcity\_data} contiene una copia de \texttt{BigCity}. La función \texttt{head()} permite observar las primeras filas y obtener una impresión rápida de la estructura de la base.

\begin{Shaded}
\begin{Highlighting}[]
\FunctionTok{head}\NormalTok{(bigcity\_data)}
\end{Highlighting}
\end{Shaded}

\begin{verbatim}
##        HHID PersonID   Stratum     PSU  Zone    Sex Age MaritalST Income
## 1 idHH00001  idPer01 idStrt001 PSU0001 Rural   Male  38   Married    555
## 2 idHH00001  idPer02 idStrt001 PSU0001 Rural Female  40   Married    555
## 3 idHH00001  idPer03 idStrt001 PSU0001 Rural Female  20    Single    555
## 4 idHH00001  idPer04 idStrt001 PSU0001 Rural   Male  19    Single    555
## 5 idHH00001  idPer05 idStrt001 PSU0001 Rural   Male  18    Single    555
## 6 idHH00002  idPer01 idStrt001 PSU0001 Rural   Male  35   Married    298
##   Expenditure Employment  Poverty
## 1         488   Employed  NotPoor
## 2         488   Employed  NotPoor
## 3         488   Inactive  NotPoor
## 4         488   Employed  NotPoor
## 5         488   Inactive  NotPoor
## 6         217   Employed Relative
\end{verbatim}

Una vez cargada la base de datos en \texttt{R}, se puede empezar a realizar los procesamientos según las necesidades de cada investigador. En este sentido, una de las primeras revisiones que se realizan al cargar una base de datos es verificar su dimensión; es decir, identificar la cantidad de filas y columnas que contiene. Lo anterior se puede hacer con la función \texttt{nrow}, que identifica el número de registros en la base de datos, y con la función \texttt{ncol}, que muestra el número de variables. Los códigos computacionales son los siguientes:

\begin{Shaded}
\begin{Highlighting}[]
\FunctionTok{nrow}\NormalTok{(bigcity\_data)}
\end{Highlighting}
\end{Shaded}

\begin{verbatim}
## [1] 150266
\end{verbatim}

\begin{Shaded}
\begin{Highlighting}[]
\FunctionTok{ncol}\NormalTok{(bigcity\_data)}
\end{Highlighting}
\end{Shaded}

\begin{verbatim}
## [1] 12
\end{verbatim}

Una forma resumida de revisar la cantidad de filas y columnas de una base de datos es mediante la función \texttt{dim}, la cual devuelve un vector cuya primera componente corresponde al número de filas y cuya segunda componente indica el número de columnas.

\begin{Shaded}
\begin{Highlighting}[]
\FunctionTok{dim}\NormalTok{(bigcity\_data)}
\end{Highlighting}
\end{Shaded}

\begin{verbatim}
## [1] 150266     12
\end{verbatim}

En algunas ocasiones, las bases de datos pueden ser extensas, ya sea porque contienen un número considerable de variables observadas o porque el número de registros es grande. Por esta razón, para visualizarlas adecuadamente, puede ser conveniente abrirlas en una ventana independiente y explorar su contenido de forma tabular. Esto se realiza mediante la función \texttt{View}, como se muestra a continuación.

\begin{Shaded}
\begin{Highlighting}[]
\FunctionTok{View}\NormalTok{(bigcity\_data)}
\end{Highlighting}
\end{Shaded}

Otra verificación importante al cargar una base de datos en \texttt{R} consiste en identificar las variables que contiene. Para ello, se puede utilizar la función \texttt{names}, la cual muestra los nombres de las variables incluidas en la base de datos.

\begin{Shaded}
\begin{Highlighting}[]
\FunctionTok{names}\NormalTok{(bigcity\_data)}
\end{Highlighting}
\end{Shaded}

\begin{verbatim}
##  [1] "HHID"        "PersonID"    "Stratum"     "PSU"         "Zone"       
##  [6] "Sex"         "Age"         "MaritalST"   "Income"      "Expenditure"
## [11] "Employment"  "Poverty"
\end{verbatim}

La salida anterior muestra los nombres de las variables contenidas en la base de datos. En este caso, la base cuenta con 12 variables. Las primeras variables corresponden a identificadores y elementos del diseño muestral, como \texttt{HHID}, \texttt{PersonID}, \texttt{Stratum}, \texttt{PSU} y \texttt{Zone}. Luego aparecen variables sociodemográficas, como \texttt{Sex}, \texttt{Age} y \texttt{MaritalST}. Finalmente, se incluyen variables económicas y de condición laboral o social, como \texttt{Income}, \texttt{Expenditure}, \texttt{Employment} y \texttt{Poverty}.

Si se desea profundizar en el tipo de información que almacena cada variable, puede utilizarse la función \texttt{str}, la cual muestra de manera compacta la estructura de un objeto y sus componentes. Para nuestra base de datos, esta función se utilizaría de la siguiente manera:

\begin{Shaded}
\begin{Highlighting}[]
\FunctionTok{str}\NormalTok{(bigcity\_data)}
\end{Highlighting}
\end{Shaded}

\begin{verbatim}
## 'data.frame':    150266 obs. of  12 variables:
##  $ HHID       : chr  "idHH00001" "idHH00001" "idHH00001" "idHH00001" ...
##  $ PersonID   : chr  "idPer01" "idPer02" "idPer03" "idPer04" ...
##  $ Stratum    : chr  "idStrt001" "idStrt001" "idStrt001" "idStrt001" ...
##  $ PSU        : chr  "PSU0001" "PSU0001" "PSU0001" "PSU0001" ...
##  $ Zone       : chr  "Rural" "Rural" "Rural" "Rural" ...
##  $ Sex        : chr  "Male" "Female" "Female" "Male" ...
##  $ Age        : int  38 40 20 19 18 35 29 14 13 6 ...
##  $ MaritalST  : Factor w/ 6 levels "Partner","Married",..: 2 2 5 5 5 2 2 5 5 NA ...
##  $ Income     : num  555 555 555 555 555 ...
##  $ Expenditure: num  488 488 488 488 488 ...
##  $ Employment : Factor w/ 3 levels "Unemployed","Inactive",..: 3 3 2 3 2 3 3 NA NA NA ...
##  $ Poverty    : Factor w/ 3 levels "NotPoor","Extreme",..: 1 1 1 1 1 3 3 3 3 3 ...
\end{verbatim}

Como se observa en la salida anterior, la función \texttt{str} permite identificar la clase de cada variable incluida en la base de datos. Por ejemplo, la variable \texttt{HHID} es de tipo carácter, al igual que la variable \texttt{Sex}, mientras que las variables \texttt{Income} y \texttt{Expenditure} son de tipo numérico. Los demás atributos asociados con las variables también pueden revisarse directamente en la salida del código. Esta función resulta especialmente útil para obtener una visión general del contenido y la estructura de una base de datos.

\section{Encadenamiento de operaciones}\label{encadenamiento-de-operaciones}

Una de las contribuciones más útiles de \texttt{R} es la incorporación de las tuberías (\emph{pipelines}), las cuales permiten construir consultas, transformaciones y nuevos objetos a partir de una base de datos de manera más intuitiva, ordenada y fácil de interpretar. El operador de encadenamiento \texttt{\%\textgreater{}\%}, proveniente del paquete \texttt{magrittr} \citep{pipe} y cargado automáticamente al utilizar la librería \texttt{tidyverse}, permite enlazar varias operaciones sucesivas, de modo que el resultado obtenido en un paso se convierte en el insumo del paso siguiente. A continuación, se presenta un ejemplo sencillo de su uso con la base de datos \texttt{BigCity}, en el cual se calcula el número total de elementos que contiene la base mediante la función \texttt{count}.

\begin{Shaded}
\begin{Highlighting}[]
\NormalTok{bigcity\_data }\SpecialCharTok{\%\textgreater{}\%}
  \FunctionTok{count}\NormalTok{()}
\end{Highlighting}
\end{Shaded}

\begin{verbatim}
##        n
## 1 150266
\end{verbatim}

La línea de código anterior puede interpretarse de forma sencilla: se toma la base de datos como punto de partida y, sobre ella, se realiza un conteo de sus registros. Además de la función \texttt{count}, existe una amplia variedad de funciones que pueden combinarse con el operador \texttt{\%\textgreater{}\%} para realizar consultas, transformaciones y resúmenes sobre los datos. El paquete \texttt{dplyr} toma ventaja de este operador, permitiendo organizar las operaciones más comunes sobre una base rectangular. En las siguientes subsecciones se presentan algunas de ellas, siguiendo un orden típico de trabajo: reducir observaciones, conservar variables, ordenar la información, crear nuevas columnas y producir resultados agregados.

\subsection{\texorpdfstring{\texttt{filter}: selección de registros}{filter: selección de registros}}\label{filter-seleccion-de-registros}

\texttt{filter()} conserva las filas que cumplen una o varias condiciones. En el siguiente ejemplo se crean dos subconjuntos de \texttt{bigcity\_data}: uno con los hombres y otro con las mujeres. El resultado son dos objetos nuevos. \texttt{male\_data} contiene únicamente registros con \texttt{Sex\ ==\ "Male"}, mientras que \texttt{female\_data} contiene registros con \texttt{Sex\ ==\ "Female"}.

\begin{Shaded}
\begin{Highlighting}[]
\NormalTok{male\_data }\OtherTok{\textless{}{-}}\NormalTok{ bigcity\_data }\SpecialCharTok{\%\textgreater{}\%}
  \FunctionTok{filter}\NormalTok{(Sex }\SpecialCharTok{==} \StringTok{"Male"}\NormalTok{)}
\NormalTok{female\_data }\OtherTok{\textless{}{-}}\NormalTok{ bigcity\_data }\SpecialCharTok{\%\textgreater{}\%}
  \FunctionTok{filter}\NormalTok{(Sex }\SpecialCharTok{==} \StringTok{"Female"}\NormalTok{)}
\end{Highlighting}
\end{Shaded}

También puede filtrarse por condición de pobreza. En este caso se conservan solo las personas clasificadas como no pobres. El objeto \texttt{nonpoor\_data} contiene el subconjunto de registros cuya variable \texttt{Poverty} toma el valor \texttt{"NotPoor"}.

\begin{Shaded}
\begin{Highlighting}[]
\NormalTok{nonpoor\_data }\OtherTok{\textless{}{-}}\NormalTok{ bigcity\_data }\SpecialCharTok{\%\textgreater{}\%}
  \FunctionTok{filter}\NormalTok{(Poverty }\SpecialCharTok{==} \StringTok{"NotPoor"}\NormalTok{)}
\end{Highlighting}
\end{Shaded}

La función \texttt{filter()} también permite conservar registros cuyos valores pertenecen a un conjunto específico. En los siguientes ejemplos se filtra por valores particulares de ingreso. Note que \texttt{\%in\%} evalúa si \texttt{Income} pertenece al conjunto indicado. Así, \texttt{working\_data} conserva ingresos iguales a 265 o 600, y \texttt{income\_filter\_data} conserva ingresos iguales a 1000 o 2000.

\begin{Shaded}
\begin{Highlighting}[]
\NormalTok{working\_data }\OtherTok{\textless{}{-}}\NormalTok{ bigcity\_data }\SpecialCharTok{\%\textgreater{}\%}
  \FunctionTok{filter}\NormalTok{(Income }\SpecialCharTok{\%in\%} \FunctionTok{c}\NormalTok{(}\DecValTok{265}\NormalTok{, }\DecValTok{600}\NormalTok{))}

\NormalTok{income\_filter\_data }\OtherTok{\textless{}{-}}\NormalTok{ bigcity\_data }\SpecialCharTok{\%\textgreater{}\%}
  \FunctionTok{filter}\NormalTok{(Income }\SpecialCharTok{\%in\%} \FunctionTok{c}\NormalTok{(}\DecValTok{1000}\NormalTok{, }\DecValTok{2000}\NormalTok{))}
\end{Highlighting}
\end{Shaded}

\subsection{\texorpdfstring{\texttt{select}: seleccionar variables}{select: seleccionar variables}}\label{select-seleccionar-variables}

La función \texttt{select()} conserva columnas específicas. Para continuar con los ejemplos, se seleccionan variables de identificación de la persona y el hogar y otras como zona, sexo, edad e ingreso. El resultado es una base más reducida en columnas. Se mantiene \texttt{Age} porque será utilizada posteriormente para ordenar registros por edad.

\begin{Shaded}
\begin{Highlighting}[]
\NormalTok{working\_data }\OtherTok{\textless{}{-}}\NormalTok{ bigcity\_data }\SpecialCharTok{\%\textgreater{}\%}
  \FunctionTok{select}\NormalTok{(PersonID, HHID, PSU, Zone, Sex, Age, Income)}
\end{Highlighting}
\end{Shaded}

Adicionalmente, \texttt{select()} también permite excluir variables. Para eliminar columnas, se antepone el signo menos (\texttt{-}) al nombre de cada variable. El objeto \texttt{compact\_data} conserva las variables de \texttt{working\_data}, excepto \texttt{HHID} y \texttt{PersonID}. Esta operación es útil cuando se desea trabajar con una base más compacta.

\begin{Shaded}
\begin{Highlighting}[]
\NormalTok{compact\_data }\OtherTok{\textless{}{-}}\NormalTok{ working\_data }\SpecialCharTok{\%\textgreater{}\%}
  \FunctionTok{select}\NormalTok{(}\SpecialCharTok{{-}}\NormalTok{HHID, }\SpecialCharTok{{-}}\NormalTok{PersonID)}
\end{Highlighting}
\end{Shaded}

\subsection{\texorpdfstring{\texttt{arrange}: ordenamiento de filas}{arrange: ordenamiento de filas}}\label{arrange-ordenamiento-de-filas}

La función \texttt{arrange()} ordena las filas de una base según una o más variables. En el primer ejemplo, \texttt{working\_data} se ordena de menor a mayor ingreso, mientras que \texttt{sorted\_data} contiene las mismas columnas que \texttt{working\_data}, pero sus filas quedan ordenadas por \texttt{Income}. La función \texttt{head()} muestra los primeros registros después del ordenamiento.

\begin{Shaded}
\begin{Highlighting}[]
\NormalTok{sorted\_data }\OtherTok{\textless{}{-}}\NormalTok{ working\_data }\SpecialCharTok{\%\textgreater{}\%}
  \FunctionTok{arrange}\NormalTok{(Income)}

\NormalTok{sorted\_data }\SpecialCharTok{\%\textgreater{}\%}
  \FunctionTok{head}\NormalTok{()}
\end{Highlighting}
\end{Shaded}

\begin{verbatim}
##   PersonID      HHID     PSU  Zone    Sex Age Income
## 1  idPer01 idHH01522 PSU0112 Rural Female  82   10.0
## 2  idPer01 idHH22167 PSU0112 Rural Female  82   10.0
## 3  idPer01 idHH10745 PSU0893 Rural   Male  68   11.9
## 4  idPer01 idHH31390 PSU0893 Rural   Male  68   11.9
## 5  idPer01 idHH17632 PSU1432 Rural   Male  58   12.1
## 6  idPer02 idHH17632 PSU1432 Rural Female  50   12.1
\end{verbatim}

El ordenamiento también puede realizarse considerando más de una variable. En este ejemplo, la base se organiza primero según la variable \texttt{Sex} y, posteriormente, dentro de cada grupo, de acuerdo con la variable \texttt{Age}. La siguiente salida permite revisar los primeros registros después de aplicar este criterio de ordenamiento.

\begin{Shaded}
\begin{Highlighting}[]
\NormalTok{working\_data }\SpecialCharTok{\%\textgreater{}\%}
  \FunctionTok{arrange}\NormalTok{(Sex, Age) }\SpecialCharTok{\%\textgreater{}\%}
  \FunctionTok{head}\NormalTok{()}
\end{Highlighting}
\end{Shaded}

\begin{verbatim}
##   PersonID      HHID     PSU  Zone    Sex Age Income
## 1  idPer04 idHH00009 PSU0001 Rural Female   0    152
## 2  idPer04 idHH20654 PSU0001 Rural Female   0    152
## 3  idPer06 idHH00042 PSU0004 Rural Female   0    528
## 4  idPer06 idHH20687 PSU0004 Rural Female   0    528
## 5  idPer04 idHH00191 PSU0014 Urban Female   0    355
## 6  idPer04 idHH20836 PSU0014 Urban Female   0    355
\end{verbatim}

Para ordenar los registros de mayor a menor, se utiliza la opción \texttt{desc()} dentro de la función \texttt{arrange()}. En el siguiente ejemplo, la base se organiza de manera que las personas con mayor edad aparezcan en los primeros registros. Este tipo de ordenamiento permite identificar rápidamente los valores más altos de una variable dentro de la base reducida.

\begin{Shaded}
\begin{Highlighting}[]
\NormalTok{working\_data }\SpecialCharTok{\%\textgreater{}\%}
  \FunctionTok{arrange}\NormalTok{(}\FunctionTok{desc}\NormalTok{(Age)) }\SpecialCharTok{\%\textgreater{}\%}
  \FunctionTok{head}\NormalTok{()}
\end{Highlighting}
\end{Shaded}

\begin{verbatim}
##   PersonID      HHID     PSU  Zone    Sex Age Income
## 1  idPer02 idHH01024 PSU0076 Urban Female  98    135
## 2  idPer02 idHH21669 PSU0076 Urban Female  98    135
## 3  idPer01 idHH02916 PSU0219 Rural Female  98    137
## 4  idPer01 idHH23561 PSU0219 Rural Female  98    137
## 5  idPer01 idHH03863 PSU0290 Rural   Male  98    245
## 6  idPer01 idHH24508 PSU0290 Rural   Male  98    245
\end{verbatim}

\subsection{\texorpdfstring{\texttt{mutate}: creación o transformación de variables}{mutate: creación o transformación de variables}}\label{mutate-creacion-o-transformacion-de-variables}

La función \texttt{mutate()} crea nuevas variables o modifica variables existentes. Dado que \texttt{BigCity} no incluye una variable de región geográfica explícita, se construye \texttt{Region} a partir de \texttt{Stratum}. Primero se extrae la parte numérica del estrato con \texttt{gsub("\textbackslash{}\textbackslash{}D",\ "",\ Stratum)}, luego se convierte a número con \texttt{as.numeric()} y finalmente se divide en cinco intervalos con \texttt{cut()}. El resultado es una nueva columna \texttt{Region} dentro de \texttt{bigcity\_data}, en la que sus categorías permiten producir resúmenes regionales en los pasos siguientes.

\begin{Shaded}
\begin{Highlighting}[]
\NormalTok{bigcity\_data }\OtherTok{\textless{}{-}}\NormalTok{ bigcity\_data }\SpecialCharTok{\%\textgreater{}\%}
  \FunctionTok{mutate}\NormalTok{(}
    \AttributeTok{Region =} \FunctionTok{cut}\NormalTok{(}
      \FunctionTok{as.numeric}\NormalTok{(}\FunctionTok{gsub}\NormalTok{(}\StringTok{"}\SpecialCharTok{\textbackslash{}\textbackslash{}}\StringTok{D"}\NormalTok{, }\StringTok{""}\NormalTok{, Stratum)),}
      \AttributeTok{breaks =} \DecValTok{5}\NormalTok{,}
      \AttributeTok{labels =} \FunctionTok{c}\NormalTok{(}\StringTok{"North"}\NormalTok{, }\StringTok{"South"}\NormalTok{, }\StringTok{"Center"}\NormalTok{, }\StringTok{"West"}\NormalTok{, }\StringTok{"East"}\NormalTok{)}
\NormalTok{    )}
\NormalTok{  )}
\end{Highlighting}
\end{Shaded}

En el análisis de bases de datos también es frecuente recodificar variables categóricas. El siguiente código crea \texttt{region\_id}, una versión codificada de \texttt{Region} que conserva el orden de las regiones y asigna códigos de dos dígitos. Esta forma de recodificación facilita la presentación o el cruce con otras fuentes de información.

\begin{Shaded}
\begin{Highlighting}[]
\NormalTok{bigcity\_data }\OtherTok{\textless{}{-}}\NormalTok{ bigcity\_data }\SpecialCharTok{\%\textgreater{}\%}
  \FunctionTok{mutate}\NormalTok{(}
    \AttributeTok{region\_id =} \FunctionTok{factor}\NormalTok{(}
\NormalTok{      Region,}
      \AttributeTok{levels =} \FunctionTok{c}\NormalTok{(}\StringTok{"North"}\NormalTok{, }\StringTok{"South"}\NormalTok{, }\StringTok{"Center"}\NormalTok{, }\StringTok{"West"}\NormalTok{, }\StringTok{"East"}\NormalTok{),}
      \AttributeTok{labels =} \FunctionTok{c}\NormalTok{(}\StringTok{"01"}\NormalTok{, }\StringTok{"02"}\NormalTok{, }\StringTok{"03"}\NormalTok{, }\StringTok{"04"}\NormalTok{, }\StringTok{"05"}\NormalTok{)}
\NormalTok{    )}
\NormalTok{  )}
\end{Highlighting}
\end{Shaded}

El siguiente ejemplo crea \texttt{log\_income}, que corresponde al logaritmo natural del ingreso (transformación de la escala del ingreso). Esta operación se usa con frecuencia cuando una variable monetaria presenta alta dispersión. Luego se muestran las primeras filas de las variables relevantes.

\begin{Shaded}
\begin{Highlighting}[]
\NormalTok{log\_income\_data }\OtherTok{\textless{}{-}}\NormalTok{ working\_data }\SpecialCharTok{\%\textgreater{}\%}
  \FunctionTok{mutate}\NormalTok{(}\AttributeTok{log\_income =} \FunctionTok{log}\NormalTok{(Income))}

\NormalTok{log\_income\_data }\SpecialCharTok{\%\textgreater{}\%}
  \FunctionTok{select}\NormalTok{(PersonID, Income, log\_income) }\SpecialCharTok{\%\textgreater{}\%}
  \FunctionTok{head}\NormalTok{()}
\end{Highlighting}
\end{Shaded}

\begin{verbatim}
##   PersonID Income log_income
## 1  idPer01    555       6.32
## 2  idPer02    555       6.32
## 3  idPer03    555       6.32
## 4  idPer04    555       6.32
## 5  idPer05    555       6.32
## 6  idPer01    298       5.70
\end{verbatim}

Note que \texttt{mutate()} puede crear más de una variable en el mismo paso. Además, las variables creadas al inicio de \texttt{mutate()} pueden ser utilizadas en cálculos posteriores dentro del mismo llamado. Por ejemplo, \texttt{income\_2} duplica el ingreso original e \texttt{income\_4} duplica \texttt{income\_2}. La siguiente salida muestra que \texttt{income\_4} equivale a cuatro veces el ingreso inicial.

\begin{Shaded}
\begin{Highlighting}[]
\NormalTok{income\_transform\_data }\OtherTok{\textless{}{-}}\NormalTok{ log\_income\_data }\SpecialCharTok{\%\textgreater{}\%}
  \FunctionTok{mutate}\NormalTok{(}
    \AttributeTok{income\_2 =} \DecValTok{2} \SpecialCharTok{*}\NormalTok{ Income,}
    \AttributeTok{income\_4 =} \DecValTok{2} \SpecialCharTok{*}\NormalTok{ income\_2}
\NormalTok{  )}

\NormalTok{income\_transform\_data }\SpecialCharTok{\%\textgreater{}\%}
  \FunctionTok{select}\NormalTok{(PersonID, income\_2, income\_4) }\SpecialCharTok{\%\textgreater{}\%}
  \FunctionTok{head}\NormalTok{()}
\end{Highlighting}
\end{Shaded}

\begin{verbatim}
##   PersonID income_2 income_4
## 1  idPer01     1110     2220
## 2  idPer02     1110     2220
## 3  idPer03     1110     2220
## 4  idPer04     1110     2220
## 5  idPer05     1110     2220
## 6  idPer01      597     1193
\end{verbatim}

También se puede crear una variable categórica mediante condiciones. En este ejemplo, la función \texttt{case\_when()} clasifica la edad en rangos y luego convierte el resultado en un factor ordenado. El resultado es \texttt{age\_group}, una variable de grupos etarios ordenados. Esta transformación facilita la producción de tablas por tramos de edad.

\begin{Shaded}
\begin{Highlighting}[]
\NormalTok{bigcity\_data }\OtherTok{\textless{}{-}}\NormalTok{ bigcity\_data }\SpecialCharTok{\%\textgreater{}\%}
  \FunctionTok{mutate}\NormalTok{(}
    \AttributeTok{age\_group =} \FunctionTok{case\_when}\NormalTok{(}
\NormalTok{      Age }\SpecialCharTok{\textless{}=} \DecValTok{5} \SpecialCharTok{\textasciitilde{}} \StringTok{"0{-}5"}\NormalTok{,}
\NormalTok{      Age }\SpecialCharTok{\textless{}=} \DecValTok{15} \SpecialCharTok{\textasciitilde{}} \StringTok{"6{-}15"}\NormalTok{,}
\NormalTok{      Age }\SpecialCharTok{\textless{}=} \DecValTok{30} \SpecialCharTok{\textasciitilde{}} \StringTok{"16{-}30"}\NormalTok{,}
\NormalTok{      Age }\SpecialCharTok{\textless{}=} \DecValTok{45} \SpecialCharTok{\textasciitilde{}} \StringTok{"31{-}45"}\NormalTok{,}
\NormalTok{      Age }\SpecialCharTok{\textless{}=} \DecValTok{60} \SpecialCharTok{\textasciitilde{}} \StringTok{"46{-}60"}\NormalTok{,}
      \ConstantTok{TRUE} \SpecialCharTok{\textasciitilde{}} \StringTok{"Over 60"}
\NormalTok{    ),}
    \AttributeTok{age\_group =} \FunctionTok{factor}\NormalTok{(}
\NormalTok{      age\_group,}
      \AttributeTok{levels =} \FunctionTok{c}\NormalTok{(}
        \StringTok{"0{-}5"}\NormalTok{, }\StringTok{"6{-}15"}\NormalTok{, }\StringTok{"16{-}30"}\NormalTok{,}
        \StringTok{"31{-}45"}\NormalTok{, }\StringTok{"46{-}60"}\NormalTok{, }\StringTok{"Over 60"}
\NormalTok{      ),}
      \AttributeTok{ordered =} \ConstantTok{TRUE}
\NormalTok{    )}
\NormalTok{  )}
\end{Highlighting}
\end{Shaded}

\subsection{\texorpdfstring{\texttt{summarise}: resúmenes descriptivos}{summarise: resúmenes descriptivos}}\label{summarise-resumenes-descriptivos}

La función \texttt{summarise()} permite construir una tabla resumen a partir de una o varias variables de la base de datos. En el siguiente ejemplo se calculan dos estadísticas descriptivas de la variable \texttt{Income}: el total y la media. Como resultado, se obtiene una tabla con una sola fila que contiene la suma de los ingresos observados y el ingreso promedio de los registros incluidos en la base.

\begin{Shaded}
\begin{Highlighting}[]
\NormalTok{bigcity\_data }\SpecialCharTok{\%\textgreater{}\%}
  \FunctionTok{summarise}\NormalTok{(}
    \AttributeTok{total\_income =} \FunctionTok{sum}\NormalTok{(Income),}
    \AttributeTok{mean\_income =} \FunctionTok{mean}\NormalTok{(Income)}
\NormalTok{  )}
\end{Highlighting}
\end{Shaded}

\begin{verbatim}
##   total_income mean_income
## 1     87893117         585
\end{verbatim}

También pueden calcularse medidas de localización para describir la posición de los valores dentro de la distribución. En este caso, se obtiene la mediana, el primer decil, el noveno decil y la diferencia entre estos dos últimos. El resultado permite resumir distintos puntos de la distribución del ingreso, mientras que el rango decílico muestra la distancia entre el 10\% inferior y el 10\% superior de los ingresos observados.

\begin{Shaded}
\begin{Highlighting}[]
\NormalTok{bigcity\_data }\SpecialCharTok{\%\textgreater{}\%}
  \FunctionTok{summarise}\NormalTok{(}
    \AttributeTok{median\_income =} \FunctionTok{median}\NormalTok{(Income),}
    \AttributeTok{decile\_1 =} \FunctionTok{quantile}\NormalTok{(Income, }\FloatTok{0.1}\NormalTok{),}
    \AttributeTok{decile\_9 =} \FunctionTok{quantile}\NormalTok{(Income, }\FloatTok{0.9}\NormalTok{),}
    \AttributeTok{decile\_range =}\NormalTok{ decile\_9 }\SpecialCharTok{{-}}\NormalTok{ decile\_1}
\NormalTok{  )}
\end{Highlighting}
\end{Shaded}

\begin{verbatim}
##   median_income decile_1 decile_9 decile_range
## 1           449      165     1126          961
\end{verbatim}

Con \texttt{summarise()} también es posible calcular medidas de dispersión, tales como la varianza, la desviación estándar, el valor mínimo, el valor máximo, el rango y el rango intercuartílico. Estas estadísticas permiten describir el grado de variabilidad de los ingresos observados en la muestra. En particular, el rango se obtiene como la diferencia entre el valor máximo y el valor mínimo, mientras que el rango intercuartílico resume la dispersión del 50\% central de la distribución.

\begin{Shaded}
\begin{Highlighting}[]
\NormalTok{bigcity\_data }\SpecialCharTok{\%\textgreater{}\%}
  \FunctionTok{summarise}\NormalTok{(}
    \AttributeTok{variance\_income =} \FunctionTok{var}\NormalTok{(Income),}
    \AttributeTok{sd\_income =} \FunctionTok{sd}\NormalTok{(Income),}
    \AttributeTok{min\_income =} \FunctionTok{min}\NormalTok{(Income),}
    \AttributeTok{max\_income =} \FunctionTok{max}\NormalTok{(Income),}
    \AttributeTok{range\_income =}\NormalTok{ max\_income }\SpecialCharTok{{-}}\NormalTok{ min\_income,}
    \AttributeTok{iqr\_income =} \FunctionTok{IQR}\NormalTok{(Income)}
\NormalTok{  )}
\end{Highlighting}
\end{Shaded}

\begin{verbatim}
##   variance_income sd_income min_income max_income range_income iqr_income
## 1          332463       577         10      32920        32910        468
\end{verbatim}

\subsection{\texorpdfstring{\texttt{group\_by}: agrupación de casos}{group\_by: agrupación de casos}}\label{group-by-agrupacion-de-casos}

La función \texttt{group\_by()} permite agrupar la base de datos de acuerdo con una o más variables. Al combinarla con \texttt{summarise()}, las estadísticas se calculan de manera independiente dentro de cada grupo definido. En el siguiente ejemplo, se cuenta el número de registros por región y luego se ordena el resultado de mayor a menor. La siguiente salida permite identificar cuántos registros están asociados con cada región de la base de datos.

\begin{Shaded}
\begin{Highlighting}[]
\NormalTok{bigcity\_data }\SpecialCharTok{\%\textgreater{}\%}
  \FunctionTok{group\_by}\NormalTok{(Region) }\SpecialCharTok{\%\textgreater{}\%}
  \FunctionTok{summarise}\NormalTok{(}\AttributeTok{n =} \FunctionTok{n}\NormalTok{()) }\SpecialCharTok{\%\textgreater{}\%}
  \FunctionTok{arrange}\NormalTok{(}\FunctionTok{desc}\NormalTok{(n))}
\end{Highlighting}
\end{Shaded}

\begin{verbatim}
## # A tibble: 5 x 2
##   Region     n
##   <fct>  <int>
## 1 East   39160
## 2 West   33868
## 3 Center 25944
## 4 South  25898
## 5 North  25396
\end{verbatim}

Para los conteos simples, \texttt{count()} ofrece una forma abreviada de escribir la combinación entre \texttt{group\_by()} y \texttt{summarise(n\ =\ n())}. En la práctica, ambas alternativas producen el mismo resultado, pues agrupan la base de datos según una variable categórica y calculan cuántos registros pertenecen a cada grupo. Por ejemplo, el siguiente código reproduce los mismos resultados que el código anterior.

\begin{Shaded}
\begin{Highlighting}[]
\NormalTok{bigcity\_data }\SpecialCharTok{\%\textgreater{}\%}
  \FunctionTok{count}\NormalTok{(Region) }\SpecialCharTok{\%\textgreater{}\%}
  \FunctionTok{arrange}\NormalTok{(}\FunctionTok{desc}\NormalTok{(n))}
\end{Highlighting}
\end{Shaded}

\begin{verbatim}
##   Region     n
## 1   East 39160
## 2   West 33868
## 3 Center 25944
## 4  South 25898
## 5  North 25396
\end{verbatim}

Para obtener el número de encuestas efectivas por sexo, se aplica la misma lógica de agrupación y conteo. En este caso, la base se agrupa según la variable \texttt{Sex} y luego se calcula el número de registros asociados con cada una de sus categorías. El resultado permite identificar cuántas observaciones corresponden a cada sexo dentro de la base de datos.

\begin{Shaded}
\begin{Highlighting}[]
\NormalTok{bigcity\_data }\SpecialCharTok{\%\textgreater{}\%}
  \FunctionTok{count}\NormalTok{(Sex) }\SpecialCharTok{\%\textgreater{}\%}
  \FunctionTok{arrange}\NormalTok{(}\FunctionTok{desc}\NormalTok{(n))}
\end{Highlighting}
\end{Shaded}

\begin{verbatim}
##      Sex     n
## 1 Female 79190
## 2   Male 71076
\end{verbatim}

El conteo también puede realizarse según la zona geográfica. En este caso, la base se agrupa por la variable \texttt{Zone} y se calcula el número de registros correspondiente a cada una de sus categorías. La siguiente salida permite conocer cómo se distribuyen las observaciones entre las distintas zonas de la base de datos.

\begin{Shaded}
\begin{Highlighting}[]
\NormalTok{bigcity\_data }\SpecialCharTok{\%\textgreater{}\%}
  \FunctionTok{count}\NormalTok{(Zone) }\SpecialCharTok{\%\textgreater{}\%}
  \FunctionTok{arrange}\NormalTok{(}\FunctionTok{desc}\NormalTok{(n))}
\end{Highlighting}
\end{Shaded}

\begin{verbatim}
##    Zone     n
## 1 Urban 78164
## 2 Rural 72102
\end{verbatim}

Además de los conteos, también es posible calcular medidas resumen por grupo, como la media. En este caso, para obtener el ingreso promedio por región junto con el número total de registros, se agrupa la base según la variable \texttt{Region} y luego se aplica \texttt{summarise()}. El resultado es una tabla regional que presenta, para cada región, el total de observaciones y el ingreso medio correspondiente.

\begin{Shaded}
\begin{Highlighting}[]
\NormalTok{bigcity\_data }\SpecialCharTok{\%\textgreater{}\%}
  \FunctionTok{group\_by}\NormalTok{(Region) }\SpecialCharTok{\%\textgreater{}\%}
  \FunctionTok{summarise}\NormalTok{(}
    \AttributeTok{n =} \FunctionTok{n}\NormalTok{(),}
    \AttributeTok{mean\_income =} \FunctionTok{mean}\NormalTok{(Income)}
\NormalTok{  )}
\end{Highlighting}
\end{Shaded}

\begin{verbatim}
## # A tibble: 5 x 3
##   Region     n mean_income
##   <fct>  <int>       <dbl>
## 1 North  25396        509.
## 2 South  25898        644.
## 3 Center 25944        701.
## 4 West   33868        571.
## 5 East   39160        530.
\end{verbatim}

El mismo procedimiento puede aplicarse para calcular el ingreso medio por sexo. En este caso, la base se agrupa según la variable \texttt{Sex} y luego se obtiene el promedio de \texttt{Income} para cada categoría. La tabla resultante permite comparar el ingreso medio observado entre los distintos grupos definidos por el sexo.

\begin{Shaded}
\begin{Highlighting}[]
\NormalTok{bigcity\_data }\SpecialCharTok{\%\textgreater{}\%}
  \FunctionTok{group\_by}\NormalTok{(Sex) }\SpecialCharTok{\%\textgreater{}\%}
  \FunctionTok{summarise}\NormalTok{(}
    \AttributeTok{n =} \FunctionTok{n}\NormalTok{(),}
    \AttributeTok{mean\_income =} \FunctionTok{mean}\NormalTok{(Income)}
\NormalTok{  )}
\end{Highlighting}
\end{Shaded}

\begin{verbatim}
## # A tibble: 2 x 3
##   Sex        n mean_income
##   <chr>  <int>       <dbl>
## 1 Female 79190        579.
## 2 Male   71076        591.
\end{verbatim}

Por último, se calculan la media, la desviación estándar y el rango intercuartílico de los ingresos según la condición de ocupación. Para ello, la base se agrupa por la variable \texttt{Employment} y luego se obtienen las estadísticas descriptivas de \texttt{Income} dentro de cada grupo. La siguiente salida presenta un resumen del comportamiento de los ingresos para cada categoría de ocupación.

\begin{Shaded}
\begin{Highlighting}[]
\NormalTok{bigcity\_data }\SpecialCharTok{\%\textgreater{}\%}
  \FunctionTok{group\_by}\NormalTok{(Employment) }\SpecialCharTok{\%\textgreater{}\%}
  \FunctionTok{summarise}\NormalTok{(}
    \AttributeTok{n =} \FunctionTok{n}\NormalTok{(),}
    \AttributeTok{mean\_income =} \FunctionTok{mean}\NormalTok{(Income),}
    \AttributeTok{sd\_income =} \FunctionTok{sd}\NormalTok{(Income),}
    \AttributeTok{iqr\_income =} \FunctionTok{IQR}\NormalTok{(Income)}
\NormalTok{  )}
\end{Highlighting}
\end{Shaded}

\begin{verbatim}
## # A tibble: 4 x 5
##   Employment     n mean_income sd_income iqr_income
##   <fct>      <int>       <dbl>     <dbl>      <dbl>
## 1 Unemployed  4630        429.      375.       392.
## 2 Inactive   44104        532.      553.       439.
## 3 Employed   62188        661.      606.       529.
## 4 <NA>       39344        541.      558.       407.
\end{verbatim}

\chapter{Inferencia en poblaciones finitas}\label{ap2}

Una cuestión fundamental en el análisis de encuestas es identificar cuál es la medida de probabilidad que rige la inferencia. En la inferencia estadística clásica se suele asumir que las observaciones disponibles son variables aleatorias independientes e idénticamente distribuidas (IID). Bajo ese enfoque, la aleatoriedad proviene del modelo que genera los datos. Sin embargo, como señala \citet{kish1987statistical}, este supuesto no describe adecuadamente la información proveniente de encuestas con diseños muestrales complejos, donde las observaciones son seleccionadas mediante estratificación, conglomeración y probabilidades de inclusión desiguales.

\section{Dos posibles inferencias}\label{dos-posibles-inferencias}

En encuestas de hogares conviene distinguir tres objetos. Primero, la población finita, denotada como \(U = \{1,2,\ldots,N\}\), la cual contiene todas las unidades de interés. Segundo, la variable de estudio \(y_k\), observada o definida para cada unidad \(k \in U\). Tercero, el diseño muestral \(p(s)\), que asigna probabilidades a las posibles muestras \(s \subset U\). En la notación utilizada a lo largo de este documento, una unidad puede identificarse como \(k\) dentro del conglomerado o UPM \(i\), perteneciente al estrato \(h\), de modo que los estimadores ponderados suelen escribirse como sumas de la forma

\[
\hat{t}_y = \sum_h \sum_i \sum_k w_{hik} \ y_{hik}
\]

donde \(w_{hik}\) representa el factor de expansión o peso muestral asociado con la unidad observada. En la teoría de muestreo, los valores \(y_k\) se consideran fijos una vez definida la población. La aleatoriedad no está en los valores de la variable, sino en los indicadores de inclusión

\[
I_k =
\begin{cases}
1, & \text{si } k \in s,\\
0, & \text{si } k \notin s,
\end{cases}
\]

que inducen la probabilidad de inclusión \(\pi_k = \Pr_p(I_k=1)\) y, a su vez, el correspondiente peso de muestreo, definido como \(w_k = \frac{1}{\pi_k}\).

Un estimador que ignora los pesos de muestreo puede no representar adecuadamente a la población objetivo, especialmente cuando la muestra proviene de un diseño complejo. En la inferencia basada en el diseño, la población se considera fija y la fuente de aleatoriedad proviene del mecanismo de selección de la muestra; por ello, sus elementos centrales son las probabilidades de inclusión \(\pi_k\) y los pesos de muestreo \(w_k\).

Por el contrario, la inferencia basada en modelos considera que los valores de la variable de interés en toda la población se conciben como realizaciones de un proceso estocástico, usualmente denotado por \(\xi\). Así, mientras la inferencia basada en los modelos estudia propiedades como \(E_{\xi}(Y_k)\) y \(Var_{\xi}(Y_k)\), la inferencia basada en el diseño se orienta a incorporar las características del plan muestral en los estimadores, para que sean insesgados \(E_{p}(\hat{t}_y) = t_y\).

Esta distinción permite entender por qué, aun cuando exista un modelo razonable para la variable de interés, el diseño muestral debe incorporarse explícitamente si se desea obtener inferencias válidas para la población objetivo. Para comprender la distinción entre inferencia basada en el modelo e inferencia basada en el diseño, partimos de un ejemplo sencillo, adaptado de \citet{binder2011estimating}. Supóngase que se generan \(N=100\) realizaciones independientes de una variable Bernoulli con parámetro \(\theta=0.3\), donde \(\theta\) representa la proporción esperada de personas desempleadas en una población generada por un modelo. Si \(Y_k \sim Bernoulli(\theta)\), entonces la esperanza bajo el modelo es \(E_{\xi}(Y_k) = \theta\), mientras que la varianza bajo el modelo es \(Var_{\xi}(Y_k) = \theta(1-\theta)\).

En la literatura especializada, el modelo \(\xi\) recibe el nombre de modelo de superpoblación, pues genera poblaciones finitas. De esta forma, para una población finita generada por ese modelo, la media poblacional es \(\bar{Y}_U = \frac{1}{N}\sum_{k \in U}Y_k\). Nótese que, bajo el modelo \(\xi\), esta media es insesgada para \(\theta\), puesto que

\[
E_{\xi}(\bar{Y}_U) =
E_{\xi}\left(\frac{1}{N}\sum_{k \in U}Y_k\right)
= \frac{1}{N}\sum_{k \in U}E_{\xi}(Y_k)
= \theta.
\]

La siguiente simulación de Monte Carlo reproduce este proceso. En cada repetición se genera una población completa de tamaño \(N=100\) y se calcula su media. El promedio de esas medias poblacionales debe aproximarse a \(\theta\).

\begin{Shaded}
\begin{Highlighting}[]
\FunctionTok{set.seed}\NormalTok{(}\DecValTok{2026}\NormalTok{)}
\NormalTok{population\_size }\OtherTok{\textless{}{-}} \DecValTok{100}
\NormalTok{theta }\OtherTok{\textless{}{-}} \FloatTok{0.3}
\NormalTok{n\_sim\_model }\OtherTok{\textless{}{-}} \DecValTok{1000}
\NormalTok{model\_estimates }\OtherTok{\textless{}{-}} \FunctionTok{rep}\NormalTok{(}\ConstantTok{NA}\NormalTok{, n\_sim\_model)}

\ControlFlowTok{for}\NormalTok{ (sim }\ControlFlowTok{in} \FunctionTok{seq\_len}\NormalTok{(n\_sim\_model)) \{}
\NormalTok{  outcome }\OtherTok{\textless{}{-}} \FunctionTok{rbinom}\NormalTok{(population\_size, }\DecValTok{1}\NormalTok{, theta)}
\NormalTok{  model\_estimates[sim] }\OtherTok{\textless{}{-}} \FunctionTok{mean}\NormalTok{(outcome)}
\NormalTok{\}}

\FunctionTok{cbind}\NormalTok{(}
\NormalTok{  theta,}
  \AttributeTok{expected\_model\_mean =} \FunctionTok{mean}\NormalTok{(model\_estimates),}
  \AttributeTok{bias =} \FunctionTok{mean}\NormalTok{(model\_estimates) }\SpecialCharTok{{-}}\NormalTok{ theta}
\NormalTok{)}
\end{Highlighting}
\end{Shaded}

\begin{verbatim}
##      theta expected_model_mean    bias
## [1,]   0.3               0.302 0.00225
\end{verbatim}

El código fija una semilla para que el ejercicio sea reproducible. Luego genera \texttt{n\_sim\_model} poblaciones independientes mediante \texttt{rbinom()}. La salida reporta el parámetro verdadero \texttt{theta}, el promedio de las medias simuladas \texttt{expected\_model\_mean} y el sesgo Monte Carlo. Como se espera bajo el modelo Bernoulli IID, \texttt{expected\_model\_mean} queda muy cerca de 0.3 y el sesgo es cercano a cero. Las pequeñas diferencias se deben a que se usan 1000 simulaciones, no un número infinito de repeticiones.

En la teoría de muestreo, en cambio, las características de interés son parámetros fijos de la población. Si una persona está desempleada, su estado se considera un valor fijo de la población finita. Lo aleatorio es el mecanismo mediante el cual se selecciona la muestra. El total poblacional se define como \(t_y = \sum_{k \in U} y_k\), mientras que la media poblacional es \(\bar{y}_U = \frac{t_y}{N}\). Si se observa una muestra \(s\), el estimador de Horvitz-Thompson del total es

\[
\hat{t}_{y} = \sum_{k \in s}\frac{y_k}{\pi_k}
= \sum_{k \in s}w_k y_k.
\]

Como se mencionó en los capítulos anteriores, cuando el tamaño poblacional \(N\) es desconocido, la media poblacional puede estimarse mediante el estimador de razón ponderado, definido como

\[
\hat{\bar{y}} =
\frac{\sum_{k \in s}w_k y_k}
{\sum_{k \in s}w_k}
\]

Estos estimadores incorporan las probabilidades de inclusión en su forma funcional y es esa, la diferencia central frente al promedio simple no ponderado de una muestra. Para ilustrar el impacto del diseño, supóngase que la población anterior se divide en \(N_I\) UPMs. La UPM \(i\) tiene tamaño \(N_i\), total \(t_{y_i}=\sum_{k \in U_i}y_{ik}\) y media \(\bar{y}_i=t_{y_i}/N_i\). Si se seleccionan \(n_I\) UPMs con probabilidades proporcionales al tamaño, entonces sus probabilidades de inclusión toman la siguiente forma:

\[
\pi_{Ii} = \frac{n_I N_i}{N}
\]

Cuando todos los miembros de las UPMs seleccionadas son observados, el estimador ponderado de la media poblacional, que es aproximadamente insesgado, puede escribirse como

\[
\hat{\bar{y}}
= \frac{1}{N}\sum_{i \in S_I}\frac{t_{y_i}}{\pi_{Ii}}
= \frac{1}{N}\sum_{i \in S_I}\frac{N_i\bar{y}_i}{n_I N_i/N}
= \frac{1}{n_I}\sum_{i \in S_I}\bar{y}_i.
\]

Esta expresión ejemplifica por qué, cuando el diseño de muestreo es complejo (en este caso selección proporcional al tamaño), el promedio de las medias de los conglomerados brinda insesgamiento. En contraste, el promedio simple de las personas observadas en las UPMs seleccionadas, sería:

\[
\bar{y}_s =
\frac{\sum_{i \in S_I}t_{y_i}}
     {\sum_{i \in S_I}N_i},
\]

el cual trata la muestra como si fuera autoponderada. Este estimador no es insesgado cuando las probabilidades de selección no son iguales o cuando el tamaño del conglomerado se relaciona con la variable de interés. La siguiente simulación fija una población finita y repite muchas veces la selección de hogares con probabilidades proporcionales al tamaño mediante \texttt{S.piPS()} del paquete \texttt{TeachingSampling} \citep{GutierrezT2020}. En cada muestra se calculan dos estimadores: \texttt{design\_estimates}, que corresponde a \(\hat{\bar{y}}\), y \texttt{simple\_estimates}, que corresponde al promedio simple de las personas observadas en las UPMs seleccionadas, \(\bar{y}_s\). Además, estos estimadores se comparan con el parámetro del modelo \texttt{theta\_population}, correspondiente a \(\theta = 0.3\).

\begin{Shaded}
\begin{Highlighting}[]
\FunctionTok{library}\NormalTok{(TeachingSampling)}

\FunctionTok{set.seed}\NormalTok{(}\DecValTok{2026}\NormalTok{)}
\NormalTok{population\_size }\OtherTok{\textless{}{-}} \DecValTok{100}
\NormalTok{theta }\OtherTok{\textless{}{-}} \FloatTok{0.3}
\NormalTok{n\_sim\_sampling }\OtherTok{\textless{}{-}} \DecValTok{1000}
\NormalTok{design\_estimates }\OtherTok{\textless{}{-}}\NormalTok{ simple\_estimates }\OtherTok{\textless{}{-}} \FunctionTok{rep}\NormalTok{(}\ConstantTok{NA}\NormalTok{, n\_sim\_sampling)}

\NormalTok{cluster\_size }\OtherTok{\textless{}{-}} \FunctionTok{rep}\NormalTok{(}\DecValTok{2}\SpecialCharTok{:}\DecValTok{6}\NormalTok{, }\AttributeTok{each =} \DecValTok{5}\NormalTok{)}
\NormalTok{n\_clusters }\OtherTok{\textless{}{-}} \FunctionTok{length}\NormalTok{(cluster\_size)}
\NormalTok{cluster\_id }\OtherTok{\textless{}{-}} \FunctionTok{rep}\NormalTok{(}\FunctionTok{seq\_len}\NormalTok{(n\_clusters), cluster\_size)}
\NormalTok{size\_center }\OtherTok{\textless{}{-}} \FunctionTok{weighted.mean}\NormalTok{(cluster\_size, cluster\_size)}
\NormalTok{prob\_person }\OtherTok{\textless{}{-}}\NormalTok{ theta }\SpecialCharTok{+} \FunctionTok{rep}\NormalTok{((cluster\_size }\SpecialCharTok{{-}}\NormalTok{ size\_center) }\SpecialCharTok{/} \DecValTok{12}\NormalTok{, cluster\_size)}
\NormalTok{outcome }\OtherTok{\textless{}{-}} \FunctionTok{rbinom}\NormalTok{(population\_size, }\DecValTok{1}\NormalTok{, prob\_person)}
\NormalTok{theta\_population }\OtherTok{\textless{}{-}} \FunctionTok{mean}\NormalTok{(outcome)}
\NormalTok{cluster\_totals }\OtherTok{\textless{}{-}} \FunctionTok{tapply}\NormalTok{(outcome, cluster\_id, sum)}
\NormalTok{cluster\_means }\OtherTok{\textless{}{-}} \FunctionTok{tapply}\NormalTok{(outcome, cluster\_id, mean)}

\NormalTok{sampled\_cluster\_count }\OtherTok{\textless{}{-}} \FunctionTok{floor}\NormalTok{(n\_clusters }\SpecialCharTok{*} \FloatTok{0.3}\NormalTok{)}
\ControlFlowTok{for}\NormalTok{ (sim }\ControlFlowTok{in} \FunctionTok{seq\_len}\NormalTok{(n\_sim\_sampling)) \{}
\NormalTok{  pps\_sample }\OtherTok{\textless{}{-}} \FunctionTok{S.piPS}\NormalTok{(sampled\_cluster\_count, cluster\_size)}
\NormalTok{  sampled\_clusters }\OtherTok{\textless{}{-}}\NormalTok{ pps\_sample[, }\DecValTok{1}\NormalTok{]}
\NormalTok{  sampled\_cluster\_size }\OtherTok{\textless{}{-}}\NormalTok{ cluster\_size[sampled\_clusters]}
\NormalTok{  sampled\_cluster\_total }\OtherTok{\textless{}{-}}\NormalTok{ cluster\_totals[sampled\_clusters]}
\NormalTok{  sampled\_cluster\_mean }\OtherTok{\textless{}{-}}\NormalTok{ cluster\_means[sampled\_clusters]}
\NormalTok{  design\_estimates[sim] }\OtherTok{\textless{}{-}} \FunctionTok{mean}\NormalTok{(sampled\_cluster\_mean)}
\NormalTok{  simple\_estimates[sim] }\OtherTok{\textless{}{-}}
    \FunctionTok{sum}\NormalTok{(sampled\_cluster\_total) }\SpecialCharTok{/} \FunctionTok{sum}\NormalTok{(sampled\_cluster\_size)}
\NormalTok{\}}

\FunctionTok{cbind}\NormalTok{(}
\NormalTok{  theta\_population,}
  \AttributeTok{expected\_pps =} \FunctionTok{mean}\NormalTok{(design\_estimates),}
  \AttributeTok{bias\_pps =} \FunctionTok{mean}\NormalTok{(design\_estimates) }\SpecialCharTok{{-}}\NormalTok{ theta\_population,}
  \AttributeTok{expected\_simple =} \FunctionTok{mean}\NormalTok{(simple\_estimates),}
  \AttributeTok{bias\_simple =} \FunctionTok{mean}\NormalTok{(simple\_estimates) }\SpecialCharTok{{-}}\NormalTok{ theta\_population}
\NormalTok{)}
\end{Highlighting}
\end{Shaded}

\begin{verbatim}
##      theta_population expected_pps bias_pps expected_simple bias_simple
## [1,]              0.3        0.304  0.00368           0.358      0.0585
\end{verbatim}

El primer bloque del código genera hogares de tamaños desiguales, entre 2 y 6 personas. Para hacer visible el efecto del diseño, la probabilidad individual de desempleo varía suavemente con el tamaño del hogar, pero se centra de manera que el promedio esperado de la población siga siendo cercano a \texttt{theta}. Luego se calcula \texttt{theta\_population}, que es la proporción real de desempleo en esa población específica. En cada repetición, \texttt{S.piPS()} selecciona UPMs con probabilidades proporcionales al tamaño; después se calculan la media de las medias de las UPMs seleccionadas y el promedio simple de las personas observadas. La salida compara ambos estimadores con \texttt{theta\_population}. El sesgo de Monte Carlo de \texttt{expected\_pps} es cercano a cero, mientras que el de \texttt{expected\_simple} no lo es, mostrando resultado de ignorar el diseño en este ejercicio.

\section{La integración de ambas inferencias}\label{la-integracion-de-ambas-inferencias}

No obstante, en muchas aplicaciones ambos marcos inferenciales no se presentan como enfoques excluyentes, sino como componentes complementarios de un mismo problema. En particular, puede considerarse que la muestra ha sido obtenida mediante un diseño probabilístico complejo y que, al mismo tiempo, los valores de la variable de interés responden a un modelo de superpoblación. Esta formulación conduce a un esquema de inferencia doble, en el cual la incertidumbre no proviene de una única fuente, sino de la combinación entre el mecanismo que genera los valores poblacionales (modelo de superpoblación) y el mecanismo que determina qué unidades son observadas (diseño de muestreo).

Formalmente, intervienen dos medidas de probabilidad: \(\xi\), asociada al proceso generador de los valores \(Y_k\), y \(p(\cdot)\), asociada al diseño muestral que selecciona la muestra \(s\). Por tanto, las propiedades de los estimadores deben evaluarse considerando simultáneamente ambas fuentes de variación. Bajo este marco, la esperanza combinada puede expresarse como

\[
E_{\xi p}(\hat{\theta})
= E_{\xi}\{E_p(\hat{\theta}\mid U)\},
\]

donde primero se evalúa el comportamiento del estimador bajo el diseño, condicionado a la población finita generada, y luego se promedia sobre el modelo de superpoblación. En este contexto, la verosimilitud conjunta completa puede resultar intratable, especialmente cuando las probabilidades de inclusión dependen de variables relacionadas con la respuesta. La solución propuesta por \citet{pfeffermann1993role} es la técnica de máxima pseudo-verosimilitud (MPV), que pondera las contribuciones de cada unidad por el inverso de su probabilidad de inclusión.

\subsection{Método de la máxima verosimilitud}\label{sec-mv}

El método de máxima verosimilitud parte de una muestra aleatoria \(y_1,\ldots,y_n\) generada por una distribución con densidad o función de probabilidad \(f(y;\theta)\). Si las observaciones son independientes e idénticamente distribuidas (IID), la función de verosimilitud se define como

\[
L(\theta)
= \prod_{k=1}^{N} f(y_k;\theta).
\]

Como los productos pueden ser difíciles de manipular, se utiliza la log-verosimilitud:

\[
\ell(\theta)
= \sum_{k=1}^{N}\log f(y_k;\theta)
\]

El estimador de máxima verosimilitud \(\hat{\theta}\) es el valor de \(\theta\) que maximiza \(\ell(\theta)\). Si la función es diferenciable, este valor satisface las ecuaciones de score:

\[
U(\theta)
= \frac{\partial \ \ell(\theta)}{\partial \ \theta}
= \sum_{k=1}^{N} u_k(\theta)
= 0
\]

donde \(u_k(\theta) = \frac{\partial}{\partial\theta}\log f(y_k;\theta)\) es la contribución de la unidad \(k\) al score total. Para una distribución Bernoulli con parámetro \(\theta\), la función de probabilidad es \(f(y_k;\theta) = \theta^{y_k}(1-\theta)^{1-y_k}\), para \(y_k \in \{0,1\}\). La log-verosimilitud es

\[
\ell(\theta)
= \sum_{k=1}^{N}
\left[
y_k\log(\theta)
+(1-y_k)\log(1-\theta)
\right]
\]

Al derivar e igualar a cero y despejar, se obtiene el estimador MV, definido por:

\[
\hat{\theta}_{MV}
= \frac{1}{N}\sum_{k=1}^{N}y_k
\]

Así, la proporción muestral es el estimador natural del parámetro Bernoulli. El punto clave es que este resultado depende del supuesto de que las observaciones tienen igual distribución y no provienen de un diseño con probabilidades desiguales.

Por otro lado, en un modelo de regresión lineal múltiple con errores normales, se tiene que \(\mathbf{y} \sim N(\mathbf{X}\boldsymbol{\beta},\sigma^2 I)\). En este caso, la log-verosimilitud, omitiendo constantes, puede escribirse como

\[
\ell(\boldsymbol{\beta},\sigma^2) = -\frac{n}{2}\log(\sigma^2) -\frac{1}{2\sigma^2}(\mathbf{y}-\mathbf{X}\boldsymbol{\beta})'(\mathbf{y}-\mathbf{X}\boldsymbol{\beta})
\]

Maximizar esta expresión respecto de \(\beta\) equivale a minimizar la suma de cuadrados de los residuos. Por ello, el estimador de máxima verosimilitud de \(\beta\) coincide con el estimador de mínimos cuadrados ordinarios \(\hat{\boldsymbol{\beta}}_{MV} = (\mathbf{X}'\mathbf{X})^{-1}\mathbf{X}'\mathbf{y}\). Este resultado es importante porque muestra que muchos procedimientos estadísticos usuales pueden interpretarse como máxima verosimilitud bajo supuestos de modelo específicos. Sin embargo, si los datos provienen de una encuesta compleja, la matriz \(\mathbf{X}'\mathbf{X}\) no refleja por sí sola el mecanismo de selección de la muestra.

\subsection{Método de máxima pseudo-verosimilitud}\label{sec-mpv-formal}

Cuando el diseño muestral es complejo, los datos observados no satisfacen necesariamente los supuestos de independencia e idéntica distribución. Algunas unidades pueden representar a muchas personas de la población y otras a pocas. Además, dos unidades seleccionadas dentro del mismo conglomerado pueden estar correlacionadas. En este contexto, usar directamente la log-verosimilitud ordinaria puede producir estimaciones que describen la muestra, pero no la población objetivo.

La máxima pseudo-verosimilitud, desarrollada a partir de los resultados de \citet{binder1983variances}, incorpora los pesos de muestreo en la log-verosimilitud. Para unidades observadas \(k \in s\), se define

\[
\ell_p(\theta)
= \sum_{k \in s} w_k \log f(y_k;\theta)
\]

El peso \(w_{hik}=1/\pi_{hik}\) indica cuántas unidades de la población representa la observación \(y_{hik}\). Por eso, la contribución de una unidad con baja probabilidad de inclusión se amplifica en la pseudo-verosimilitud. Las ecuaciones de estimación resultantes son

\[
U_p(\theta)
= \sum_{k \in s} w_k \ u_k(\theta)
= 0
\]

La estimación MPV conserva la lógica de máxima verosimilitud, pero reemplaza el \emph{score} ordinario por un \emph{score} ponderado. La solución \(\hat{\theta}_{MPV}\) es el valor que hace que la suma ponderada de contribuciones individuales sea igual a cero. Para una distribución Bernoulli, la pseudo-log-verosimilitud es

\[
\ell_p(\theta)
= \sum_{k \in s} w_k
\left[
y_k\log(\theta)
+(1-y_k)\log(1-\theta)
\right].
\]

Al derivar e igualar a cero, se obtiene que:

\[
\frac{\partial l_p(\theta)}{\partial \theta}
=
\sum_{k \in s} w_k
\left[
\frac{y_k}{\theta}
-\frac{1-y_k}{1-\theta}
\right]
=0
\]

Cuya solución es:

\[
\hat{\theta}_{MPV}
=
\frac{\sum_{k \in s}w_k y_k}
     {\sum_{k \in s}w_k}
= \hat{p}_d
\]

Por tanto, para una variable binaria, la MPV conduce al estimador ponderado de la proporción, equivalente al estimador de Hájek. Este resultado conecta directamente la teoría de pseudo-verosimilitud con los estimadores de proporciones presentados en los capítulos de variables categóricas.

Asimismo, para un modelo de regresión lineal múltiple, la pseudo-log-verosimilitud ponderada implica minimizar una suma de cuadrados ponderada:

\[
(\mathbf{y}-\mathbf{X}\beta)'\mathbf{W}(\mathbf{y}-\mathbf{X}\boldsymbol{\beta}),
\]

donde

\[
\mathbf{W}=
\begin{pmatrix}
w_1 & 0 & \cdots & 0\\
0 & w_2 & \cdots & 0\\
\vdots & \vdots & \ddots & \vdots\\
0 & 0 & \cdots & w_n
\end{pmatrix}.
\]

La solución es el estimador de mínimos cuadrados ponderados:

\[
\hat{\boldsymbol{\beta}}_{MPV}
= (\mathbf{X}'\mathbf{W}\mathbf{X})^{-1}\mathbf{X}'\mathbf{W}\mathbf{y}
\]

Este resultado extiende de manera natural el estimador ordinario al contexto de encuestas complejas: cada observación contribuye al ajuste del modelo en proporción a su peso muestral. No obstante, para la inferencia no basta con ponderar la estimación puntual; también se requiere estimar correctamente la varianza considerando estratos, conglomerados y pesos.

La siguiente simulación de Monte Carlo anidada ilustra la inferencia doble. En el ciclo externo se genera una nueva población bajo el modelo Bernoulli. En el ciclo interno se seleccionan muchas muestras por diseño y se comparan el estimador que incorpora el diseño y el promedio simple de la muestra.

\begin{Shaded}
\begin{Highlighting}[]
\FunctionTok{set.seed}\NormalTok{(}\DecValTok{2026}\NormalTok{)}
\NormalTok{population\_size }\OtherTok{\textless{}{-}} \DecValTok{100}
\NormalTok{theta }\OtherTok{\textless{}{-}} \FloatTok{0.3}
\NormalTok{n\_sim\_population }\OtherTok{\textless{}{-}} \DecValTok{100}
\NormalTok{expected\_design\_estimates }\OtherTok{\textless{}{-}}\NormalTok{ expected\_simple\_estimates }\OtherTok{\textless{}{-}}\NormalTok{ population\_means }\OtherTok{\textless{}{-}}
  \FunctionTok{rep}\NormalTok{(}\ConstantTok{NA}\NormalTok{, n\_sim\_population)}

\NormalTok{cluster\_size }\OtherTok{\textless{}{-}} \FunctionTok{rep}\NormalTok{(}\DecValTok{2}\SpecialCharTok{:}\DecValTok{6}\NormalTok{, }\AttributeTok{each =} \DecValTok{5}\NormalTok{)}
\NormalTok{n\_clusters }\OtherTok{\textless{}{-}} \FunctionTok{length}\NormalTok{(cluster\_size)}
\NormalTok{cluster\_id }\OtherTok{\textless{}{-}} \FunctionTok{rep}\NormalTok{(}\FunctionTok{seq\_len}\NormalTok{(n\_clusters), cluster\_size)}
\NormalTok{size\_center }\OtherTok{\textless{}{-}} \FunctionTok{weighted.mean}\NormalTok{(cluster\_size, cluster\_size)}
\NormalTok{prob\_person }\OtherTok{\textless{}{-}}\NormalTok{ theta }\SpecialCharTok{+} \FunctionTok{rep}\NormalTok{((cluster\_size }\SpecialCharTok{{-}}\NormalTok{ size\_center) }\SpecialCharTok{/} \DecValTok{12}\NormalTok{, cluster\_size)}
\NormalTok{sampled\_cluster\_count }\OtherTok{\textless{}{-}} \FunctionTok{floor}\NormalTok{(n\_clusters }\SpecialCharTok{*} \FloatTok{0.3}\NormalTok{)}

\ControlFlowTok{for}\NormalTok{ (population\_sim }\ControlFlowTok{in} \FunctionTok{seq\_len}\NormalTok{(n\_sim\_population)) \{}
\NormalTok{  outcome }\OtherTok{\textless{}{-}} \FunctionTok{rbinom}\NormalTok{(population\_size, }\DecValTok{1}\NormalTok{, prob\_person)}
\NormalTok{  population\_means[population\_sim] }\OtherTok{\textless{}{-}} \FunctionTok{mean}\NormalTok{(outcome)}
\NormalTok{  cluster\_totals }\OtherTok{\textless{}{-}} \FunctionTok{tapply}\NormalTok{(outcome, cluster\_id, sum)}
\NormalTok{  cluster\_means }\OtherTok{\textless{}{-}} \FunctionTok{tapply}\NormalTok{(outcome, cluster\_id, mean)}

\NormalTok{  design\_estimates }\OtherTok{\textless{}{-}}\NormalTok{ simple\_estimates }\OtherTok{\textless{}{-}} \FunctionTok{rep}\NormalTok{(}\ConstantTok{NA}\NormalTok{, }\DecValTok{100}\NormalTok{)}
  \ControlFlowTok{for}\NormalTok{ (sample\_sim }\ControlFlowTok{in} \FunctionTok{seq\_len}\NormalTok{(}\DecValTok{100}\NormalTok{)) \{}
\NormalTok{    pps\_sample }\OtherTok{\textless{}{-}} \FunctionTok{S.piPS}\NormalTok{(sampled\_cluster\_count, cluster\_size)}
\NormalTok{    sampled\_clusters }\OtherTok{\textless{}{-}}\NormalTok{ pps\_sample[, }\DecValTok{1}\NormalTok{]}
\NormalTok{    sampled\_cluster\_size }\OtherTok{\textless{}{-}}\NormalTok{ cluster\_size[sampled\_clusters]}
\NormalTok{    sampled\_cluster\_total }\OtherTok{\textless{}{-}}\NormalTok{ cluster\_totals[sampled\_clusters]}
\NormalTok{    sampled\_cluster\_mean }\OtherTok{\textless{}{-}}\NormalTok{ cluster\_means[sampled\_clusters]}
\NormalTok{    design\_estimates[sample\_sim] }\OtherTok{\textless{}{-}} \FunctionTok{mean}\NormalTok{(sampled\_cluster\_mean)}
\NormalTok{    simple\_estimates[sample\_sim] }\OtherTok{\textless{}{-}}
      \FunctionTok{sum}\NormalTok{(sampled\_cluster\_total) }\SpecialCharTok{/} \FunctionTok{sum}\NormalTok{(sampled\_cluster\_size)}
\NormalTok{  \}}
\NormalTok{  expected\_design\_estimates[population\_sim] }\OtherTok{\textless{}{-}} \FunctionTok{mean}\NormalTok{(design\_estimates)}
\NormalTok{  expected\_simple\_estimates[population\_sim] }\OtherTok{\textless{}{-}} \FunctionTok{mean}\NormalTok{(simple\_estimates)}
\NormalTok{\}}

\FunctionTok{cbind}\NormalTok{(}
\NormalTok{  theta,}
  \AttributeTok{mean\_population =} \FunctionTok{mean}\NormalTok{(population\_means),}
  \AttributeTok{expected\_pps =} \FunctionTok{mean}\NormalTok{(expected\_design\_estimates),}
  \AttributeTok{bias\_pps =} \FunctionTok{mean}\NormalTok{(expected\_design\_estimates) }\SpecialCharTok{{-}}\NormalTok{ theta,}
  \AttributeTok{expected\_simple =} \FunctionTok{mean}\NormalTok{(expected\_simple\_estimates),}
  \AttributeTok{bias\_simple =} \FunctionTok{mean}\NormalTok{(expected\_simple\_estimates) }\SpecialCharTok{{-}}\NormalTok{ theta}
\NormalTok{)}
\end{Highlighting}
\end{Shaded}

\begin{verbatim}
##      theta mean_population expected_pps bias_pps expected_simple bias_simple
## [1,]   0.3           0.304        0.302  0.00248           0.332      0.0317
\end{verbatim}

El código reproduce las dos fuentes de aleatoriedad. El ciclo externo representa el modelo \(\xi\), porque genera nuevas poblaciones finitas con \texttt{rbinom()}. El ciclo interno representa el diseño \(p\), porque para cada población se seleccionan múltiples muestras con \texttt{S.piPS()}. La salida incluye \texttt{mean\_population}, que es el promedio de las medias poblacionales generadas; \texttt{expected\_pps}, que resume el comportamiento del estimador que incorpora el diseño; y \texttt{expected\_simple}, que resume el promedio simple de la muestra seleccionada. Cuando el diseño se incorpora correctamente, el sesgo Monte Carlo respecto de \texttt{theta} debe ser pequeño. La comparación con \texttt{expected\_simple} muestra por qué la inferencia basada en encuestas requiere pesos y varianzas de diseño, incluso cuando se parte de un modelo estadístico sencillo.

En síntesis, la máxima verosimilitud ordinaria es adecuada cuando las observaciones pueden tratarse como IID bajo un modelo probabilístico. En encuestas complejas, la máxima pseudo-verosimilitud ofrece una extensión natural: conserva la estructura de las ecuaciones de estimación, pero pondera cada contribución individual por su peso muestral. Así, los modelos estadísticos se ajustan de manera más coherente con la población que la encuesta busca representar.

\chapter*{Referencias}\label{referencias}
\addcontentsline{toc}{chapter}{Referencias}

\bibliography{book.bib}

\end{document}
